\documentclass[11pt,a4paper]{article}
\usepackage[margin=25mm]{geometry}
\usepackage{fontspec}
\setmainfont[BoldFont=lmroman10-bold.otf,ItalicFont=lmroman10-italic.otf,
 BoldItalicFont=lmroman10-bolditalic.otf]{lmroman10-regular.otf}
\setsansfont{lmsans10-regular.otf}
\setmonofont[Scale=0.87]{lmmono10-regular.otf}
\usepackage{amsmath,amssymb,amsthm,mathtools}
\usepackage{booktabs,longtable,array,enumitem,microtype}
\usepackage{xcolor}
\definecolor{inkblue}{HTML}{173B52}
\usepackage[colorlinks=true,linkcolor=inkblue,urlcolor=inkblue]{hyperref}
\usepackage{xurl}
\usepackage{bookmark}
\usepackage{fancyhdr}
\pagestyle{fancy}\fancyhf{}
\fancyhead[L]{\small GRH\_NIGHT: consequences and proof gates}
\fancyfoot[C]{\thepage}
\setlength{\headheight}{14pt}
\setlength{\emergencystretch}{3em}
\setlist{nosep}
\newtheorem{proposition}{Proposition}
\newcommand{\A}{A_D}
\newcommand{\M}{M_H}
\title{\textbf{GRH\_NIGHT}\\Consequences, methods, and the Grothendieck programme}
\author{GPT-6 Astra Pro research run\\Model effort setting not exposed}
\date{21 November 2026, 03:\allowbreak{}32 CET}
\begin{document}
\maketitle
\begin{abstract}
This finite catalogue records 199 consequences, refinements, counterexamples and explicit gates.\allowbreak{}The input $A_D$ is the Alp\"oge RH/Dirichlet-GRH theorem designated in the task. The input $M_H$ is Masullo\textquotesingle s claimed universal unitary-Hecke GRH and is always conditional. Classical proofs are located rather than reproduced; new formulations receive complete proofs. ``Proved here means established in this run, not a claim of publication novelty. The source formalization was not rebuilt. The full 940-page draft was not certified. Explicit source errors are distinguished from the conditional endpoint assumption. IN PROGRESS.
\end{abstract}
\tableofcontents
\newpage
\section*{How to read the evidence}
\addcontentsline{toc}{section}{How to read the evidence}
The six-column \path{CONSEQUENCES.tsv} is the authoritative statement and
dependency catalogue. \path{QUESTIONS.tsv} records the state of every
question treated. \path{CHECK_EVENTS.tsv} records actual check events,
including who performed them and whether a proof was read, independently
derived, checked by a finite computation, or carried from the supplied
record. A citation verified at the statement level is not a claim that the
whole classical proof was rechecked. All new proof texts in the final
checked part have a complete-reading event; computational examples are
not substituted for a general proof.

All critical-line assertions concern the stated nontrivial divisor, with
Tate shifts, trivial zeros, poles and omitted Euler factors handled as
specified. A signed Brauer factorization gives zero/pole locations; it
does not supply the nonnegative divisor orders required by Artin
holomorphy. The notation $A$, $A_D$ refers to the same adopted theorem,
and $M$, $M_H$ to the same conditional draft claim.

The original prompt asks for an open-ended exhaustive investigation.
This is a finite research record. Unresolved conjectures and unavailable
proofs are displayed as such rather than given invented resolutions.
Archive locators distinguish PDF page, repository footer and manuscript
numbering where they differ. The old papers absent from the supplied ZIP
are cited through their attached return records, without a claim of fresh
whole-proof verification.

\part{Checked findings and located consequences}

% BEGIN number_theory
\section{Number-theory consequences with the hypotheses separated}
\label{sec:number-theory}

Throughout this section $A$ denotes the user-designated Alp\"oge theorem
for the Riemann zeta function and Dirichlet L-functions.  It is an adopted
cited input, not a claim that this run rebuilt the unseen Lean dependencies.
$M$ denotes the conditional all-number-field unitary Hecke-GRH assertion
of Masullo, \S1.2.  Every use of $M$ below remains conditional.  The
accompanying catalogue is \path{work/number_theory/rows.json}; its
literature records include retrieval evidence and distinguish a checked
statement from a reread of an entire classical proof.

There are three different geometric translations.  Rational primes are
the closed points of $\operatorname{Spec}\mathbb Z$; prime ideals are the
closed points of $\operatorname{Spec}\mathcal O_K$; and a conjugacy condition
on Frobenius is a condition on a finite \'{e}tale cover of an open arithmetic
curve.  The last dictionary is the Galois-category equivalence in SGA~1,
Expos\'e~V, \S7.  The class group is
$\operatorname{Pic}(\operatorname{Spec}\mathcal O_K)$ and its class fields
are finite \'{e}tale covers.  These translations identify the objects to
which the estimates apply.  They do not turn number-field GRH into a
consequence of finite-field purity.  The inherited cote-1 trace programme
is a comparison of explicit formulas; no new archive scan is certified by
this section.  Exact archive work belongs to the geometry section.

\subsection{The classical arithmetic bounds being used}

\paragraph{NT-P01--NT-P04: rational primes and gaps.}
Schoenfeld, \emph{Sharper bounds for the Chebyshev functions
$\theta(x)$ and $\psi(x)$. II} (1976), Theorem~10, printed p.~337,
equations (6.2) and (6.3), gives, under RH,
\[
 |\psi(x)-x|<\frac{\sqrt{x}\log^2x}{8\pi}\quad(x\ge73.2),\qquad
 |\theta(x)-x|<\frac{\sqrt{x}\log^2x}{8\pi}\quad(x\ge599).
\]
Equation (6.1) replaces $\log^2x$ by $(\log x-2)\log x$ for
$x\ge23\cdot10^8$.  Corollary~1 gives
$|\pi(x)-\operatorname{li}(x)|<\sqrt{x}\log x/(8\pi)$ for
$x\ge2657$.  Here $\operatorname{li}$ has its usual normalization.
The constants in Theorem~10 were checked in the primary indexed text.
An independently retrieved primary proof locator for the same
prime-counting bound is Johnston (2022), Theorem~3.2, equation (3.6),
\emph{The Ramanujan Journal} \textbf{59}, 1307--1321.
It requires RH only to a height $T$ satisfying
$9.06\sqrt{x/\log x}/\log\log x\le T$; under $A$ one may choose
such a $T$ for each $x$.  Its \S2 explicitly says its main theorem applies
also to the ordinary, non-half-weight counting functions.  Equation (3.3)
there also improves the threshold for the displayed $\psi$ inequality
to $x\ge59$.  Thus the inaccessible original Corollary~1 page is no
longer a proof-source gate for NT-P03.

Dudek--Greni\'e--Molteni, \emph{Primes in explicit short intervals on RH}
(2016), Theorem~1.1, author PDF p.~1, gives a prime in
\[
 (x-c\sqrt{x}\log x,\ x+c\sqrt{x}\log x),
 \qquad c=\tfrac12+2/\log x,\qquad x\ge2.
\]
The same theorem gives at least $\sqrt{x}$ primes after replacing $c$ by
$c+1$.  These results imply the classical
$p_{n+1}-p_n=O(\sqrt{p_n}\log p_n)$ consequence of RH.  Neither this
bound nor the assertion of RH supplies a proof here of Cram\'er's
$O(\log^2p_n)$ conjecture, twin primes, or the binary Goldbach conjecture.
Those problems require additional information about prime correlations.

\paragraph{NT-P05--NT-P08: progressions and nonresidues.}
Greni\'e--Molteni--Perelli (2016), Theorem~1, author PDF p.~2, proves that
absolute $x_0,c_1,c_2>0$ work in
\[
 \pi(x+h;q,a)-\pi(x;q,a)\ge
 c_2\frac{h}{\varphi(q)\log x},\qquad
 c_1\varphi(q)\sqrt{x}\log x\le h\le x,\quad x\ge x_0,
\]
assuming the Dirichlet GRH modulo $q$ and $(a,q)=1$.  The following
sentence in that source gives the asymptotic if the ratio to the lower
scale tends to infinity.  We have not converted the unspecified absolute
constants into numerical constants.

Lamzouri--Li--Soundararajan (LLS), corrected
\path{arXiv:1309.3595v3}, is used for the sharp explicit conductor
statements.  Corollary~1.2, author p.~3, gives the least prime in a reduced
class modulo $q>3$ at most $(\varphi(q)\log q)^2$.
Corollary~1.1, p.~2, gives a quadratic nonresidue below $(\log q)^2$ for
prime $q\ge5$.  Theorem~1.1(1), p.~2, gives the proper-subgroup bound
recorded with its explicit functions $A(q),B(q)$ in NT-P08.
These are the already carried IA02 consequences, with the corrected
version and hypotheses rechecked; they are not new findings of this run.
Geometrically the progression is a Frobenius class in a cyclotomic cover,
and the nonresidue is detected by a finite rank-one character.

\paragraph{NT-P09: the least primitive root is a different quantifier.}
Shoup, \emph{Searching for primitive roots in finite fields} (1992),
Theorem~3, author PDF p.~4, gives under Dirichlet GRH
\[
 g(p)\ll r^4(\log r+1)^4(\log p)^2,
 \qquad r=\omega(p-1).
\]
The printed theorem was checked on the page image.  Here is the short
simplification to $g(p)\ll\log^6p$.  If $q_1<\cdots<q_r$ are the distinct
prime divisors of $p-1$, then $q_j\ge j+1$ and therefore
$p-1\ge\prod_{j=1}^r(j+1)$.  For $r\ge4$ the last at least $r/2$ factors
are at least $r/2$, so
$\log(p-1)\ge(r/2)\log(r/2)$.  Consequently
$r(\log r+1)\ll\log p$; enlarging the absolute constant handles $r<4$.
Substitution proves the assertion.  This finds a small generator for a
fixed $\mathbb F_p^\times$.  The density theorem for one fixed integer
as $p$ varies uses different, generally nonabelian, fields (NT-A01).

\subsection{A quantitative Frobenius trace transfer}

\paragraph{NT-C01 and NT-C02: the precise cited Chebotarev input.}
Let $L/K$ be a finite Galois extension, $G=\operatorname{Gal}(L/K)$,
$n_L=[L:\mathbb Q]$, and $\Delta_L=|\operatorname{disc}L|$.
For a conjugacy class $C$ put $\delta_C=|C|/|G|$.  For unramified
primes only, define
\[
 \psi_C(x)=\sum_{\substack{\mathfrak p\text{ unramified},\ m\ge1\\
                    (N\mathfrak p)^m\le x}}
 \mathbf1_C(\operatorname{Frob}_{\mathfrak p}^{m})\log N\mathfrak p,
 \qquad
 \pi_C(x)=\sum_{\substack{\mathfrak p\text{ unramified}\\N\mathfrak p\le x}}
             \mathbf1_C(\operatorname{Frob}_{\mathfrak p}).
\]
Set
\begin{align*}
 E_L(x)&=\sqrt{x}\left[
       \left(\frac{\log x}{2\pi}+2\right)\log\Delta_L+
       \left(\frac{\log^2x}{8\pi}+2\right)n_L\right],\qquad x\ge1,\\
 B_L(x)&=\sqrt{x}\left[
       \left(\frac1{2\pi}+\frac3{\log x}\right)\log\Delta_L+
       \left(\frac{\log x}{8\pi}+\frac1{4\pi}+\frac6{\log x}\right)n_L
                         \right],\qquad x\ge2,\\
 \operatorname{Li}_2(x)&=\int_2^x\frac{dt}{\log t}.
\end{align*}
Greni\'e--Molteni, \emph{An explicit Chebotarev density theorem under GRH}
(2019), Theorem~1.1 and Corollary~1.2, author PDF pp.~2--3, give
\[
 |\psi_C(x)-\delta_Cx|\le\delta_C E_L(x),\qquad
 |\pi_C(x)-\delta_C\operatorname{Li}_2(x)|\le\delta_C B_L(x),
\tag{NT-C}
\]
under RH for $\zeta_L$.  The theorem explicitly includes the unramified
version of $\psi_C$.  Its constants and ranges were checked on the page
images.  In particular Artin holomorphy is not an additional hypothesis.
$M$ supplies this input for every $L$; $A$ supplies it if $L/\mathbb Q$
is abelian by the factorization into primitive Dirichlet functions.

\paragraph{NT-F01. Statement and full proof.}
For a complex class function $f$ on $G$ put
\[
 \mu_f=\frac1{|G|}\sum_{g\in G}f(g),\qquad
 \|f\|_{1,G}=\frac1{|G|}\sum_{g\in G}|f(g)|,
\]
and define $\Psi_f$ and $\Pi_f$ by replacing $\mathbf1_C$ in the two
sums above by $f$.  Then
\[
 |\Psi_f(x)-\mu_fx|\le\|f\|_{1,G}E_L(x),\qquad
 |\Pi_f(x)-\mu_f\operatorname{Li}_2(x)|\le\|f\|_{1,G}B_L(x).
\tag{NT-F01}
\]
If $f=\chi_V$ is the character of a finite-dimensional complex
representation, then $\mu_f=\dim V^G$ and
$\|f\|_{1,G}\le\sqrt{\langle\chi_V,\chi_V\rangle_G}$.

To prove the estimates, choose one representative $g_C$ from each
conjugacy class.  The finite identities
\[
 f=\sum_C f(g_C)\mathbf1_C,
 \qquad \mu_f=\sum_C\delta_Cf(g_C)
\]
give $\Psi_f-\mu_fx=\sum_Cf(g_C)(\psi_C-\delta_Cx)$.
The triangle inequality and (NT-C) bound its modulus by
$E_L(x)\sum_C\delta_C|f(g_C)|=E_L(x)\|f\|_{1,G}$.
The same finite argument proves the assertion for $\Pi_f$.
The averaging operator $|G|^{-1}\sum_g\rho(g)$ is the projection to $V^G$,
so its trace is $\mu_{\chi_V}=\dim V^G$.  Cauchy--Schwarz for the
probability measure assigning mass $1/|G|$ to each group element yields
\[
 \left(\frac1{|G|}\sum_g|\chi_V(g)|\right)^2
 \le\frac1{|G|}\sum_g|\chi_V(g)|^2
 =\langle\chi_V,\chi_V\rangle_G.
\]
This proves every assertion.  This is a linear functional estimate on the
complexified representation ring; virtual differences need not be
nonnegative, so their error is controlled by the displayed norm.

\paragraph{NT-F02. Finite \texorpdfstring{\'{e}}{e}tale point counts, with the boundary.}
Let $U$ be a nonempty open subscheme of
$\operatorname{Spec}\mathcal O_K$, and let $X\to U$ be finite \'{e}tale
of degree $d$.  Let $S_X$ be the set of geometric points of its generic
fibre, let $L/K$ be its finite Galois splitting field, and let $r$ be the
number of $G$-orbits on $S_X$, equivalently the number of generic connected
components.  Let $S$ be the finite set of primes not in $U$ which are
unramified in $L$.  Under RH for $\zeta_L$,
\[
 \left|\sum_{\substack{\mathfrak p\in U,\ m\ge1\\(N\mathfrak p)^m\le x}}
       \#X(k_{\mathfrak p^m})\log N\mathfrak p-rx\right|
 \le rE_L(x)+d|S|\log x,\qquad x\ge1.
\tag{NT-F02}
\]
Here $k_{\mathfrak p^m}$ denotes the degree-$m$ extension of
$k_{\mathfrak p}$.  In particular, when no additional unramified prime is
removed, the error is $rE_L(x)$, independent of the degree $d$.

For the proof, finite \'{e}taleness implies that inertia at a prime of $U$
acts trivially on the generic geometric fibre.  The action of
$\operatorname{Gal}(L/K)$ on this fibre is faithful by the definition of
the splitting field, so $L/K$ is unramified at every prime of $U$.
Let $\chi(g)=\#S_X^g$ be the permutation character.  The elementary
Frobenius fixed-point formula for a finite \'{e}tale scheme over a finite
field is
$\#X(k_{\mathfrak p^m})=\chi(\operatorname{Frob}_{\mathfrak p}^{m})$:
its geometric points rational over that field are precisely the points
fixed by the $m$th Frobenius power.  Also $0\le\chi(g)\le d$.
Double counting pairs $(g,s)$ with $gs=s$ gives
\[
 \sum_{g\in G}\chi(g)=\sum_{s\in S_X}|G_s|=r|G|.
\]
Thus both $\mu_\chi$ and $\|\chi\|_{1,G}$ equal $r$.  NT-F01 bounds
the sum over every unramified prime by $rx$ with error $rE_L(x)$.
For each omitted prime, the nonnegative lost contribution is at most
\[
 d\left\lfloor\frac{\log x}{\log N\mathfrak p}\right\rfloor
          \log N\mathfrak p\le d\log x.
\]
If $N\mathfrak p>x$ the contribution is zero.  Summing this bound over
$S$ and applying the triangle inequality proves (NT-F02), including
$x=1$.  Notice that replacing the nonnegative character norm by the square
root of its second moment would give a weaker bound; the orbit count is
the relevant parameter here.

\paragraph{NT-F03. Splitting types.}
For a partition $\lambda$ of $d=|S_X|$, let $D_\lambda\subseteq G$ be
the elements with cycle partition $\lambda$ on $S_X$, and put
$\delta_\lambda=|D_\lambda|/|G|$.  This is a union of conjugacy classes.
Sum the second inequality (NT-C) over these classes and use the triangle
inequality.  Their densities sum to $\delta_\lambda$, proving
\[
 \left|\#\{\mathfrak p:\ N\mathfrak p\le x,\ \mathfrak p
 \text{ unramified},\ \operatorname{Frob}_{\mathfrak p}\in D_\lambda\}
 -\delta_\lambda\operatorname{Li}_2(x)\right|
 \le\delta_\lambda B_L(x).
\]
If $D_\lambda$ is empty both sides before the absolute value are zero.
For the canonical finite \'{e}tale model of the generic algebra over the
unramified locus, each Frobenius orbit is one closed point of its fibre,
and its orbit length is that point's residue degree.  Hence this cycle
partition is exactly its splitting type.  For the originally specified
model $X$ over $U$, restrict the counted primes to $U$; at most $|S|$
terms have been removed, so the same estimate holds with the additional
error $|S|$.  This proves both interpretations without assigning a fibre
of $X$ outside its specified base.

\paragraph{NT-F04. Avoiding a prescribed finite set.}
Let $D\subseteq G$ be a nonempty conjugacy-stable set of density
$\delta_D$, and let $S$ be any finite set of primes of $K$.  The same
finite-union argument gives
$\pi_D(x)\ge\delta_D(\operatorname{Li}_2(x)-B_L(x))$.
At most $|S|$ of the counted primes belong to $S$.  If
\[
 \delta_D\bigl(\operatorname{Li}_2(x)-B_L(x)\bigr)>|S|,
\]
the count outside $S$ is strictly positive.  It is an integer, so an
eligible prime exists.  This proves the avoidance assertion without adding
an unsupported arbitrary-exclusion clause to Bach--Sorenson's least-prime
theorem.  That theorem is NT-C03, Theorem~5.1, printed p.~1729, with
RH hypotheses specified in its \S2, p.~1720.

\subsection{Class numbers, Kummer towers, and Euclidean algorithms}

\paragraph{NT-H01--NT-H04: the regulator normalization.}
Write
\begin{align*}
 U(q)&=2e^\gamma\left(\log\log q-\log2+\tfrac12+
                                      \frac1{\log\log q}\right),\\
 V(q)&=\frac{12e^\gamma}{\pi^2}\left(\log\log q-\log2+\tfrac12+
                   \frac1{\log\log q}+\frac{14\log\log q}{\log q}\right).
\end{align*}
LLS Theorem~1.5, author pp.~3--4, gives
$V(q)^{-1}\le|L(1,\chi)|\le U(q)$ for primitive conductor
$q\ge10^{10}$.  Its Corollary~1.3 is the already carried imaginary
quadratic bound in NT-H02.  The inequality is not a two-sided estimate of
order $\log\log q$: the reciprocal lower bound must be retained.

Here is the full deduction of NT-H03, including its normalization.
For $D>0$ fundamental put $K=\mathbb Q(\sqrt D)$ and let $R_D$ be the
regulator of the full unit group modulo its torsion, so
$R_D=\log\varepsilon$ for a fundamental unit $\varepsilon>1$, regardless
of its norm's sign.  The analytic class-number formula is
\[
 \mathop{\rm Res}_{s=1}\zeta_K(s)=\frac{2h_DR_D}{\sqrt D}.
\]
An exact independent locator is Milne, \emph{Class Field Theory}, v4.03,
Chapter~VI, Corollary~2.12 and Example~2.13, printed p.~191.
Since $\zeta_K(s)=\zeta(s)L(s,\chi_D)$ and $\zeta$ has residue one,
$L(1,\chi_D)=2h_DR_D/\sqrt D>0$.
Multiplying the two inequalities in NT-H01 by $\sqrt D/2$ proves
\[
 \frac{\sqrt D}{2V(D)}\le h_DR_D\le\frac{\sqrt D}{2}U(D),
 \qquad D\ge10^{10}.
\]
This completes the short proof; no bound on $R_D$ was suppressed.

For completeness, the already inherited abelian-field formulation
NT-H04 is the following precise product calculation.  If $K/\mathbb Q$
is abelian with Dirichlet-character group $X_K$, then
\[
 h_KR_K=\frac{w_K\sqrt{\Delta_K}}{2^{r_1}(2\pi)^{r_2}}
              \prod_{\substack{\chi\in X_K\\\chi\ne1}}L(1,\chi).
\tag{NT-H04}
\]
Indeed Artin factorization gives the product for $\zeta_K/\zeta$,
and taking the residue at one gives this formula by the same class-number
formula.  Its left side and the prefactor are positive, so taking moduli
turns the product into $\prod|L(1,\chi)|$.  For any factors whose
conductors are at least $10^{10}$ insert $U$ and $1/V$; retain the finite
remaining product exactly or bound it separately.  This proves every
claimed product bound.  Dividing by an unestimated regulator does not
prove the corresponding bound on $h_K$.  This is the correction already
made in RETURN\_ZETA\_TWO IA02 and IA09, not a new result of this run.

\paragraph{NT-A01 and NT-A02: the nonabelian input in Artin's primitive-root problem.}
For a fixed integer $g\ne-1$ that is not a square, put
$K_n=\mathbb Q(\zeta_n,g^{1/n})$.  Hooley's classical theorem gives a
positive density
\[
 \delta(g)=\sum_{n\ge1}\frac{\mu(n)}{[K_n:\mathbb Q]}
\]
of primes where $g$ is a primitive root, assuming RH for these Kummer
fields.  We cite Hooley, \emph{On Artin's conjecture} (1967), the
unnumbered theorem on pp.~219--220.  The original internal page has not
been freshly obtained here; a fully inspected precise proof locator for
the density implication is Lenstra, \emph{On Artin's conjecture and
Euclid's algorithm in global fields} (1977), Theorem~(3.1), printed
p.~206, specialized to $K=F=\mathbb Q$, $W=\langle g\rangle$, $k=1$,
with density (2.15), p.~206, and positivity criterion (4.1), p.~207.
The series uses actual field degrees; replacing it by the universal
Artin constant for every $g$ loses the entanglement correction.

Lenstra's version also imposes a conjugacy-stable Frobenius condition
$C\subseteq\operatorname{Gal}(F/K)$ and asks that the residue image of
a finitely generated positive-rank $W\subseteq K^\times$ have index
dividing $k$.  Its auxiliary fields are formed from
$K(\zeta_{q(\ell)},W^{1/q(\ell)})$ with
$q(\ell)=\ell^{v_\ell(k)+1}$, taking the compositum for primes dividing
a squarefree $n$.  Theorem~(3.1) assumes RH for these fields' zeta
functions.  Theorem~(4.1) gives the finite obstruction criterion for
positive density.  $M$ supplies the hypotheses; $A$ alone does not
supply them because the Kummer fields need not be abelian over
$\mathbb Q$.  In geometric terms this is a tower of Kummer torsors of
$\mathbb G_m$, not merely the cyclotomic torsion cover.

\paragraph{NT-A03: a PID becomes Euclidean under the stated Kummer GRH.}
Lenstra's Theorem~(9.1), printed p.~218, states the result recorded in
NT-A03 for the ring of $S$-integers with $|S|\ge2$ and trivial ideal
class group.  Equation (9.2) gives its least Euclidean function.
The original page image confirms the weak inequality $|S|\ge2$;
the OCR's apparent strict inequality was not adopted.  The result is a
classical cited theorem with an exact proof locator, not a new proof of
the Euclidean algorithm.  The independent PID hypothesis and the
distinction from a norm-Euclidean algorithm remain in the catalogue.

\subsection{Growth, cancellation, and zero statistics}

\paragraph{NT-L02 and NT-L03. The fixed-character Lindel\"of deduction.}
The classical input is Simoni\v c, \emph{Estimates for L-functions in the
critical strip under GRH with effective applications} (2023), Theorem~1,
equations (4)--(6), proved in its \S3.  For a function $F$ in the Selberg
class with polynomial Euler product, zero-free to the right of the
critical line, it implies, for fixed $\delta>0$ and large $|t|$,
\[
 \log|F(\sigma+it)|,\qquad -\log|F(\sigma+it)|
 \ll_{F,\delta}(\log|t|)^{1-2\delta}\log\log|t|,
 \quad \tfrac12+\delta\le\sigma\le1.
\tag{NT-L-log}
\]
For $\delta\ge1/2$ one uses instead a smaller positive $\delta$ and
the absolutely convergent Euler product on $\sigma>1$.
The modulus bound for $\log F$ in the cited theorem gives both signs of
its real part.  Dividing the right side of (NT-L-log) by $\log|t|$
shows that it tends to zero for fixed $0<\delta<1/2$.  Hence, for every
$\epsilon>0$,
\[
 F(\sigma+it),\quad F(\sigma+it)^{-1}
 \ll_{F,\delta,\epsilon}(1+|t|)^\epsilon
 \qquad(\sigma\ge\tfrac12+\delta),
\tag{NT-L-right}
\]
with pole neighbourhoods omitted if $F$ has a pole.  The finite range of
$t$ is absorbed into the constant on any closed set avoiding poles and
zeros.  On $\sigma\ge1+\delta$ the Euler product gives the required
uniform bound directly; compact intermediate strips are already covered
by (NT-L-log) or the cited theorem's adjacent range.

Here is the passage to the critical line, often suppressed in invoking
Lindel\"of.  Fix $0<\delta<1/2$.  If $F$ has degree $d$, its functional
equation and Stirling's formula imply
\[
 |F(\tfrac12-\delta+it)|
 \ll_{F,\delta}(1+|t|)^{d\delta}
       |F^\vee(\tfrac12+\delta-it)|.
\]
The dual has the same zero-free hypothesis, so NT-L-right bounds this by
$O((1+|t|)^{d\delta+\eta})$ for every $\eta>0$.  The right boundary
$\sigma=1/2+\delta$ has exponent $\eta$.
The Phragm\'en--Lindel\"of principle on this strip gives at its midpoint
\[
 |F(\tfrac12+it)|\ll_{F,\delta,\eta}
                       (1+|t|)^{d\delta/2+\eta}.
\]
Its hypotheses hold because the standard functional equation and
finite-order analytic continuation give polynomial growth in each fixed
vertical strip; if there is a pole, apply the principle above and below
a fixed height containing it.  Choose $\delta$ and $\eta$ so that
$d\delta/2+\eta<\epsilon$.  The remaining bounded interval is compact,
and the critical line contains no pole for these L-functions.  This
proves $F(1/2+it)\ll_{F,\epsilon}(1+|t|)^\epsilon$.

For primitive Dirichlet functions the hypotheses are classical, and $A$
is the zero-free input; this proves NT-L02.  For a fixed non-pure-norm
unitary Hecke character, the Hecke L-function is entire of finite order,
its gamma shifts have nonnegative real parts, and its degree is
$[K:\mathbb Q]$.  On grouping its ideal Euler factors by rational primes,
one obtains a polynomial Euler product of that degree with local roots
of modulus at most one.  The Dirichlet coefficients are bounded by
$d_{[K:\mathbb Q]}(n)=O_\epsilon(n^\epsilon)$: each local coefficient
is bounded by the corresponding coefficient of $(1-p^{-s})^{-[K:\mathbb Q]}$.
Thus it meets the explicit axioms in Simoni\v c \S2.1, and $M$ supplies
zero-freeness.  If the character is a pure unitary norm power, its
L-function is a vertical translate of $\zeta_K$, to which the same
argument applies before translation.  Imprimitive Euler factors, when
included in the convention, form a fixed finite product; they and their
inverses are bounded on each $\sigma\ge1/2+\delta$, since their roots
have modulus at most one.  This proves NT-L03 in all the conventions just
specified.  Constants may depend on the field, character, infinity type,
and $\epsilon$; this proof makes no uniform conductor-aspect claim.

\paragraph{NT-L05 and NT-L06. Mertens, Liouville, and reciprocal tails.}
Titchmarsh, second edition, Theorem~14.25(C), printed p.~370, is the exact
classical equivalence
$\mathrm{RH}\Longleftrightarrow M(x)=O_\epsilon(x^{1/2+\epsilon})$
for every $\epsilon>0$.  We use that theorem rather than reproduce its
contour proof.  The following short deductions are included in full.
For a prime power $p^a$,
\[
 \sum_{d^2\mid p^a}\mu(p^a/d^2)=
 \begin{cases}1,&a\text{ even},\\-1,&a\text{ odd},\end{cases}
\]
because only the term with remaining exponent zero or one is nonzero.
Both sides are multiplicative, hence
$\lambda(n)=\sum_{d^2\mid n}\mu(n/d^2)$ for every $n$.
Consequently
\[
 \left|\sum_{n\le x}\lambda(n)\right|
 \le\sum_{d\le\sqrt x}|M(x/d^2)|
 \ll_\epsilon x^{1/2+\epsilon}
                \sum_{d\ge1}d^{-1-2\epsilon}
 \ll_\epsilon x^{1/2+\epsilon}.
\]
This proves Liouville cancellation, with the dependence on $\epsilon$
retained in the convergent sum.

Let $s=\sigma+it$ and choose $\eta>0$ with $1/2+\eta<\sigma$.
Partial summation for $y>x$ gives
\[
 \sum_{x<n\le y}\frac{\mu(n)}{n^s}
 =M(y)y^{-s}-M(x)x^{-s}
     +s\int_x^y M(u)u^{-s-1}\,du.
\]
The first term tends to zero as $y\to\infty$, and the integral converges
absolutely after using $|M(u)|\ll_\eta u^{1/2+\eta}$.  Its modulus is at
most a constant times
$|s|x^{1/2+\eta-\sigma}/(\sigma-1/2-\eta)$.
Together with the boundary term this proves the stated tail estimate.
It also proves local uniform convergence of the reciprocal Dirichlet
series on $\operatorname{Re}s>1/2$ by choosing $\eta$ below the distance
of a compact set from the boundary.  It agrees with $1/\zeta(s)$ on
$\operatorname{Re}s>1$ by the absolutely convergent Euler product and
therefore throughout that half-plane by analytic continuation.  At
$s=1$, $1/\zeta$ has value zero.  Hence $\sum\mu(n)/n=0$, as asserted.

\paragraph{NT-L04, NT-L07, NT-L08: what moments and zero density say.}
Harper (2013), Theorem~1, author p.~1, gives
$\int_T^{2T}|\zeta(1/2+it)|^{2k}dt\ll_kT(\log T)^{k^2}$ for fixed
$k\ge0$ under RH.  Its proof does not supply the conjectural leading
constant for every $k$.  Simoni\v c (2023), Theorem~2, equations (7),(8),
provides the explicit effective Mertens constants in NT-L07.
For NT-L08, if $\sigma>1/2$, every nontrivial zero covered by $A$ or
conditionally $M$ has real part $1/2<\sigma$, so the indicated counting
set is empty.  At $\sigma=1/2$ this reasoning says nothing about
multiplicity.  Trivial zeros from gamma factors and zeros of removed
finite Euler factors are not silently included in this definition.

\paragraph{NT-Z01 and NT-Z02: the exact pair-correlation range.}
For zeta zero ordinates $0<\gamma,\gamma'\le T$, counted with
multiplicity, put
\[
 F(\alpha,T)=\left(\frac{T}{2\pi}\log T\right)^{-1}
 \sum_{\gamma,\gamma'}T^{i\alpha(\gamma-\gamma')}
                    \frac4{4+(\gamma-\gamma')^2}.
\]
Montgomery (1973), \S1, unnumbered theorem, printed p.~181, proves under
RH the formula in NT-Z01 uniformly on each $0\le\alpha\le1-\epsilon$.
Corollary~2, p.~182, gives the proportion $2/3+o(1)$ of simple zeros.
The extension of the pair-correlation formula beyond that range is
explicitly a conjecture in that paper.  Adopting RH removes the theorem's
hypothesis; it does not replace its range by all $\alpha$.

\paragraph{NT-Z03--NT-Z05: low zeros of the prime-modulus Dirichlet family.}
For prime $q$ average over the $q-2$ nonprincipal characters, with scaled
ordinates $x_{\chi,j}=\gamma_{\chi,j}\log q/(2\pi)$.
Hughes--Rudnick, \emph{Linear statistics of low-lying zeros of
L-functions} (2003), Theorem~3.1, author p.~5, gives the unitary mean for
admissible even test functions with Fourier support in $[-2,2]$.
Theorem~4.1, p.~7, gives variance
$\int_{-1}^1|u|\widehat f(u)^2du$ with support in $[-1,1]$.
Theorem~5.1, p.~9, gives the Gaussian moments only for
$m<2/\alpha$ when the support is in $[-\alpha,\alpha]$.
These identities are unconditional with complex $\gamma$ allowed;
$A$ supplies the real-ordinate interpretation.  They are not an
unrestricted all-moment Gaussian theorem for one fixed test function.

In that source's ordering on p.~1, $\gamma_{\chi,1}$ is the first
\emph{nonnegative} ordinate, so a central zero is included with value zero.
With exactly this convention, Corollary~8.2, p.~20, gives
$\liminf_q\min_{\chi\ne\chi_0}x_{\chi,1}\le1/4$ under $A$.
Theorem~8.3, pp.~20--21, gives a positive proportion for
$x_{\chi,1}<\beta$ if $\beta\ge0.633$.
No strictly-positive-zero variant was inferred by discarding central
zeros.  The geometric comparison is with finite-field family monodromy
and its symmetry group; it does not identify the unrestricted
number-field limiting distribution.

\paragraph{NT-Z06 and NT-Z07: the additional gates.}
No argument in the cited sources proves full GUE, simplicity of every
zero, or central nonvanishing for every quadratic Dirichlet function
from zero location alone.  We record the specific missing correlation
or nonvanishing estimates, without claiming a logical independence
theorem.  For Iwaniec--Luo--Sarnak, \emph{Low lying zeros of families of
L-functions} (2000), the standing assumption on printed p.~56 explicitly
includes RH for the GL$_2$ functions in the family as well as Dirichlet
GRH.  Theorem~1.1, pp.~58--59, gives support $(-2,2)$ for fixed even
weight and squarefree level tending to infinity.  $A$ supplies its
Dirichlet part.  A universal Hecke statement does not identify a general
non-CM GL$_2$ newform with a rank-one Hecke character.  Thus the GL$_2$
hypothesis remains an additional input in this invocation.

\paragraph{NT-Z08 and NT-Z09: two exact classical equivalents.}
Rodgers--Tao, \path{arXiv:1801.05914v5}, Theorem~1, author p.~3,
proves $\Lambda\ge0$ for the classical de Bruijn--Newman constant.
The definition and RH equivalence in its \S1 give
$A\Rightarrow\Lambda\le0$, hence $\Lambda=0$.
For the deformation as normalized there, real-rootedness holds exactly
for $t\ge\Lambda$, yielding the two assertions in NT-Z08.
Lagarias, \path{arXiv:math/0008177v2}, Theorem~1.1, author p.~1,
is the exact locator for NT-Z09's divisor-sum inequality.
These are cited classical consequences, not additional new proofs of RH.

\subsection{Fixed-character cancellation and the local zero count}

\paragraph{NT-L09 and NT-L10. A full Perron transfer.}
Here is the elementary analytic form used for both rational and ideal
M\"obius sums.  Suppose
$G(s)=\sum_{n\ge1}a_nn^{-s}$ is absolutely convergent for
$\operatorname{Re}s>1$, $|a_n|\le d_D(n)$ for a fixed positive integer
$D$, and $G$ extends holomorphically to $\operatorname{Re}s>1/2$.
Assume that for every $\delta,\eta>0$,
\[
 G(\sigma+it)\ll_{\delta,\eta}(1+|t|)^\eta
 \quad\text{uniformly for }\tfrac12+\delta\le\sigma\le2.
\tag{NT-Perron-input}
\]
Then $\sum_{n\le x}a_n\ll_\epsilon x^{1/2+\epsilon}$ for every
$\epsilon>0$.

We prove the assertion, recording the truncation error as well as the
contour displacement.  It suffices to take $0<\epsilon<1/4$ and
half-integral $x\ge4$, because replacing any real $y$ by
$\lfloor y\rfloor+1/2$ leaves the partial sum unchanged and changes its
size by a bounded factor; the bounded remaining range is absorbed into
the constant.  Put
$c=1+1/\log x$, $T=x^2$, $\delta=\epsilon/4$ and
$\eta=\epsilon/8$.  The classical truncated Perron kernel gives
\[
 \sum_{n\le x}a_n=
 \frac1{2\pi i}\int_{c-iT}^{c+iT}G(s)\frac{x^s}{s}\,ds+R,
\]
with
\[
 |R|\ll\sum_{n\ge1}|a_n|(x/n)^c
           \min\left(1,\frac1{T|\log(x/n)|}\right).
\tag{NT-Perron-error}
\]
This is the usual Perron formula; an exact classical locator is
Titchmarsh, Lemma~3.12 and its proof, printed pp.~61--62.  The kernel
estimate follows by moving the integral of $y^s/s$ to the left or right
according to $y>1$ or $y<1$, taking the residue at zero in the first
case, and integrating the horizontal tails; the alternative trivial
bound handles $|\log y|\le T^{-1}$.

For $n\le x/2$ or $n\ge2x$ the denominator $|\log(x/n)|$ is bounded
below by $\log2$, so that part of (NT-Perron-error) is
\[
 \ll x^cT^{-1}\sum_{n\ge1}d_D(n)n^{-c}
 =x^cT^{-1}\zeta(c)^D\ll xT^{-1}(\log x)^D.
\]
For $x/2<n<2x$, use $d_D(n)\ll_{D,\eta}x^\eta$ and
$|\log(x/n)|\gg |x-n|/x$.  Since $x$ is half-integral,
$\sum_{x/2<n<2x}|x-n|^{-1}\ll\log x$.  Also $(x/n)^c\ll1$,
uniformly for these $x,n$.  This part of the error is therefore
$O(x^{1+\eta}T^{-1}\log x)$.  Both parts are smaller than
$x^{1/2+\epsilon}$.

Move the vertical line to $\sigma_0=1/2+\delta$.  The integrand has no
pole in the intervening rectangle: $G$ is holomorphic there and
$\sigma_0>0$.  On either horizontal segment,
(NT-Perron-input) bounds the integral by
$O(x^cT^{\eta-1})=O(x^{-1+2\eta})$.
On the new vertical segment it gives
\[
 \ll x^{\sigma_0}\left(1+\int_1^Tt^{\eta-1}\,dt\right)
 \ll_\eta x^{1/2+\delta}T^\eta
 =O(x^{1/2+\epsilon/2}).
\]
The bounded part $|t|\le1$ is harmless because $\sigma_0>0$ and $G$
is bounded on its compact segment.  This proves the asserted partial-sum
bound with room in the exponent.

Apply this lemma first to $G(s)=1/L(s,\chi)$.
Its coefficients are $a_n=\mu(n)\chi(n)$, including the usual zero
values at integers not coprime to the conductor.  Under $A$, the
reciprocal is holomorphic on $\operatorname{Re}s>1/2$; the principal
L-function's pole produces a zero of the reciprocal, not a pole.
NT-L02 supplies (NT-Perron-input), so the twisted bound follows.
Character orthogonality, for $(a,q)=1$, is the finite identity
\[
 \sum_{\substack{n\le x\\n\equiv a\ ({\rm mod}\ q)}}\mu(n)
 =\frac1{\varphi(q)}\sum_{\chi\ ({\rm mod}\ q)}
       \overline{\chi(a)}\sum_{n\le x}\mu(n)\chi(n).
\]
Imprimitive functions differ by finitely many Euler factors whose zeros
lie on $\operatorname{Re}s=0$, and their reciprocals are bounded in any
closed half-strip to the right of $1/2$.  Thus each term satisfies the
same fixed-character argument.  A finite sum proves the fixed-$q$
progression bound, with no uniformity in $q$ claimed.

For a unitary Hecke character, expand
\[
 \frac1{L(s,\eta)}
 =\sum_{(\mathfrak a,\mathfrak f)=1}
         \mu_K(\mathfrak a)\eta(\mathfrak a)(N\mathfrak a)^{-s}
 =\sum_{n\ge1}a_nn^{-s}.
\]
For each $n$, $|a_n|$ is at most the number of ideals of norm $n$.
If $[K:\mathbb Q]=D$, comparison of the rational-prime Euler factors
of $\zeta_K$ with $(1-p^{-s})^{-D}$ gives this number at most $d_D(n)$.
Indeed there are at most $D$ prime ideals above $p$, each of residue
degree at least one, and replacing every local factor by a degree-one
factor can only increase its nonnegative coefficients.  Holomorphy and
(NT-Perron-input) follow conditionally on $M$ from NT-L03, including
pure norm twists.  The lemma proves NT-L10.  Averaging over the finite
character group of any fixed ray class group then projects onto one ray
class exactly as above; all sums are over ideals prime to its modulus.
This completes every claimed cancellation statement.

\paragraph{NT-Z10. The fixed infinite-order Hecke zero count.}
Let $\eta$ be any fixed unitary Hecke character of a number field $K$.
Its number of nontrivial zeros, counted with multiplicity, in
\[
 0\le\operatorname{Re}s\le1,\qquad |\operatorname{Im}s-T|\le1
\]
is $O_{K,\eta}(\log(|T|+3))$, unconditionally.  This assertion includes
infinite-order characters; its constant is not asserted uniform in
their infinity types.

Here is a full local argument.  On any fixed vertical strip the finite
Hecke L-function has a polynomial upper bound in $|t|$, away from its
possible pole.  In the present argument this follows directly from the
unconditional first part of Simoni\v c's Theorem~1, equation (2), on
$\sigma\ge1/2$, and from the functional equation and Stirling's formula
on the reflected strip.  The Selberg-class hypotheses for a non-pure-norm
character were verified in NT-L03 without using GRH; a pure norm twist
is a translate of the Dedekind function.  Consequently constants $C,A>0$
exist such that, for large $|T|$,
\[
 \max_{|s-(2+iT)|\le4}|L(s,\eta)|\le C(|T|+3)^A.
\tag{NT-J-upper}
\]
For a pure norm twist the only pole has a fixed imaginary coordinate, so
choose $|T|$ sufficiently large that this disk does not contain it.
In all other cases the function is entire.

The Euler product at the centre gives a lower bound that a general
Schwartz Mellin transform need not have:
\[
 |L(2+iT,\eta)|
 \ge\prod_{\mathfrak p\nmid\mathfrak f}(1+(N\mathfrak p)^{-2})^{-1}
 \ge\prod_{\mathfrak p}(1-(N\mathfrak p)^{-2})
 =\zeta_K(2)^{-1}>0.
\tag{NT-J-lower}
\]
These comparisons follow term by term from $|\eta(\mathfrak p)|=1$
and $1/(1+u)\ge1-u$ for $0\le u<1$, and persist on taking limits of
the absolutely convergent products.

Let $n_T(3)$ count zeros in the disk of radius $3$ about $2+iT$.
Jensen's formula on the radius-$4$ disk, or its limiting-radius version
if a zero lies on the outer circle, implies
\[
 n_T(3)\log(4/3)
 \le\log\max_{|s-(2+iT)|\le4}|L(s,\eta)|
                -\log|L(2+iT,\eta)|.
\]
The right side is $O_{K,\eta}(\log(|T|+3))$ by (NT-J-upper) and
(NT-J-lower).  Every point with $0\le\operatorname{Re}s\le1$ and
$|\operatorname{Im}s-T|\le1$ is at distance at most $\sqrt5<3$
from the centre, so its zeros are included in $n_T(3)$.  The remaining
bounded interval of $T$ meets only finitely many zeros, by meromorphic
analyticity and the absence of an identically zero function; increasing
the constant covers it.  This proves the claim.

The finite-order conductor-explicit counterpart is
Greni\'e--Molteni2019, Appendix~A, Theorem~A.1, author p.~26.
NT-Z10 instead proves the fixed-character form needed for the local
arithmetic divisor, and supplies the infinite-order version of root
MT-M04.  The argument uses the lower bound at its Jensen centre and
polynomial vertical growth.  It does not repair a statement quantified
over every function in an arbitrary Schwartz Mellin-transform space.

\subsection{An explicit uniform cyclotomic specialization}

\paragraph{NT-P11 and NT-P12.}
For $q\ge3$ and $(a,q)=1$, define
\[
 \psi(x;q,a)=\sum_{\substack{p^m\le x\\p^m\equiv a\ ({\rm mod}\ q)}}\log p,
 \qquad
 \pi(x;q,a)=\#\{p\le x:p\equiv a\ ({\rm mod}\ q)\}.
\]
Put $e_q=1$ if $q\equiv2\pmod4$ and $e_q=0$ otherwise.
Then $A$ implies
\begin{align*}
 \left|\psi(x;q,a)-\frac{x}{\varphi(q)}\right|
 &\le\sqrt{x}\left[
 \left(\frac{\log x}{2\pi}+2\right)\log q+
 \frac{\log^2x}{8\pi}+2\right]+e_q\log x,\qquad x\ge1,\\
 \left|\pi(x;q,a)-\frac{\operatorname{Li}_2(x)}{\varphi(q)}\right|
 &\le\sqrt{x}\left[
 \left(\frac1{2\pi}+\frac3{\log x}\right)\log q+
 \frac{\log x}{8\pi}+\frac1{4\pi}+\frac6{\log x}\right]+e_q,
 \qquad x\ge2.
\end{align*}

For the full proof, take $L=\mathbb Q(\zeta_q)$ and $K=\mathbb Q$.
The Galois group is $(\mathbb Z/q\mathbb Z)^\times$, of order
$n_L=\varphi(q)$, and its conjugacy classes are singletons.
The cyclotomic discriminant formula gives
\[
 \Delta_L=q^{\varphi(q)}
    \prod_{p\mid q}p^{-\varphi(q)/(p-1)},\qquad
 \log\Delta_L\le\varphi(q)\log q.
\]
At every prime $p\nmid q$, arithmetic Frobenius acts on $\zeta_q$ by
$\zeta_q\mapsto\zeta_q^p$.  Thus its $m$th power is the element
represented by $p^m$ modulo $q$.  The use of arithmetic or geometric
Frobenius merely replaces the labelled class by its inverse; the estimate
is identical after that relabelling.  The Dirichlet factorization of
$\zeta_L$ makes RH for it a consequence of $A$.
Apply NT-C01 and NT-C02 with $\delta_C=1/\varphi(q)$ and insert the
discriminant upper bound; all displayed coefficients multiplying
$\log\Delta_L$ are positive on the stated ranges.

If $q$ is odd or $4\mid q$, the primes unramified in $L$ are exactly
those not dividing $q$, so the resulting sums are the progression sums
just defined.  If $q\equiv2\pmod4$, write $q=2m$ with $m$ odd.
Then $L=\mathbb Q(\zeta_m)$ and the prime $2$ is unramified, whereas it
is excluded from a reduced class modulo $q$.  For odd primes the class
condition modulo $m$ is equivalent to that modulo $q$, since reduction
$(\mathbb Z/q\mathbb Z)^\times\to(\mathbb Z/m\mathbb Z)^\times$ is
an isomorphism.  The unramified Chebotarev sum can therefore differ only
by the powers of $2$ whose Frobenius powers lie in the designated class.
Their total weight is at most
$\lfloor\log x/\log2\rfloor\log2\le\log x$.
For the prime-counting sum the discrepancy is at most one.  Applying the
triangle inequality proves both formulas, including this degeneration.

The small moduli are transparent separately:
\begin{align*}
 \psi(x;1,0)&=\psi(x),& \pi(x;1,0)&=\pi(x),\\
 \psi(x;2,1)&=\psi(x)-\lfloor\log x/\log2\rfloor\log2,&
 \pi(x;2,1)&=\pi(x)-\mathbf1_{x\ge2}.
\end{align*}
For $1\le x<2$ the floor is zero.  Hence NT-P01 and NT-P03 give the
corresponding bounds in their stated ranges, with the explicit powers-of-two
correction retained for modulus two.  These are estimates on the open
cyclotomic arithmetic curve, and the additional term is exactly its
removed-prime contribution.

\subsection{A specified quadratic family and central nonvanishing}

\paragraph{NT-Z11. The exact one-level-density input.}
Let $\mathcal D$ consist of the nonzero odd squarefree integers, of both
signs.  For $d\in\mathcal D$, $8d$ is a fundamental discriminant and
$\chi_{8d}$ is its primitive quadratic character.  Put
$\ell_X=\log(X/(2\pi e))$.  Let $w$ be an even, nonnegative, nonzero
Schwartz function and write
\[
 W_w(X)=\sum_{d\in\mathcal D}w(d/X),\qquad
 \mathcal Z_{w,\phi}(X)=\frac1{W_w(X)}
  \sum_{d\in\mathcal D}w(d/X)
       \sum_{\rho=1/2+i\gamma}\phi(\gamma\ell_X/(2\pi)),
\]
where the inner sum counts every nontrivial zero with multiplicity.
Use the Fourier convention
$\widehat\phi(u)=\int_{\mathbb R}\phi(t)e^{-2\pi iut}\,dt$.
For a real even Schwartz function $\phi$ with
$\operatorname{supp}\widehat\phi\subset(-2,2)$, $A$ implies
\[
 \lim_{X\to\infty}\mathcal Z_{w,\phi}(X)
 =\widehat\phi(0)-\frac12\int_{-1}^1\widehat\phi(u)\,du.
\tag{NT-Z11}
\]
This is the leading term of Fiorilli--Parks--S\"odergren,
\path{arXiv:1601.06833}, Theorem~1.1, author p.~2, with definitions
(1.1)--(1.2) on the same page.  Its proof uses Theorem~3.5, author p.~12, and is concluded on p.~14.
The GRH input in the prime-sum step is Dirichlet GRH
(Lemma~2.5, author p.~7), which $A$ supplies.  No unrestricted Fourier
support, individual central nonvanishing, or nonabelian GRH is being
asserted here.

\paragraph{NT-Z12. A full nonvanishing deduction, including deweighting.}
Under $A$, at least $15/16$ of this specified family have a nonzero
central value in the following precise sense:
\[
 \liminf_{X\to\infty}
 \frac{\#\{d\in\mathcal D:|d|\le X,\ L(1/2,\chi_{8d})\ne0\}}
      {\#\{d\in\mathcal D:|d|\le X\}}\ge\frac{15}{16}.
\tag{NT-Z12}
\]
The same lower bound holds for the weighted proportions defined by any
of the fixed weights $w$ in NT-Z11.  This is a deduction from the cited
classical density theorem, not a claim of a new nonvanishing theorem.
Here are all approximation and counting details.

Let $m_d$ be the order of vanishing at $1/2$.  The completed primitive
quadratic function has root number $+1$.  One way to see the sign without
choosing a Gauss-sum square root is to divide the symmetric completed
Dedekind functional equation of $\mathbb Q(\sqrt{8d})$ by that of
$\mathbb Q$: the factorization
$\zeta_{\mathbb Q(\sqrt{8d})}=\zeta L(s,\chi_{8d})$ and the real or
complex gamma factors identify the quotient, up to a nonzero constant,
with the completed quadratic function.  The gamma duplication formula
handles the imaginary quadratic case.  Thus this entire function is
even about $s=1/2$.  Its gamma factor is finite and nonzero at the
centre, so $m_d$ is an even nonnegative integer.  In particular
\[
 2\mathbf1_{L(1/2,\chi_{8d})=0}\le m_d.
\tag{NT-central-parity}
\]
These functional equations are the standard Hecke equations; a precise
source for the completed convention is Ramakrishnan--Valenza,
Theorem~7-19, printed p.~279.

Choose real even smooth functions $g_j$ compactly supported in $(-1,1)$
such that $\int g_j=1$ and
$g_j\to g=\tfrac12\mathbf1_{[-1,1]}$ in $L^2(\mathbb R)$.
For example, smooth even cutoffs of $[-1+1/j,1-1/j]$, normalized by
their integrals, have these properties.  Define
\[
 f_j(t)=\int g_j(u)e^{2\pi iut}\,du,\qquad
 \phi_j(t)=|f_j(t)|^2.
\]
Then $\phi_j$ is real, even, nonnegative and Schwartz,
$\phi_j(0)=1$, and
$\operatorname{supp}\widehat\phi_j\subset(-2,2)$.
The support assertion follows because
$\widehat\phi_j=g_j*\widetilde g_j$, where
$\widetilde g_j(u)=\overline{g_j(-u)}$.
Plancherel gives $f_j\to f$ in $L^2$, with
\[
 f(t)=\frac{\sin(2\pi t)}{2\pi t}.
\]
Moreover, Cauchy--Schwarz yields
\[
 \|\phi_j-|f|^2\|_1
 \le\|f_j-f\|_2(\|f_j\|_2+\|f\|_2)\longrightarrow0.
\]
Consequently their Fourier transforms converge uniformly.  Since
\[
 \widehat{|f|^2}(u)=
 \begin{cases}\tfrac12-\tfrac14|u|,&|u|\le2,\\0,&|u|>2,
 \end{cases}
\]
we obtain
\[
 \lim_{j\to\infty}
 \left(\widehat\phi_j(0)-\frac12
                   \int_{-1}^1\widehat\phi_j(u)\,du\right)
 =\frac12-\frac12\cdot\frac34=\frac18.
\tag{NT-test-limit}
\]
Under $A$ every ordinate in NT-Z11 is real.  The nonnegativity of
$\phi_j$ and $\phi_j(0)=1$ therefore imply
\[
 \frac1{W_w(X)}\sum_{d\in\mathcal D}w(d/X)m_d
 \le\mathcal Z_{w,\phi_j}(X).
\]
First take $\limsup_{X\to\infty}$ for a fixed $j$, use NT-Z11, and
then let $j\to\infty$.  The result and (NT-central-parity) give
\[
 \limsup_{X\to\infty}
 \frac{\sum_{d\in\mathcal D}w(d/X)
               \mathbf1_{L(1/2,\chi_{8d})=0}}{W_w(X)}\le\frac1{16}.
\tag{NT-weighted-vanishing}
\]
This proves the weighted assertion without applying the density theorem
directly to a nonsmooth Fourier transform or to support touching $2$.

For completeness, the ordinary count follows by a separate elementary
argument.  If $A_+(Y)$ counts positive odd squarefree integers up to $Y$,
the squarefree indicator gives
\begin{align*}
 A_+(Y)&=\sum_{\substack{k\le\sqrt Y\\k\ {\rm odd}}}\mu(k)
       \#\{n\le Y/k^2:n\ {\rm odd}\}\\
 &=\frac Y2\sum_{k\ {\rm odd}}\frac{\mu(k)}{k^2}+O(\sqrt Y)
 =\frac4{\pi^2}Y+O(\sqrt Y).
\end{align*}
Here the omitted tail is $O(Y\sum_{k>\sqrt Y}k^{-2})=O(\sqrt Y)$,
and the Euler product is
$\sum_{k\ {\rm odd}}\mu(k)/k^2=8/\pi^2$.
Thus $N(X):=\#\{d\in\mathcal D:|d|\le X\}
 =8X/\pi^2+O(\sqrt X)$.
Partial summation of this counting formula gives, for every fixed
smooth compactly supported $w$,
\[
 W_w(X)=\frac4{\pi^2}X\int_{\mathbb R}w(t)\,dt+O_w(\sqrt X).
\]
Indeed apply Stieltjes integration to the positive count $A_+$ and
its reflected negative count; the integrated error is bounded by a
constant depending on $\int_0^\infty\sqrt t\,|w'(t)|dt$.

Fix $\delta>0$.  Take such an even smooth $w$ equal to $1$ on
$[-1,1]$, between $0$ and $1$ everywhere, and supported in
$[-1-\delta,1+\delta]$.  The number of vanishing members with
$|d|\le X$ is bounded above by its weighted counterpart.  Dividing by
$N(X)$ and using (NT-weighted-vanishing) gives a $\limsup$ of at most
$\frac1{16}\cdot\frac12\int w\le(1+\delta)/16$.
Letting $\delta\downarrow0$ proves (NT-Z12).

The characters in this statement are quadratic rank-one Artin objects.
Their varying conductors describe a family of double covers of open
arithmetic curves.  The symplectic expression in NT-Z11 matches the
finite-field hyperelliptic-family comparison discussed immediately after
Theorem~1.1 of the cited paper.  This comparison does not enlarge the
analytic support range or turn the $15/16$ bound into nonvanishing for
each member.  The ordinary proportion proved here concerns the
$\chi_{8d}$ family exactly; it has not been silently relabelled as all
fundamental discriminants.

\subsection{Growing fields, class numbers, and restricted unramified towers}

\paragraph{NT-H05 and NT-H06. Exact growing-field inputs.}
For pairwise nonisomorphic number fields $K_i\ne\mathbb Q$, write
$g_i=\tfrac12\log\Delta_{K_i}$ and $n_i=[K_i:\mathbb Q]$.
Let $N_q(K_i)$ count prime ideals of norm $q$, where $q$ ranges over
rational prime powers.  The family is \emph{asymptotically exact} when
the following limits exist for every indicated index:
\[
 \phi_q=\lim_i\frac{N_q(K_i)}{g_i},\qquad
 \phi_{\mathbb R}=\lim_i\frac{r_1(K_i)}{g_i},\qquad
 \phi_{\mathbb C}=\lim_i\frac{r_2(K_i)}{g_i}.
\]
Let $h_i,R_i$ be the ordinary class number and full unit regulator, and
let $\kappa_i=\operatorname{Res}_{s=1}\zeta_{K_i}(s)$.
Tsfasman--Vl\u adu\c t, \path{arXiv:math/0205129},
GRH Theorem~7.2, author p.~44 (proof pp.~44--49), gives
\begin{align*}
 \lim_i\frac{\log(h_iR_i)}{g_i}
 &=1+\sum_q\phi_q\log\frac q{q-1}
      -\phi_{\mathbb R}\log2-\phi_{\mathbb C}\log(2\pi),\\
 \lim_i\frac{\log\kappa_i}{g_i}
 &=\sum_q\phi_q\log\frac q{q-1}.
\end{align*}
Only RH for the family of Dedekind functions is needed: $M$ supplies it;
$A$ suffices when every $K_i/\mathbb Q$ is abelian.  Neither normality
nor a tower is required in this theorem.  If $n_i/g_i\to0$, then
$N_q(K_i)\le n_i$ makes all the displayed invariants zero, so the
first limit is $1$.  This proves the usual Brauer--Siegel asymptotic
in this range without adding a normality hypothesis.

The same paper's GRH Theorem~3.1, author p.~19 (proof pp.~19--21),
supplies the basic inequality
\[
 \sum_q\frac{\phi_q\log q}{\sqrt q-1}
 +\alpha_{\mathbb R}\phi_{\mathbb R}
 +\alpha_{\mathbb C}\phi_{\mathbb C}\le1,
\tag{NT-basic}
\]
where
\[
 \alpha_{\mathbb R}=\tfrac12\log(8\pi)+\tfrac\pi4+\tfrac\gamma2,
 \qquad
 \alpha_{\mathbb C}=\log(8\pi)+\gamma,
\]
and $\gamma$ is Euler's constant.  Both original theorem pages were
inspected as images.  These are cited classical theorems.  The next
two consequences and their exact tower interpretations are proved below.

\paragraph{NT-H07. A root-discriminant bound with prescribed splitting.}
Suppose $n_i\to\infty$,
\[
 \Delta_{K_i}^{1/n_i}\longrightarrow\varrho\in(1,\infty),
 \qquad r_1(K_i)/n_i\longrightarrow\tau\in[0,1].
\]
Let $P$ be any set of rational primes such that every $p\in P$ splits
completely in every $K_i$.  Under RH for all $\zeta_{K_i}$,
\[
 \log\varrho\ge\log(8\pi)+\gamma+\frac\pi2\tau
          +2\sum_{p\in P}\frac{\log p}{\sqrt p-1}.
\tag{NT-H07}
\]
If $P$ is infinite, the inequality includes convergence of this
nonnegative series.  No normality assumption on $K_i/\mathbb Q$ is
needed.

To prove the claim, first take a subsequence with strictly increasing
degrees; its fields are pairwise nonisomorphic and its stated limits
are unchanged.  Note that
$g_i/n_i\to\tfrac12\log\varrho>0$.
For every fixed prime power $q$, the numbers $N_q(K_i)/g_i$ are
bounded since $N_q(K_i)\le n_i$.  Successive subsequence extraction,
followed by a diagonal subsequence over the countable prime powers,
makes all of them converge.  The archimedean ratios already converge;
on the resulting asymptotically exact subfamily they are
\[
 \phi_{\mathbb R}=\frac{2\tau}{\log\varrho},\qquad
 \phi_{\mathbb C}=\frac{1-\tau}{\log\varrho}.
\]
Complete splitting gives $N_p(K_i)=n_i$, hence
$\phi_p=2/\log\varrho$ for $p\in P$.
Apply (NT-basic) and discard the nonnegative terms from other norms.
For every finite $P_0\subset P$, multiplication by $\log\varrho$
yields
\[
 2\sum_{p\in P_0}\frac{\log p}{\sqrt p-1}
 +2\tau\alpha_{\mathbb R}+(1-\tau)\alpha_{\mathbb C}
 \le\log\varrho.
\]
The identity
$2\alpha_{\mathbb R}-\alpha_{\mathbb C}=\pi/2$
gives the claimed expression for finite $P_0$.  Taking the supremum
over finite subsets proves convergence and the same bound for all of
$P$.  The limiting quantities $\varrho,\tau$ are unchanged by the
subsequence, so this proves the original assertion.

\paragraph{NT-H08. Finiteness of a restricted arithmetic fundamental group.}
Fix a number field $K$ and a finite set $P$ of rational primes that
split completely in $K$.  Let $K^{\rm ur,P}$ be the compositum of all
finite extensions of $K$ unramified at every finite prime, in which
every prime of $K$ above $P$ splits completely.  Put
\[
 C(P)=8\pi\exp\left(\gamma+
                 2\sum_{p\in P}\frac{\log p}{\sqrt p-1}\right).
\]
If $\Delta_K^{1/[K:\mathbb Q]}<C(P)$, then $M$ implies
$[K^{\rm ur,P}:K]<\infty$.
If $K$ is totally real and only totally real extensions are allowed,
the corresponding conclusion holds under the weaker inequality
$\Delta_K^{1/[K:\mathbb Q]}<e^{\pi/2}C(P)$.
The inequalities are strict: no equality case is claimed.

Here is the full proof.  Being unramified at finite primes and splitting
completely at specified primes are both preserved by conjugation over
$K$ and by finite composita.  Thus the displayed compositum is Galois
over $K$, and all of its finite subextensions have these properties.
If it had infinite degree, it would contain a sequence of finite
Galois subextensions $K_i/K$ with degrees tending to infinity.
The relative discriminant formula and finite unramifiedness give
\[
 \Delta_{K_i}=\Delta_K^{[K_i:K]},\qquad
 \Delta_{K_i}^{1/[K_i:\mathbb Q]}=
                  \Delta_K^{1/[K:\mathbb Q]}=: \varrho.
\]
If $K=\mathbb Q$, this equality gives discriminant one for any such
$K_i$; Minkowski's discriminant bound already rules out a nontrivial
extension.  Otherwise $\varrho>1$.
After taking a subsequence, $r_1(K_i)/[K_i:\mathbb Q]$ tends to some
$\tau\in[0,1]$.  Every $p\in P$ splits completely in $K_i/\mathbb Q$
because it splits in $K$ and each prime above it splits in $K_i/K$.
$M$ supplies the Dedekind RH hypotheses of NT-H07.  That result yields
$\varrho\ge e^{\pi\tau/2}C(P)\ge C(P)$, a contradiction.
In the totally real version all $K_i$ are totally real, so $\tau=1$
and the same proof gives the second threshold.  This completes the proof.

For $P=\varnothing$, the thresholds are
$8\pi e^\gamma\simeq44.763232$ and
$8\pi e^{\gamma+\pi/2}\simeq215.332516$.
The group $\operatorname{Gal}(K^{\rm ur,P}/K)$ is the quotient of
$\pi_1^{\text{ét}}(\operatorname{Spec}\mathcal O_K)$ obtained by
killing the closed normal subgroup generated by Frobenius elements at
the primes above $P$.  Requiring total reality also kills the
archimedean complex-conjugation classes.  This is the finite-\'{e}tale
Galois-category dictionary, SGA~1, Expos\'e~V, \S7.  The discriminant
criterion therefore gives a concrete conditional finiteness statement
for this quotient.  It is a restriction on arithmetic towers in the
cote-149/GT comparison programme, not a claim of a new anabelian
reconstruction theorem.

\subsection{Prime races: zero location versus relations among ordinates}

\paragraph{NT-R01 and NT-R02. What $A$ alone supplies.}
Fix distinct reduced classes $a_1,\ldots,a_r$ modulo $q$, and set
\[
 E_{q;\boldsymbol a}(x)=\frac{\log x}{\sqrt x}
 \bigl(\varphi(q)\pi(x;q,a_1)-\pi(x),\ldots,
       \varphi(q)\pi(x;q,a_r)-\pi(x)\bigr).
\]
Rubinstein--Sarnak, \emph{Chebyshev's Bias} (1994), Theorem~1.1,
printed p.~174, proof \S2.1, gives under $A$ a probability measure
$\mu_{q;\boldsymbol a}$ such that, for every bounded continuous $f$,
\[
 \lim_{X\to\infty}\frac1{\log X}
       \int_2^X f(E_{q;\boldsymbol a}(x))\frac{dx}{x}
   =\int_{\mathbb R^r}f\,d\mu_{q;\boldsymbol a}.
\tag{NT-R01}
\]
Its normalization uses $\pi(x)$ in each coordinate.  Theorem~1.2,
printed p.~175, proof \S2.2, and the discussion following it imply
\[
 \underline\delta\{x:\pi(x;4,3)>\pi(x;4,1)\}>0,
 \qquad
 \underline\delta\{x:\pi(x;4,1)>\pi(x;4,3)\}>0,
\tag{NT-R02}
\]
where $\underline\delta S=\liminf_{X\to\infty}
(\log X)^{-1}\int_{S\cap[2,X]}dx/x$.
The same conclusion holds for the two classes modulo three.
These statements concern logarithmic sampling.  The existence of the
limiting distribution for continuous tests does not by itself justify
substituting an ordering-cone indicator: its boundary could have mass.

\paragraph{NT-R03 and NT-R04. The separate ordinate-independence input.}
Use the source's Grand Simplicity Hypothesis (GSH): the nonnegative
ordinates of zeros of all primitive Dirichlet functions, listed with
multiplicity, are linearly independent over $\mathbb Q$.  In particular
this excludes central zeros and repeated ordinates.  GSH remains an
additional hypothesis here.

Under $A$ and GSH, for $r\ge2$ every strict ordering race has a
logarithmic density in $(0,1)$; the source's formula (1.2) and
Remark~1.3, printed p.~176, give the requisite regularity and positive
mass.  For fixed $r$,
\[
 \max_{\substack{a_1,\ldots,a_r\pmod q\\
                   \text{distinct reduced classes}}}
 \left|\delta\{x:\pi(x;q,a_1)>\cdots>\pi(x;q,a_r)\}
                   -\frac1{r!}\right|\longrightarrow0
       \quad(q\to\infty).
\tag{NT-R04}
\]
This is Theorem~1.5, printed p.~177, proof \S3.2.
For the exact symmetry classification, put
$c(q,a)=-1+\#\{b\pmod q:b^2\equiv a\pmod q\}$.
Theorem~1.4, printed p.~177, proof Proposition~3.1, says that the
limiting law is invariant under every coordinate permutation precisely
in the following cases: $r=2$ and $c(q,a_1)=c(q,a_2)$; or $r=3$
and, after labelling, the classes are
$a_1,a_1\rho,a_1\rho^2$ with $\rho\ne1$ and $\rho^3=1\pmod q$.
This is symmetry of the full law.  It must not be conflated with a
converse assertion about one numerical race density.  The discussion
on p.~177 gives, in particular, the strict bias
$1/2<\delta\{x:\pi(x;4,3)>\pi(x;4,1)\}<1$ under these extra
hypotheses.  The relevant original pp.~174--178 were visually inspected
because their PDF text encoding is defective.

Geometrically these are races between Frobenius classes in cyclotomic
covers of $\operatorname{Spec}\mathbb Z[1/q]$.  Their limiting
distribution is an explicit-formula statistic of the corresponding
rank-one objects.  The additional GSH gate concerns relations among
zero ordinates, so the zero-location inputs $A$ and $M$ do not discharge
it in this argument.  A finite-field monodromy comparison likewise
does not remove that number-field hypothesis.

\subsection{Multiplicity, gaps, and the first zero: quantitative limits}

\paragraph{NT-Z13. An argument bound with its actual archimedean hypothesis.}
Carneiro--Chandee--Milinovich, \emph{A note on the zeros of zeta and
$L$-functions} (2015), author PDF Theorem~1, p.~2, proves under RH
\[
 |S(t)|\le \frac14\frac{\log t}{\log\log t}
       +O\left(\frac{\log t\log\log\log t}{(\log\log t)^2}\right).
\tag{NT-Z13a}
\]
Here $S(t)=\pi^{-1}\arg\zeta(1/2+it)$, with the argument continued
from $2$ through $2+it$ and the mean of its two one-sided values used
at a zero.  The proof is in \S3, pp.~3--5.

We record the precise Hecke specialization of their Theorem~5,
p.~7, equations (4.4)--(4.5), whose proof occupies pp.~7--12.
Let $\eta$ be a primitive finite-order Hecke character of a number
field $K$ of degree $n$.  Write its completion, up to a nonzero
constant independent of $s$, as
\[
 \Lambda(s,\eta)=G_\eta(s)L(s,\eta),\qquad
 G_\eta(s)=Q_\eta^{s/2}\prod_{j=1}^n\Gamma_{\mathbb R}(s+\mu_j),
 \qquad Q_\eta=\Delta_K\operatorname N\mathfrak f_\eta,
\]
where $\Gamma_{\mathbb R}(s)=\pi^{-s/2}\Gamma(s/2)$.
Each real place contributes $\mu_j=0$ or $1$ according to the sign
character.  Each complex place contributes the pair $0,1$.
Set
\[
 C_\eta=Q_\eta\prod_{j=1}^n(\mu_j+3),\quad
 C_\eta(t)=C_\eta(|t|+1)^n,\quad
 D_\eta(t)=\log\log\bigl(C_\eta(t)^{3/n}\bigr).
\]
For $t>0$ define
$S_\eta(t)=-\pi^{-1}\int_{1/2}^{\infty}
\operatorname{Im}(L'/L)(\sigma+it,\eta)\,d\sigma$ away from zero
ordinates, and use one-sided means at them.  Under RH for this
particular $L$-function,
\[
 |S_\eta(t)|\le\frac{\log C_\eta(t)}{4D_\eta(t)}
   +O\left(\frac{\log C_\eta(t)\log D_\eta(t)}{D_\eta(t)^2}\right).
\tag{NT-Z13b}
\]
The implied constant is absolute.  In particular, for fixed $K,\eta$,
\[
 |S_\eta(t)|\le\frac n4\frac{\log t}{\log\log t}
   +O_{K,\eta}\left(
       \frac{\log t\log\log\log t}{(\log\log t)^2}\right).
\tag{NT-Z13c}
\]

Here is the check of the source hypotheses, which matters for $M$.
After grouping prime ideals by their underlying rational prime, the
Hecke Euler factor is a product of at most $n$ linear reciprocal
factors with local parameters of modulus at most one: factor each
$1-\eta(\mathfrak p)p^{-fs}$ into its $f$ roots and pad missing
factors by zero parameters.  Thus the source's parameter $\vartheta$
is zero.  The completion has order one, the standard functional
equation with the conjugate character, and is entire except for the
trivial character's poles at $0,1$.  Its listed gamma parameters are
real and nonnegative.  These are precisely the analytic conditions
in \S4, p.~6, supplied by classical Hecke theory (RV1999,
Theorem~7-19).  $M$ supplies the remaining zero-location assumption;
$A$ suffices for Dirichlet characters over $\mathbb Q$.
In particular, the source explicitly requires gamma parameters to be
real or to occur in conjugate pairs.  We do not silently use it for
an arbitrary infinite-order unitary Hecke character.  The paired
$\pm it_k$ family in root MT-H01--H03 can be checked separately
against this same hypothesis.

\paragraph{NT-Z14. Full deduction of multiplicity and gap bounds.}
For each fixed $K,\eta$ as in NT-Z13, under its GRH, a zero
$1/2+i\gamma$ with $\gamma\to+\infty$ has multiplicity
\[
 m_\eta(\gamma)\le\frac n2\frac{\log\gamma}{\log\log\gamma}
   +O_{K,\eta}\left(
    \frac{\log\gamma\log\log\log\gamma}{(\log\log\gamma)^2}\right).
\tag{NT-mult}
\]
For consecutive distinct sufficiently large positive ordinates
$\gamma<\gamma'$,
\[
 \gamma'-\gamma\le\frac{\pi}{\log\log\gamma}
   +O_{K,\eta}\left(
    \frac{\log\log\log\gamma}{(\log\log\gamma)^2}\right).
\tag{NT-gap}
\]
Consequently every sufficiently large real $T$ lies within
$\pi/(2\log\log T)+O_{K,\eta}(
\log\log\log T/(\log\log T)^2)$ of a zero ordinate.  The analogous
negative-height statements follow by conjugating $\eta$.  These are
classical consequences in a fixed-character formulation; they assert
neither simplicity nor pair correlation.

We give the entire deduction, including the jump and endpoint issues.
In a sufficiently high horizontal region $L(s,\eta)$ has no pole;
GRH and the boundary nonvanishing theorem give no zero to the right
of the critical line.  Continue $\log L$ from its Euler-product
normalization there.  Near $\rho=1/2+i\gamma$ write
$L(s,\eta)=(s-\rho)^m U(s)$ with $U$ holomorphic and nonzero.
Approaching the critical line from its right, the argument of
$s-\rho$ changes from $-\pi/2$ to $\pi/2$ across the ordinate.
Hence
\[
 S_\eta(\gamma+)-S_\eta(\gamma-)=m.
\]
Taking one-sided limits in NT-Z13c bounds this jump by twice its
right-hand side and proves (NT-mult).

Let $\Theta_\eta(t)=\operatorname{Im}\log G_\eta(1/2+it)$, with
any continuous determination at large positive $t$.
The functional equation and conjugation give
$\Lambda(1/2+it,\eta)=\varepsilon_\eta
\overline{\Lambda(1/2+it,\eta)}$.
Therefore, on an interval containing no zero, the argument of the
completed function is constant modulo $\pi$, and
\[
 S_\eta'(t)=-\frac1\pi\Theta_\eta'(t),\qquad
 \frac1\pi\Theta_\eta'(t)
    =\frac n{2\pi}\log t+O_{K,\eta}(1).
\tag{NT-phase}
\]
The derivative formula is Stirling's formula applied to each of the
$n$ displayed $\Gamma_{\mathbb R}$ factors; the conductor term is
constant.  The same identity and NT-Z13c show that positive ordinates
are unbounded: if there were no zeros beyond a height, integration
would force $S_\eta(t)$ to have size comparable to $t\log t$.

For large $\gamma$, a zero-free gap cannot have length at least one.
Otherwise it contains an interval of length $1/2$ whose endpoints
are between $\gamma$ and $\gamma+1$.  Integrating (NT-phase) over it
gives a decrease of size $(n/(4\pi)+o(1))\log\gamma$, whereas
NT-Z13c bounds the difference of its endpoint values by
$O_{K,\eta}(\log\gamma/\log\log\gamma)$, a contradiction.
Thus $h=\gamma'-\gamma<1$ eventually.  Integrate (NT-phase) from
just above $\gamma$ to just below $\gamma'$ and then take endpoint
limits.  Uniformly on this interval,
\[
 \left(\frac n{2\pi}\log\gamma+O_{K,\eta}(1)\right)h
   =S_\eta(\gamma+)-S_\eta(\gamma'-)
   \le\frac n2\frac{\log\gamma}{\log\log\gamma}
      +O_{K,\eta}\left(
       \frac{\log\gamma\log\log\log\gamma}
            {(\log\log\gamma)^2}\right).
\]
Division proves (NT-gap).  Any $T$ between consecutive ordinates is
at distance at most half their gap from one endpoint.  Since that
gap is less than one, replacing $\gamma$ by $T$ in the slowly varying
bound changes only the error term.  This proves the last assertion
and completes the proof.  For $K=\mathbb Q$ and the trivial
character, $n=1$ gives the usual zeta multiplicity and gap bounds
under $A$ alone.

\paragraph{NT-Z15 and NT-Z16. Conductor aspect and the first ordinate.}
In the same finite-order Hecke setting put
$D_\eta=\log\log(C_\eta^{3/n})$.
Theorem~5, equation (4.5), gives
\[
 \operatorname{ord}_{s=1/2}L(s,\eta)
 \le\frac{\log C_\eta}{2D_\eta}
   +O\left(\frac{\log C_\eta\log D_\eta}{D_\eta^2}\right).
\tag{NT-Z15}
\]
There is no gamma pole or zero at $1/2$, so the order equals that
of the completion in the source.  Theorem~7, p.~13, with its full
proof on that page, gives at least one zero $1/2+i\gamma$ satisfying
\[
 |\gamma|\le\frac{\pi}{2D_\eta}
          +O\left(\frac{\log D_\eta}{D_\eta^2}\right)
 \quad\text{as } C_\eta^{3/n}\longrightarrow\infty.
\tag{NT-Z16}
\]
Both implied constants are universal.  A possible central zero is
included: this is not a bound for the first strictly positive
ordinate.  Keeping $n$ inside $D_\eta$ is essential when the degree
varies.  For primitive Dirichlet characters of conductor $q>2$, the
specializations are
\[
 \operatorname{ord}_{1/2}L(s,\chi)
       \le(1/2+o(1))\frac{\log q}{\log\log q},\qquad
 \min_\rho|\operatorname{Im}\rho|
       \le\frac{\pi/2+o(1)}{\log\log q}.
\]
Here $A$ supplies every required GRH hypothesis.  For finite-order
Hecke characters over general fields the corresponding input is
the finite-order part of $M$; it is not supplied by $A$ merely
because the base field happens to be abelian over $\mathbb Q$.

These bounds constrain the zero divisor of rank-one arithmetic
local systems and their Artin inductions.  In the Grothendieck
comparison language they are quantitative restrictions on the
archimedean explicit-formula spectrum.  They supply neither a
finite-field purity argument over $\operatorname{Spec}\mathbb Z$
nor a general geometric multiplicity-one theorem.  They refine the
specific open gates in NT-Z06 while leaving them open.

\subsection{Further prime constants and quadratic-family moments}

\paragraph{NT-P13. Explicit and asymptotic prime-gap constants.}
Under the RH part of $A$, every $x\ge4$ has a prime in
\[
 \left[x,x+\frac{22}{25}\sqrt x\log x\right].
\]
This is Carneiro--Milinovich--Soundararajan,
\emph{Fourier optimization and prime gaps} (2019), author v2
Theorem~5, p.~5; the explicit proof is \S6, pp.~21--26.
Consequently, for every consecutive pair with $p_n>3$,
$p_{n+1}-p_n\le(22/25)\sqrt{p_n}\log p_n$.
For clarity about the closed left endpoint, apply the interval
statement to $x=p_n+\varepsilon$ with
$0<\varepsilon<p_{n+1}-p_n$, then let $\varepsilon\downarrow0$.
Corollary~4 on the same page, following Theorem~3 and its proof in
\S5, strengthens the asymptotic constant:
\[
 \limsup_{n\to\infty}
       \frac{p_{n+1}-p_n}{\sqrt{p_n}\log p_n}<\frac{21}{25}.
\]
The explicit all-$x$ statement and the sharper asymptotic statement
have different quantifiers.  No claim is made that these are the
latest possible constants.

\paragraph{NT-P14--NT-P16. Fourier improvements in the conductor aspect.}
Carneiro--Milinovich--Quesada-Herrera--Ramos,
\emph{Fourier optimization, the least quadratic non-residue, and
the least prime in an arithmetic progression}, arXiv v2 dated
August 2025, Theorems~1 and 2, author p.~2, gives under $A$
\[
 \limsup_{p\to\infty}\frac{n_p}{(\log p)^2}<0.7615,
 \qquad
 \limsup_{p\to\infty}\frac{r_p}{(\log p)^2}<0.7615.
\tag{NT-P14}
\]
Here $p$ runs over primes, $n_p$ is the least positive quadratic
nonresidue modulo $p$, and $r_p$ is the least prime $r$ with
$(r/p)=1$.  The nonresidue is itself prime.  Theorem~1 follows from
Theorem~6, whose proof is \S4, pp.~13--17; the modification proving
Theorem~2 is \S4.4, pp.~16--17.  These bounds improve the
asymptotic constant; they do not replace NT-P07's explicit threshold
$p\ge5$ by an unstated numerical threshold.

More generally, for a nonprincipal character $\chi$ modulo $q$ of
fixed order $\ell\ge2$, define
$n_\chi=\min\{n\ge1:\chi(n)\notin\{0,1\}\}$.
Theorem~6, p.~5, and Corollary~7, p.~6, give
\[
 n_\chi\le(c_\ell+o_\ell(1))(\log q)^2,
 \qquad
 \begin{array}{c|rrrr}
   \ell&2&3&4&5\\ \hline
   c_\ell&0.7615&0.6707&0.6131&0.5765
 \end{array}
\tag{NT-P15}
\]
as $q\to\infty$.  For fixed $\ell\ge6$ the same corollary permits
\[
 c_\ell=\frac14\left(
  1-\frac{\frac{\ell}{\ell-1}\log\bigl((3\ell-4)/\ell\bigr)}
             {\log(\ell-1)}+\frac1{\ell-1}\right)^{-2}.
\]
The character need not be primitive.  Its values zero at primes
dividing $q$ are explicitly excluded in the definition.  If
$n_\chi$ were composite, every prime divisor would be smaller and
coprime to $q$, hence would have character value one by minimality;
multiplicativity would give $\chi(n_\chi)=1$, a contradiction.
Thus these are also bounds for the least prime outside the kernel
of the cyclotomic character.  The limit $c_\ell\to1/4$ is a
statement about these constants; we do not assert uniformity when
$\ell$ varies with $q$.

For the least prime $P(a,q)\equiv a\pmod q$, the same paper's
Theorem~3, author p.~3, proof \S5, pp.~18--22, gives under $A$
\[
 \limsup_{q\to\infty}
  \max_{(a,q)=1}\frac{P(a,q)}{(\varphi(q)\log q)^2}<\frac89.
\tag{NT-P16}
\]
The uniformity in the reduced class follows from the proof: its
errors in (5.15)--(5.16) depend on the fixed Fourier test function
and not on $a$.  The numerical optimization certification in the
paper is cited, not rerun here.  These results concern closed points
and selected Frobenius classes of cyclotomic covers.  They sharpen
the effective detection bounds in the SGA~1 Galois-category
dictionary; they do not enlarge the class of $L$-functions furnished
by $A$.

\paragraph{NT-L11. Quadratic central moments with the full family specified.}
For each fixed real $k\ge0$, $A$ and Harper,
\emph{Sharp conditional bounds for moments of the Riemann zeta
function}, Theorem~2, author p.~15, give
\[
 \sum_{\substack{X\le |d|\le2X\\d\ {
m fundamental}}}
       |L(1/2,\chi_d)|^k
     \ll_k X(\log X)^{k(k+1)/2}.
\tag{NT-L11}
\]
The family includes positive and negative fundamental discriminants;
$\chi_d$ is the associated primitive quadratic character.  The
argument is in \S5, pp.~15--17, using Propositions~3 and 4.
The input is only GRH for the quadratic functions in question.
This is an upper bound of the predicted order, not an asymptotic
formula with its leading constant.  In particular it asserts no
bound for negative moments or reciprocals of central values.

\paragraph{NT-L12. Full fixed-threshold large-value deductions.}
For every fixed $\alpha>0$, $A$ implies
\[
 \operatorname{meas}\{t\in[T,2T]:
            |\zeta(1/2+it)|\ge(\log T)^\alpha\}
       \ll_\alpha \frac{T}{(\log T)^{\alpha^2}}.
\tag{NT-L12a}
\]
For every fixed $\alpha>1/2$ it also implies
\[
 \#\{d\text{ fundamental}:X\le|d|\le2X,
          |L(1/2,\chi_d)|\ge(\log X)^\alpha\}
       \ll_\alpha
          \frac{X}{(\log X)^{(\alpha-1/2)^2/2}}.
\tag{NT-L12b}
\]
Here is the complete proof.  For any fixed $k>0$, Markov's
inequality and Harper's Theorem~1 (NT-L04) bound the first measure by
\[
 (\log T)^{-2k\alpha}
       \int_T^{2T}|\zeta(1/2+it)|^{2k}\,dt
       \ll_k T(\log T)^{k^2-2k\alpha}.
\]
Choose the fixed value $k=\alpha$ to obtain (NT-L12a).
Similarly NT-L11 bounds the second count by
$O_k(X(\log X)^{k(k+1)/2-k\alpha})$.
For $\alpha>1/2$, choose the positive fixed value
$k=\alpha-1/2$; completing the square gives
$k(k+1)/2-k\alpha=-(\alpha-1/2)^2/2$.
This proves (NT-L12b).  All constants may depend on the fixed
threshold exponent; no choice of $k$ growing with $T$ or $X$ has
been made.  Central zeros cause no difficulty because only positive
moments occur.  The result is a direct optimization of the cited
moment bounds, not a new moment theorem.

The quadratic family consists of rank-one Artin objects coming from
quadratic covers with growing discriminant.  Its central-value
moment exponent is an arithmetic family statement.  A geometric
symmetry-type comparison is useful context, but is not substituted
for Harper's analytic proof or the exact GRH input.

\subsection{Closed-point Frobenius relations in unramified towers}

\paragraph{NT-H09. A relative splitting-degree budget.}
Let $K\ne\mathbb Q$ be a number field of degree $n$, discriminant
$\Delta_K$, and root discriminant $\varrho=\Delta_K^{1/n}$.
Let $S$ be a finite set of prime ideals of $K$, with
$q_{\mathfrak p}=\operatorname N\mathfrak p$, and assign to each
$\mathfrak p\in S$ a positive integer $m_{\mathfrak p}$.
Suppose that $K_i/K$ are finite Galois extensions unramified at all
finite primes, $d_i=[K_i:K]\to\infty$, and the common residue
degree of primes above $\mathfrak p$ in $K_i/K$ is at most
$m_{\mathfrak p}$ for every $i$.  Assume
$r_1(K_i)/[K_i:\mathbb Q]\to\tau$ and RH for every
$\zeta_{K_i}$.  Then
\[
 \log\varrho\ge\log(8\pi)+\gamma+\frac\pi2\tau
  +\frac2n\sum_{\mathfrak p\in S}
       \frac{\log q_{\mathfrak p}}
            {q_{\mathfrak p}^{m_{\mathfrak p}/2}-1}.
\tag{NT-H09}
\]
Thus $M$ supplies the needed zero hypotheses.  No nesting of the
extensions is required.  Complete splitting is the case
$m_{\mathfrak p}=1$ and does not require the rational prime below
$\mathfrak p$ to split in $K$.

Here is the full proof.  Pass to strictly increasing degrees so that
the fields are pairwise nonisomorphic.  Finite unramifiedness gives
$\Delta_{K_i}=\Delta_K^{d_i}$ and
$g_i=\tfrac12 d_i\log\Delta_K$.
There are finitely many possible residue degrees at each point of
the finite set $S$, so pass to a subsequence on which these degrees
are fixed, say $f_{\mathfrak p}\le m_{\mathfrak p}$.
As in NT-H07, boundedness of $N_Q(K_i)/g_i$ and a diagonal
subsequence over rational prime powers $Q$ give an asymptotically
exact family.  Neither the limiting signature nor the displayed
residue degrees change.  Since the extension is Galois and
unramified, precisely $d_i/f_{\mathfrak p}$ primes lie above
$\mathfrak p$, each of norm $q_{\mathfrak p}^{f_{\mathfrak p}}$.
Different base prime ideals give disjoint sets of primes upstairs,
even when their absolute norms coincide.  Therefore
\[
 \phi_Q\ge\frac2{\log\Delta_K}
       \sum_{\substack{\mathfrak p\in S\\
                    q_{\mathfrak p}^{f_{\mathfrak p}}=Q}}
           \frac1{f_{\mathfrak p}}.
\]
The archimedean invariants are
$\phi_{\mathbb R}=2n\tau/\log\Delta_K$ and
$\phi_{\mathbb C}=n(1-\tau)/\log\Delta_K$.
Substitute these lower bounds in the TV inequality (NT-basic),
discard the other nonnegative terms, and multiply by
$\log\Delta_K/n$.  Since
$\log(q_{\mathfrak p}^{f_{\mathfrak p}})/f_{\mathfrak p}
=\log q_{\mathfrak p}$, the finite-place contribution is at least
\[
 \frac2n\sum_{\mathfrak p\in S}
      \frac{\log q_{\mathfrak p}}
           {q_{\mathfrak p}^{f_{\mathfrak p}/2}-1}
 \ge\frac2n\sum_{\mathfrak p\in S}
      \frac{\log q_{\mathfrak p}}
           {q_{\mathfrak p}^{m_{\mathfrak p}/2}-1}.
\]
The archimedean part simplifies using
$2\alpha_{\mathbb R}-\alpha_{\mathbb C}=\pi/2$ exactly as in
NT-H07.  This proves (NT-H09), including the possibility of
coincident prime-power norms.

\paragraph{NT-H10. A finite quotient from bounded-order Frobenius relations.}
Let
$G_K=\pi_1^{\text{ét}}(\operatorname{Spec}\mathcal O_K)$, and
for the data $S,\boldsymbol m$ above let $N$ be the closed normal
subgroup generated by
$\operatorname{Frob}_{\mathfrak p}^{m_{\mathfrak p}}$ for
$\mathfrak p\in S$.  The normal subgroup is independent of the
choices inside the Frobenius conjugacy classes.  Put
\[
 C_K(S,\boldsymbol m)=8\pi\exp\left(
       \gamma+\frac2n\sum_{\mathfrak p\in S}
       \frac{\log q_{\mathfrak p}}
            {q_{\mathfrak p}^{m_{\mathfrak p}/2}-1}\right).
\]
Under $M$, if $\varrho<C_K(S,\boldsymbol m)$, then $G_K/N$ is
finite; equivalently, $N$ is open.  For totally real $K$, first
quotient $G_K$ by the closed normal subgroup generated by complex
conjugations.  In that totally real quotient the weaker condition
$\varrho<e^{\pi/2}C_K(S,\boldsymbol m)$ suffices.  No equality
case is asserted.

To prove this, suppose the stated quotient is infinite.  Its finite
Galois subextensions $K_i/K$ can then be chosen with degrees tending
to infinity.  They are unramified at every finite prime.  In each
finite Galois group the order of
$\operatorname{Frob}_{\mathfrak p}$ divides
$m_{\mathfrak p}$, and that order equals its residue degree.
Choose a subsequence whose signature ratios tend to
$\tau\in[0,1]$.  NT-H09 gives
$\varrho\ge e^{\pi\tau/2}C_K(S,\boldsymbol m)
\ge C_K(S,\boldsymbol m)$, contradicting the strict inequality.
In the totally real quotient every $K_i$ is totally real, so
$\tau=1$ and the second conclusion follows.  If $K=\mathbb Q$,
Minkowski's discriminant bound already makes the whole unramified
group trivial, so this omitted case satisfies the conclusion
unconditionally.  This completes the proof.

For $m_{\mathfrak p}=1$ this is the unrestricted closed-point
version of NT-H08.  If $S$ consists of all primes above rational
primes that split completely in $K$, then each such rational prime
contributes $n$ equal summands and NT-H08's constant is recovered.
For larger $m_{\mathfrak p}$ the imposed relations only bound the
orders of Frobenius elements.  The conclusion concerns the index
of their closed normal closure; it does not assert that this normal
subgroup is generated by finitely many elements as an ordinary
topological group.  This is a quantitative Galois-category
consequence of the TV explicit-formula inequality, with its precise
arithmetic input recorded, rather than a general anabelian
reconstruction statement.

% END number_theory

% BEGIN geometry_programme
\section{Geometry, motives, and the Grothendieck programme}
\label{sec:geometry}

The letters A and M have the dependency meaning fixed in the introduction.
In particular, accepting the stipulated endpoint A is not a certification of
every intermediate statement in the note, and assuming M does not certify
the functional analysis in its proposed proof. The findings below are new
formulations proved in this run; no claim of priority over the literature is
intended. Earlier findings are used with their original identifiers.

\subsection{The transfer results already established}
\label{ag:transfer}

\paragraph{AG-AT (the exact Artin--Tate domain).}
For a finite-type scheme over $\mathbb Z$ with exact all-prime point counts
\[
 \#X(\mathbb F_{p^r})=
 \sum_j p^{jr}\operatorname{Tr}(F_p^r\mid V_j^{I_p}),
\]
where the $V_j$ are finitely many virtual finite-image representations,
the identity
\[
 \zeta_X(s)=\prod_j L(V_j,s-j)
\]
is unconditional. This is ZETA\_TWO ZT01, in
\path{RETURN_ZETA_TWO_final_IMPLICATIONS.md}, Section 1,
extending BRIDGES F1. If the representations have abelian image, A gives
noninteger global divisor support on the lines $\Re s=j+1/2$.
For arbitrary finite image, M gives the corresponding meromorphic
divisor-support assertion through integer Brauer induction (ZT02).
Negative exponents can produce poles; Artin holomorphy is an extra question.

\paragraph{AG-DED and AG-LOCAL (ramification and models).}
The permutation representation on the embeddings of a number field has
inertia invariants giving the Euler factors of $\operatorname{Spec}\mathcal
O_K$ at ramified primes as well. Thus every Dedekind zeta function is covered
by M, and abelian fields over $\mathbb Q$ are covered by A (ZT03).
If a factorization holds only away from a finite set $S$, the actual factors
at $S$ must be restored (ZT04). The counterexample ZT05,
\[
 \zeta_{\operatorname{Spec}\mathbb Z[1/p]}(s)
       =(1-p^{-s})\zeta(s),
\]
has additional zeros $2\pi i k/\log p$, $k\ne0$. Such local divisor lines
are not shifted global critical lines. This is a statement about the
integral model, not a failure of the generic-fibre motivic normalization.

\paragraph{AG-CM (the potentially CM domain).}
ZETA\_TWO CM01 supplies the all-weight extension to pure objects in the
absolute-Hodge category generated by Artin motives, Tate motives, and
potentially CM abelian varieties. Its cited CM descent input gives a signed
product of Hecke functions of one weight; norm twisting makes these
characters unitary. The required M therefore includes infinite-order
unitary characters. The result is meromorphic divisor RH, and does not
declare absolute-Hodge projectors algebraic. The exact prior locator is
\path{RETURN_ZETA_TWO_final_IMPLICATIONS.md}, Section 2, CM01,
with its reported Anderson Section 5.7, Corollary 5.7.2 and Remark 5.7.3.
The original source has now been recovered and checked: G.~W.~Anderson,
\emph{Cyclotomy and an extension of the Taniyama group}, Compositio
Math.~57 (1986), 153--217. The category is defined in Section 5.3,
printed p.~178 (PDF p.~27); Theorem 6, Section 5.6, printed pp.~179--180
(PDF pp.~28--29), identifies its one-dimensional characters with algebraic
Hecke characters, including the inertia and Frobenius comparisons.
Corollary 5.7.2 and Remark 5.7.3 are both on printed p.~181 (PDF p.~30).
The former gives integer exponents, with no positivity claim; the latter
gives the common weight. Its internal reference to ``Corollary 5.8.2'' is
a printed cross-reference error, not a new result used here.

\paragraph{AG-CM-REC and AG-CM-PERIOD (classical conclusions already
stronger than zero location).}
In this precise potentially CM absolute-Hodge category, equality of the
whole $L$-functions implies isomorphism of objects: Anderson,
Proposition 5.7.5, printed pp.~181--182 (PDF pp.~30--31). This is
unconditional. The proof uses semisimplicity and Corollary 5.7.4's
identification, for a weight-zero object, of the pole order at $1$ with
the dimension of its invariant part. Equality of the local factors makes
the tensor products with a dual agree away from finitely many primes;
their local factors are nonzero and nonpolar at $1$. A pole of
$L(M\otimes M^\vee,s)$ therefore gives a nonzero morphism
$M\longrightarrow N$, and induction on the common rank proves the
claim. The proof was read in full, rather than inferred from GRH.
This is a recognition theorem for semisimple motives; it supplies no
effective finite bound for general infinite CM monodromy.

Anderson's Section 5.8, Theorem 7, printed p.~182 (PDF p.~31), also
proves Deligne's period conclusion
$L(M,0)/c^+(M)\in\mathbb Q$ for the stated \emph{tractible, critical}
objects: those for which the $L$-function is a single algebraic Hecke
$L$-function over a totally real or CM field. This is a separate
classical special-value theorem, using Siegel and Blasius, and is not
enlarged to arbitrary motives by A or M. It complements, rather than
replaces, the cycle/height and integral comparison gates below.

\paragraph{AG-WEIGHT and AG-NONCM (the residual geometric factors).}
With geometric Frobenius,
\[
 \operatorname{wt}N(n)=\operatorname{wt}N-2n,
 \qquad L(N(n),s)=L(N,s+n).
\]
The line for weight $w$ is $\Re s=(w+1)/2$. Positive weight is not the
remaining obstruction: CM objects of positive weight are covered by M,
whereas $\operatorname{Sym}^2h^1(E)(1)$ for a non-CM elliptic curve still
has a non-CM factor despite having weight zero. ZT06 reduces a finite-type
scheme to finitely many smooth projective geometric factors and actual
exceptional fibres. ZT07 proves that, under A, an off-line elliptic zero
in $1/2<\Re s<3/2$ cannot be cancelled by the two Tate factors in the
good-model formula
\[
 \zeta_{\mathcal E}(s)
   =\zeta^S(s)\zeta^S(s-1)/L^S(E,s).
\]
These are retained reductions, not a proof of RH for the non-CM factor.
Framework locators carried from ZETA\_TWO are SGA 4 XVII Sections 5.1--5.3;
SGA 5 XV Section 3, no.~2, unnumbered theorem, equation (1);
EGA IV$_3$ (8.8.2); and Deligne, Weil I, Theorem 1.6, and Weil II,
Theorem 3.3.1.

\subsection{AG-FROB: a finite recognition theorem for Artin--Tate objects}
\label{ag:frob}

Let $L/\mathbb Q$ be a nontrivial finite Galois extension, with group $G$,
absolute discriminant $D_L$, and degree $n_L$. Put
\[
 B_L=(4\log D_L+\tfrac52 n_L+5)^2.
\]
Let
\[
 U=\bigoplus_j V_j(-j),\qquad W=\bigoplus_j W_j(-j),
\]
where all sums are finite and $V_j,W_j$ are actual finite-dimensional
characteristic-zero representations of $G$. Their Euler polynomials at
unramified primes are understood as those of the associated compatible
Artin--Tate systems; a single auxiliary prime may instead be chosen greater
than $B_L$.

\textbf{Statement.} Assume RH for $\zeta_L$ and $\zeta_{\mathbb Q}$.
If
\[
 \det(1-TF_p\mid U)=\det(1-TF_p\mid W)
 \quad\text{for every }p\le B_L\text{ unramified in }L,
\]
then $V_j\simeq W_j$ for every $j$, and hence $U\simeq W$.
For rational coefficients this is also an isomorphism in the rational
Artin--Tate category. A supplies the hypothesis when $L/\mathbb Q$ is
abelian; M supplies it in general. When $L=\mathbb Q$, the one prime $p=2$
suffices.

\begin{proof}
At an unramified prime, every eigenvalue in $V_j(-j)$ has the form
$p^j\omega$, where $\omega$ is a root of unity. Therefore its complex
absolute value is $p^j$. Distinct integers $j$ give distinct absolute
values, so equality of the total Euler polynomials separates uniquely into
equality of the polynomials in each weight. It follows that
$\operatorname{Tr}(F_p\mid V_j)=\operatorname{Tr}(F_p\mid W_j)$ for every
$j$ and every prime in the hypothesis.

Bach--Sorenson, \emph{Explicit bounds for primes in residue classes},
Math. Comp. 65 (1996), Theorem 5.1, printed p.~1729 (PDF p.~13), gives
an unramified prime $p\le B_L$ in every conjugacy class of $G$.
The precise analytic assumptions are stated in their Section 2,
printed p.~1720: zero-freeness of $\zeta_L$ and $\zeta_{\mathbb Q}$ in
$\Re s>1/2$. The reduction from cyclic to general Galois extensions is
Corollary 3.3, p.~1721. Thus the characters of $V_j$ and $W_j$ agree on
every conjugacy class. In characteristic zero finite-group representations
are semisimple, and equal characters give isomorphic representations.

For completeness, this last assertion also descends to a non-algebraically
closed characteristic-zero coefficient field $E$. The space of
$G$-equivariant maps is the solution space over $E$ of finitely many linear
equations. If it contains an invertible map after scalar extension, its
determinant polynomial is nonzero on that $E$-vector space. Since $E$ is
infinite it has a nonzero value at an $E$-point. That point is an invertible
$G$-equivariant map over $E$. For $L=\mathbb Q$, the eigenvalues $2^j$
already recover all multiplicities from the polynomial at $2$.
\end{proof}

In Grothendieck's language this is effective recognition inside a specified
semisimple Artin--Tate realization category. It extends
ALPOGE\_BRIDGES AQ-FINITE, whose sharper abelian conductor bound remains
available. It does not construct algebraic projectors for an arbitrary
geometric realization. Moreover, recognition of an Artin motive is weaker
than recognition of a finite cover as a scheme; the connected Gassmann
example in CX-GEO-004's companion finding makes this distinction explicit.
The forbidden-prime variant is handled there by a quadratic compositum.

\paragraph{AG-COVER: equal Artin motives need not give equal covers.}
Put $L=\mathbb Q(\sqrt2,\sqrt3)$ and let $K_1,K_2,K_3$ be its three
quadratic subfields. The finite schemes
\[
 X=\operatorname{Spec}L\amalg\operatorname{Spec}\mathbb Q
       \amalg\operatorname{Spec}\mathbb Q,
 \qquad
 Y=\coprod_{i=1}^3\operatorname{Spec}K_i
\]
are not isomorphic, but their rational Artin motives, and all their local
Artin Euler factors, are equal. Consequently
\[
 \zeta_L(s)\zeta_{\mathbb Q}(s)^2
       =\zeta_{K_1}(s)\zeta_{K_2}(s)\zeta_{K_3}(s).
\]
This identity includes the factors at ramified primes.

\begin{proof}
Write $G=\operatorname{Gal}(L/\mathbb Q)=C_2\times C_2$, with its three
subgroups $H_i$ of order two. The character of
$\mathbb Q[G]\oplus\mathbb Q^2$ takes value $6$ at the identity and
$2$ at every other element. Each nonidentity element belongs to exactly
one $H_i$, so the character of $\bigoplus_i\mathbb Q[G/H_i]$ has the
same values. Characteristic-zero semisimplicity and the descent argument
in AG-FROB give an isomorphism of rational representations. Isomorphic
representations have isomorphic inertia invariants with their residual
Frobenius action, proving equality of every local factor and hence the
zeta identity in the initial half-plane and by continuation. On the other
hand the degrees of the connected components are $4,1,1$ for $X$ and
$2,2,2$ for $Y$, precluding a scheme isomorphism.
\end{proof}

The same example over $\mathbb Z$ uses the respective rings of integers;
over $\mathbb Z[1/6]$ these are finite \emph{\'etale} covers. Thus even
complete Euler data forget the multiplication on an Artin algebra. A
applies to all fields here because they are abelian over $\mathbb Q$,
but does not restore the forgotten multiplication. The independent
transitive seven-element example CX-GEO-005 treats the connected
$G$-set distinction without asserting an arithmetic realization.

\subsection{AG-LAMBDA: the precise closure supplied by tensor operations}
\label{ag:lambda}

\textbf{Statement.} The finite direct sums of abelian Artin characters with
integer Tate twists are closed under tensor products, duals, exterior and
symmetric powers, and direct summands in characteristic zero. All their
$L$-functions have the A divisor-support conclusion. Replacing abelian
finite image by arbitrary finite image gives the corresponding M conclusion.
The same operations preserve the potentially CM absolute-Hodge category
used in CM01, so its M conclusion applies there as well.

\begin{proof}
Over an algebraic closure of the coefficient field an abelian Artin--Tate
object is a sum of lines $\chi_a(-j_a)$. A tensor product or a monomial
appearing in a symmetric or exterior power is a line
$(\prod_a\chi_a^{e_a})(-\sum_a e_aj_a)$. The character still has finite
abelian image and the Tate index is an integer. Duality replaces the
character and index by their inverses and negatives. Direct summands remain
in the same semisimple category, and descent preserves the asserted
factorization. The nonabelian assertion follows because all these
operations still factor through a finite group and commute with Tate
twists; apply ZT02 to their resulting Artin representations. Finally, the
potentially CM category is the tensor, dual, and summand closure in CM01.
Symmetric and exterior powers are the images of the usual idempotents in
the symmetric-group action on tensor powers, so belong to it too.
\end{proof}

The placement is SGA 6 V Section 2 and Section 7, especially (7.1.1),
and the determinant/$\lambda$-ring discussion carried in ZETA\_TWO ZT08
and SV01. These closures do not generate an arbitrary non-CM $GL_2$
motive from CM objects; enlarging a list of tensor operations is not the
missing analytic theorem.

\subsection{AG-SIGN: central signs with the necessary normalization}
\label{ag:sign}

\textbf{Statement.} Let $F$ be a nonzero holomorphic function on a
neighbourhood of $[c,\infty)$, real-valued there. Suppose $F$ has no zeros
on $(c,\infty)$ and is positive at one sufficiently large real point.
If $r=\operatorname{ord}_{s=c}F$, then
\[
 \frac{F^{(r)}(c)}{r!}>0.
\]
If, moreover, $F(c+z)=\varepsilon F(c-z)$ near $0$, with
$\varepsilon\in\{1,-1\}$, then $(-1)^r=\varepsilon$.
For a real, standardly normalized, holomorphic Hecke or CM $L$-function,
the absence of real zeros to the right of the centre follows from M
together with its usual analytic continuation and Euler-product theory.
The analogous Dirichlet statement uses A whenever no pole is present.

\begin{proof}
A continuous real nonvanishing function has constant sign on the connected
interval $(c,\infty)$. The normalization makes that sign positive.
Write $F(c+t)=a_rt^r+O(t^{r+1})$, with $a_r\ne0$. For $t>0$ the
left-hand side is positive, so $a_r>0$ after division by $t^r$ and passage
to the limit. Comparing the leading terms in the functional equation gives
$a_r=\varepsilon(-1)^r a_r$, hence the parity assertion.
\end{proof}

\paragraph{AG-HEIGHT (what this proves for a supplied height formula).}
If $F(c)=0$, then $F'(c)\ge0$: it is positive when $r=1$ and zero when
$r>1$. Consequently, an independently established identity
$H(Z)=C F'(c)$ with a specified constant $C>0$ yields $H(Z)\ge0$.
The elementary deduction is complete; the existence, domain, and
normalization of the height identity are extra inputs. It does not assert
that every homologically trivial cycle has such a formula, or that its
$L$-function has simple central vanishing. For an identity involving the
first nonzero derivative, the sign is strictly positive. For a formula with
an additional sign, that sign must be retained.

\paragraph{AG-SIGN-NORM (a counterexample to the unnormalized inference).}
Zero locations and functional-equation parity alone do not determine the
sign: $F(s)=s-c$ and $-F(s)$ have exactly the same zeros and the same odd
functional equation, but their central derivatives have opposite signs.
This is a counterexample to the stated analytic inference, not an example
of two standardly normalized arithmetic $L$-functions.

These findings refine \path{OURS_HODGE_BRIDGE_SKETCH_20261121.md},
Section 3.5. Their archive placement is cote 44, scans 2, 4 and 6,
Theorems 4--5, and scans 8--9, the letter to Serre of 12 August 1964,
as reported in that attached sketch. No original scan has been inferred
from this sign argument. The remaining arithmetic Hodge-index gate is
the cycle/height comparison for the stated class of cycles, not a sign
convention suppressed in a formula.

\subsection{AG-CRITICAL: why off-centre critical integers gain little}
\label{ag:critical}

\textbf{Statement.} Suppose $N$ is a pure, integral-weight, explicitly
Hecke-factored motive of weight $w$, with its usual functional equation
and archimedean factors. At an endpoint of the unitary strip, require each
factor in the specified signed Hecke factorization (or its dual, when the
functional equation is used) to be nonpolar. This factorwise requirement
prevents a denominator pole from producing a zero. Let $m$ be critical in
Deligne's sense: $L_\infty(N,s)$ and $L_\infty(N^\vee,1-s)$
are both regular at $s=m$. If $m\ne(w+1)/2$, then the standard
zero-free half-plane and boundary for unitary Hecke $L$-functions,
together with the functional equation, already give $L(N,m)\ne0$.
M is not needed for this conclusion.

The precise unitary boundary theorem used here is
D.~Ramakrishnan--R.~J.~Valenza, \emph{Fourier Analysis on Number Fields}
(1999), Section 7.7, Theorem 7-28, printed p.~289: for every number
field and every unitary idele class character, there is no zero at
$1+it$. Section 7.4, Theorem 7-19, printed p.~279, supplies the completed
continuation and functional equation, with its pure norm-character pole
exceptions retained. These are classical inputs independent of M.

\begin{proof}
After unitary normalization set $z=m-w/2$. If $w$ is even, $z$ is an
integer; if $w$ is odd, $z$ is a half-integer. An integer $m$ with
$0<z<1$ therefore has $z=1/2$, and is precisely the excluded centre.
For every other $m$ one has $z\ge1$ or $z\le0$. In the first case
standard unitary Hecke nonvanishing in $\Re z\ge1$, away from a pole,
gives the result for each factor and hence its finite signed product.
In the second case the functional equation relates the completed value
to its dual at $1-z\ge1$. Criticality ensures that both archimedean
factors are finite; gamma functions have no zeros. The nonzero dual
value therefore gives a nonzero original value. The factorwise nonpolar
condition excludes a Tate pole or its inverse being misread as a finite
nonzero special value. For a product with nonnegative exponents, the
usual no-pole condition implies this condition.
\end{proof}

Thus the strongest new real-axis information from M lies between the
off-centre integer points and at signs of central leading coefficients;
it does not evaluate regulators, periods, central ranks, or integral
torsion. This is the special-value distinction in ZETA\_TWO Section 4
and in the Hodge sketch's edge-versus-centre table. The relevant motivic
conventions are Deligne, \emph{Valeurs de fonctions $L$ et p\'eriodes
d'int\'egrales}, Sections 1.1--1.3; the period conjectures numbered 1.8
and 2.8 are separate inputs.

\subsection{AG-DET: Hecke versions of the square-variable determinant}
\label{ag:det}

The following formulation records exactly which analytic data are needed.
Let $\Lambda$ be a nonzero entire function of order at most one, with
\[
 \Lambda(c+z)=\varepsilon\Lambda(c-z),\qquad
 \varepsilon\in\{1,-1\},
\]
and let $r=\operatorname{ord}_{s=c}\Lambda$. Assume that every other zero
is $c+i\gamma$, with $\gamma$ real and nonzero. Choose one positive
$\gamma$ from each pair, with multiplicity, and write the sequence as
$\gamma_1,\gamma_2,\ldots$ in nondecreasing order, so $\gamma_1$
denotes the least positive zero when the sequence is nonempty.
Define the removable entire function
\[
 H(t)=t^{-r}\Lambda(c+it),\qquad
 D(w)=H(\sqrt w)/H(0).
\]
No choice of square-root branch affects $D$.

\textbf{Statement.} The diagonal operator
$B=\operatorname{diag}(\gamma_j^{-2})$ is positive and trace class, and
\[
 D(w)=\prod_j(1-w\gamma_j^{-2})=\det(I-wB).
\]
In particular this applies under M to standard entire selfdual completed
Hecke functions and to the entire completed $L$-function of a CM elliptic
curve, with its centre $c=1$. The usual order, functional equation,
and completion are additional classical inputs, not assertions supplied
by the word ``GRH''. Poles must first be removed in a separately specified
entire completion; arbitrary meromorphic Artin quotients are not included.

\begin{proof}
Comparison of the first nonzero Taylor coefficients shows
$\varepsilon=(-1)^r$. Hence $H(-t)=H(t)$, so $D$ is entire by the even
power series for $H$. Its order is at most one half: a bound of order
$1+\delta$ for $H(t)$ gives a bound of order $(1+\delta)/2$ for $D(w)$.
Its zeros, with multiplicity, are exactly $\gamma_j^2$.
The convergence exponent of the zeros of an entire function of order at
most one is at most one. In particular
$\sum_j|\gamma_j|^{-2}<\infty$. Equivalently the diagonal positive
operator $B$ has finite trace norm. Hadamard factorization for $D$, whose
order is strictly less than one, has no nonconstant exponential factor
and uses genus-zero factors. Since $D(0)=1$, its constant factor is one.
This gives the displayed product. The Fredholm determinant of a diagonal
trace-class operator is that same absolutely locally uniformly convergent
product, proving the identity. An exact analytic locator for the
factorization used here is C.~McMullen, \emph{Advanced Complex Analysis},
Math 213a course notes, version 16 November 2026, Theorem 3.16,
printed p.~92 (PDF p.~93); Theorems 3.10--3.12 give the zero-exponent
bound used just above. This is classical background, not a new claim
dependent on that revision date.
\end{proof}

This extends the prior zeta realization BR-TR-01 to a precisely specified
Hecke domain. Its relationship to cote 1 scans 79--80 and 111--112 is
the trace/Fredholm theorem reported in
\path{OURS_INTAKE_RESEARCH_SA23_20261120.md}, Section 4.1,
geo-H-lidskii and geo-H-genus. Those sheets are not themselves an
arithmetic RH claim. The operator here is defined from the zeros; the
construction has not supplied a cohomology theory or an operator defined
independently from prime data.

\subsection{AG-MOMENT and AG-HANKEL: recoverable positive spectral data}
\label{ag:moment}

Assume the hypotheses of AG-DET and infinitely many distinct positive
zeros. Put $\lambda_j=\gamma_j^{-2}$ and
\[
 s_k=\operatorname{Tr}(B^k)=\sum_j\lambda_j^k\quad(k\ge1).
\]
The logarithmic determinant expansion is
\[
 -\log D(w)=\sum_{k\ge1}\frac{s_k}{k}w^k
 \qquad (|w|<\gamma_1^2).
\]

\textbf{AG-MOMENT.} For every $k\ge1$,
\[
 s_k^{-1/k}\le\gamma_1^2\le\frac{s_k}{s_{k+1}}.
\]
The left bounds increase to $\gamma_1^2$; the right bounds decrease to it.
If $m_1$ is the multiplicity of the least positive zero, then
$\lim_{k\to\infty}s_k\gamma_1^{2k}=m_1$.

\begin{proof}
Let $a=\max_j\lambda_j=\gamma_1^{-2}$. Then
$s_k\ge a^k$ and $s_{k+1}\le a s_k$, proving the two inequalities.
The $\ell^p$ norm monotonicity gives
$s_{k+1}^{1/(k+1)}\le s_k^{1/k}$. One direct proof is to divide by
$s_k^{1/k}$: the resulting nonnegative numbers have $k$-th powers
summing to one, are each at most one, and thus have $(k+1)$-th powers
summing to at most one. Hence the left bounds increase.
Cauchy--Schwarz applied to
$\lambda_j^{k/2}$ and $\lambda_j^{(k+2)/2}$ gives
$s_{k+1}^2\le s_ks_{k+2}$, so the right bounds decrease.

There are finitely many indices with $\lambda_j=a$, since $\sum\lambda_j$
is finite. For $k\ge1$,
$0\le(\lambda_j/a)^k\le\lambda_j/a$, whose sum is finite.
Dominated convergence shows
$s_k/a^k\to m_1$. It follows that both $s_k^{1/k}$ and
$s_{k+1}/s_k$ tend to $a$, which proves both limits and the
multiplicity formula.
\end{proof}

\textbf{AG-HANKEL.} For all integers $h\ge1$ and $d\ge0$, the matrix
\[
 (s_{h+i+j})_{0\le i,j\le d}
\]
is positive definite. Its determinant is therefore strictly positive.

\begin{proof}
For a nonzero vector $a=(a_0,\ldots,a_d)$ and
$P(X)=\sum_{i=0}^d a_iX^i$, its Hermitian quadratic form equals
\[
 \sum_{j\ge1}\lambda_j^h|P(\lambda_j)|^2.
\]
This converges because the $\lambda_j$ are bounded and
$\sum\lambda_j^h<\infty$. All terms are nonnegative. A nonzero
polynomial has only finitely many distinct zeros, whereas the positive
$\lambda_j$ have infinitely many distinct values. At least one term
is positive. This proves strict positive definiteness.
\end{proof}

These consequences strengthen the two-by-two moment inequalities carried
in BR-TR-02. They describe actual determinant coefficients available in
the cote 1 analytic framework, while leaving the arithmetic/geometric
realization gate unchanged. They assert no simplicity and no
random-matrix statistics.

\subsection{AG-SIGNED: a finite moment witness for a surviving pole}
\label{ag:signed}

Consider a finite signed product of specified entire selfdual functions
of order at most one, with a common centre $c$, each satisfying the
zero-location hypotheses of AG-DET. Let $r\in\mathbb Z$ be the net
central order. Divide out the central Laurent factor and normalize its
nonzero constant term. Its square-variable quotient has the form
\[
 D(w)=\prod_{\lambda\in\mathcal L}(1-w\lambda)^{m_\lambda},
 \qquad m_\lambda\in\mathbb Z\setminus\{0\},\quad
 \lambda>0,\quad \sum_\lambda |m_\lambda|\lambda<\infty,
 \qquad D(0)=1.                                      \tag{AG-S1}
\]
Here $\mathcal L$ has no positive accumulation point, and the signed
multiplicities are taken \emph{after} cancellation. The form (AG-S1)
follows by applying AG-DET to each factor. More generally the proposition
below applies whenever (AG-S1) is independently supplied. Put
\[
 s_k=\sum_\lambda m_\lambda\lambda^k\quad(k\ge1),\qquad
 H_d=(s_{1+i+j})_{0\le i,j\le d}.
\]
These moments are the local logarithmic coefficients
$-\log D(w)=\sum_{k\ge1}s_kw^k/k$ near $0$.

\textbf{Statement.} All the matrices $H_d$ are positive semidefinite
if and only if $m_\lambda>0$ for every $\lambda\in\mathcal L$.
Equivalently, the normalized quotient has no noncentral poles. The
central order condition $r\ge0$ must be checked separately. If a
negative multiplicity survives, a finite matrix $H_d$ already detects it,
with the explicit witness described below.

\begin{proof}
Set $a=\sup\mathcal L$ and introduce the finite signed measure
\[
 \mu=\sum_\lambda m_\lambda\lambda\,\delta_\lambda
 \quad\hbox{on }[0,a].
\]
The empty support case is immediate. For every complex polynomial
$P(X)=\sum_{i=0}^d b_iX^i$, absolute convergence gives
\[
 \sum_{i,j=0}^d\overline b_i s_{1+i+j}b_j
       =\int |P(x)|^2\,d\mu(x).
\]
Positive multiplicities therefore imply all the required inequalities.

Conversely, suppose $m_{\lambda_0}<0$. The point $\lambda_0>0$ is
isolated in the compact set $\{0\}\cup\mathcal L$. Let
\[
 \delta=\operatorname{dist}
   (\lambda_0,(\{0\}\cup\mathcal L)\setminus\{\lambda_0\})>0,
 \qquad R>a+\lambda_0,
 \qquad q=1-\delta^2/R^2<1.
\]
For each nonnegative integer $N$, the polynomial
\[
 P_N(x)=\left(1-\frac{(x-\lambda_0)^2}{R^2}\right)^N
\]
has degree at most $2N$, takes the value $1$ at $\lambda_0$, and has
$|P_N(\lambda)|\le q^N$ at every other point of $\mathcal L$.
Writing $W=\sum_{\lambda\ne\lambda_0}|m_\lambda|\lambda$, we obtain
\[
 \int |P_N|^2\,d\mu
   \le -|m_{\lambda_0}|\lambda_0+Wq^{2N}.
\]
Choose $N$ with $Wq^{2N}<|m_{\lambda_0}|\lambda_0$; take $N=0$
if $W=0$. Then the coefficient vector of $P_N$ is a negative
quadratic-form witness for the matrix $H_{2N}$, of size $2N+1$.
This proves the converse and its finite-witness assertion. In (AG-S1)
a negative net multiplicity is exactly a pole at $w=\lambda^{-1}$,
which proves the analytic equivalence. Removing the central Laurent
factor erased its order, explaining the separate condition on $r$.
\end{proof}

Geometrically this strengthens the cote 1 determinant bridge in the
meromorphic setting. A or M can first localize an identified divisor;
positivity of these finite moment matrices is then an exact additional
condition preventing cancellation from leaving poles. It is not supplied
by zero location alone. In particular, universal Hecke GRH plus Brauer
induction does not prove this positivity, and hence does not prove Artin
holomorphy. The criterion is for the specified quotient: symmetrizing
an arbitrary nonselfdual function with its dual can cancel a pole between
the two factors, so holomorphy of that product must not be attributed to
each unsymmetrized factor. The local character-matrix test in MT-ART02
is a complementary criterion at an individual Artin divisor point.

\subsection{AG-EXT: the precise information that the determinant forgets}
\label{ag:ext}

The earlier findings BR-MIX-01 and ZETA\_TWO SV01--SV02 are retained.
At a good prime an adapted matrix for a mixed object is block triangular,
so its determinant equals that of its weight-graded object. The nonsplit
Kummer extension from the unit $2$ over $\mathbb Z[1/2]$ and its split
Tate object therefore have the same incomplete $L$-function. The omission
of ramified factors is essential: inertia invariants need not preserve a
short exact sequence. These findings refute recovery of extensions from
semisimplified good Euler data; they do not assert that all richer
arithmetic data agree.

The exact independent background locator is Deligne--Goncharov,
\emph{Groupes fondamentaux motiviques de Tate mixte},
arXiv:math/0302267v2 (27 February 2003), (2.1.3), printed p.~14
(PDF p.~16):
\[
 \operatorname{Ext}^1(\mathbb Q(0),\mathbb Q(n))=
 \begin{cases}
 0&n\le0,\\
 \mathcal O_S^\times\otimes\mathbb Q&n=1,\\
 K_{2n-1}(k)\otimes\mathbb Q&n\ge2.
 \end{cases}
\]
Propositions 2.2--2.3 on the same page identify the corresponding graded
free motivic Lie algebra, while Section 2.4 explicitly distinguishes a
free presentation from a canonical set of generators. These statements
are unconditional background. They do not calculate integral torsion or
the pro-$\ell$ kernels in the prompt's cote 149 programme.

\subsection{AG-KUMMER: an integral obstruction on the actual cote 149 sheets}
\label{ag:kummer}

The original Montpellier PDF was inspected here, not merely the intake.
It has a cover page: PDF p.~80 has repository footer 79, and so on.
At PDF p.~80/footer 79, equations (13)--(17) relate an inertia-normalizer
extension, its $H^2$ class, and Kummer cohomology. PDF p.~81/footer 80,
(18)--(19), identifies the relevant line bundle as the normal bundle of
the section. PDF p.~82/footer 81, (20)--(21), and PDF p.~85/footer 84,
(22), describe the $H^1$ and completed-unit ambiguity. These are the
actual geometric objects on the cited sheets. Their link to LM6
A.44--A.45 is supplied by ORDER item 31; no theorem about an irregular
prime or $K_8(\mathbb Z)$ is printed on these pages.

Here is a precise integral consequence, independent of A and M. Let $S$
be a scheme on which $\ell$ is invertible, put $P=\operatorname{Pic}(S)$,
and let $\operatorname{Br}'(S)=H^2_{\text{ét}}(S,\mathbb G_m)$.
Write $\widehat P_\ell=\varprojlim_nP/\ell^nP$ and
$T_\ell A=\varprojlim_n A[\ell^n]$, with transition maps multiplication
by $\ell$. There are exact sequences
\begin{align*}
 0&\longrightarrow\widehat P_\ell
   \longrightarrow\varprojlim_nH^2_{\text{ét}}(S,\mu_{\ell^n})
   \longrightarrow T_\ell\operatorname{Br}'(S)\longrightarrow0,\\
 0&\longrightarrow\Gamma(S,\mathcal O_S^*)^{\wedge}_\ell
   \longrightarrow\varprojlim_nH^1_{\text{ét}}(S,\mu_{\ell^n})
   \longrightarrow T_\ell P\longrightarrow0.
\end{align*}
We deliberately use inverse-limit cohomology: identifying it with a
chosen definition of continuous cohomology requires its own
$\varprojlim{}^1$ check. If $P$ is finitely generated, the second right
term is zero. A line bundle $\mathcal L$ has zero compatible integral
$\ell$-adic Chern class if and only if its class in $P$ is torsion of
order prime to $\ell$.

\begin{proof}
The Kummer exact sequence gives
\[
 0\longrightarrow P/\ell^nP
 \longrightarrow H^2(S,\mu_{\ell^n})
 \longrightarrow\operatorname{Br}'(S)[\ell^n]\longrightarrow0,
\]
and the analogous degree-one sequence with units modulo $\ell^n$th
powers and $P[\ell^n]$. The exact source is Stacks Project,
Lemma 59.28.1, tag 03PK, and its immediately following cohomology
sequence. The transition maps on the left terms are surjective.
Consequently a compatible sequence on the right lifts compatibly in
the middle: after choosing a lift at level $n$, correct any lift at
level $n+1$ by a lift of its discrepancy in the left term. This recursive
construction proves right-surjectivity after taking inverse limits;
left-exactness is immediate. It proves both displayed sequences.

Now write the finitely generated group $P$ as
$\mathbb Z^a\oplus P[\ell^\infty]\oplus P_{\ell'}$, with finite torsion
summands. Its $\ell$-adic completion is
$\mathbb Z_\ell^a\oplus P[\ell^\infty]$, and the kernel of completion
is exactly $P_{\ell'}$. Also no nonzero compatible infinitely
$\ell$-divisible torsion sequence can occur in a finitely generated
group, so $T_\ell P=0$. The injection in degree two then proves the
Chern-class criterion.
\end{proof}

\paragraph{AG-KUMMER-2: a fully explicit integral/rational separation.}
Let $K=\mathbb Q(\sqrt{-35})$, $\omega=(1+\sqrt{-35})/2$, and
$S=\operatorname{Spec}\mathbb Z[\omega,1/2]$. The invertible ideal
$\mathfrak p=(3,\omega)$ determines a line bundle $\mathcal L$ whose
compatible integral $2$-adic Chern class is nonzero, although its image
after tensoring with $\mathbb Q_2$ is zero.

\begin{proof}
The ring of integers is $\mathbb Z[\omega]$ and its discriminant is
$-35$. The minimal polynomial $X^2-X+9$ is irreducible modulo $2$,
so $(2)$ is the only prime over $2$, and is principal. Modulo $3$ it is
$X(X-1)$; thus the primes over $3$ are
$\mathfrak p=(3,\omega)$ and $\overline{\mathfrak p}=(3,\omega-1)$.
Since $N(\omega)=9$ and $\omega$ is a unit modulo
$\overline{\mathfrak p}$, unique ideal factorization gives
$\mathfrak p^2=(\omega)$. The ideal $\mathfrak p$ is not principal:
an integral generator would have norm $3$, whereas
\[
 N(a+b\omega)=a^2+ab+9b^2=(a+b/2)^2+35b^2/4
\]
cannot be $3$ for integers $a,b$.

Minkowski's bound is $2\sqrt{35}/\pi<4$, so every ideal class has a
representative of norm $1$, $2$, or $3$. There is no ideal of norm $2$;
the two of norm $3$ have the same class of order two. Hence the class
group is $\mathbb Z/2$. The precise classical bound is Milne,
\emph{Algebraic Number Theory}, v3.08 (19 July 2020), Theorem 4.3,
PDF p.~72. Inverting $2$ quotients the class group by the classes of
the primes over $2$, which here are principal, so
$\operatorname{Pic}(S)=\mathbb Z/2$ as well. To see this last assertion
directly, fractional ideals of the localization are obtained by forgetting
the valuation at $(2)$; a relation introduced in the localization differs
from a principal ideal only by a power of that principal prime.

AG-KUMMER now injects this nonzero class into the inverse limit of
$H^2(S,\mu_{2^n})$. It has order two there, so rationalization kills it.
The total space of $\mathcal L$ has a zero section with normal bundle
$\mathcal L$, making the example a literal normal-bundle obstruction of
the kind appearing in cote 149. No claim identifying this particular
scheme's group-cohomology extension class is needed.
\end{proof}

This proves a concrete limit of rational information: a rationally zero
normal Chern class need not be integrally zero. A supplies Dirichlet GRH
for the quadratic field here, but changes none of this calculation.
The finite-level Kummer test is therefore a genuine extra gate for an
integral inertia-splitting argument. It does not refute the original
LM6 statements, whose comparison maps must still be supplied.

\subsection{AG-C34-TYPE and AG-C34-REPAIR: resolving the commutator notation}
\label{ag:c34repair}

The official cote 34 PDF has a cover: PDF p.~53 is manuscript/footer
p.~52, and PDF p.~64 is p.~63. Theorem 1.1 on the former page fixes a
field $k$ finitely generated over its prime field, a normal geometrically
connected $k$-scheme $X$ locally of finite type, a prime
$\ell\ne\operatorname{char}k$, and an open subgroup
$\pi\subseteq\pi_1(X)$. It defines $\pi_0$ to be the image in
$G_k$ and $I=\pi\cap\pi_1(X_{\bar k})$, in
\[
 1\longrightarrow I\longrightarrow\pi\longrightarrow\pi_0
 \longrightarrow1.
\]
Its $I(\ell)$ is the maximal abelian pro-$\ell$ quotient of $I$.
The theorem states that $I(\ell)$ is finitely generated over
$\mathbb Z_\ell$ and that its $\pi_0$-coinvariants are finite.
The proof was located and read in Sections 2--3, PDF pp.~70--78
(manuscript pp.~69--77), rather than guessed from the intervening
applications. Section 2.4, PDF p.~72, makes the dual $H^1$-invariant
reduction; Section 2.5, PDF pp.~73--74, uses a curve compactification,
its cyclotomic boundary, and the finite-field curve theorem after
spreading out. Sections 3.1--3.3, PDF pp.~75--78, reduce normal $X$ to
a generic hyperplane curve using surjectivity on fundamental groups.
These are classical geometric inputs; A and M do not enter them.

In the crossed-out Remark 1.17 on PDF p.~64, $H$ is a quotient of $I$
whose kernel is $\pi$-stable, and the readable final hypothesis is that
$H$ is a solvable pro-$\ell$ group. The preceding proposed
$\ell$-adic Lie-group language is itself struck out. The remark writes
\[
 1\longrightarrow H\longrightarrow\pi'\longrightarrow\pi_0
 \longrightarrow1
\]
and places the closure of $[\pi_0,H]$ as an open subgroup of $\pi'$.
This is a tentative, struck manuscript assertion, not a theorem.

\paragraph{AG-C34-TYPE (the literal ambient-group obstruction).}
For an infinite $\pi_0$, every commutator of a lift of an element of
$\pi_0$ with an element of $H$ lies in $H$, and its closure cannot be
open in $\pi'$. Indeed the quotient by $H$ is the infinite group
$\pi_0$, whereas a subgroup open in a profinite group has finite index.
For a concrete example satisfying the manuscript's geometric hypotheses,
take $X=\mathbb G_m/\mathbb F_q$ and $\ell\nmid q$. Its geometric
pro-$\ell$ quotient $H=\mathbb Z_\ell(1)$ is abelian and hence
solvable; the rational point $1$ gives a section. Arithmetic Frobenius
acts as multiplication by $q$, so the action commutators generate
$(q-1)\mathbb Z_\ell(1)$. This is open in $H$ but is not open in
$\pi'$, whose quotient is $G_{\mathbb F_q}\simeq\widehat{\mathbb Z}$.
Thus the literal ambient group in the inherited Q03 cannot be retained.
For nonabelian $H$ there is a second issue: $\pi_0$ acts canonically
only by outer automorphisms unless a section or lifts are specified.

\paragraph{AG-C34-REPAIR (the intrinsic corrected statement).}
Let $H$ be \emph{any} pro-$\ell$ quotient of $I$ with $\pi$-stable
kernel, and let $P$ be the resulting quotient extension of $\pi_0$ by
$H$. Define $K$ to be the closed subgroup generated by all
$[p,h]$, $p\in P$, $h\in H$. Then $K$ is open in $H$, and
\[
 H/K\simeq (H^{\mathrm{ab}})_{\pi_0},
 \qquad [H:K]\le |I(\ell)_{\pi_0}|<\infty.       \tag{AG-C34.1}
\]
In particular the solvability and Lie-group hypotheses are unnecessary
for this precise corrected conclusion.

\begin{proof}
Because $H$ is normal in $P$, conjugating a commutator $[p,h]$ by
$a\in P$ gives $[apa^{-1},aha^{-1}]$. Hence $K$ is normal in $P$.
It contains the closed commutator subgroup of $H$, so $H/K$ is abelian;
all conjugation by $P$ is trivial on it. Conjugation on
$H^{\mathrm{ab}}$ factors canonically through $\pi_0$, irrespective
of any section. Killing its action differences is exactly the additional
quotient by the images of the commutators $[p,h]$. This proves the
isomorphism in (AG-C34.1).

The surjection $I\twoheadrightarrow H$ induces an equivariant
surjection $I(\ell)\twoheadrightarrow H^{\mathrm{ab}}$ by the
universal property of the maximal abelian pro-$\ell$ quotient.
Taking coinvariants remains surjective: it is simply quotienting both
compact abelian groups by the closed subgroups generated by their
action differences. Theorem 1.1 therefore supplies a surjection from
the finite group $I(\ell)_{\pi_0}$ to $H/K$. This proves the bound.
Since $K$ is closed, finite index in the compact group $H$ implies
openness (the finite Hausdorff quotient is discrete).
\end{proof}

This repairs the archived Q03 as a theorem already implied by its
stated abelian predecessor. It does not identify the unqualified
symbol $[\pi_0,H]$ in every possible intended meaning. A smaller
commutator subgroup formed from a chosen section, or openness in a
different ambient group, needs a separately stated question. In
particular the repair is not described as an automatic proof of all
interpretations of the crossed-out remark.

\paragraph{AG-C34-INDEX (an explicit Frobenius refinement).}
More generally, let $M=H^{\mathrm{ab}}$ be a finitely generated
$\mathbb Z_\ell$-module with continuous $\Gamma$-action, let $T$ be
its finite torsion subgroup, and let $\Lambda=M/T$. Suppose some
$F\in\Gamma$ has no eigenvalue $1$ on
$\Lambda\otimes\mathbb Q_\ell$. Then
\[
 |\operatorname{coker}(F-1\mid M)|
   =|T^F|\ell^{v_\ell\det(1-F\mid\Lambda\otimes\mathbb Q_\ell)},
 \qquad
 |M_\Gamma|\le
 |T^F|\ell^{v_\ell\det(1-F\mid\Lambda\otimes\mathbb Q_\ell)}.
                                                               \tag{AG-C34.2}
\]
Equality in the second assertion holds if $\Gamma$ is topologically
generated by $F$.

\begin{proof}
The assumption gives $\Lambda^F=0$. Apply the snake lemma to
$0\to T\to M\to\Lambda\to0$ with vertical maps $F-1$.
It yields
\[
 0\longrightarrow T/(F-1)T\longrightarrow M/(F-1)M
 \longrightarrow\Lambda/(F-1)\Lambda\longrightarrow0.
\]
For an endomorphism of the finite group $T$, kernel and cokernel have
the same cardinality. Smith normal form over $\mathbb Z_\ell$ gives
the order of the last group as the displayed power of $\ell$.
This proves the equality. The continuous image $(F-1)M$ is compact
and hence closed; the full coinvariant quotient is a quotient of it.
If the powers of $F$ are dense in $\Gamma$, every action difference
is a limit of differences of powers, all of which belong to
$(F-1)M$, so the two quotients agree.
\end{proof}

For a smooth proper geometrically connected curve $C/\mathbb F_q$,
let $H$ be its full geometric pro-$\ell$ fundamental group and $J$ its
Jacobian. Then $H^{\mathrm{ab}}=T_\ell J$, compatibly with Galois
action, and the corrected commutator quotient is
\[
 H/\overline{[P,H]}
       \simeq T_\ell J/(F-1)T_\ell J
       \simeq J(\mathbb F_q)[\ell^\infty].             \tag{AG-C34.3}
\]
For the fundamental-group comparison, an exact modern primary locator
is J.~S.~Milne, \emph{Jacobian Varieties}, Proposition 9.1,
pp.~30--32, together with \emph{Abelian Varieties}, Remarks 15.3
and 15.5, pp.~28--29. To check the last isomorphism directly, use
$0\to T_\ell J\to V_\ell J\to J(\overline{\mathbb F}_q)
[\ell^\infty]\to0$. The fixed subgroup of $T_\ell J$ is zero:
$J(\mathbb F_q)$ is finite, so a compatible indefinitely divisible
sequence of its $\ell$-power torsion points vanishes. Thus $F-1$ is
invertible on $V_\ell J$, and the snake lemma gives (AG-C34.3).
Consequently the index is exactly $\ell^{v_\ell\#J(\mathbb F_q)}$.
This finite-field calculation even uses less than the full Weil bound.
In higher-dimensional smooth proper situations purity supplies the
corresponding no-eigenvalue-$1$ test, but finite-field purity is already
classical and no new global GRH input is involved.

\subsection{AG-C34-SECTION: the solvable quotient with a chosen section}
\label{ag:c34section}

There is a second precise interpretation of the struck Remark 1.17 that
can now be settled. Keep the exact geometric hypotheses and notation of
AG-C34-REPAIR. Suppose the quotient extension
$1\to H\to P\to\pi_0\to1$ has a specified continuous section, so
$P=H\rtimes\pi_0$. Assume that $H$ is solvable of finite derived
length, using closed derived subgroups. This is stronger than merely
being an inverse limit of finite solvable groups; every pro-$\ell$
group already satisfies that latter condition. For the action $\alpha$
given by the section put
\[
 N=\overline{\langle h^{-1}\alpha_g(h):h\in H,
                                      \ g\in\pi_0\rangle}.
\]
\textbf{Statement.} The subgroup $N$ is normal in $P$ and open in
$H$. Thus the chosen-section action commutators have finite index in
the geometric group. They need not equal the intrinsic subgroup
$\overline{[P,H]}$ of AG-C34-REPAIR.

\begin{proof}
For a fixed automorphism $\alpha$, write
$c_\alpha(h)=h^{-1}\alpha(h)$. The identity
\[
 c_\alpha(hx)=x^{-1}c_\alpha(h)x\,c_\alpha(x)
\]
shows that the subgroup generated by all the $c_\alpha(h)$ is normal
in $H$. Applying $\alpha_g$ to $c_{\alpha_f}(h)$ gives
$c_{\alpha_{gfg^{-1}}}(\alpha_g(h))$, so it is also stable under
$\pi_0$. Its closure $N$ is consequently normal in $P$. Set
$Q=H/N$. The induced action of $\pi_0$ on $Q$ is trivial by
construction, whence $P/N=Q\times\pi_0$.

We claim that every open subgroup $U$ of $Q$ has finite abelianization.
Take the inverse image of $U\times\pi_0$ first in $P$ and then in
the original open arithmetic subgroup $\pi$. This gives an open
subgroup $\pi_U$ of $\pi_1(X)$ with the same arithmetic image
$\pi_0$. Its geometric kernel $I_U$ maps onto $U$. Conjugation on
$U^{\mathrm{ab}}$ is trivial: the $U$-part acts by inner
automorphisms and the $\pi_0$-part acts trivially. Therefore
\[
 I_U(\ell)_{\pi_0}\twoheadrightarrow U^{\mathrm{ab}}.
\]
The left group is finite by Theorem 1.1 applied to the open subgroup
$\pi_U$, proving the claim. This use of the theorem on \emph{all
the relevant finite covers}, rather than on $H$ alone, is essential.

Let $Q^{(0)}=Q$ and
$Q^{(i+1)}=\overline{[Q^{(i)},Q^{(i)}]}$. Starting with the open
subgroup $Q$, the claim implies inductively that $Q^{(i+1)}$ has
finite index in the open subgroup $Q^{(i)}$, and hence is open in
$Q$. Since $Q$ is a quotient of the solvable group $H$, this finite
derived series ends at $1$. Thus $1$ is open in $Q$, so the compact
group $Q$ is finite. Equivalently $N$ is open in $H$.

For the non-equality assertion, take a finite nonabelian $\ell$-group
$H$ and let $\pi_0$ act trivially. Then $N=1$ whereas
$\overline{[P,H]}=[H,H]\ne1$. Both are open in $H$. This last
example concerns the subgroup identity, and asserts no realization of
that abstract example by a particular arithmetic variety.
\end{proof}

This is additional movement on BRIDGES Q03: under the explicitly stated
section hypothesis, its solvability condition is sufficient for the
action-commutator conclusion in the corrected ambient group $H$.
No section is inferred from purity or from A/M. Without a section the
intrinsic assertion AG-C34-REPAIR remains the canonical formulation.

\subsection{AG-C34-NONSOLV: why compact analytic groups alone do not suffice}
\label{ag:c34nonsolv}

The finite-derived-length condition just used cannot be replaced, at
the level of this group argument, by finite generation or compact
$\ell$-adic analyticity. The following counterexample is deliberately
abstract; no geometric realization is asserted.

Fix an odd prime $\ell$ and put
\[
 H_n=\ker\bigl(\mathrm{SL}_2(\mathbb Z_\ell)
            \longrightarrow\mathrm{SL}_2(\mathbb Z/\ell^n)\bigr),
 \qquad H=H_1.
\]
Then $H$ is an infinite, topologically finitely generated compact
pro-$\ell$ group, and every open subgroup of $H$ has finite
abelianization. Nevertheless, for the trivial action of any profinite
group $\Gamma$, the action-commutator subgroup is $1$, which is not
open in $H$. In particular the absence of eigenvalue $1$ on
$H^{\mathrm{ab}}\otimes\mathbb Q_\ell$ does not repair that
inference: here this vector space is zero.

\begin{proof}
For each $n\ge1$, the map
$1+\ell^n A\mapsto A\bmod\ell$ identifies
$H_n/H_{n+1}$ with the additive group
$\mathfrak{sl}_2(\mathbb F_\ell)$. The determinant equation gives
trace zero. Surjectivity follows from the upper and lower unipotent
matrices and the diagonal matrix
$\operatorname{diag}(1+\ell^n,(1+\ell^n)^{-1})$, whose leading terms
are the basis
\[
 E=\begin{pmatrix}0&1\\0&0\end{pmatrix},\quad
 F=\begin{pmatrix}0&0\\1&0\end{pmatrix},\quad
 T=\begin{pmatrix}1&0\\0&-1\end{pmatrix}.
\]
Thus $H$ is an inverse limit of finite $\ell$-groups and is compact;
the upper unipotent matrices with entry in $\ell\mathbb Z_\ell$
already show that it is infinite.
It is explicitly a three-dimensional $\ell$-adic analytic group:
write its matrices as
$\bigl(\begin{smallmatrix}a&b\\c&(1+bc)/a\end{smallmatrix}\bigr)$
with $a\in1+\ell\mathbb Z_\ell$ and
$b,c\in\ell\mathbb Z_\ell$. These three parameters are analytic
coordinates, and multiplication and inversion are analytic in them.

The three displayed lifts at level $1$ topologically generate $H$.
Indeed, for $n\ge1$ and odd $\ell$,
\[
 (1+\ell^n A)^\ell\equiv1+\ell^{n+1}A
                       \pmod{\ell^{n+2}}.
\]
Their successive $\ell$-power iterates therefore give all three basis
vectors in every $H_n/H_{n+1}$. Any element of $H$ can be matched
successively modulo $H_2,H_3,\ldots$ by products of elements in the
closed group they generate. The products converge to the given element,
so that closed group is all of $H$.

For $n,m\ge1$, direct matrix multiplication gives
\[
 [1+\ell^n A,1+\ell^m B]
       \equiv1+\ell^{n+m}(AB-BA)
                      \pmod{\ell^{n+m+1}}.
\]
The Lie brackets $[E,F]=T$, $[T,E]=2E$, and $[T,F]=-2F$
span $\mathfrak{sl}_2(\mathbb F_\ell)$ because $\ell$ is odd.
Let $C_n=\overline{[H_n,H_n]}$. For every $k\ge2n$, commutators
of $H_n$ with $H_{k-n}\subseteq H_n$ therefore show that
$C_n\cap H_k$ surjects onto $H_k/H_{k+1}$. Successively lifting a
given element of $H_{2n}$ through these quotients and taking the
convergent product proves
\[
 H_{2n}\subseteq C_n.
\]
Every open subgroup $U\subseteq H$ contains some $H_n$. Its closed
commutator subgroup contains $C_n$, and hence $H_{2n}$, so
$U^{\mathrm{ab}}$ is finite. Taking $U=H$ also proves
$H^{\mathrm{ab}}\otimes\mathbb Q_\ell=0$.
For a trivial $\Gamma$-action, each action commutator is $1$,
which has infinite index in $H$. This proves the claimed failure.
\end{proof}

The example pinpoints what AG-C34-SECTION uses beyond abelian
coinvariant control: its induction must terminate. It does not challenge
the intrinsic commutator theorem, since $\overline{[H,H]}$ is open
here, nor any additional geometric constraints that could exclude this
trivial-action example in a proposed realization.

\subsection{Cote 149: the fundamental-group class and the section obstruction}
\label{ag:c149-comparison}

The comparison omitted by a bare Kummer calculation is available under
precise homotopy hypotheses. Borne--Emsalem--Stix,
\emph{Lifting Galois sections along torsors}, Math. Nachr.~288 (2015),
1939--1951, arXiv:1301.4429v2, Proposition 3.1 and Proposition 3.4
(PDF pp.~5--6), and Proposition 4.4 (PDF pp.~8--9), supply the
following classical statement. The complete proofs and the continuous
cohomology argument in Proposition 4.2 (PDF p.~7) have been read.

\paragraph{AG-C149-COMP (the exact comparison).}
Let $S$ be a geometrically connected, geometrically unibranch variety
over a characteristic-zero field $k$, with
$\pi^{\text{ét}}_2(S,\bar s)=0$. Let $\mathcal L$ be a line
bundle and $\mathcal L^\times$ its complement of the zero section.
For compatible base points there is an exact sequence
\[
 1\longrightarrow\widehat{\mathbb Z}(1)
 \longrightarrow\pi_1(\mathcal L^\times)
 \longrightarrow\pi_1(S)\longrightarrow1.
\]
The conjugation action on its kernel factors through $G_k$ and is
cyclotomic. Quotienting the kernel by its prime-to-$\ell$ factors gives
an extension
\begin{equation}\label{eq:ag-c149-extension}
 1\longrightarrow\mathbb Z_\ell(1)
 \longrightarrow\mathcal E_\ell(\mathcal L)
 \longrightarrow\pi_1(S)\longrightarrow1.
\end{equation}
The comparison isomorphism
\[
 H^2_{\mathrm{cont}}(\pi_1(S),\mathbb Z_\ell(1))
 \simeq H^2_{\mathrm{cont}}(S,\mathbb Z_\ell(1))
\]
sends its extension class $e_\ell(\mathcal L)$ to the arithmetic
Kummer class $c_{1,\ell}(\mathcal L)$. Continuous cohomology here is
the cohomology of the inverse system $(\mu_{\ell^n})$; it is not
silently identified with an inverse limit in arbitrary situations.

The source proves exactness from the homotopy sequence of the torus
torsor and proves the class comparison by compatible transgressions:
the identity of the Tate module maps to the negative of each of the
two classes. The hypothesis on $\pi_2$ is substantive: over an
algebraically closed field, the torsor
$\mathbb A^{d+1}\setminus\{0\}\to\mathbb P^d$, $d\ge1$, has trivial
total-space fundamental group, so its fibre inertia does not inject
(source Remark 3.2, PDF p.~5). Smooth geometrically connected curves
other than the geometric projective line are algebraic $K(\pi,1)$
spaces (source Section 3.1, PDF p.~4).

Geometrically, $\mathcal L$ is the normal bundle of the zero section in
$\operatorname{Tot}(\mathcal L)$, so this is an actual punctured
normal-bundle extension. It instantiates the $K(\pi,1)$ comparison
explicitly contemplated on cote 149, original PDF p.~80, manuscript
p.~79, equations (13)--(17), followed by the normal-bundle discussion
on PDF p.~81, manuscript p.~80. Identifying a separately given
hyperbolic inertia normalizer with this group still requires the
corresponding geometric comparison; the result does not identify
the global tripod kernel in LM6 A.44--A.45.

\begin{proposition}[AG-C149-LIFT: torsion classes vanish on every section]
Under the hypotheses of AG-C149-COMP, a continuous section
$s:G_k\to\pi_1(S)$ of the arithmetic projection lifts to
$\mathcal E_\ell(\mathcal L)$ if and only if
$s^*e_\ell(\mathcal L)=0$. If $\mathcal L$ is torsion in
$\operatorname{Pic}(S)$, every such section lifts, although the
extension \eqref{eq:ag-c149-extension} need not split over
$\pi_1(S)$. When lifts exist, their classes modulo conjugation by
$\mathbb Z_\ell(1)$ form a torsor under
\[
 H^1_{\mathrm{cont}}(k,\mathbb Z_\ell(1))
   \simeq\varprojlim_n k^\times/(k^\times)^{\ell^n}.
\]
\end{proposition}
\begin{proof}
Pulling back \eqref{eq:ag-c149-extension} by $s$ gives an extension of
$G_k$ with class $s^*e_\ell(\mathcal L)$. A lift is precisely a
continuous splitting of this pullback, so the class criterion follows
from extension classification, equivalently source Lemma 4.5
(PDF p.~9) and its preceding comparison.

We verify the continuous-cohomology issue directly for the field.
Hilbert 90 and the Kummer sequence give
\[
 H^1(k,\mu_{\ell^n})=k^\times/(k^\times)^{\ell^n},\qquad
 H^2(k,\mu_{\ell^n})=\operatorname{Br}(k)[\ell^n].
\]
The maps on the first groups are the surjective quotient maps, and
those on the second groups are multiplication by $\ell$. Thus
$\varprojlim^1 H^1(k,\mu_{\ell^n})=0$, and the defining exact
sequence for continuous cohomology gives
\[
 H^2_{\mathrm{cont}}(k,\mathbb Z_\ell(1))
   \simeq T_\ell\operatorname{Br}(k).
\]
For any abelian group $A$, its Tate module
$T_\ell A=\varprojlim(A[\ell^n],\ell)$ is torsion-free over
$\mathbb Z_\ell$: if $\ell t=0$, then compatibility gives
$t_n=\ell t_{n+1}=0$ for every $n$, and induction excludes all
$\ell$-power torsion. Multiplication by an integer prime to $\ell$
is invertible. If $m[\mathcal L]=0$, additivity of the Kummer class
therefore gives
$m s^*e_\ell(\mathcal L)=0$, whence $s^*e_\ell(\mathcal L)=0$.
This is the $\mathbb G_m$, $\ell$-primary specialization of the
source's Proposition 4.2 and Lemma 4.6.

Finally, once one splitting is fixed, the difference of a second
splitting from it is a continuous $1$-cocycle. Conjugation by the
kernel changes that cocycle by a coboundary. This proves the torsor
claim. The finite groups $H^0(k,\mu_{\ell^n})$ form a
Mittag--Leffler inverse system; hence their $\varprojlim^1$ vanishes
and the degree-one continuous-cohomology sequence gives the displayed
Kummer identification.
\end{proof}

\begin{proposition}[AG-C149-ELLIPTIC: non-splitting does not obstruct all sections]
Let $E/\mathbb Q$ be the elliptic curve
\[
 y^2=x^3-x,
 \qquad O=\text{the point at infinity},\qquad P=(0,0),
\]
and let $\mathcal L=\mathcal O_E(P-O)$.
The class of the extension $\mathcal E_2(\mathcal L)\to\pi_1(E)$
is nonzero of order $2$, so this extension does not split. Its image
with $\mathbb Q_2(1)$ coefficients is zero. Nevertheless, every
Galois section $G_{\mathbb Q}\to\pi_1(E)$ lifts through this
extension. Thus vanishing of the obstruction after pullback by
every Galois section does not force the global integral class to
vanish.
\end{proposition}
\begin{proof}
The point $P$ is a nonidentity point of order $2$, since inversion on
this Weierstrass curve sends $(x,y)$ to $(x,-y)$. The standard
identification $E(\mathbb Q)=\operatorname{Pic}^0(E)$ sends $P$ to
$[\mathcal L]$, so $\mathcal L$ has order $2$.

Its discriminant is $64$, so the displayed integral model has good
reduction at $3$. Direct enumeration gives
\[
 E(\mathbb F_3)=\{O,(0,0),(1,0),(2,0)\}
                 \simeq(\mathbb Z/2\mathbb Z)^2.
\]
Properness and good reduction give a group homomorphism
$E(\mathbb Q)\to E(\mathbb F_3)$. If $P=2Q$ for a rational point
$Q$, reduction would yield $\bar P=2\bar Q=0$, whereas
$\bar P=(0,0)\ne O$. Therefore $P\notin2E(\mathbb Q)$.
If $\mathcal L=\mathcal M^{\otimes2}$ in $\operatorname{Pic}(E)$,
degrees give $\deg\mathcal M=0$, and the same Picard identification
would express $P$ as twice a rational point. Thus
$[\mathcal L]\ne0$ in $\operatorname{Pic}(E)/2$.

The Kummer injection
\[
 \operatorname{Pic}(E)/2\hookrightarrow H^2(E,\mu_2)
\]
shows that the reduction of $c_{1,2}(\mathcal L)$ modulo $2$ is
nonzero. Consequently the continuous integral class is nonzero;
its order is exactly $2$ by additivity and $\mathcal L^{\otimes2}
\simeq\mathcal O_E$. The curve has genus $1$ and is an algebraic
$K(\pi,1)$, so AG-C149-COMP identifies this with the nonzero
extension class. A torsion class maps to zero in cohomology with
$\mathbb Q_2(1)$ coefficients. Finally AG-C149-LIFT applies to
every Galois section, and such sections do exist because
$O\in E(\mathbb Q)$. This proves all the claims.
\end{proof}

This supplies the group/cohomology comparison missing from the purely
arithmetic example AG-KUMMER-2, under explicit hypotheses and for a
literal normal bundle. It also separates three questions on the cote:
nonvanishing of a global integral extension class, its rational image,
and the obstruction for a specified section. None of these distinctions
uses A or M, and none identifies an irregular-prime class-group
component or proves the remaining LM6 tripod-kernel assertion.

\subsection{LM6 A.44--A.45: the exact integral generator gate}
\label{ag:lm6}

\paragraph{AG-LM6-REG (the known kernel-field theorem).}
For a prime $p$, let $\Omega_p^*$ be the fixed field of the kernel
of the outer Galois action on
$\pi_1^{\mathrm{pro}\text{-}p}
(\mathbb P^1_{\overline{\mathbb Q}}\setminus\{0,1,\infty\})$.
Let $\Omega_p$ be the maximal pro-$p$ extension of
$\mathbb Q(\mu_p)$ unramified away from $p$. Then
$\Omega_p^*\subseteq\Omega_p$, and equality is classical for odd
regular $p$. An exact primary statement is Hoshi,
\emph{On the Kernels of the Pro-$l$ Outer Galois Representations
Associated to Hyperbolic Curves over Number Fields}, revised March
2014 manuscript, Definition 1.8, Remark 1.8.1, and Theorem 1.9,
PDF p.~11. Its proof cites Sharifi's Theorem 1.1 and Brown's theorem.
This formulation corrects the informal phrase ``the kernel is the
Galois group of $\Omega_p$'': the kernel is
$\operatorname{Gal}(\overline{\mathbb Q}/\Omega_p^*)$.

Brown, \emph{Mixed Tate motives over $\mathbb Z$}, Ann. Math.~175
(2012), 949--976, Theorem 1.1 and the following paragraph, printed
p.~950, prove the motivic-generation assertion formulated as
Conjecture 1 on p.~949. Its Deligne--Ihara consequence concerns the
graded Lie algebra after tensoring with $\mathbb Q_p$. It does not
discard integral indices. Brown's theorem is a classical input here;
its entire proof has not been re-refereed in this run.

\paragraph{AG-LM6-IRREG (where the regular-prime argument uses more).}
Write $K=\mathbb Q(\mu_{p^\infty})$ and
$G=\operatorname{Gal}(\Omega_p^*/K)$, with its outer-action
filtration $F^mG$ and graded Lie algebra $\mathfrak g$.
Choose the usual maximal-image Soulé elements $\sigma_m\in F^mG$
for odd $m\ge3$. They define maps $\Psi:S\to G$ from the free
pro-$p$ group on these generators and
$\psi:\mathfrak s\to\mathfrak g$ from the free graded
$\mathbb Z_p$-Lie algebra. Sharifi,
\emph{Relationships between conjectures on the structure of pro-$p$
Galois groups unramified outside $p$}, Proc. Sympos. Pure Math.~70
(2002), 275--284, author PDF Theorem 1.1 p.~2 and proof pp.~8--9,
uses the freeness of
$\operatorname{Gal}(\Omega_p/\mathbb Q(\mu_p))$ at regular primes.
The full proof and its preceding group lemmas have been read.

For irregular $p$, his Theorem 1.3 (author PDF p.~3; proof p.~10),
together with Brown's rational freeness, says that $\Psi$ and
$\psi$ are not surjective if Greenberg's generalized conjecture
holds for $\mathbb Q(\mu_p)$. Here that conjecture says that the
unramified abelian pro-$p$ Iwasawa module over the compositum of
all $\mathbb Z_p$-extensions is pseudo-null. This is a conditional
obstruction to the proposed generator argument, not a disproof of
$\Omega_p^*=\Omega_p$: the latter could hold with additional
integral generators and relations.

\begin{proposition}[AG-LM6-NORM: the filtered Soulé normalization]
For odd $p$ and odd $m\ge3$, in the integral normalization of
Sharifi's equation (2.10), author PDF p.~7,
\[
 \kappa_m(F^mG)\subseteq
 p^{\sum_{j\ge1}\lfloor(m-1)/p^j\rfloor}\mathbb Z_p(m).
\]
In particular $\kappa_5(F^5G)\subseteq3\mathbb Z_3(5)$,
although $3$ is regular. Thus surjectivity of these filtered
characters onto the full displayed lattice is false in every
degree with $m-1\ge p$.
\end{proposition}
\begin{proof}
The cited identity is
\[
 \kappa_m=(p^{m-1}-1)(m-1)!\lambda_m^{(p)}
       \quad\hbox{on }F^mG,
\]
where $\lambda_m^{(p)}$ is integral, obtained from the integral
stable-derivation algebra. The first factor is a $p$-adic unit.
Counting the multiples of $p^j$ among $1,\ldots,m-1$ gives
$v_p((m-1)!)=\sum_{j\ge1}\lfloor(m-1)/p^j\rfloor$.
The inclusion follows, and for $(p,m)=(3,5)$ this valuation is $1$.
\end{proof}

This is a precise test for the suggestion in
\path{OURS_FABLE_PROGRAM_SUGGESTIONS_20261120_EXPLAINED.md},
Section 5, linked there to LM6 A.44--A.45, bridge B3, and cote 149
scans 79--85. It does not assert nonsurjectivity on a different,
larger Galois domain, and it does not replace a normalized character
by an unnormalized one. A proposed extension of the regular-prime
proof must identify the exact integral generator quotient, its
relations, and the cohomological comparison producing them.
Neither endpoint A nor endpoint M supplies these data. The
$p=2$ and irregular-prime kernel-field questions in the intake
therefore remain unresolved by this argument; they are not listed
as refuted.

\subsection{The formal-height mechanism in Alpöge step 2}
\label{ag:formal-height}

The supplied Fable suggestion, Section 2, gives the height-dependent
derivative estimate but passes through a Lubin--Tate normal form.
The following proof works directly with the given formal group.
It strengthens the estimate to an explicit bound and requires no
identification of a general deformation with a particular
Lubin--Tate group.

\begin{proposition}[AG-FG-EXP: an exact exponential coefficient bound]
Let $R$ be a complete unramified discrete valuation ring of mixed
characteristic $(0,p)$, with fraction field $K$, perfect residue
field, and valuation $v(p)=1$, so $v(K^\times)=\mathbb Z$;
set $v(0)=+\infty$.
Let $F(X,Y)$ be a one-dimensional commutative formal group over
$R$. Suppose its reduction has finite height $h\ge1$, and put
$q=p^h$. Write
\[
 [p]_F(t)=pt+\sum_{j\ge2}a_jt^j.
\]
Thus $a_j\in pR$ for $2\le j<q$, $a_q\in R^\times$, and
all $a_j\in R$. Normalize the formal exponential by
\[
 e(z)=\exp_F(z)=z+\sum_{n\ge2}e_nz^n.
\]
Then
\[
 v(e_n)\ge-\frac{n-1}{q-1},\qquad
 v(e_n)\ge-\left\lfloor\frac{n-1}{q-1}\right\rfloor.
\]
The second inequality records that the valuation group is integral.
\end{proposition}
\begin{proof}
We include the logarithm construction to specify all normalizations.
Put $b(t)=\partial_2F(t,0)\in1+tR[[t]]$ and
\[
 \omega_F=\frac{dt}{b(t)},\qquad D=b(t)\frac{d}{dt}.
\]
The differential is integral, and $D$ preserves $R[[t]]$ with
$(Dt)(0)=1$. Associativity, differentiated in the third variable
at zero, gives
\[
 b(F(x,y))=\partial_2F(x,y)b(y).
\]
Define $\lambda(t)=\int_0^t b(u)^{-1}\,du\in K[[t]]$.
The displayed identity implies
$\partial_y\lambda(F(x,y))=\lambda'(y)$.
Evaluating the resulting constant at $y=0$ yields
$\lambda(F(x,y))=\lambda(x)+\lambda(y)$.
Its linear coefficient is $1$, so it has a formal inverse $e$,
with $e(0)=0$ and $e'(0)=1$. In particular
\[
 e(pz)=[p]_F(e(z)).
\]

We prove the first inequality by induction; it is equality for
$n=1$. Comparing the coefficient of $z^n$ in this identity gives
\[
 (p^n-p)e_n=
 \sum_{j=2}^n a_j
  \sum_{\substack{n_1+\cdots+n_j=n\\n_i\ge1}}
       e_{n_1}\cdots e_{n_j}.
\]
Every $n_i<n$, and the induction hypothesis bounds the valuation
of each product below by $-(n-j)/(q-1)$.
For $2\le j<q$, multiplication by $a_j$ adds at least $1$;
division by $p^n-p$, whose valuation is $1$, therefore leaves
a bound $-(n-j)/(q-1)\ge-(n-1)/(q-1)$.
For $j\ge q$, multiplication adds at least $0$, and division
leaves
\[
 -1-\frac{n-j}{q-1}\ge-\frac{n-1}{q-1}.
\]
The ultrametric inequality proves the induction. Since each
$e_n\in K$ has integral valuation, rounding up gives the second
bound. No convergence assertion or special normal form was used.
\end{proof}

\begin{proposition}[AG-FG-DER: integral derivatives with the exact error]
With $R,F,D,h,q$ as above, for every $f\in R[[t]]$ and integer
$a\ge1$,
\begin{equation}\label{eq:ag-height-derivative}
 v\bigl((D^af)(0)\bigr)\ge
 \max\left\{0,\ v(a!)-
                 \left\lfloor\frac{a-1}{q-1}\right\rfloor\right\}.
\end{equation}
In particular, writing $s_p(a)$ for the base-$p$ digit sum,
\[
 v\bigl((D^af)(0)\bigr)\ge
 a\left(\frac1{p-1}-\frac1{p^h-1}\right)
 -\frac{s_p(a)}{p-1}+\frac1{p^h-1}.
\]
For height $2$ the coefficient of $a$ is
\[
 \frac{p}{p^2-1}=\frac1p+\frac1{p^3}+\frac1{p^5}+\cdots.
\]
\end{proposition}
\begin{proof}
Differentiating $\lambda(e(z))=z$ gives $e'(z)=b(e(z))$.
Consequently, for each $a$,
\[
 (D^af)(0)=a!\,[z^a]f(e(z)).
\]
In a term of $[z^a]e(z)^j$, with $1\le j\le a$, the
coefficient is a product of $j$ coefficients whose indices sum
to $a$. AG-FG-EXP bounds its valuation by $-(a-j)/(q-1)$,
which is at least $-(a-1)/(q-1)$. The integral coefficients of
$f$ and the ultrametric inequality give the same bound for
$[z^a]f(e(z))$. Rounding to integral valuations and multiplying
by $a!$ gives the second entry in the maximum. The first follows
because $D$ preserves $R[[t]]$.

Finally, counting multiples of $p^j$ among $1,\ldots,a$ gives
$v(a!)=\sum_{j\ge1}\lfloor a/p^j\rfloor
=(a-s_p(a))/(p-1)$. Subtracting $(a-1)/(q-1)$ gives the
displayed digit-sum bound. The height-$2$ identity is the
geometric-series sum. Since
$s_p(a)\le(p-1)(1+\lfloor\log_pa\rfloor)$, it also gives the
claimed linear gain with an explicit logarithmic error.
\end{proof}

\begin{proposition}[AG-FG-PPOWER: the exact restricted derivation]
With the same coordinate and normalization, reduction modulo $p$
satisfies
\[
 \overline D^{p}=\overline{a_p}\,\overline D
             \quad\hbox{on }(R/p)[[t]].
\]
Thus the $p$-th power is zero for $h\ge2$ and nonzero for $h=1$.
\end{proposition}
\begin{proof}
In characteristic $p$, the iterated Leibniz formula shows that
the $p$-th power of a derivation is a derivation. It is invariant
under the formal-group translations because $\overline D$ is.
Every invariant derivation is a scalar multiple of
$\overline D$: an invariant vector field $c(t)\partial_t$
satisfies $c(F(t,u))=\partial_1F(t,u)c(t)$, and setting $t=0$
gives $c(u)=b(u)c(0)$. It remains to compute the scalar
$(D^pt)(0)$ modulo $p$.

The coefficients $e_i$ are integral for $1\le i<p$ by
AG-FG-EXP, since $q\ge p$. In the coefficient recursion for
$n=p$, every summand with $2\le j<p$ lies in $pR$, and the
$j=p$ term is $a_p$. Hence
\[
 pe_p\equiv-a_p\pmod p.
\]
As $(D^pt)(0)=p!e_p$ and $(p-1)!\equiv-1\pmod p$,
the scalar is $a_p$ modulo $p$. The factorial congruence follows
by pairing inverse elements in $\mathbb F_p^\times$ and holds
also at $p=2$. The height definition gives the last assertions.
\end{proof}

\begin{proposition}[AG-FG-PRODUCT: independent formal factors]
Suppose the integral formal group is a product of one-dimensional
groups as above, of heights $h_1,\ldots,h_r$, with commuting
normalized invariant derivations $D_1,\ldots,D_r$. For
$f\in R[[t_1,\ldots,t_r]]$ and integers $a_i\ge0$, put
\[
 B_i(0)=0,\qquad B_i(a)=
 \max\left\{0,v(a!)-
       \left\lfloor\frac{a-1}{p^{h_i}-1}\right\rfloor\right\}
          \quad(a\ge1).
\]
Then
\[
 v\bigl((D_1^{a_1}\cdots D_r^{a_r}f)(0)\bigr)
       \ge\sum_{i=1}^r B_i(a_i).
\]
\end{proposition}
\begin{proof}
Regard $f$ as a power series in all variables other than $t_1$,
with coefficients in $R[[t_1]]$. Apply AG-FG-DER to each
coefficient and evaluate at $t_1=0$. The result lies in
$p^{B_1(a_1)}R[[t_2,\ldots,t_r]]$. Divide out this scalar and
repeat for each remaining coordinate. Evaluation and the other
coordinate derivations commute, proving the bound. If $a_i=0$,
the corresponding operation is simply evaluation and preserves
integrality. An isogeny decomposition over $K$ alone is not an
integral product decomposition and does not meet this hypothesis.
\end{proof}

\paragraph{AG-FG-THETA (what transfers to a theta calculation).}
Alpöge's note, step 2, PDF p.~3, asserts the height-$2$ exponent
for the normalized derivatives of a CM theta value. AG-FG-DER
supplies that exponent whenever the theta section, in a specified
local trivialization and integral parameter, is an element of
$R[[t]]$, and the operator is the integral invariant derivation
above. The period and trivialization comparison must be supplied
for that theta object. More generally, for constants
$u,b\in K^\times$ and $\widetilde D=uD$,
\[
 v\bigl((\widetilde D^{a}(bf))(0)\bigr)
  =v(b)+a\,v(u)+v\bigl((D^af)(0)\bigr),
\]
so their effect on the bound is explicit. A nonconstant change of
trivialization or a connection requires its own comparison.
The supplied Lean map, step-2 row, reports only a low-confidence
match to CM-curve/invariant-derivation declarations; no new exact
Lean lemma identifier is asserted here.

This proves the formal-group mechanism in Fable Section 2 and
shrinks its general-height gate to the actual theta integrality
and normalization data. Its archive map is inherited explicitly
from that record: cote 54 scans 5--7, the restricted $p$-Lie
algebra and Cartier-duality discussion; cotes 4--9 for crystals;
cotes 136--137, cote 16 scan 14, and cote 47 scan 9 for the
divided-power line. These additional sheets have not been newly
image-checked here. The theorem does not by itself place a theta
section on the crystalline site or replace a theta-lift chain.
It uses neither A nor M.

\subsection{Archive and run consequences, with the original scope retained}
\label{ag:archive}

\paragraph{AG-C1.} Cote 1/79--80 and 111--112 gives the trace/Fredholm
framework, and SA23 geo-H-lidskii/geo-H-genus is its reported modern
reconstruction. A yields the earlier zeta positivity realization;
M yields AG-DET for the stated Hecke functions. AG-MOMENT and AG-HANKEL
give further consequences in that framework. A geometric operator defined
from primes remains a separate construction.

\paragraph{AG-C16.} Cote 16, scans 14--18, no.~7, concerns the integral
Lefschetz correspondence ring. Scans 44--50 contain the Coates letter's
conditions A, C and index positivity; the attached Hodge sketch uses a scan
numbering one higher than the GCT\_NSRH intake. These are algebraicity and
integrality problems. Neither A nor M constructs the required
correspondence or bounds its denominator. The sketch's explicit route is
the rank comparison on $X\times X$, not a claim that zero location
reverses Grothendieck's implication to the Weil conjectures.

\paragraph{AG-C34.} Original PDF p.~74, manuscript p.~73, (2.5.5),
has now been visually checked: its multiplicative $1$ denotes the trivial
invariant group, equivalently additive $0$. The finite-field curve theorem
is the stated input, so inherited Q04 is unchanged. The original
definitions and full proof route for Theorem 1.1 have also been read.
AG-C34-TYPE identifies the literal obstruction in Q03, and
AG-C34-REPAIR proves the coherent intrinsic correction. AG-C34-INDEX
gives its finite Frobenius/Jacobian index. AG-C34-SECTION also proves
the formulation using a specified section and solvable $H$;
AG-C34-NONSOLV gives its abstract group-theoretic boundary.
No assertion of new purity or
of a consequence of A/M is made.

\paragraph{AG-C44.} Cote 44 Theorems 4--5 and the Serre letter, at the
locators in AG-SIGN, concern polarization and heights. AG-SIGN corrects
the central-derivative bridge but leaves cycle existence and the
height formula explicit. The supplied Hodge sketch also cites a separate
2026 CM-Hodge theorem for its abelian positivity consequences; those
claims are not consequences of either A or M and are carried with
their original dependency rather than independently certified here.

\paragraph{AG-C149 and AG-GT.} ORDER\_tonight, item 31, points to cote
149 scans 79--85, LM6 A.44--A.45, integral bridge B3, and the
irregular-prime/$\ell=2$ questions. GRH does not remove the integral
class-group and generator-relation gates in AG-LM6-IRREG/NORM.
For an actual punctured normal bundle, AG-C149-COMP now gives the
classical comparison under the explicit $\pi_2=0$ hypothesis, and
AG-C149-ELLIPTIC separates its global class from every section's
obstruction. A separately given inertia normalizer still requires
its identification with that geometric extension.
The semi-graded Euler information in AG-EXT explains why the endpoint
zero theorems do not evaluate that extension data. In SA22 run 1,
A6-node-character is the quadratic Artin character times the cyclotomic
character; A gives effective detection of that finite quadratic sign.
M additionally covers finite nonabelian splitting fields in effective
Frobenius tests such as AG-FROB. These improvements do not reverse the
recorded refutation of run 2 A:inertia-maximality: finite-index inertia
subgroups have the same weights. Nor do they identify the profinite GT
kernel from the cyclotomic quotient. The existing supplied claim about
rational Deligne--Drinfeld freeness (CC-GT03, external 2026 input) is
carried separately from the full profinite Ihara question R-GTHAT.

\paragraph{AG-PURITY and AG-WM.} Finite-field purity and the usual
Grothendieck--Lefschetz trace formula are already classical. The full
geometric weight--monodromy statement remains represented in these
materials by BR-WM-02's finite exterior-power trace criterion; its
remaining inequalities are geometric hypotheses, not consequences of
global Dirichlet or Hecke zero locations. Likewise the supplied archive
questions about standard conjectures require algebraic cycles or
correspondences. No formal logical independence from A or M is asserted;
the precise point is that the displayed deductions do not supply those
inputs.

% END geometry_programme

% BEGIN complexity
\section{Complexity, effective algebra, and finite geometric tests}
\label{sec:complexity}

Here $A_D$ denotes the rational Dirichlet-GRH input designated by the task as
Alp\"oge's theorem. The letter $M$ denotes Masullo's all-number-field Hecke-GRH
claim, and every use of $M$ remains conditional. The input size of an arithmetic
problem includes the circuits, explicitly listed variables, rational constants
in binary, and any promised field representation. A geometric object without
an encoding has no associated complexity bound.

\subsection{Classical consequences and inherited results}

The following catalogue entries preserve the original finding identifiers of
the earlier runs. They are not advertised as new findings of this session.
An exact locator to an earlier proof is distinguished from a fresh reading of
that proof. In particular, the earlier complete file \path{CX_GRH.tex} is
not in tonight's attachment: its \path{IMPLICATIONS}, \path{LEDGER},
\path{QUESTIONS}, and targeted intake are available. Statements that depend
on its elaborate new finite-field constructions are carried from those records;
this session does not manufacture an independent verification of the absent
proof text.

\paragraph{CX-ALG-001 and CX-ALG-002: primality.}
Bach, \emph{Explicit bounds for primality testing and related problems}
(1990), Theorem 2, p.~372, and its application on p.~373 give the cutoff
$2(\log N)^2$ for a compositeness witness. Schoof, \emph{Four primality testing
algorithms} (2008), Theorem 2.2, pp.~105--106, states and proves the Miller test
with this cutoff using only quadratic Dirichlet characters. Therefore $A_D$
discharges its analytic input. Writing $N-1=2^s m$ with $m$ odd, test the usual
conditions $a^m=1$ or $a^{2^i m}=-1$ modulo $N$ for every base through the
cutoff, also checking nontrivial gcds. Even inputs are handled separately.
There are $O((\log N)^2)$ bases and each modular exponentiation takes a
polynomial number of bit operations. With $b=\lceil\log_2N\rceil$ and
multiplication cost $\mathcal M(b)$, the elementary implementation uses
$O(b^3\mathcal M(b))$ bit operations. This is the transfer already recorded
as CX\_GRH CX-N01 and ALPOGE\_GCT C-GRH-02.

The statement $\mathrm{PRIMES}\in\mathrm P$ was already unconditional:
Agrawal--Kayal--Saxena, \emph{PRIMES is in P}, Annals 160 (2004), Theorems 4.1
and 5.1, with their 2019 erratum. The new input validates this particular
short-witness test; it does not first place primality in $\mathrm P$.
Geometrically these are computations in $\mathbf G_m$ over $\mathbf Z/N\mathbf Z$
and, for AKS, finite cyclotomic algebras. The archival companion is cote 28,
scan 3 (group-action axioms), as mapped in the CX\_GRH return, not freshly
inspected from an archive scan in this session.

\paragraph{CX-ALG-003--CX-ALG-004: generators and nonresidues.}
Bach's Theorem 3, p.~372, gives units below $3(\log N)^2$ generating
$(\mathbf Z/N\mathbf Z)^\times$. Applied to the proper subgroup of $r$th
powers in $\mathbf F_p^\times$, it gives a polynomial search for a specified
prime $r\mid p-1$. Membership is tested by exponentiation to $(p-1)/r$;
the other factors of $p-1$ need not be known. This is CX\_GRH CX-N02.
The original Bach PDF was not freshly recovered here (the publisher returned
403); its exact locator is carried from the preceding checked run. Schoof's
primary source independently supplies the quadratic special case used above.
In the Kummer sequence this is finding a nonzero class in
$\mathbf F_p^\times/(\mathbf F_p^\times)^r$, or a nontrivial
$\boldsymbol\mu_r$-torsor. The established companion locator is SGA 4 IX 3.2.

\paragraph{CX-ALG-006--CX-ALG-008: primitive roots and extension fields.}
CX\_GRH CX-N03 constructs a certified generator of $\mathbf F_p^\times$ when
the full factorization of $p-1$ is supplied, by constructing each primary
component and multiplying them. Shoup, \emph{Searching for Primitive Roots
in Finite Fields} (1992), author version Theorem 3, p.~4, and Section 6,
p.~13, gives a least primitive root of order
$O(r^4(\log r+1)^4(\log p)^2)$, where $r=\omega(p-1)$,
and hence $O((\log p)^6)$ under the
rational Dirichlet hypothesis. This is a guaranteed short candidate list
without a supplied factorization, as CX\_GRH CX-N04 explicitly records.
A candidate list does not identify or certify a particular primitive root.

The stronger extension-field construction CX\_GRH CX-N13 supplies a specified
$r$th nonresidue in a supplied $\mathbf F_{p^m}$ in
$\operatorname{poly}(m\log p)$ time when $r\mid p-1$ is prime. Its quadratic
case with even $m$ already has an unconditional construction. The theorem is
inherited from the full targeted check in
\path{OURS_INTAKE_CX_GRH_20261121.md}, Section 1. A fresh source audit
found a nonvanishing gap in one cited auxiliary formula; CX-ALG-016 gives
the exact counterexample, and CX-ALG-017--CX-ALG-018 give a fully checked
repair and re-proof of the needed construction. The defining
irreducible polynomial is part of the input. These are computations with
Kummer covers inside an explicitly represented finite \`etale algebra.

\paragraph{CX-ALG-009--CX-ALG-010: polynomial factorization over finite fields.}
The inherited theorem CX\_GRH CX-N15 has bound
\[
                 (d^{\log d}\log q)^{O(1)}
\]
for a degree-$d$ polynomial over a supplied $\mathbf F_q$. Under $M$ its
analytic input is satisfied by $\zeta_{\mathbf Q}$ and the particular
$\zeta_{E_\nu}$ used in the construction, where
$E_\nu=\mathbf Q(\zeta_{2\nu},p^{1/\nu})$ and $\nu\leq d^2$.
Its exact component locators are Adleman--Lenstra (1986), Theorem 1, p.~350,
Algorithm A, pp.~350--351; Bach--von zur Gathen--Lenstra (2001), Theorem 2,
p.~7; and the 2008 Ivanyos--Karpinski--R\'onyai--Saxena version, Main Theorem,
p.~3, Lemma 4.2, p.~14, Proposition 4.6, pp.~16--17.
The fresh primary published IKRS article confirms the quasipolynomial
factorization-versus-automorphism distinction (Theorems 1.1 and 1.3), and
its unconditional noncommutative zero-divisor result (Theorem 1.7).

The displayed bound remains quasipolynomial in $d$. The residual $A_D$-only
problem in CX\_GRH CX-U02 is the high-valuation branch
$p>d$, $r\nmid p(p-1)$ and
$v_r(p^{\operatorname{ord}_r(p)}-1)>1$: construct the maximal coefficient
Teichm\"uller unit needed by that descent, or replace the descent. The new
CX-ALG-023--CX-ALG-026 below replace the degree-$sr$ field construction by
a prime-degree field, cover the additional case $r\mid m$ for the supplied
base field $\mathbf F_{p^m}$, and reduce the remaining analytic family to
pure prime-degree number fields. They do not solve the unrestricted
$A_D$-only gate. Geometrically factorization computes
primitive idempotents, hence connected components, of a finite \`etale
algebra (SGA 1 V Section 7). It is a different computational problem from
factoring an integer with unknown prime divisors.

\paragraph{CX-ALG-011--CX-ALG-013: arithmetic Picard groups.}
Bach (1990), Theorem 4, p.~375, and its corollary, p.~376, give the familiar
$12(\log|D_K|)^2$ class-group generator bound, using class-character Hecke
$L$-functions over $K$. Thus $M$ supplies the premise. This is CX\_GRH
CX-N07 and ALPOGE\_BRIDGES CC-F03/F03b. If the Hilbert class field is
abelian over $\mathbf Q$, $A_D$ suffices by Dirichlet factorization; the
weaker condition that $K/\mathbf Q$ be abelian does not establish this.
A factor base is not the full lattice of relations between its ideal
classes. Relation collection, saturation, regulator information and their
running times remain separate inputs in CX\_GRH CX-N08 and CX-U03.
In Grothendieck's language the class group is
$\operatorname{Pic}(\operatorname{Spec}\mathcal O_K)$ and principal divisors
supply the relation lattice. The established companion locators are SGA 4
IX Theorem 3.3, SGA 6 V 3.10 and VI 3.2--3.3, and cote 2 scan 5 as recorded
by CX\_GRH.

\paragraph{CX-ALG-019: a complete class-group algorithm retains randomness.}
For a supplied negative fundamental discriminant $-d$, Biasse--Erukulangara,
\emph{A proof of the conjectured run time of the Hafner--McCurley class
group algorithm}, Main Theorem, author PDF p.~4, Algorithm 4, p.~15,
and Corollary 1, p.~16, give a Las Vegas computation of
$\operatorname{Cl}(\mathcal O_{\mathbf Q(\sqrt{-d})})$ with expected bound
\[
 \exp\!\left(\left(\frac3{\sqrt8}+o(1)\right)
                         \sqrt{\log d\log\log d}\right).
\]
Their Conjecture 1, p.~4, explicitly states number-field Hecke-GRH; hence
$M$ supplies the hypothesis. The proof also uses the class-group Cayley
graph expansion of Jao--Miller--Venkatesan (2009), Theorem 1.1 and
Corollary 1.3. Rational quadratic Dirichlet-GRH alone does not supply all
these class-character estimates. This is the complete algorithmic transfer
already recorded as CX\_GRH CX-N12, not a new algorithm of this session.

The source's bounded trial returns a certified answer or FAILURE, with
success probability $1-d^{-1+o(1)}$ and the displayed time bound. Repeating
on detected failure gives the stated expected Las Vegas bound. The
certification in Algorithm 4 compares the order $h'$ of the group presented
by the collected relations with $2h^*$, where $h^*\le h<2h^*$ and $h$
is the actual class number. Since $h\mid h'$ and $h'>0$, the condition
$h'<2h^*$ forces $h'=h$. Thus the restart does not accept an incorrect
relation lattice. The full earlier proof and its detailed source check are
credited to CX-N12; the primary theorem, analytic-input definition,
Algorithm 4 and Corollary 1 were freshly inspected here.
The runtime is subexponential in $\log d$, not polynomial in that input
length, and the algorithm remains randomized. Its geometric output is the
finite Picard group, with actual principal-divisor relations rather than
only a factor base. The prior companions are SGA 4 IX Theorem 3.3,
SGA 6 V 3.10 and VI 3.2--3.3, and cote 2 scan 5.

\paragraph{CX-ALG-020: small-prime walks in ray class groups.}
Jao--Miller--Venkatesan (2009), Theorem 1.1, author PDF p.~3,
Remark 1.2(a) and Corollary 1.3, p.~4, give the following exact transfer.
For degree-$n$ $K$, modulus $\mathfrak m$, and
$Q=|D_K|N\mathfrak m$, take the narrow ray class group $G$ and the
multiset of degree-one prime ideals of norm at most $x$, coprime to
$\mathfrak m$, together with their inverses. Conditional on $M$, its
Cayley adjacency eigenvalues satisfy
\[
 \lambda_1=2\operatorname{li}(x)+O(n\sqrt x\log(xQ)),\qquad
 |\lambda_\chi|=O(n\sqrt x\log(xQ))\quad(\chi\ne1).
\]
The constants here are absolute. For fixed $n$ and $B>2$,
$x\ge(\log Q)^B$, and sufficiently large $Q$, a walk of length
$t\ge C_{n,B}\log|G|/\log\log Q$ reaches any fixed subset $T$ with
probability at least $|T|/(2|G|)$. Quotients inherit the bound.
The precise hypothesis is GRH for the Hecke characters actually occurring;
for $K=\mathbf Q$, $A_D$ suffices. Random walks remain randomized.
Geometrically these are bounded-prime correspondences on arithmetic
Picard/ray-class quotients, complementing CX\_GRH CX-N12; the Picard
companions are SGA 4 IX Theorem 3.3 and cote 2 scan 5.

\paragraph{CX-ALG-022: the ordinary elliptic isogeny graph.}
The same paper's Theorem 1.5, p.~5, applies conditionally on $M$ to ordinary
elliptic curves over a fixed $\mathbf F_q$, with fixed point count and
the same geometric endomorphism order. The graph of isogenies of prime
degree less than $(\log(4q))^B$, $B>2$ fixed, has the stated expansion
bound. Theorem 1.6, pp.~6--7, transfers a discrete-logarithm solver working
on a fraction $\mu$ of that level to every curve on it with an expected
$\mu^{-1}\operatorname{poly}(\log q)$ overhead. These are conditional
classical transfers; they do not supply such a solver. The Hecke input is
for imaginary-quadratic ring-class characters, rather than an assertion
about a non-CM elliptic curve's global L-function. The geometric link is
the CM ideal-class action (Proposition 3.1, p.~11) and Deuring lifting
(Theorem 3.3, p.~12); ORDER item 20 supplies the abelian-scheme/isogeny
programme, alongside the inherited Picard comparison of cote 2 scan 5.

\paragraph{CX-HN-001--CX-HN-004: the exact Nullstellensatz input.}
Koiran's DIMACS 96-27, Section 2, Theorem 2, pp.~5--6, proves
$\mathrm{HN}\in\mathrm{AM}\subseteq\Pi_2^{\mathrm P}$ under its GRH input.
The August 8, 1997 erratum corrects the lower estimate in Theorem 8 and a
prime-product expression in Theorem 10; it explicitly preserves the main
statements. The primary author-hosted report and the entire erratum were
retrieved in this session. Section 5, Theorem 9 and Corollary 1, printed p.~12 (PDF page~13),
show that the analytic estimate is for the bounded primitive solution
field itself. A Galois closure is not required. These are the earlier
CX\_GRH CX-K01--K02; the circuit and homogeneous nonzero forms are also
stated in GCT V Theorem 2.10, p.~25.

Consequently $M$ gives the unrestricted transfer, whereas $A_D$ gives the
bounded abelian-witness promise of CX\_GRH CX-K03. All the arithmetic size
bounds in that promise matter: field degree, logarithmic discriminant and
denominators have the stated single-exponential bounds in the input length.
The universal replacement by an abelian solution field was already refuted
by $x^3-2$ in CX\_GRH CX-K04 and ALPOGE\_GCT C-GRH-03. This refutes that
specific replacement, not every possible future $A_D$-only proof.
The geometric object is a finite \`etale spread of the solution field over
an open of $\operatorname{Spec}\mathbf Z$; degree-one Frobenius orbits give
rational reductions. The companion locators are SGA 1 V Section 7 and
cote 98 scans 11--12 as carried in CX\_GRH.

\paragraph{CX-NN-001--CX-NN-004: normalization maps.}
The preceding CX\_GRH CX-K05 refines GCT V Theorem 3.4, pp.~27--28, to
$\mathrm{FP}^{\Sigma_2^{\mathrm P}}$ under $M$ for equation-given,
standard-graded algebras with explicit ambient variables and polynomial
output bits. The source proof's prefix-existence queries already have
$\Sigma_2$ form. Its unconditional polynomial-space statement remains
available. The resulting map
$\operatorname{Spec}R\to\mathbf A^{\dim R}$ is finite; the companion is
EGA II 6.1.4--6.1.5, and cote 98 scan 44 in the prior crosswalk.

The historical EXPH transfers for succinct small separating systems are
GCT V Theorems 4.10, p.~34, and 5.12, carried in CX\_GRH CX-K06. The
optimal-cardinality h.s.o.p. bounds of Theorems 4.1, p.~30, and 5.5 retain
their efficient defining-equation premise. Their output conventions cannot
be silently changed. Earlier unconditional improvements already exist:
Forbes--Shpilka, ECCC TR17-163, Theorem 7.1, p.~21, constructs the relevant
hitting set in polynomial space for unary $n,s,r$; its introductory
Theorems 1.2--1.3 record normalization consequences. Guo--Saxena--Sinhababu
(2019), Theorems 1.5--1.6, are recorded in CX\_GRH CX-K07. These are not
new consequences of tonight's analytic inputs.

\subsection{New formulations with complete proofs}

\paragraph{CX-ALG-005: deterministic prime-field square roots.}
\emph{Statement.} Under $A_D$, given an odd prime $p$ and $a\in\mathbf F_p$,
one can find all roots of $X^2-a$ or certify their absence in deterministic
time polynomial in $\log p$.

\emph{Proof.} The case $a=0$ is immediate. Euler's criterion decides whether
$a$ is a square. If it is not, return no roots. Otherwise find a quadratic
nonresidue by the deterministic search of CX-ALG-004. Apply the classical
Tonelli--Shanks algorithm with this nonresidue. More specifically, Algorithm
3.34 of Menezes--van Oorschot--Vanstone, \emph{Handbook of Applied Cryptography},
p.~100, has only one random step, its choice of the nonresidue in Step 2.
All its remaining steps are modular exponentiation, gcd inversion and a
loop of length at most $v_2(p-1)\leq\log_2p$; Fact 3.35 gives a polynomial
bit bound for those operations. Replacing Step 2 by the bounded deterministic
search preserves polynomial time and the algorithm's proved correctness.
The resulting roots are $r,-r$. Completing the square factors every quadratic
with nonzero leading coefficient; in characteristic two the separate field
case is trivial for $\mathbf F_2$. This does not provide a square-root oracle
modulo an integer whose factors are unknown.\hfill$\square$

\paragraph{CX-ALG-014: a compositeness witness need not give the usual gcd factor.}
\emph{Statement.} The rule ``take a Miller witness and compute
$\gcd(a,N)$ or $\gcd(a^m\pm1,N)$'' does not always output a nontrivial divisor.

\emph{Proof.} Take $N=35$, so $N-1=2\cdot17$, and $a=2$. Since
$2^{12}=1\pmod{35}$, one has $2^{17}=32\pmod{35}$. This is neither $1$ nor
$-1$ modulo $35$, so $2$ is a Miller witness. Yet
\[
 \gcd(2,35)=\gcd(32-1,35)=\gcd(32+1,35)=1.
\]
This is a counterexample to precisely the displayed extraction rule.
It is not a lower bound for general integer factorization and does not
rule out other algorithms.\hfill$\square$

\paragraph{CX-ALG-015: integer factorization with a supplied exponent multiple.}
\emph{Statement.} Given integers $N\ge2$ and $T\ge1$, with the promise
\[
       \lambda'(N):=\operatorname{lcm}_{p\mid N}(p-1)\mid T,
\]
the complete factorization of $N$ is computable deterministically in time
polynomial in $\log N+\log T$ under $A_D$. This is the input-length form
of Miller (1976), Lemma 5 and Theorem 3, p.~310. In particular, a supplied
value of $\varphi(N)$ or $\lambda(N)$ meets the promise. Neither value is
asserted to be available for an arbitrary $N$.

\emph{Proof of the stated input-length form.} Remove powers of two, test
primality, and detect perfect powers by integer-root computations. These
are polynomial-time operations. For a perfect power, recurse on its base
and multiply the resulting exponents. It remains to split an odd composite
with at least two distinct prime divisors. Write $T=2^s t$ with $t$ odd.
For every $a\leq\lceil3(\log N)^2\rceil$, first compute $\gcd(a,N)$;
return a proper nontrivial gcd if found, and otherwise test, for units $a$,
\[
             \gcd\bigl(a^{2^j t}-1,N\bigr),\qquad 0\leq j<s.
\]
Powers are computed modulo $N$ by repeated squaring.

To prove that some test splits $N$, put
$e_p=v_2(p-1)$ and $e=\max_{p\mid N}e_p$. If $e_p=e$ and $e_q<e$ for
two divisors, the quadratic character $a\mapsto(a/p)$ is a nontrivial
character of $(\mathbf Z/N\mathbf Z)^\times$. If all $e_p=e$, use
$a\mapsto(a/p)(a/q)$ for distinct divisors $p,q$ instead. This too is
nontrivial, by the Chinese remainder theorem. Bach's generating bound
CX-ALG-003 implies that some tested unit has character $-1$.
In the first case its reduction modulo $p$ has multiplicative order with
2-part $2^e$, whereas the 2-part modulo $q$ is smaller. In the second
case one Legendre symbol is $-1$ and the other is $1$, with the same
consequence after interchanging $p,q$. Since $t$ is a multiple of the odd
part of each $p-1$, one obtains
\[
       a^{2^{e-1}t}=-1\pmod p,\qquad
       a^{2^{e-1}t}=1\pmod q.
\]
The corresponding gcd is divisible by $q$ and not by $p$, so is proper.
The promise gives $s\ge e$, hence this test occurs in the list. This
argument uses only reductions modulo the distinct primes; repeated prime
factors in $N$ cause no problem.

Every divisor $D$ encountered in the recursion still satisfies
$\lambda'(D)\mid T$. A recursion tree has $O(\log N)$ leaves and nodes
after perfect-power compression. Each node uses polynomially many bases
and at most $\log_2T$ squarings per base, so the total bit cost is
polynomial in the displayed input length. Geometrically the known exponent
annihilates the multiplicative groups of the reduced residue fibers;
quadratic characters separate two unknown fibers and a gcd recovers their
different supports. This is a precise promise case of the integer-factoring
gate in CX-ALG-014, not an algorithm for producing the exponent promise.
\hfill$\square$

\paragraph{CX-ALG-016: the fixed-generator resolvent needs a nonzero test.}
\emph{Statement.} The literal fixed-generator formula in
Bhargava--Ivanyos--Mittal--Saxena (2017), Theorem 3.1, author PDF p.~3,
can be zero. This does not refute their nonresidue-construction conclusion;
CX-ALG-017 supplies the missing nonzero choice.

\emph{Proof.} Let $q=7$, $r=3$, $d=3$, and
$E=\mathbf F_7[X]/(X^3-2)$, with $\alpha=X\bmod(X^3-2)$.
The nonzero cubes in $\mathbf F_7$ are $1,-1$, so this cubic has no root
and is irreducible. Take the primitive cube root $\zeta=4$. Then
$\alpha^7=4\alpha$ and $\alpha^{49}=2\alpha$, whence
\[
 L=\alpha+4\alpha^7+2\alpha^{49}
   =(1+16+4)\alpha=0.
\]
Consequently $L^{(7^3-1)/3}=0$, whereas the stated right-hand side is
$\zeta^{-1}=2$. The displayed formula and its hypotheses were checked
against the rendered original PDF page, not merely its text extraction.
The invalid step is division by $L$ when inferring
$L^{q^k-1}=\zeta^{-1}$ from $L^{q^k}=\zeta^{-1}L$.
\hfill$\square$

\paragraph{CX-ALG-017: a basis-search repair for the resolvent.}
\emph{Statement.} Let $E=\mathbf F_q[X]/(f)$ be a supplied field of degree
$d=rk$, where $r$ is prime. Suppose a primitive $r$th root of unity
$\zeta\in\mathbf F_q$ is supplied, and assume $r\mid q-1$ for odd $r$,
or $4\mid q-1$ when $r=2$. Then an $r$th nonresidue in $E$, and its
nonresidue norm in $\mathbf F_q$, can be found deterministically in time
polynomial in $d\log q$. No analytic hypothesis is needed for this statement.

\emph{Proof.} Write $Q=q^k$ and define the $\mathbf F_q$-linear map
\[
       P(z)=\sum_{i=0}^{r-1}\zeta^i z^{Q^i}.
\]
As a polynomial in $z$, it is nonzero and has degree
$Q^{r-1}<Q^r=|E|$. Thus it cannot vanish at every element of $E$.
Since it is linear, it cannot vanish on every element of the power basis
$1,\alpha,\ldots,\alpha^{d-1}$. Scan this basis and choose
$\beta=P(\alpha^j)\ne0$. Cyclically shifting the sum and using
$z^{Q^r}=z$ gives $\beta^Q=\zeta^{-1}\beta$, hence
$\beta^{Q-1}=\zeta^{-1}$.

Put $S=1+Q+\cdots+Q^{r-1}$. For odd $r$, write $Q=1+rc$ and expand
modulo $r^2$ to get
\[
 S\equiv r+rc\frac{r(r-1)}2\equiv r\pmod{r^2}.
\]
For $r=2$, the hypothesis gives $S=1+Q\equiv2\pmod4$.
In both cases $S/r\equiv1\pmod r$. Therefore
\[
 \beta^{(q^d-1)/r}
   =\left(\beta^{Q-1}\right)^{S/r}=\zeta^{-1}\ne1.
\]
This proves that $\beta$ is an $r$th nonresidue. Its norm satisfies
\[
 \left(N_{E/\mathbf F_q}\beta\right)^{(q-1)/r}
      =\beta^{(q^d-1)/r}=\zeta^{-1},
\]
so is also a nonresidue. There are $d$ basis candidates and $r\le d$
terms per test; all exponent bit lengths are $O(d\log q)$.
Field arithmetic in the supplied representation and these exponentiations
therefore take polynomial time.\hfill$\square$

For the inherited CX\_GRH CX-N13, obtain $\zeta$ under $A_D$ from a
prime-field nonresidue. If odd $r\mid m$, apply the theorem with
$q=p,d=m$. If $r\nmid m$, a prime-field nonresidue $a$ remains a
nonresidue upstairs, since
\[
 a^{(p^m-1)/r}
   =\left(a^{(p-1)/r}\right)^{1+p+\cdots+p^{m-1}}
   =\left(a^{(p-1)/r}\right)^m\ne1.
\]
The remaining quadratic-even case is proved next. Thus the prior
construction is preserved with an explicit nonvanishing justification.
In the Kummer-cover language, the repair chooses a nonzero vector in the
specified Frobenius eigenspace before passing to its torsor class; a
particular chosen generator need not have a nonzero projection.

\paragraph{CX-ALG-018: the even-degree quadratic case, with no GRH.}
\emph{Statement.} For odd $p$ and even $m$, a supplied
$E=\mathbf F_p[X]/(f)$ of degree $m$ admits a deterministic
$\operatorname{poly}(m\log p)$ construction of a quadratic nonresidue.
This is a fresh complete re-proof of the even-degree clause of inherited
CX\_GRH CX-N13, not a new claim of priority over BIMS17 Corollary 1.3.

\emph{Proof.} Put $Q=p^{m/2}$ and $\alpha=X\bmod f$.
If $p\equiv1\pmod4$ or $4\mid m$, then $Q\equiv1\pmod4$.
The element $u=\alpha^Q-\alpha$ is nonzero, since $\alpha$ has degree
$m$ and is not in the subfield of degree $m/2$. It satisfies $u^Q=-u$.
Consequently
\[
       u^{(p^m-1)/2}=(-1)^{(Q+1)/2}=-1,
\]
so it is the required nonresidue.

It remains that $p\equiv3\pmod4$ and $m\equiv2\pmod4$.
Compute the subfield $K=\mathbf F_{p^2}\subseteq E$ as the kernel of
the $\mathbf F_p$-linear map $z\mapsto z^{p^2}-z$.
Choose $v\in K\setminus\mathbf F_p$ from a computed basis and set
$u=v^p-v\ne0$. Then $u^p=-u$ and $c=u^2\in\mathbf F_p$ is a
nonresidue: a square $c$ would force $u\in\mathbf F_p$.
Because $-c$ is a square, the exponent formula
$t=(-c)^{(p+1)/4}$ gives $t^2=-c$. Thus $i=u/t$ satisfies $i^2=-1$,
and $K=\mathbf F_p(i)$ is explicitly represented.

Start with $z=i$, of order four. Its norm to $\mathbf F_p$ is one.
Whenever the current $z=a+bi$ has norm one, its order is a power of two
at least four; hence $b\ne0$, as the 2-part of $\mathbf F_p^\times$ has
order two. The nonzero numbers
\[
                 \delta_+=(a+1)/2,\qquad \delta_-=(a-1)/2
\]
have product $-b^2/4$, a nonsquare. Exactly one is a square in
$\mathbf F_p$. Determine which by Euler's criterion, choose its square
root $x=\delta^{(p+1)/4}$, and set $y=b/(2x)$.
Using $a^2+b^2=1$ verifies directly that
$(x+yi)^2=a+bi=z$. The norm of this root is $1$ if
$\delta=\delta_+$ and $-1$ if $\delta=\delta_-$.
Replace $z$ by this root, thereby doubling its order, until the norm is
$-1$.

This loop terminates after $O(\log p)$ steps. Indeed the norm-one subgroup
has order $p+1$, whereas the 2-part of $K^\times$ has order
$2^{v_2(p^2-1)}=2^{1+v_2(p+1)}$. An element's order cannot continue
doubling inside the former subgroup. At termination
\[
        z^{(p^2-1)/2}
          =N_{K/\mathbf F_p}(z)^{(p-1)/2}=-1.
\]
The extension $E/K$ has odd degree $m/2$, so this $z$ remains a nonresidue
in $E$. All operations are modular exponentiation or linear algebra in
the supplied $m$-dimensional field, with polynomial bit cost.
\hfill$\square$

The quadratic square-root algorithm then applies in $E$ with this known
nonresidue. In characteristic two the separate square-root problem is
inverse Frobenius, $a\mapsto a^{2^{m-1}}$, and there is no quadratic
nonresidue. These are finite-\'etale and Kummer computations; SGA 1 V
Section 7 and SGA 4 IX 3.2 are the same companions as for CX-ALG-008.

\paragraph{CX-ALG-021: the genus quotient only needs rational Dirichlet-GRH.}
\emph{Statement.} For an imaginary quadratic field $K$, put
$C=\operatorname{Cl}(\mathcal O_K)$ and $G=C/C^2$. The small-prime
spectral and mixing bounds in CX-ALG-020 hold on $G$ under $A_D$, with
$n=2$, $\mathfrak m=(1)$ and $Q=|D_K|$. No Hecke-GRH assumption for
the remaining, non-genus characters of $C$ is needed. If $C$ has exponent
two, this covers the whole class-group graph.
Moreover, the images of prime ideals of prime norm at most
\[
                         (8\log|D_K|+15)^2
\]
generate $C/C^2$ under $A_D$.

\emph{Proof.} Let $H$ be the Hilbert class field of $K$.
By class field theory $\operatorname{Gal}(H/K)=C$. Complex conjugation
preserves $H$: the conjugate of an unramified abelian extension of $K$
is again such an extension, and the maximal one is unique. Its action on
$C$ is inversion, since
$\mathfrak a\overline{\mathfrak a}=(N\mathfrak a)$ is principal.
Actual complex conjugation has order two and restricts nontrivially to
$K$, so
\[
       \operatorname{Gal}(H/\mathbf Q)=C\rtimes\langle c\rangle,
                    \qquad cac^{-1}=a^{-1}.
\]
The subgroup $C^2$ is normal. Hence $F=H^{C^2}$ is Galois over
$\mathbf Q$, and its Galois group is
$(C/C^2)\times\langle c\rangle$, since inversion is trivial on a group
of exponent two. In particular $F/\mathbf Q$ is abelian.

Every character $\chi$ of $C/C^2$ extends to two one-dimensional
characters $\chi_+,\chi_-$ of $\operatorname{Gal}(F/\mathbf Q)$, by
assigning $c$ the values $1,-1$. Directly on the two cosets of
$\operatorname{Gal}(F/K)$,
\[
       \operatorname{Ind}_{\operatorname{Gal}(F/K)}^{
                       \operatorname{Gal}(F/\mathbf Q)}\chi
                         =\chi_+\oplus\chi_-.
\]
Artin induction of L-functions and abelian class field theory over
$\mathbf Q$ therefore express the class-character Hecke function
$L_K(s,\chi)$ as a product of the two associated Dirichlet L-functions.
This includes $\chi=1$, for which the product is $\zeta_K$.
The input $A_D$ gives precisely the zero-location premises for all these
functions.

For the explicit generator bound, fix any nontrivial character $\chi$ of
$C/C^2$. Its kernel in $C$ defines a quadratic unramified extension
$F_\chi/K$ inside $H$. The preceding conjugation argument shows that
$F_\chi/\mathbf Q$ is biquadratic. Its absolute discriminant is
$|D_K|^2$ by the relative-discriminant tower formula and unramifiedness.
Apply Bach--Sorenson Theorem 5.1 to the nonidentity element of
$\operatorname{Gal}(F_\chi/K)$, regarded inside
$\operatorname{Gal}(F_\chi/\mathbf Q)$. Its zeta hypothesis follows from
$A_D$, since this degree-four field is abelian over $\mathbf Q$.
The resulting unramified rational prime satisfies
\[
            p\le(4\log(|D_K|^2)+2.5\cdot4+5)^2
                    =(8\log|D_K|+15)^2.
\]
Its Frobenius is trivial on $K$, so it splits there, and either prime
ideal above it has $\chi$-value $-1$. If the prime ideals in the displayed
range generated a proper subgroup of $C/C^2$, a nontrivial character of
that elementary two-group would be trivial on the subgroup. The prime just
constructed contradicts this. This proves the generator assertion without
paying for the degree of the entire genus field.

The character-by-character proof of JMV Theorem 1.1, author p.~9,
uses only the L-function of the individual character being estimated.
The spectrum of the quotient graph consists exactly of characters
factoring through $C/C^2$ (Remark 1.2(a), p.~4). Thus all its eigenvalue
bounds hold under the premises just established. Lemma 2.1, pp.~7--8,
and the proof of Corollary 1.3, p.~10, give the claimed mixing bound
with $|G|=|C/C^2|$. If $G$ is trivial, the statement is immediate.
\hfill$\square$

This is the genus-level part of the arithmetic Picard action. It sharpens
the input split in CX-ALG-013: the full Hilbert class field may be
nonabelian over $\mathbf Q$, while its genus quotient is abelian and
already covered by $A_D$. It does not give the missing estimates for the
other class characters or a general deterministic class-group algorithm.

\paragraph{CX-GEO-001: existential constructible feasibility.}
\emph{Statement.} Conditional on $M$, deciding whether a quantifier-free
Boolean circuit of rational polynomial equalities and inequations has a
complex solution belongs to $\mathrm{AM}$. Its emptiness problem belongs to
$\mathrm{coAM}\subseteq\Sigma_2^{\mathrm P}$. Polynomial expressions are
given by arithmetic circuits and all original variables are explicitly listed.

\emph{Proof.} Represent each atomic equality $f(x)=0$ by a new Boolean
variable $b$ and a new variable $u$, with equations
\[
 b(b-1)=0,\qquad bf(x)=0,\qquad (1-b)(u f(x)-1)=0.
\]
If $b=1$, these require $f(x)=0$. If $b=0$, they require $f(x)\ne0$
and an inverse $u=f(x)^{-1}$. Conversely any assignment to $x$ extends
with $b$ equal to the truth value of $f(x)=0$; when $b=1$, any $u$ works.
An inequation is the negation of its equality atom.
For each Boolean gate, introduce its output variable $t$ and use
\[
 t-uv=0\quad(\mathrm{AND}),\qquad
 t-u-v+uv=0\quad(\mathrm{OR}),\qquad
 t+u-1=0\quad(\mathrm{NOT}),
\]
where $u,v$ here denote that gate's Boolean inputs. Impose $t(t-1)=0$
for each output if desired, and impose that the final output equals $1$.
Induction through the acyclic Boolean circuit proves that the resulting
polynomial system is feasible exactly when the original formula is true
at some $x$. The construction shares each original arithmetic circuit;
it does not expand a Boolean formula into disjunctive normal form.
Its size and rational-constant bit lengths increase polynomially.
Apply CX-HN-001 to the resulting system and take complements for emptiness.
The reduction concerns equations and inequations over $\mathbf C$, not
ordered-field inequalities. Geometrically it realizes the given constructible
locus as the image of an explicitly presented affine scheme and preserves
nonemptiness.\hfill$\square$

\paragraph{CX-GEO-002: radical membership.}
\emph{Statement.} Conditional on $M$, for
$I=(f_1,\ldots,f_s)\subseteq\mathbf Q[x_1,\ldots,x_n]$ and $g$ given
by rational arithmetic circuits, radical membership $g\in\sqrt I$ is in
$\mathrm{coAM}$ and nonmembership is in $\mathrm{AM}$.

\emph{Proof.} By the Hilbert Nullstellensatz over $\mathbf C$,
$g\notin\sqrt{I\mathbf C[x]}$ exactly when there is a point $x$ with
$f_i(x)=0$ for every $i$ and $g(x)\ne0$. This is equivalent to feasibility of
\[
               f_i(x)=0\quad(1\leq i\leq s),\qquad y g(x)-1=0.
\]
For completeness, the rational and complex membership tests agree. If
$g^m\in I\mathbf C[x]$, a representation by a finite list of polynomial
multiples gives a finite linear system with rational coefficients for those
multipliers' coefficients. A rational linear system having a complex solution
has a rational solution, by row reduction. Thus $g^m\in I$. The converse
follows by extension of scalars. The new system has polynomial size, so
CX-HN-001 gives the asserted classes. In scheme language the question is
whether the principal open $D(g)$ in $\operatorname{Spec}(\mathbf Q[x]/I)$
has a geometric point. Nilpotents matter only through the radical in this
particular problem.\hfill$\square$

\paragraph{CX-GEO-003: reduced containment and equality.}
\emph{Statement.} Conditional on $M$, inclusion and equality of affine
algebraic sets specified by rational arithmetic-circuit equations are in
$\mathrm{coAM}$. This statement does not assert the same bound for
scheme-theoretic ideal inclusion.

\emph{Proof.} Put $X=V(f_1,\ldots,f_s)$ and
$Y=V(g_1,\ldots,g_t)$. Then $X\nsubseteq Y$ exactly when
\[
        f_i(x)=0\quad(1\leq i\leq s),\qquad
        \sum_{j=1}^t u_jg_j(x)=1                              \tag{CX.1}
\]
is feasible. Indeed, the last equation is impossible if all $g_j(x)$ vanish;
if one does not vanish, its inverse supplies a solution with all other
$u_j=0$. For $t=0$ the sum is zero, making (CX.1) impossible, correctly
expressing $Y=\mathbf A^n$. For an empty $f$-list there is no restriction
on $x$, also correctly. Thus noninclusion is in $\mathrm{AM}$ and
inclusion is in $\mathrm{coAM}$.

Equality can be handled by the same single-reduction method. Introduce
$b$ with $b(b-1)=0$. Impose
\[
\begin{aligned}
 b f_i(x)&=0,& (1-b)g_j(x)&=0,\\
 b\left(\sum_j u_jg_j(x)-1\right)&=0,&
 (1-b)\left(\sum_i v_if_i(x)-1\right)&=0.
\end{aligned}
\]
For $b=1$ this is a point of $X\setminus Y$; for $b=0$ it is a point
of $Y\setminus X$. Hence inequality is one HN instance of polynomial size.
Finally, the ideals $(x)$ and $(x^2)$ have the same complex zero set but
are unequal. They show exactly why this reduction tests underlying reduced
sets rather than the structure sheaves or ideal inclusions.
For GCT this test applies when defining equations of the orbit closures
are actually supplied; producing such equations remains a separate gate.
\hfill$\square$

\paragraph{CX-GEO-004: a finite Frobenius test with prescribed excluded primes.}
\emph{Statement.} Let $L/\mathbf Q$ be finite Galois, let $G$ be its Galois
group, and let $E_i=L^{H_i}$ for subgroups $H_i\leq G$ ($i=1,2$).
Let $\chi_i$ be the permutation character of $G/H_i$. Set $n=[L:\mathbf Q]$
and $D=|D_L|$. Suppose the right halves of the critical strips for
$\zeta_L$ and $\zeta_{\mathbf Q}$ are zero-free, as supplied by $M$, or by
$A_D$ if $L/\mathbf Q$ is abelian. If $\chi_1\ne\chi_2$, there is a prime
unramified in $L$ such that
\[
       p\leq B_L=(4\log D+\tfrac52 n+5)^2                  \tag{CX.2}
\]
and the numbers of degree-one prime ideals of $E_1$ and $E_2$ above $p$
differ. Equivalently the numbers of $\mathbf F_p$-points of the two integer
models differ.

For any nonempty finite set $S$ of rational primes put
$d=\prod_{p\in S}p$, $F=\mathbf Q(\sqrt d)$, and $L'=LF$.
The same assertion holds with $p\notin S$ and bound $B_{L'}$, under the
corresponding hypotheses for $\zeta_{L'}$ and $\zeta_{\mathbf Q}$.
If $n'=[L':\mathbf Q]$ and $D'=|D_{L'}|$, an explicit upper estimate is
\[
 n'\leq2n,\qquad \log D'\leq2\log D+n\log|D_F|
                     \leq2\log D+n\log(4d),               \tag{CX.3}
\]
so the entirely explicit larger bound
\[
        \bigl(8\log D+4n\log(4d)+5n+5\bigr)^2            \tag{CX.4}
\]
also suffices. If $S=\varnothing$, use (CX.2) directly.

\emph{Proof.} Choose $\sigma\in G$ such that
$\chi_1(\sigma)\ne\chi_2(\sigma)$. Bach--Sorenson,
\emph{Explicit Bounds for Primes in Residue Classes} (1996), Theorem 5.1,
p.~1729, gives an unramified prime whose Frobenius conjugacy class is that
of $\sigma$, with bound (CX.2). Their Section 2, p.~1720, specifies the
needed zero-free right-half assumption for the top-field zeta function and
$\zeta_{\mathbf Q}$; no Artin-holomorphy conjecture is required.
For an unramified prime the Frobenius orbits on $G/H_i$ correspond to the
prime ideals above $p$, and the orbit lengths are their residue degrees.
Its fixed points therefore count precisely the degree-one primes.
This proves the first assertion.

For the exclusion refinement, every prime of $S$ ramifies in $F$. For odd
$p\mid d$ this follows from the quadratic discriminant formula; if $2\in S$,
then $d$ is even and $D_F=4d$, so 2 ramifies as well. Consequently every
prime in $S$ ramifies in $L'$, since its ramification index in the subextension
$F/\mathbf Q$ divides its index upstairs. Choose any lift of $\sigma$ to
$\operatorname{Gal}(L'/\mathbf Q)$ and apply the same theorem there. The
prime obtained is unramified in $L'$, hence avoids $S$, and its restriction
to $L$ has the required character difference.

We justify (CX.3) without an unstated discriminant estimate. If $F\subseteq L$,
then $L'=L$ and the inequalities are immediate. Otherwise $L$ and $F$ are
linearly disjoint because $F$ is quadratic. If $\omega$ is an integral basis
generator for $\mathcal O_F$, the order $\mathcal O_L[\omega]\subseteq
\mathcal O_{L'}$ has relative discriminant generated by $D_F$. The relative
field discriminant divides the order discriminant. The tower formula gives
$D'\leq D^2|D_F|^n$. Also $n'=2n$ and $|D_F|\leq4d$.
Substitution in (CX.2) yields (CX.4). If $L/\mathbf Q$ is abelian, $LF$ is
also abelian, so the stated $A_D$ alternative survives the exclusion step.
The exceptional source case $L=\mathbf Q$ cannot have unequal characters
for the two subfields in this statement, so causes no gap.\hfill$\square$

\emph{Consequences and limits.} Agreement at all the tested primes implies
$\chi_1=\chi_2$. Hence the corresponding characteristic-zero permutation
representations are isomorphic and their Artin $L$-functions, including
ramified factors formed on inertia invariants, agree. In particular
$\zeta_{E_1}=\zeta_{E_2}$. This is a finite trace certificate, not necessarily
an efficient algorithm in a short presentation of $E_i$: a common Galois
closure may have large degree. Nor is equality of permutation characters
an isomorphism theorem for finite covers, as the next result proves.
The Grothendieck framework is SGA 1 V Section 7 and the Frobenius trace
formalism; cote 98 scans 11--12 are the finite-$G$-set archival companion
carried from the CX\_GRH crosswalk.

\paragraph{CX-GEO-005: equal characters do not identify transitive $G$-sets.}
\emph{Statement.} Let $V=\mathbf F_2^3$, $G=\operatorname{GL}(V)$,
$X=V\setminus\{0\}$ and $Y=V^*\setminus\{0\}$ with the natural and dual
actions. Then $X$ and $Y$ are transitive seven-element $G$-sets with identical
permutation characters but are not isomorphic as $G$-sets.

\emph{Proof.} For $g\in G$, the number of fixed nonzero vectors is
$2^{\dim\ker(g-I)}-1$. The corresponding number of fixed nonzero covectors
is $2^{\dim\ker((g^{-1})^{\mathsf T}-I)}-1$. The matrices
$(g^{-1})^{\mathsf T}-I$ and $g^{\mathsf T}-I$ have the same rank, since
they differ by multiplication by an invertible matrix and sign. Transpose
preserves rank, so both fixed-point counts are equal.

Suppose the two $G$-sets were isomorphic. The stabilizer $H$ of a nonzero
vector, say $e_1$, would stabilize some nonzero covector and therefore its
kernel, a plane $W\subset V$. Such an invariant plane does not exist.
If $e_1\in W$, its image $W/\langle e_1\rangle$ is a line in
$V/\langle e_1\rangle$, invariant under all of
$\operatorname{GL}_2(\mathbf F_2)$ because $H$ induces that full group.
No line is invariant under that group. If $e_1\notin W$, projection identifies
$W$ with the two-dimensional quotient, so $W$ is a complement to
$\langle e_1\rangle$. But $H$ contains all transvections
$(a,v)\mapsto(a+u(v),v)$ with $u$ a linear functional on that quotient.
Choosing $u(v)=1$ for a nonzero projected vector in $W$ moves its lift by
$e_1$, which cannot stay in $W$. Both cases contradict invariance.
Thus no $G$-set isomorphism exists.\hfill$\square$

This is a representation-theoretic counterexample, needing no GRH and no
inverse-Galois assertion. For any connected finite \`etale $G$-torsor that
realizes this group, its two associated connected covers exhibit the same
limitation. No particular unverified realization over $\mathbf Q$ is used.
The example distinguishes the Galois category from its linearized
representation category; the latter is what Frobenius traces recognize.

\paragraph{CX-GEO-006: orbit membership with an explicit action.}
\emph{Statement.} Conditional on $M$, suppose an algebraic group $G$ and its
action on an affine space have polynomial-length rational circuit descriptions,
with every inverted denominator explicitly listed. Given rational points
$v,w$, deciding whether $w\in G(\mathbf C)v$ belongs to $\mathrm{AM}$.

\emph{Proof.} Introduce coordinates for $g$, impose the supplied equations
and inequations defining $G$, and impose the coordinate equations
$g\cdot v=w$. For example, in $\operatorname{GL}_n$ an additional variable
$z$ and equation $z\det(g)-1=0$ enforce invertibility; the determinant has
a polynomial-size arithmetic circuit. Inverted denominators of a rational
action are handled by the same inverse-variable equations. The resulting
constructible-feasibility instance has polynomial size and is equivalent to
orbit membership by definition of the action. Apply CX-GEO-001.
Orbit-closure membership is not represented by this system: points in the
boundary need not equal $g\cdot v$ for any $g$. If closure equations are
provided, CX-HN-001 and CX-GEO-003 apply to those equations; without them
the equation-access problem persists.\hfill$\square$

\paragraph{CX-GEO-007: polynomially parameterized orbit closures already admit a PSPACE test.}
\emph{Statement.} Let the ambient coordinates of $V$ be explicitly listed.
Suppose the action on a supplied rational vector $v$ is given by
division-free arithmetic circuits with rational coefficients for all its
coordinates. In either of the following
cases, membership of a rational point $w$ in the complex orbit closure is
in $\mathrm{PSPACE}$, without either GRH input: (a) the group is
$\operatorname{GL}_n$ and the map $g\mapsto g\cdot v$ extends by the supplied
circuits to a polynomial map on all $\operatorname{Mat}_n$; (b) the group is
$\operatorname{SL}_n$ and the supplied action coordinates are polynomial
in its matrix entries. This statement charges the explicit ambient
coordinate count in the input size.

\emph{Proof.} Guo--Saxena--Sinhababu (2019), Theorem 1.3, p.~5, together
with Lemma 4.1 and Theorem 4.2, pp.~15--16, gives a polynomial-space test
for whether zero lies in the Zariski closure of the image of a polynomial
map specified by arithmetic circuits: this is their approximate polynomial
satisfiability problem. In case (a), $\operatorname{GL}_n$ is dense in
$\operatorname{Mat}_n$. For a polynomial map $F$ on the latter, every closed
set containing $F(\operatorname{GL}_n)$ has closed inverse image containing
a dense subset, hence contains $F(\operatorname{Mat}_n)$. Their image
closures are equal. Apply the cited test to $F-w$.

For case (b), construct a polynomial map from an affine space into
$\operatorname{SL}_n$ whose image contains its Gaussian-elimination big
cell. Products of elementary unipotent matrices parameterize arbitrary
lower and upper unitriangular matrices. A diagonal matrix of determinant
one is a product of adjacent diagonal blocks
$\operatorname{diag}(t,t^{-1})$: for diagonal entries $d_1,\ldots,d_n$,
use $t_i=d_1\cdots d_i$ in positions $(i,i+1)$ for $1\leq i<n$.
Such a block belongs to the image of a
six-factor elementary-unipotent word, because, in its $2$ by $2$ block,
\[
 E_{12}(t)E_{21}(-t^{-1})E_{12}(t)
 E_{12}(-1)E_{21}(1)E_{12}(-1)
       =\operatorname{diag}(t,t^{-1}).
\]
Use independent parameters in each of the six positions of this fixed
word. The resulting map is polynomial and always has determinant one;
the displayed specializations show that it contains every such diagonal
block in its image. Place these words between the lower and upper
unitriangular words. Its image contains every $LDU$ decomposition in the
big cell, which is a dense open subset of $\operatorname{SL}_n$.
The number of elementary factors and parameters is $O(n^2)$ and the
matrix product has polynomial-size circuits. Composing with the supplied
action on $v$ gives a polynomial map with the same image closure as the
orbit map, by the same closed-inverse-image argument. Subtract $w$ and
apply the cited PSPACE algorithm. For $n=1$, the special-linear action is
trivial and membership reduces to $w=v$.\hfill$\square$

This gives an equation-free route for this exact orbit-closure encoding.
It does not give an AM bound from $M$ for approximate satisfiability, and
it does not turn an exponentially long list of ambient coefficients into
a polynomial-length input. Those are separate questions from the theorem.

\paragraph{CX-GEO-008: a finite Frobenius generating certificate.}
\emph{Statement.} Let $L/\mathbf Q$ be finite Galois with group $G$ and
assume the zeta hypotheses of CX-GEO-004. For each unramified rational
prime $p\leq B_L$, choose any representative $g_p$ of its Frobenius class.
The elements $g_p$ generate $G$. More strongly, a subgroup meeting every
Frobenius class observed in this range equals $G$. A sublist of at most
$\lfloor\log_2|G|\rfloor$ representatives suffices to generate $G$.
For a prescribed excluded set $S$, replace the bound by $B_{L'}$ from
CX-GEO-004 and use only primes outside $S$. As before, $A_D$ suffices for
abelian $L/\mathbf Q$, and $M$ supplies the general hypothesis.

\emph{Proof.} Let $H<G$ be proper and let $X=G/H$. Double counting gives
\[
 \sum_{g\in G}|\operatorname{Fix}_X(g)|
      =\sum_{x\in X}|\operatorname{Stab}_G(x)|=|G|.
\]
The identity fixes $|X|>1$ points. If every element fixed at least one,
the left side would exceed $|G|$. Hence there is $g$ with no fixed coset.
Since $g$ fixes $aH$ exactly when $a^{-1}ga\in H$, its conjugacy class
misses $H$. The explicit Chebotarev theorem used in CX-GEO-004 produces
that class at an unramified prime $p\leq B_L$. With excluded primes,
choose a lift to $\operatorname{Gal}(L'/\mathbf Q)$ and apply the same
argument at bound $B_{L'}$; every prime of $S$ ramifies in $L'$ and is
therefore avoided. Thus no proper subgroup can meet every observed class,
and in particular no proper subgroup contains all the chosen $g_p$.

Starting from the trivial subgroup, choose a listed $g_p$ outside the
subgroup generated so far until it equals $G$. Each proper enlargement
has index at least two, so after $k$ selections the generated subgroup has
order at least $2^k$. Consequently $k\leq\lfloor\log_2|G|\rfloor$.
For $G=1$ the empty list suffices.\hfill$\square$

In SGA 1 V Section 7 this is a finite quotient of the arithmetic
fundamental group, with a finite list of Frobenius conjugacy classes
certifying that no proper subgroup contains the image. The bound is
relative to the specified extension and its discriminant; it is not a
uniform presentation of the entire profinite fundamental group. The
archival companion is the finite-$G$-set discussion of cote 98 scans
11--12 recorded in the preceding CX\_GRH crosswalk, and the work target
is ORDER item 30's effective finite Galois-image question.

\paragraph{CX-GEO-009: bounded Laurent degenerations give an AM promise.}
\emph{Statement.} Let $F:\mathbf A^n\to\mathbf A^m$ and
$w\in\mathbf Q^m$ be specified by division-free rational arithmetic
circuits of total size $s$. Supply in unary an upper bound $D\ge1$ for
their syntactic degrees and an integer $B\ge0$.
Conditional on $M$, it is in $\mathrm{AM}$ to decide whether there are
Laurent polynomials
\[
 \beta_i(\varepsilon)=\sum_{j=-B}^{B}c_{i,j}\varepsilon^j,
 \qquad c_{i,j}\in\mathbf C,
\]
such that every coordinate of $F(\beta)-w$ belongs to
$\varepsilon\mathbf C[\varepsilon]$.
Consequently image-closure membership is in $\mathrm{AM}$ on the promise
that every positive instance admits such an arc with $D$ and $B$ bounded
polynomially in the original encoding length. This includes the polynomial
orbit parameterizations of CX-GEO-007 with the same arc promise.

\emph{Proof.} Use the coefficients $c_{i,j}$ as unknowns. At each circuit
gate $g$ of syntactic degree $d_g$, maintain arithmetic circuits for the
coefficients of its substituted Laurent polynomial in the range
$[-Bd_g,Bd_g]$. Input-variable gates have their prescribed coefficient
unknowns; a rational constant has only a constant coefficient. At an
addition or subtraction gate, add or subtract coefficient circuits.
At a multiplication gate, take the usual coefficient convolutions.
Syntactic degrees add at multiplication, so no exponent outside the stated
range can occur. Discard gates unused by the outputs; then each remaining
degree is at most $D$.

There are at most $2BD+1$ coefficient positions per gate, and each product
uses at most $(2BD+1)^2$ pairs. Sharing all earlier coefficient circuits
therefore gives a polynomial-size arithmetic-circuit construction in
$s,n,m,B,D$. For every output impose that all coefficients of exponents
$-BD,\ldots,0$ in $F(\beta)-w$ vanish. These are polynomial equations in
the $n(2B+1)$ unknown coefficients. The resulting HN instance has polynomial
encoding size and is feasible exactly when the displayed bounded Laurent
arc exists. Koiran's theorem gives the claimed AM bound under $M$.
The construction also covers $B=0$, when it is ordinary exact feasibility.

Any such arc puts $w$ in the image closure by GSS Lemma 4.1 and Theorem
4.2, pp.~15--16. Conversely the promise provides a bounded arc for every
point of that closure under consideration, proving the promise corollary.
All oracle queries concern this explicit coefficient system; its bounded
solution field is the exact number-field zeta input. An $A_D$ alternative
requires the bounded abelian-witness promise for this coefficient system,
not merely for the original rational points $v,w$.
\hfill$\square$

The generic GSS bound on Laurent precision can be exponential; writing
that bound in binary would not justify the polynomial-size coefficient
construction. This result therefore isolates a concrete smaller gate for
the promised families: a polynomially bounded degeneration, rather than
access to all equations of the orbit closure. It does not assert that
arbitrary GCT orbit closures meet that promise. The geometric objects are
finite Laurent jets of a parameterized degeneration; the group-action
framework is SGA 3 I 4.6.3.1 and 4.7.1--4.7.2, with cote 28 scan 3 as the
archival companion inherited from CX\_GRH. The work target is ORDER item
28's Koiran/GCT route.

\subsection{Gates retained rather than silently discharged}

\paragraph{CX-LIM-001: general derandomization and GCT.}
The exact remaining premises of the cited routes are uniform hitting sets
or independent circuit-hardness hypotheses. GCT V Theorem 4.7, p.~33,
assumes a specific permanent border-hardness bound. Its Lemma 4.8, on the
same page, extracts a hitting set from a sufficiently explicit separating
system. Neither construction is furnished by a statement about zeros of
Dirichlet or Hecke $L$-functions. This records the boundary of these proofs;
it is not a proof that no logical implication from GRH to general
derandomization could ever exist. No $\mathrm{BPP}=\mathrm P$, integer
factorization in $\mathrm P$, or $\mathrm{VP}\ne\mathrm{VNP}$ is claimed.

\paragraph{CX-LIM-002: positivity and purity.}
The earlier ALPOGE\_GCT P-F7 and P-F10 and PLETH PLE-F5/F7 already show
why pure geometry, nonnegative Hilbert data, efficient counting and actual
orbit-closure multiplicities are different inputs. Their established
counterexamples and corrected positivity-index implications are carried,
not rebranded. The new tests CX-GEO-002--CX-GEO-003 can verify specified
reduced-locus properties once equations are given. They do not construct a
uniform positive multiplicity formula or a strict asymptotic quotient-rank
inequality. Those are the surviving GCT gates.

\paragraph{CX-LIM-003: the separate counting premise.}
The attached CX\_GRH return also assumed an independent theorem denoted
$E$ and drew ETH/counting consequences in CX-E01--CX-E05. Tonight's two
analytic inputs $A_D$ and $M$ do not contain $E$. Those rows therefore do
not become consequences of these two inputs by being present in the same
archive. This is an input classification, not a new assertion about ETH.

\subsection{A smaller field-construction gate for deterministic factorization}

\paragraph{CX-ALG-023: prime-degree witnesses and maximal coefficient units.}
\emph{Statement.} Let $p,r$ be distinct primes, put
$s=\operatorname{ord}_r(p)$, $R=r^{v_r(p^s-1)}$, and
\[
 D_r=\mathbf F_p[Z]/(1+Z+\cdots+Z^{r-1}),\qquad \zeta=Z\bmod\Phi_r.
\]
Write $\rho_a(\zeta)=\zeta^a$ for $a\in\mathbf F_r^\times$, and
$\omega(a)$ for its multiplicative Teichm\"uller lift to the relevant
$r$-power modulus. An explicitly supplied field $F=\mathbf F_{p^r}$
permits, unconditionally in $\operatorname{poly}(r\log p)$ time,
construction of $\eta\in D_r^\times$ satisfying
\[
          \operatorname{ord}(\eta)=R,\qquad
          \rho_a(\eta)=\eta^{\omega(a)}\quad(a\in\mathbf F_r^\times).
\]
Conversely, such an $\eta$ permits construction of an explicit
$\mathbf F_{p^r}$ in the same polynomial bound. In particular, a supplied
$\mathbf F_{p^m}$ already supplies this coefficient unit whenever
$r\mid m$, regardless of the value of $v_r(p^s-1)$.

The equivalence is classical: Lenstra, \emph{Finding isomorphisms between
finite fields} (1991), Theorems (5.1)--(5.2), pp.~338--340, with its
normal-basis algorithm, Theorem (3.1), pp.~334--336. The following full
proof makes the exact constructive input and the new application to the
remaining CX\_GRH gate explicit; no priority for the equivalence is claimed.

\emph{Proof.} Compute a normal-basis element $\alpha$ for
$F/\mathbf F_p$ by the cited unconditional linear-algebra algorithm.
Form $C=F\otimes_{\mathbf F_p}D_r$, let $\sigma$ act by $p^s$-power
Frobenius on $F$ and trivially on $D_r$, and set
\[
             \beta=\sum_{i=0}^{r-1}\zeta^i\sigma^i(\alpha).
\]
The factors of $\Phi_r$ over $\mathbf F_p$ all have degree $s$.
Since $s\mid r-1$, the fields $F$ and $\mathbf F_{p^s}$ are linearly
disjoint. Thus each component of $C$ is a field, and the elements
$\sigma^i(\alpha)$ remain linearly independent over its coefficient
field $\mathbf F_{p^s}$. The displayed sum is consequently nonzero in
every component: $\beta$ is a unit. This use of a normal basis prevents
the vanishing defect in CX-ALG-016.

On $C$, $\sigma$ equals $p^s$-power Frobenius. Shifting the sum gives
$\sigma(\beta)=\zeta^{-1}\beta$. With $u=(p^s-1)/R$ and
$\delta=\beta^u$ it follows that
\[
                  \delta^R=\zeta^{-1},\qquad \delta^{rR}=1.
\]
Compute the lifts $\omega(a)$ modulo $rR$ and their inverses, and put
\[
          y=\prod_{a\in\mathbf F_r^\times}
                        \rho_a(\delta)^{\omega(a)^{-1}}.
\]
Reindexing the product by $a\mapsto ba$ proves
$\rho_b(y)=y^{\omega(b)}$. Moreover,
\[
          y^R=\prod_a\zeta^{-a\omega(a)^{-1}}
                   =\zeta^{-(r-1)}=\zeta.
\]
Hence $y$ has order $rR$. The element $\eta=y^r$ has order $R$,
and $\eta^{p^s-1}=1$. Therefore $\sigma(\eta)=\eta$.
Since $\gcd(s,r)=1$, the fixed field of $\sigma$ on $F$ is
$\mathbf F_p$, so $C^\sigma=D_r$. This proves $\eta\in D_r$ and
the required equivariance.

For the converse, equivariance and transitivity of the $\rho_a$ on the
components of $D_r$ imply that the image of $\eta$ in every component
$\mathbf F_{p^s}$ has order $R$. It is an $r$th nonresidue there.
The polynomial $Y^r-\eta$ is irreducible in each component: in a field
containing $\boldsymbol\mu_r$, its splitting group is a subgroup of
the prime-order group $\boldsymbol\mu_r$, and reducibility would force
a root and hence make $\eta$ an $r$th power. Form
\[
                       C'=D_r[Y]/(Y^r-\eta).
\]
Extend $\rho_a$ by $\rho_a(Y)=Y^{\omega(a)}$, taking lifts modulo
$rR$. The relations are preserved, and uniqueness and multiplicativity
of the lifts give an action of $\mathbf F_r^\times$.
This group acts transitively on the field components of $C'$, so its
fixed algebra is a field. It has degree $r$ over $\mathbf F_p$:
after extending scalars to an algebraic closure, $C'$ has $r(r-1)$
points, the action on them is free because its action on $\zeta$ is
free, and there are exactly $r$ orbits. Invariants commute with this
scalar extension, being kernels of explicit linear maps. The fixed
field is therefore an explicit $\mathbf F_{p^r}$.

All algebras in this proof have dimension at most $r(r-1)$.
The integer $R$ has $O(r\log p)$ bits; the lifts can be calculated by
Hensel lifting the simple roots of $X^{r-1}-1$ modulo $r$, or by
modular exponentiation. Fixed algebras and their multiplication tables
are obtained by linear algebra. These operations give the stated bound.
If $r\mid m$, the subfield $F\subseteq\mathbf F_{p^m}$ is the kernel
of $z\mapsto z^{p^r}-z$, computable in time polynomial in the supplied
field size in bits. Applying the first construction finishes the proof.
\hfill$\square$

Geometrically $D_r$ is a finite \`etale coefficient algebra; adjoining
an $r$th root is a Kummer cover, and taking the displayed invariants
is descent of a finite \`etale algebra. The companions are SGA~1,
Expos\'e~V, \S7 and SGA~4, Expos\'e~IX, 3.2, as in CX\_GRH
CX-N13--N15. The new use is the additional $r\mid m$ branch of
CX\_GRH CX-U02, not a claim that its older proof was absent or wrong.

\paragraph{CX-ALG-024: a pure prime-degree zeta function supplies the auxiliary prime.}
\emph{Statement.} Let $r$ be an odd prime and $p\ne r$ any prime, and write
\[
 F_r=\mathbf Q(p^{1/r}),\quad K_r=\mathbf Q(\zeta_r),\quad
 E_r=\mathbf Q(\zeta_r,p^{1/r}).
\]
Under $A_D$, RH for $\zeta_{F_r}$ is equivalent to RH for
$\zeta_{E_r}$. It implies a rational prime $q\ne p,r$ with
\[
 q\equiv1\pmod r,\qquad p^{(q-1)/r}\not\equiv1\pmod q,
 \qquad q\le B_{p,r},
\]
where the explicit sufficient bound is
\[
 B_{p,r}=
 \left(4r(2r-3)\log r+4(r-1)^2\log p
                       +\tfrac52r(r-1)+5\right)^2.
\]
Consequently an explicit $\mathbf F_{p^r}$ is constructible in
$\operatorname{poly}(r\log p)$ time from this analytic input.
Under the task convention $M$ supplies RH for the degree-$r$ field
$F_r$; its role here is conditional.

\emph{Proof.} The discriminant of the order $\mathbf Z[\zeta_r]$ is
$\pm r^{r-2}$, as follows by taking the product of the derivatives
of $\Phi_r$ at its roots. Thus $p$ is unramified in $K_r$.
The polynomial $X^r-p$ is Eisenstein at every prime of $K_r$ above $p$,
so $[E_r:K_r]=r$. The group of the splitting field $E_r/\mathbf Q$
is the affine group
\[
             G=\mathbf F_r\rtimes\mathbf F_r^\times,
       \qquad (a,b):t\longmapsto bt+a
\]
acting on the $r$ roots of $X^r-p$.
Let $P$ be this permutation character, and let $Q$ be the regular
character of the quotient $G/\mathbf F_r$, inflated to $G$.
There is an identity of characters
\[
                    \operatorname{reg}_G=(r-1)(P-1)+Q.
\]
At the identity both sides equal $r(r-1)$; at a nonzero translation,
$P=0$ and $Q=r-1$; and at an affine transformation with $b\ne1$,
$P=1$ and $Q=0$. These cases exhaust $G$, proving the identity.
Artin formalism, including its ramified inertia-invariant local factors,
now gives
\[
          \zeta_{E_r}(s)=\zeta_{K_r}(s)
                    \left(\frac{\zeta_{F_r}(s)}{\zeta_{\mathbf Q}(s)}
                    \right)^{r-1}.
\]
The abelian field $K_r$ and $\mathbf Q$ have their zero-location
assertions under $A_D$. Thus in the open critical strip off the middle
line the displayed identity proves the asserted equivalence. This uses
no general Artin-holomorphy conjecture.

For a discriminant bound, the relative order
$\mathcal O_{K_r}[p^{1/r}]$ has discriminant generated, up to sign, by
$r^r p^{r-1}$. The field discriminant divides this order discriminant.
The discriminant tower formula therefore yields
\[
      |D_{E_r}|\le
                  r^{\,r(2r-3)}p^{(r-1)^2},
      \qquad [E_r:\mathbf Q]=r(r-1).
\]
Apply Bach--Sorenson (1996), Theorem~5.1, p.~1729, to the conjugacy
class of the automorphism fixing $K_r$ and sending $p^{1/r}$ to
$\zeta_r p^{1/r}$. Its conjugates send that root to
$\zeta_r^a p^{1/r}$ for $a\in\mathbf F_r^\times$.
The theorem's precise zeta input (\S2, p.~1720) is supplied by the
preceding factorization and $A_D$. Substitution of the degree and
discriminant bounds into $(4\log|D_{E_r}|+2.5[E_r:\mathbf Q]+5)^2$
gives $B_{p,r}$.

The resulting unramified prime splits in $K_r$, hence $q\equiv1\pmod r$.
At a prime of $E_r$ over $q$, its Frobenius relation is
\[
       (p^{1/r})^q\equiv\zeta_r^a p^{1/r},\qquad
              p^{(q-1)/r}\equiv\zeta_r^a\ne1.
\]
Unramifiedness excludes $p$ and $r$, and the root remains nonzero.
For prime $r$, the last inequality is equivalent to
$r\nmid(q-1)/\operatorname{ord}_q(p)$.
The image of $p$ consequently generates the order-$r$ quotient of
$(\mathbf Z/q\mathbf Z)^\times$. Equivalently, $p$ is inert in the
unique degree-$r$ subfield of $\mathbf Q(\zeta_q)$.

This is exactly the algebraic input to Adleman--Lenstra (1986),
Algorithm~A, steps (2)--(3), pp.~350--351, and Lemma~6, p.~352.
That algorithm constructs an irreducible degree-$r$ polynomial by
Gaussian periods using $\operatorname{poly}(q\log p)$ operations;
the primality and power-congruence tests for $q$ use the same bound.
Enumerating primes until the first suitable $q$ halts by $B_{p,r}$,
which is $O(r^4\log^2(rp))$. No additional Hecke or radical oracle
enters this construction. This proves the polynomial bound.
\hfill$\square$

The geometric interpretation is a relation among the permutation
objects of an affine finite Galois cover: the degree-$r$ cover and its
cyclotomic quotient determine the required zeta input of the full cover.
It belongs to the Artin-motive and finite-monodromy framework of SGA~1
V \S7 and to the representation-ring operations of SGA~6 V.
It sharpens the auxiliary-prime input of CX\_GRH CX-N11--N15;
the earlier sufficient family remains valid.

\paragraph{CX-ALG-025: a complete factoring composition with a smaller analytic family.}
\emph{Statement.} Let $f\in\mathbf F_{p^m}[X]$ have degree $d\ge2$,
with the base field explicitly supplied. If $p\le d$, deterministic
factorization is unconditional in $\operatorname{poly}(d\log(p^m))$
time. Suppose $p>d$, and define
\[
 \mathcal B(p,d,m)=\{r\le d:\ r\text{ prime},\ r\nmid(p-1)m,
                 \ v_r(p^{\operatorname{ord}_r(p)}-1)>1\}.
\]
Under $A_D$ and RH for the finite family
\[
             \zeta_{\mathbf Q(p^{1/r})},\qquad r\in\mathcal B(p,d,m),
\]
one can factor $f$ deterministically in
$(d^{\log d}\log(p^m))^{O(1)}$ time. Only the primes from this family
actually encountered need their zeta hypothesis. If $\mathcal B$ is
empty, the result is $A_D$-only. In general $M$ supplies the remaining
input, conditionally. Every number field in the displayed new family
has prime degree at most $d$. If every prime $r\le d$ either divides
$m$ or satisfies $v_r(p^{\operatorname{ord}_r(p)}-1)=1$, the same
quasipolynomial bound is unconditional.

\emph{Proof.} First reduce to a square-free polynomial by gcds and
inverse Frobenius; repeated-factor multiplicities are retained.
For a square-free factor $h$, put
$A=\mathbf F_{p^m}[X]/(h)$ and compute the $\mathbf F_p$-linear fixed
algebra
\[
                         B=\{x\in A:x^p=x\}.
\]
If $h$ has $t$ irreducible factors, its Chinese-remainder decomposition
shows that $B\simeq\mathbf F_p^t$. This identity is used to prove
correctness, not as supplied factorization data. A basis and a
multiplication table for $B$ are found by linear algebra in dimension
at most $m\deg h$. If $t=1$, $h$ is irreducible and no split is needed.

If $p\le d$, take a nonscalar basis element $b\in B$ and test
$b-a$ for $a\in\mathbf F_p$. One of these is a nonzero zero divisor:
in the unknown split coordinates $b$ has at least two different values.
Thus at most $p$ tests of multiplication matrices supply a proper
idempotent, and recursive splitting has polynomial cost in this case.

Assume henceforth $p>d$. Apply the unconditional IKRS algorithm,
published Theorem~5.8, author PDF pp.~28--29, to the explicit algebra
$B$ and its scalar subalgebra $\mathbf F_p$. It either returns a
nonzero zero divisor or a semiregular automorphism $\sigma$ of order
$t$. Choose a prime $r\mid t$ and put $\tau=\sigma^{t/r}$;
this is a semiregular automorphism of prime order $r\le d<p$.
The group in question acts freely on the split components; in this
setting its powers have the required semiregularity directly.

There are four ways to obtain the coefficient unit $\eta$ of
CX-ALG-023 for this $r$.
If $r\mid m$, use the degree-$r$ subfield of the supplied base field.
If $r\mid p-1$, CX-ALG-004 finds an $r$th nonresidue $a\in\mathbf F_p$
under $A_D$; the prime-degree Kummer polynomial $X^r-a$ is irreducible,
so CX-ALG-023 applies. If $v_r(p^{\operatorname{ord}_r(p)}-1)=1$,
the formal root $\zeta\in D_r$ itself already has the required order
and equivariance. Otherwise $r\in\mathcal B$, and CX-ALG-024 followed
by CX-ALG-023 supplies $\eta$ from its one additional zeta hypothesis.
All these operations take polynomial time in $d\log(p^m)$.

For completeness, here is why that coefficient unit closes the
automorphism step. IKRS Proposition~4.6, published author PDF
pp.~18--19, constructs a Teichm\"uller element
$x\in(B\otimes D_r)^\times$ with
$\tau(x)=\zeta x$. Every component of $B\otimes D_r$ is
$\mathbf F_{p^s}$, where $s=\operatorname{ord}_r(p)$, so its
$r$-primary unit group has exponent $R$.
If $\langle x,\eta\rangle$ were cyclic, the element $\eta$ of
order $R$ would generate it. Then $x$ would be a power of the
coefficient element $\eta$ and hence fixed by $\tau$, contradicting
$\tau(x)=\zeta x$ and invertibility of $x$.
IKRS Lemma~4.2, published author PDF p.~16, therefore supplies a
nonzero zero divisor in $B$ in polynomial time. Its proof also gives
the descended idempotent: the annihilator of the difference of two
powers of these Teichm\"uller elements is invariant under every
$\rho_a$, so its identity lies in $B$.

From any nonzero zero divisor $z\in B$, compute the identity $e$ of
the nonzero proper ideal $Bz$ by linear algebra; it is a proper
idempotent. Such an $e\in B\subseteq A$ has values zero or one in
the factor fields. Hence $\gcd(h,e)$, with $e$ represented by its
polynomial residue, is a proper factor of $h$. Recurse on the two
factors. There are at most $2d-1$ nodes; at each node the IKRS cost
is bounded by $(d^{\log d}\log(p^m))^{O(1)}$, and all the added
operations are polynomial. If each $r\le d$ is covered by the
supplied-field or valuation-one branch, choose those branches first;
then no analytic input is used. This proves both stated bounds.
\hfill$\square$

This is a new composition and finite-input refinement using classical
algorithms. It gives a strictly smaller sufficient analytic family than
the earlier CX\_GRH CX-N15 construction through degrees $sr$ and
auxiliary primes dividing $r-1$. It also strictly enlarges the earlier
$A_D$ cases by exploiting $r\mid m$. The unrestricted $A_D$ gate is
now to construct a degree-$r$ field, equivalently its maximal coefficient
unit, only when
\[
       p>d,\qquad r\nmid p(p-1)m,\qquad
             v_r(p^{\operatorname{ord}_r(p)}-1)>1.
\]
This is a sufficient construction problem, not a claim that those RH
assumptions are necessary. The geometric task is still computation of
primitive idempotents of a finite \`etale algebra, SGA~1 V \S7;
the Kummer descent is SGA~4 IX 3.2. It changes CX\_GRH CX-U02 and
ORDER~28 at their actual coefficient-field gate.

\paragraph{CX-ALG-026: exact field extensions have the same prime-degree gate.}
\emph{Statement.} Given $\mathbf F_{p^m}$ and a positive integer $n$
charged in unary, an explicit extension of degree $n$ is constructible
in $\operatorname{poly}(n\log(p^m))$ time under $A_D$ and RH for
$\zeta_{\mathbf Q(p^{1/r})}$ only for prime divisors $r\mid n$ such that
\[
       r\ne p,\qquad r\nmid(p-1)m,\qquad
               v_r(p^{\operatorname{ord}_r(p)}-1)>1.
\]
If every prime divisor $r$ of $n$ either divides $pm$ or satisfies
$r\ne p$ and $v_r(p^{\operatorname{ord}_r(p)}-1)=1$, no analytic
hypothesis is needed. In particular, this holds if every such $r$ divides
$pm$. The output size, and hence the unary or equivalent dense-field
convention on $n$, is essential.

\emph{Proof.} Use Lenstra (1991), Theorem~(9.1), pp.~344--346.
It requires an irreducible degree-$r$ polynomial over the supplied base
field only for primes $r\mid n$ not dividing $m$.
For $r=p$, the Artin--Schreier construction of its Theorem~(7.1),
p.~342, is unconditional and polynomial in $p\log(p^m)$; here
$p\le n$ is charged to the output size. For $r\ne p$, obtain an
explicit $\mathbf F_{p^r}$ by the appropriate branch of
CX-ALG-023--CX-ALG-024. In the valuation-one branch use the converse
of CX-ALG-023 with $\eta=\zeta$; in the $r\mid p-1$ branch use
the prime-field nonresidue under $A_D$.
Convert this explicit field to an irreducible polynomial by Lenstra's
Theorem~(1.1), p.~330. Since $r\nmid m$, that polynomial stays
irreducible over $\mathbf F_{p^m}$: the two finite extensions of
$\mathbf F_p$ have coprime degrees. These are all inputs to
Theorem~(9.1), which supplies the claimed extension and polynomial
bound. Under the stated unconditional promise, use the supplied-field,
Artin--Schreier, or valuation-one branch before considering the
$r\mid p-1$ nonresidue branch. Thus no analytic branch is called.\hfill$\square$

The output is a connected degree-$n$ finite \`etale cover of the
supplied finite-field point. This is the explicit construction counterpart
to the isomorphism/uniqueness theorem for finite fields, not a replacement
for general polynomial factorization. Its run locators are CX\_GRH
CX-N15/CX-U02 and the construction interface in ORDER~28; the same
SGA~1 V \S7 finite-\'etale companion applies.

\subsection{Binary cocharacters remove a Laurent-precision promise}

\paragraph{CX-GEO-010: a polynomial-size equation system for one-parameter limits.}
\emph{Input convention.} Let $G/\mathbf Q$ be a connected split reductive
group acting linearly on $V=\mathbf Q^N$, with $N\ge1$.
The input supplies equations and inverted denominators defining $G$,
rational arithmetic circuits for its action, and a weight basis for a
split maximal torus $T\simeq\mathbf G_m^r$. Its cocharacter lattice is
identified with $\mathbf Z^r$, and the complete list of coordinate
weights
\[
                     a_1,\ldots,a_N\in\mathbf Z^r
\]
is supplied in binary. These data are promised to define the indicated
group, action, and full cocharacter lattice. All dimensions, circuits,
equations, denominators, and coefficient bit lengths are charged to the
input size. The points $v,w$ have rational coordinates. Rank zero is the
trivial connected group and can be handled directly; assume $r\ge1$ below.

\emph{Statement.} The relation
\[
 \mathcal R(v,w):\quad
 \exists g,h\in G(\mathbf C),\ b\in\mathbf Z^r:\quad
       \lim_{t\longrightarrow0}\lambda_b(t)\,g v=h w,
 \qquad \lambda_b(t)=(t^{b_1},\ldots,t^{b_r}),
\]
has a polynomial-time many-one reduction to Hilbert Nullstellensatz over
$\mathbf Q$. In particular it is in AM under $M$. No bound on the
absolute exponents is supplied in unary. The equations remain a
polynomial-size system when $w$ is an additional unknown vector instead
of a fixed input.

\emph{Proof: bounded sign patterns.} Put
$H=\max(1,|a_{ji}|)$ and $B=N!H^N$.
Every sign pattern of the numbers $a_j\cdot b$ attained by an integer
$b$ is attained by another integer vector $b'$ with $|b'_i|\le B$.
Here is a complete size argument. Orient each row having nonzero sign
so that its dot product is positive, and require that dot product to be
at least one; retain a zero row as an equality. Write the unknown vector
as $x^+-x^-$, with both vectors nonnegative, and introduce a nonnegative
slack variable for each inequality. This gives a feasible system
\[
                       Cx=c,\qquad x\ge0,
\]
having $N$ rows, integral coefficients of absolute value at most $H$,
and entries of $c$ equal to zero or one. It is feasible because the
original integer vector supplies one such solution.

Choose a solution with the fewest positive coordinates. Its supported
columns of $C$ are linearly independent: otherwise a linear dependence,
oriented to have a positive coefficient, could be subtracted up to the
first vanishing coordinate, preserving nonnegativity and all equations.
Let $k\le N$ be the number of supported columns. Choose $k$ rows making
their square submatrix nonsingular. Cramer's rule gives the supported
coordinates with common positive denominator the absolute value of that
integer determinant. Multiplying by this denominator gives nonnegative
integer coordinates at most $k!H^k\le B$, since both the denominator and
each replaced-column determinant have that bound. The resulting
$x^+-x^-$ has coordinates of absolute value at most $B$. The zero
equalities remain zero, and every positive inequality remains positive,
because its right hand side has been multiplied by a positive integer.
Thus the entire prescribed sign pattern is preserved. If $k=0$, the
zero vector supplies the same conclusion.

\emph{Proof: algebraic encoding.} Choose
$L=\lceil\log_2(B+1)\rceil$. Represent $b_i$ as
\[
           b_i=\sum_{\ell=0}^{L-1}2^\ell
                         (u_{i\ell}-v_{i\ell}),
        \qquad u_{i\ell}^2=u_{i\ell},\quad
               v_{i\ell}^2=v_{i\ell}.
\]
Every complex solution of these equations assigns actual zero-or-one
values to the bits, so the represented vector is an integer vector.
Write $d_j=a_j\cdot b$ and give each coordinate three Boolean selectors
$s_j^+,s_j^0,s_j^-$ whose sum is one. Choose
\[
       J=L+\lceil\log_2(rH)\rceil+1.
\]
Let $P_j,Q_j$ be nonnegative integers represented by $J$ Boolean bits
each, and impose
\[
 s_j^+(d_j-1-P_j)=0,\qquad
 s_j^0d_j=0,\qquad s_j^-(d_j+1+Q_j)=0.
\]
The selectors now describe exactly the signs of $d_j$. All necessary
values of $|d_j|-1$ fit, since $|d_j|\le rH(2^L-1)$.
No ordering predicate over $\mathbf C$ has been inserted: the signs
are enforced by bounded integers represented by Boolean variables.

Impose the supplied equations for $g,h\in G$, including inverse
variables for all inverted denominators. Write $u=gv$ and $z=hw$ using
the supplied action circuits. For each coordinate impose
\[
              s_j^-u_j=0,\qquad
              (s_j^++s_j^-)z_j=0,\qquad
              s_j^0(u_j-z_j)=0.
\]
These are exactly the conditions that $t^{d_j}u_j$ have finite limit
$z_j$: a negative weight requires $u_j=z_j=0$, a zero weight requires
$u_j=z_j$, and a positive weight requires $z_j=0$.

Any solution therefore proves $\mathcal R(v,w)$. Conversely, from any
one-parameter limit, replace $b$ by the bounded vector having the same
signs. The limit has precisely the same coordinates, and all Boolean
and slack variables can be filled in. Hence the equation system is
feasible exactly when $\mathcal R(v,w)$ holds.
Finally $\log B=O(N\log N+N\log H)$, so $L,J$ and every coefficient
bit length are polynomial in the declared input size. Sharing the action
circuits gives a polynomial-size HN system. If $w$ is unknown, the action
circuits remain polynomial-size since the representation is linear in
that vector. Koiran's theorem gives the claimed AM consequence under
$M$.\hfill$\square$

This is an explicit encoding lemma proved here, not a claim of priority
for bounded integral witnesses in rational polyhedra. The geometric
object is a cocharacter filtration and its associated one-parameter
specialization. The inherited companions are SGA~3, Expos\'e~I,
4.6.3.1 and 4.7.1--4.7.2 for actions, and Expos\'e~VIII, 1.1--1.2 for
diagonalizable groups and their character lattices. Its run target is
the equation-access gate in CX\_GRH CX-K02/CX-K05 and ORDER~28; it
refines the bounded Laurent-arc route CX-GEO-009.

\paragraph{CX-GEO-011: orbit-closure intersection and the closed-target promise.}
\emph{Statement.} Under the explicit representation convention of
CX-GEO-010, deciding
\[
             \overline{Gv}\cap\overline{Gw}\ne\varnothing
\]
reduces in polynomial time to HN over $\mathbf Q$, and is consequently
in AM under $M$; disjointness is in coAM. If $Gw$ is promised closed,
the same upper bound holds for $w\in\overline{Gv}$. The nullcone test
$0\in\overline{Gv}$ is included without a separate promise, since the
orbit of zero is closed. These assertions concern complex Zariski
closures. They require neither defining equations for those closures nor
a unary bound on degeneration exponents.

\emph{Classical geometric input.} Every orbit closure of a reductive
group on an affine variety contains a unique closed orbit, and a
one-parameter subgroup has a limit in that orbit. An exact modern primary
locator is Bate--Herpel--Martin--R\"ohrle,
\emph{Cocharacter-closure and the rational Hilbert--Mumford Theorem},
arXiv:1411.7849v5, introduction p.~1 and Theorem~4.3, pp.~12--13,
specialized to the algebraically closed field $\mathbf C$.
Their introduction identifies the classical closed-subset formulation
as Kempf's Theorem~1.4 (1978). The current source check uses the displayed
primary theorem and its full proof, not a fresh audit of Kempf's
unretrieved original proof. Every cocharacter of a connected reductive
group is conjugate into the chosen maximal torus.

\emph{Proof.} First suppose $Gw$ is closed and
$w\in\overline{Gv}$. The unique closed orbit in that closure is then
$Gw$. The cited theorem gives a cocharacter with a limit in $Gw$.
Conjugate this cocharacter into $T$, moving the source point by a group
element and moving the limit within $Gw$. Thus $\mathcal R(v,w)$
holds. Conversely, if $\mathcal R(v,w)$ holds, its limit is in
$\overline{Gv}$; applying $h^{-1}$ and using invariance of the closure
gives $w\in\overline{Gv}$. This converse needs no closedness promise.
Applying CX-GEO-010 proves the promised membership statement.

For the intersection assertion, a nonempty intersection of the two
closures is a closed $G$-stable affine subvariety. It contains a closed
orbit: choose an orbit of minimal dimension in it, whose boundary, if
nonempty, would consist of smaller-dimensional orbits. That orbit is
closed also in $V$, since the intersection is closed. Uniqueness shows
that it is the closed orbit of both $\overline{Gv}$ and
$\overline{Gw}$. Choose a point $z$ on it. The preceding paragraph
then gives both $\mathcal R(v,z)$ and $\mathcal R(w,z)$.
Conversely, these two relations imply that $z$ belongs to both closures.
We have proved the exact existential equivalence
\[
 \overline{Gv}\cap\overline{Gw}\ne\varnothing
       \quad\Longleftrightarrow\quad
       \exists z\in\mathbf C^N:\
                  \mathcal R(v,z)\ \text{and}\ \mathcal R(w,z).
\]
Take two copies of the polynomial-size HN system in CX-GEO-010 and share
the unknown vector $z$. This is still one polynomial-size HN instance.
Koiran's theorem and complementation give the stated AM and coAM bounds.
\hfill$\square$

These are full reductions written here from classical geometry; no claim
that this is the first complexity formulation is made. The exact analytic
input is the Dedekind zeta function of a bounded algebraic solution field
of the constructed HN system, as in CX-HN-002. $M$ supplies it
conditionally. An $A_D$ alternative needs the bounded abelian-witness
promise for this new system. Rational input points and integral
cocharacters do not force the group elements or the shared limit point
to be defined over an abelian field.

The closed-target promise is substantive for membership. This argument
does not assert that an arbitrary point in an orbit closure is reachable
by a single cocharacter. Nor does deciding intersection construct a
separating invariant, a homogeneous system of parameters, or a
representation-theoretic multiplicity inequality. Such search and
lower-bound inputs remain the separate gates already recorded in
CX-NN-001--004 and CX-LIM-001--002. The geometric companion is the
reductive group action and its torus filtrations in SGA~3; the work
locators are CX\_GRH CX-K02/CX-K05, the GCT V orbit-intersection
interface (Theorem~9.8, p.~68), and ORDER~28.

% END complexity

% BEGIN root
\section{The two methods: precise transfers and repairs}
\label{sec:methods}

\subsection*{MT-A01. What the adopted Dirichlet theorem says adelically}
The supplied declarations \path{riemannHypothesis},
\path{LFunction_ne_zero_of_half_lt_re}, and
\path{LFunction_eq_zero_re} in \path{ALPOGE_Main.lean} have the
definitions printed in \path{ALPOGE_Statements.lean}.  We use these as the
cited input $A_D$, as requested.  The supplied archive contains these exported
declarations and the scan report, but not the imported tree
\path{Supplements.Kernels.MainTheoremsAll}; a kernel build of that tree is
not a check performed in this run.  This distinction concerns our evidence,
not the stipulated status of the input theorem.

Every unitary idele-class character of $\mathbb Q$ has the form
$\chi_0|\cdot|^{it}$, where $\chi_0$ is a finite-order Dirichlet character
and $t\in\mathbb R$.  Consequently $A_D$ gives the rank-one Hecke GRH over
$\mathbb Q$ for these characters, not just its finite-order specialization.
Indeed the norm map has a continuous section and compact kernel, which for
$\mathbb Q$ is $\widehat{\mathbb Z}^{\times}$ after the sign convention at
infinity.  A continuous character of this profinite group into the circle has
finite image: its image is a compact quotient of a profinite group, hence
totally disconnected, whereas an infinite closed subgroup of the circle is
the circle.  It therefore factors through some
$(\mathbb Z/N\mathbb Z)^\times$.  The remaining unitary character of
$\mathbb R_{>0}$ is $x^{it}$.  In the convention
$L(s,\chi_0|\cdot|^{it})=L(s+it,\chi_0)$, zeros are translated only in
their imaginary coordinate.  This proves the assertion.  Geometrically,
the finite-order part is a rank-one representation of the arithmetic
fundamental group; the continuous norm twist is an analytic normalization,
not an additional finite-monodromy motive.  The result does not cover
arbitrary characters over a nontrivial number field.

\subsection*{MT-M01. A nonzero counterexample to the draft's general zero-count lemma}
Masullo, Lemma 9.3.2.4, PDF p.~425, claims a local $O(\log |T|)$ zero count
for the Mellin transform of every $\varphi\in\mathcal S_E(G)$.
We checked the image of the original PDF page, including the displayed
Jensen inequality and its last inference.  The quantifier is present in the
PDF; this is not an inference from the model transcription.

Take $G=\mathbb R_{>0}$, trivial character, Haar measure $dt/t$, and set
\[
 g_a(u)=\frac{1}{\sqrt{4\pi a}}\exp\left(-\frac{u^2}{4a}\right),
 \qquad a=1+i,\qquad
 \varphi(e^u)=\frac{g_{1+i}(u)-g_{1-i}(u)}{2i}.
\]
The square roots use the branch with positive real part.  This is a real,
nonzero smooth function.  For every real $b$ and every nonnegative integer
$j$, the $j$th derivative of $e^{bu}g_{1\pm i}(u)$ is a polynomial in $u$
times $e^{bu-u^2/(4(1\pm i))}$.  Its absolute value is bounded by a
polynomial times $e^{bu-u^2/8}$, hence is rapidly decreasing.  Thus
$t^b\varphi(t)$ is Schwartz in the logarithmic coordinate for every $b$;
this is exactly the condition of Definition D.3.3.1, PDF p.~804 (also
\S1.1.2, PDF p.~11).  In particular $\varphi\in\mathcal S_E(G)$.

For $\operatorname{Re}a>0$, the Gaussian integral and analytic continuation
in $a$ give
\[
 \int_{\mathbb R}g_a(u)e^{su}\,du=e^{as^2},\qquad s\in\mathbb C.
\]
For completeness, the identity first follows by completing the square for
positive real $a$ and real $s$; both sides are entire in $s$ by Gaussian
domination on compact sets, and holomorphic in $a$ on
$\operatorname{Re}a>0$, where $\operatorname{Re}(1/a)>0$.  The identity
theorem proves the claimed continuation.  Therefore
\[
 \mathcal M\varphi(s)=e^{s^2}\sin(s^2).
\]
Its zeros on the positive imaginary axis are
$i\sqrt{k\pi}$, $k=1,2,\ldots$, and are simple.  For $I=[-1,1]$ and
$T\ge2$, their number in $I+i[T,T+1]$ is
\[
 \#\{k:T^2/\pi\le k\le(T+1)^2/\pi\}
   =\frac{2T+1}{\pi}+O(1).
\]
This is not $O(\log T)$.  The function is nonzero, so removing the zero
function from the lemma does not repair it.  The proof on p.~425 needs a
\emph{lower} bound for the modulus at a Jensen centre, whereas rapid decay
supplies upper bounds.  No statement about the zeros of a Hecke L-function
is contradicted by this example.

\subsection*{MT-M02. The displayed approximate identities are proper kernel projections}
Theorem 9.3.2.6, equation (9.3.2.10), PDF p.~426, defines subspaces by
intersecting kernels of all jet evaluations up to height $T$, then claims
their orthogonal projections tend strongly to the identity.  The printed
formula and the asserted inclusion in its proof have both been checked on
the page image.

Here is a counterexample even with a single zero, so the zero-count issue
in MT-M01 is irrelevant.  Put
\[
 g(u)=(4\pi)^{-1/2}e^{-u^2/4},\qquad
 \psi(e^u)=-g'(u).
\]
Then $\psi\in\mathcal S_E(\mathbb R_{>0})$, and integration by parts gives
$\mathcal M\psi(z)=z e^{z^2}$.  Its only zero is $0$, with multiplicity one,
strictly inside the strip for $I=[-1,1]$.  Choose Sobolev index $0$.
The draft's Hilbert norm is then
\[
 \|f\|_H^2=\int_{\mathbb R}|f(e^u)|^2e^{2|u|}\,du.
\]
The functional $\delta_0(f)=\int_{\mathbb R}f(e^u)\,du$ is continuous,
because Cauchy--Schwarz bounds its modulus by
$\|f\|_H(\int e^{-2|u|}\,du)^{1/2}=\|f\|_H$.
It is nonzero, since $\delta_0(g)=1$.  For every $T>0$ the subspace displayed
in (9.3.2.10) is the same proper closed subspace $\ker\delta_0$.
Its orthogonal projector cannot converge to the identity.  In particular,
the Riesz vector $r(e^u)=e^{-2|u|}$ has $\|r\|_H=1$ and is killed by every
one of these projectors.  It belongs to the indicated completion because
$C_c^\infty(\mathbb R)$ is contained in $\mathcal S_E$ in logarithmic
coordinates and is dense in this weighted $L^2$ space: truncate first and
then smooth on compact intervals, where the positive weight is bounded
above and below.  The ranges decrease when more evaluation
conditions are added, whereas the proof on p.~426 uses the opposite
inclusion.

\subsection*{MT-M03. The exact projection repair}
Let $H$ be a Hilbert space, let $(\ell_j)_{j\ge1}$ be continuous linear
functionals, and put
\[
 V_N=\bigcap_{j\le N}\ker\ell_j,\qquad V=\bigcap_{j\ge1}\ker\ell_j.
\]
Then $P_{V_N}\to P_V$ strongly.  The complementary projectors converge
strongly to $P_{V^\perp}$, the projector onto the closed span of the Riesz
vectors of the $\ell_j$.  Thus they give an approximate identity on the
spectral quotient $H/V\simeq V^\perp$.

Here is a full proof.  For fixed $x\in H$, put $y_N=P_{V_N}x$.  For
$M\ge N$, nestedness gives $P_{V_M}y_N=y_M$ and hence
\[
 \|y_N-y_M\|^2=\|y_N\|^2-\|y_M\|^2.
\]
The nonincreasing nonnegative sequence $\|y_N\|^2$ converges, so $(y_N)$ is
Cauchy.  Write its limit as $y$.  For each fixed $N$, all sufficiently late
$y_M$ lie in the closed space $V_N$, so $y\in V$.  Also
$x-y_N\perp V$ for every $N$; passing to the limit gives $x-y\perp V$.
Therefore $y=P_Vx$.  Taking complements proves the second assertion.
Finally the orthogonal complement of the closed span of the Riesz vectors
is precisely $V$, which proves the last identification.  Exhaustion by
height instead of an integer enumeration gives the same proof.

This repairs a functional-analytic direction error.  It does not identify
the common kernel with the closed adelic theta image, prove a Schauder
basis, establish nuclearity of the actual convolution operators, or justify
the geometric-kernel decomposition.  Those are separate statements and
must be checked in the quotient on which they are used.  In Grothendieck's
trace language this is the distinction between an annihilator and its
spectral quotient, before any trace or Fredholm determinant is taken.

\subsection*{MT-M04. An arithmetic replacement for the false general lemma}
For every fixed unitary Hecke character, the number of nontrivial zeros
in a unit-height window is $O_{K,\eta}(\log(|T|+3))$, unconditionally.
The full proof for arbitrary infinity type is NT-Z10 above: polynomial
vertical growth and a uniform lower bound at the Euler-product centre
give Jensen's estimate.  It includes pure norm twists, whose single
translated pole is excluded before applying Jensen at large height.
For a fixed primitive finite-order character a conductor-explicit
primary locator is Greni\'e--Molteni,
\emph{An explicit Chebotarev density theorem under GRH}, Appendix A,
Theorem A.1, author PDF p.~26: that appendix's zero count is unconditional.
Subtract the counts at consecutive heights.  The main term has derivative
$O(\log(|T|+3))$ for a fixed field and character and each remainder has the
same bound.  The arithmetic restriction therefore closes the local
zero-count gate, including infinite-order characters, without using GRH.
It supplies a route around MT-M01 for arithmetic divisors, but does not
establish Lemma 9.3.2.5 for arbitrary Schwartz convolution ideals.  The
constant in the infinite-order statement is fixed-character, not a
claimed uniform archimedean-conductor constant.

\subsection*{MT-ART01. A nonmonomial GL\texorpdfstring{$_2$}{2} consequence: odd Artin forms}
Assume $M_H$.  Let $\rho:G_{\mathbb Q}\to\operatorname{GL}_2(\mathbb C)$
be continuous, irreducible, odd and of finite image.  Then its Artin
L-function is entire and all of its nontrivial zeros lie on
$\operatorname{Re}s=1/2$.

The holomorphy input is already a theorem: Khare--Wintenberger,
\emph{Serre's modularity conjecture (I)}, Corollary 10.2(ii), author PDF
p.~21 (with Theorem 10.1(ii), pp.~20--21), associates $\rho$ to a weight-one
cuspidal newform.  Its standard L-function is entire.  Let $L$ be the finite
Galois extension cut out by $\rho$.  Brauer induction writes its character
as an integer linear combination of characters induced from
one-dimensional characters of subgroups of $\operatorname{Gal}(L/\mathbb Q)$.
Artin induction and class field theory therefore express $L(s,\rho)$ as
a finite product of integer powers of finite-order Hecke L-functions over
the corresponding fixed fields.  In the open critical strip, $M_H$ makes
each factor holomorphic and nonzero away from the critical line.  The
product, including its negative exponents, is consequently holomorphic and
nonzero there.  Combining this with the independent holomorphy theorem
proves GRH for $L(s,\rho)$.

This includes projective image $A_5$: no monomial expression with only
nonnegative exponents is required.  It is a new transfer in this catalogue,
not a new proof of the classical modularity theorem.  In motivic terms it
concerns odd rank-two Artin motives and their weight-one automorphic
realizations.  It gives no GRH for the weight-two L-function of an arbitrary
non-CM elliptic curve.  Nor does it prove holomorphy of arbitrary even
icosahedral or higher-dimensional Artin representations.  The inherited
Brauer divisor-support statements in RETURN\_ZETA\_TWO are the general
input behind this specialization.

\subsection*{MT-ART02. The missing positivity is a finite matrix at each zero}
Let $L/K$ be a finite Galois extension with group $G$, and let
$0<\operatorname{Re}s_0<1$.  By Brauer meromorphy define
\[
 n_\chi=\operatorname{ord}_{s=s_0}L(s,\chi,L/K)\in\mathbb Z,
 \qquad\Theta_{s_0}=\sum_{\chi\in\operatorname{Irr}(G)}n_\chi\chi.
\]
Form the $|G|$ by $|G|$ matrix
\[
 K_{g,h}=\Theta_{s_0}(g h^{-1}).
\]
Then all Artin L-functions for $L/K$ are holomorphic at $s_0$ if and only
if $K$ is positive semidefinite.  The matrix is Hermitian because
$\chi(g^{-1})=\overline{\chi(g)}$ and $n_\chi$ is real.

To prove the equivalence, consider the central convolution operator
$T=|G|^{-1}K$ on $\mathbb C[G]$.  On the matrix-coefficient space of an
irreducible character $\chi$, Schur orthogonality makes $T$ scalar with
eigenvalue $n_\chi/\chi(1)$ (interchanging $\chi$ and its dual under the
opposite convolution convention only permutes the list).  Each such space
has dimension $\chi(1)^2$ and together they exhaust the regular
representation.  Thus $K\succeq0$ is equivalent to $n_\chi\ge0$ for every
$\chi$, which is exactly absence of a pole of every factor.

The restrictions of $\Theta_{s_0}$ to cyclic subgroups are already actual
characters.  Indeed, for a cyclic subgroup $C$ and a one-dimensional
character $\eta$ of $C$, Frobenius reciprocity and Artin induction give
\[
 \langle\Theta_{s_0}|_C,\eta\rangle_C
   =\operatorname{ord}_{s=s_0}L(s,\eta,L/L^C)\ge0,
\]
because the right-hand side is an order of a holomorphic Hecke function in
the strip.  This proves positivity on every such restriction; it does not
prove the full matrix is positive.  Also
\[
 \Theta_{s_0}(1)=\sum_\chi n_\chi\chi(1)
   =\operatorname{ord}_{s=s_0}\zeta_L(s).
\]
The point of the criterion is the exact remaining object: a central
positive kernel in the Artin character ring.  $M_H$ removes all tests off
the critical line, while neither zero location nor cyclic positivity
supplies the full critical-line kernel.  An Alp\"oge-type positivity
argument would have to produce this kernel or an equivalent positive
decomposition, with the actual Artin orders as coefficients.

\subsection*{MT-ART04. All induced-linear positivity tests still leave a gap}
There is a virtual character which passes every positivity inequality
obtained by inducing a one-dimensional character of a subgroup, but is not
an actual character.  This refutes an abstract reconstruction rule; we do
not claim that this virtual character occurs as the orders of actual
Artin L-functions at a point.

Let $G=A_5$, and let $\chi$ be its four-dimensional augmentation character
for the natural action on five letters.  Define
\[
             \Theta=\operatorname{reg}_G-5\chi.
\]
The natural permutation action is doubly transitive, so its permutation
character has inner-product norm squared two.  Its invariant subspace is
one-dimensional; the orthogonal complement therefore has character
$\chi$ with norm squared one, proving that $\chi$ is irreducible.
The coefficient of $\chi$ in $\Theta$ is $4-5=-1$, so $\Theta$ is not an
actual character and its full matrix from MT-ART02 is not positive.

Nevertheless, for every $H\le G$ and every linear character $\eta$ of $H$,
\[
       \left\langle\Theta,\operatorname{Ind}_H^G\eta\right\rangle_G
                    \ge0.                                      \tag{MT.1}
\]
Here are all the details.  The conjugacy classes of $A_5$ have sizes
$1,15,20,12,12$, as follows by counting the double transpositions,
three-cycles and the two five-cycle classes.  A normal subgroup is a union
of these classes containing the identity, and its size divides $60$.
No proper nontrivial sum of these class sizes containing $1$ divides $60$;
thus $A_5$ is simple.  An action on cosets of a proper subgroup of index
$2,3,$ or $4$ would be an injective homomorphism into the corresponding
symmetric group, which is impossible by order.  Lagrange's theorem leaves
only the possible proper indices
\[
                    k\in\{5,6,10,12,15,20,30,60\}.
\]
Put $m=\langle\chi|_H,\eta\rangle_H$.  Since $\chi|_H$ has dimension four,
$0\le m\le4$.  Frobenius reciprocity says that
$\operatorname{Ind}_H^G\eta$, of dimension $k$, contains $m$ copies of
$\chi$, so $4m\le k$.  These two bounds imply $k-5m\ge0$ for every index
in the list except possibly $k=12$.  At $k=12$, $|H|=5$, and $H$ acts
regularly on the five letters.  The restriction of $\chi$ is the sum of
the four distinct nontrivial characters of $C_5$, so $m\le1$ and again
$k-5m\ge0$.  Finally
\[
 \left\langle\operatorname{reg}_G,
       \operatorname{Ind}_H^G\eta\right\rangle=k,
\]
proving (MT.1) for proper $H$.  For $H=G$, simplicity and nonabelianity
make the trivial character the only linear character, and its coefficient
in $\Theta$ is one.  This proves (MT.1) in all cases.

For comparison with the methods, Artin induction and Hecke holomorphy
give exactly inequalities of the form (MT.1) for an actual local Artin
order character.  The example proves that these inequalities alone do not
force a positive character, even before any question of zero location.
Thus a transfer of determinant positivity from the Alp\"oge method would
need additional arithmetic information, not just all the existing
induced-linear nonnegative orders.  The elementary exact computation
\path{checks/artin_cone_certificate.py} checks the negative matrix
quadratic form: with $v(g)=\chi(g)$ one has $v^*Kv=-900$.

\subsection*{MT-ART03. The classical simple-zero reduction}
Stark, \emph{Some effective cases of the Brauer--Siegel theorem},
Inventiones Mathematicae \textbf{23} (1974), Theorem 3, p.~144, proves
that if $\operatorname{ord}_{s_0}\zeta_L\le1$ and $s_0\ne1$, then all
Artin factors for $L/K$ are holomorphic at $s_0$.  We cite this classical
result rather than count another proof as a finding.  Applied with the
divisor-support theorem under $M_H$, it confines any still-possible pole
inside the strip to a multiple zero of the top-field zeta on the critical
line.  No simplicity assertion for these Dedekind zetas is being made.

\subsection*{MT-Q01. Exact path quadrics in every characteristic}
For $k\ge2$ let
\[
 V_k=\left\{(x_1,\ldots,x_k):\sum_{i=1}^{k-1}x_ix_{i+1}=0\right\}
       \subset\mathbb A^k_{\mathbb Z},\qquad
 r=\lfloor k/2\rfloor,\quad d=\lceil k/2\rceil.
\]
Then for every finite field $\mathbb F_q$, including characteristic two,
\[
 \#V_k(\mathbb F_q)=q^{k-1}+(q-1)q^{d-1}.                 \tag{MT.2}
\]
In the relative Grothendieck ring of varieties over $\mathbb Z$ one has
\[
                 [V_k]=\mathbb L^{k-1}+\mathbb L^d-\mathbb L^{d-1}.
                                                               \tag{MT.3}
\]
Here $\mathbb L=[\mathbb A^1_{\mathbb Z}]$; the identity is an identity of
classes, not a claim of a direct-sum splitting of all integral realizations.

To prove these formulas, group the odd coordinates into $x\in\mathbb A^d$
and the even coordinates into $y\in\mathbb A^r$.  The quadratic form
becomes $x^{\mathsf T}Ay$, where the integral $d$ by $r$ matrix $A$ has an
$r$ by $r$ triangular minor with diagonal entries one.  It has a left
inverse over $\mathbb Z$ and its columns extend to a unimodular basis of
$\mathbb Z^d$.  Thus after an invertible integral change of the $x$
coordinates the equation is $\sum_{j=1}^r x_jy_j=0$, with $d-r$ free
$x$ coordinates.  Over $y=0$ its fibre is $\mathbb A^d$.  Over
$\mathbb A^r\setminus\{0\}$, the map of vector bundles
$\mathcal O^d\to\mathcal O$ given by $(y_1,\ldots,y_r,0,\ldots,0)$ is
surjective, and its kernel is a vector bundle of rank $d-1$.
It is Zariski locally trivial on the opens where a coordinate $y_j$ is
invertible.  Additivity and the vector-bundle relation give
\[
 [V_k]=\mathbb L^d+(\mathbb L^r-1)\mathbb L^{d-1},
\]
which is (MT.3).  Counting the same fibres over $\mathbb F_q$ proves
(MT.2) directly, without a purity theorem.

Consequently the exact error term in the unweighted orthogonality identity
of the supplied Fable sketch is
\[
             q\,\#V_k(\mathbb F_q)=q^k+(q-1)q^d.
\]
For a fixed nonzero additive-character parameter its quadratic sum is
$q^d$, by summing first over $x$.  This gives square-root size when $k$ is
even; for odd $k$ the extra factor $q^{1/2}$ reflects the one-dimensional
radical.  It explains the sign and the exact constant for this family.
The general weighted operator estimate in
\path{OURS_FABLE_PROGRAM_SUGGESTIONS_20261120_EXPLAINED.md},
\S3, remains a separate inherited result.  This calculation does not prove
the sign or norm assertions for arbitrary sheaves in Alp\"oge's steps 3
and 5, nor identify this elementary model with all the constructions in
the omitted Lean dependency tree.

\subsection*{MT-Q02. Its arithmetic zeta function and shifted zero lines}
The exact all-prime formula in MT-Q01 gives
\[
 Z(V_k/\mathbb F_q,T)
     =\frac{1-q^{d-1}T}{(1-q^{k-1}T)(1-q^dT)},
 \qquad
 \zeta_{V_k}(z)
     =\frac{\zeta(z-k+1)\zeta(z-d)}{\zeta(z-d+1)}.       \tag{MT.4}
\]
Indeed apply (MT.2) over every $\mathbb F_{q^m}$ in
$Z=\exp(\sum_{m\ge1}\#V_k(\mathbb F_{q^m})T^m/m)$ and use
$-\log(1-aT)=\sum_{m\ge1}a^mT^m/m$.  Multiplying the resulting factors
over rational primes proves the second identity first in a right
half-plane of absolute convergence, and then meromorphically.
Under $A_D$ the nontrivial divisor of this function is supported on
\[
    \operatorname{Re}z=k-\tfrac12,\qquad
    \operatorname{Re}z=d+\tfrac12,\qquad
    \operatorname{Re}z=d-\tfrac12.
\]
Divisor orders are signed and cancellations are kept exactly as in
(MT.4); the denominator's zeros are possible poles, not additional zeros.
The shifted pole of $\zeta$ at $1$ and shifted trivial zeros are recorded
by the same explicit formula, separately from this nontrivial-divisor
assertion.  This is an all-characteristic, concrete Artin--Tate example
linking the path model of the method with arithmetic zeta functions and
the Grothendieck ring.  It illustrates why a mixed-weight scheme has
several shifted lines rather than a single universal critical line.

\subsection*{MT-L01. Damped Euler products and Adams operations}
Let $\rho$ be a finite-image complex representation of $G_K$ with no
invariant vectors, of dimension $d$.  At each finite place $v$ put
$A_v=\rho(\operatorname{Frob}_v)|_{V^{I_v}}$.  Its eigenvalues have
absolute value one.  For $|c|\le1$ define, initially for $\operatorname{Re}s>1$,
\[
 D_c(s,\rho)=\prod_v\det(1-cA_v(Nv)^{-s})^{-1}.
\]
If $L(s,\rho)$ is holomorphic and nonzero on
$H=\{\operatorname{Re}s>1/2\}$, then $D_c$ has a holomorphic nonvanishing
continuation to $H$, given by
\begin{align*}
 D_c(s,\rho)&=\exp\{c\log L(s,\rho)+E_c(s)\},\\
 E_c(s)&=\sum_v\sum_{k\ge2}
   \frac{c^k-c}{k}\operatorname{tr}(A_v^k)(Nv)^{-ks}.
\end{align*}
The branch of $\log L$ is fixed by the Euler series in
$\operatorname{Re}s>1$.

For a full proof, expand $-\log\det(1-X)$ in powers of $X$ in the initial
half-plane.  The terms with $k=1$ cancel in
$\log D_c-c\log L$, giving the displayed $E_c$.  On a compact subset with
$\operatorname{Re}s\ge1/2+\varepsilon$, its terms in absolute value sum to
at most
\[
 2d\sum_v\sum_{k\ge2}(Nv)^{-k(1/2+\varepsilon)}
 \le \frac{2d}{1-2^{-1/2-\varepsilon}}
       \sum_v(Nv)^{-1-2\varepsilon}<\infty.
\]
Normal convergence, also uniform in $|c|\le1$, proves that $E_c$ is
holomorphic in $s$.  A nonvanishing holomorphic function on the simply
connected half-plane $H$ has a holomorphic logarithm, uniquely normalized
by its existing branch on the right.  The exponential above is therefore
holomorphic, nonzero, and agrees with the product where it converges.
This proves the continuation.

Under $M_H$, Brauer induction supplies the needed holomorphic nonvanishing
on $H$ for $\rho$ with no invariants: all Hecke factors have no zeros
there, and their orders at $1$ combine to
$-\dim V^{G_K}=0$.  This does not require Artin holomorphy on the boundary
line.  Under $A_D$ the same argument applies to sums of nonprincipal
Dirichlet characters over $\mathbb Q$.

The traces $\operatorname{tr}(A_v^k)$ are the local Adams-power characters;
the proof separates the first prime term from all higher Adams terms,
whose series converge in the larger half-plane.  This is a concrete
$\lambda$-ring link, rather than a claim that positivity of coefficients
follows from the formal operations.  For $c=1/m$ the phenomenon is closely
related to Gun--Hazra--Sahu, \emph{On holomorphy and non-vanishing of Artin
L-functions} (2025), Theorem 5.  The displayed formula covers all complex
$|c|\le1$ and makes its hypotheses explicit.

\subsection*{MT-L02. One irrational parameter detects the zero-free half-plane}
Fix an irrational real number $0<c<1$ and a nontrivial irreducible Artin
representation $\rho$.  Then the following are equivalent:
\begin{enumerate}
 \item $L(s,\rho)$ is holomorphic and nonzero on $H$;
 \item the Euler product $D_c(s,\rho)$ admits a single-valued meromorphic
       continuation to $H$.
\end{enumerate}
In the second assertion meromorphic is a global, single-valued condition;
a branch on a slit domain is not enough.

MT-L01 proves the forward implication.  Conversely $L$ is meromorphic by
Brauer induction and $E_c$ is holomorphic on $H$ by the same normally
convergent series without any hypothesis about the zeros of $L$.
Logarithmic differentiation in $\operatorname{Re}s>1$ gives
\[
 \frac{D_c'}{D_c}=c\frac{L'}L+E_c'.
\]
Both sides are meromorphic on connected $H$ if the asserted continuation
of $D_c$ exists, so the identity holds on $H$.  Taking residues at any
$s_0\in H$ gives
\[
 \operatorname{ord}_{s_0}D_c
      =c\operatorname{ord}_{s_0}L.
\]
Both orders are integers, and $c$ is irrational.  Both must be zero.
Thus $L$ is holomorphic and nonvanishing on $H$, proving the equivalence.
If instead $c=a/b$ in lowest terms, the same equality proves
$b\mid\operatorname{ord}_{s_0}L$ everywhere in $H$.

GRH in the meromorphic-divisor sense for the nontrivial irreducible Artin
factors is therefore equivalent, by the functional equation, to these
irrational-parameter continuation tests.  To include the trivial factor one
must separately add RH for $\zeta_K$; that input is already part of $M_H$.
This is an analytic multiplicity test.  It neither constructs the
continuation from arithmetic lattices nor removes poles on the critical
line.  With a trivial summand of multiplicity $r$, the known pole at $1$
contributes local monodromy $e^{-2\pi i c r}$; this is why the no-invariants
hypothesis was retained rather than suppressed.

\subsection*{MT-ART05. The exact induced-linear cone for \texorpdfstring{$A_5$}{A5}}
The coarse counterexample MT-ART04 can be sharpened to a degree-three
virtual character. More usefully, the same calculation gives explicit
lower bounds for any pole of an actual $A_5$ Artin function. These are
group-theoretic deductions, not realizations of the virtual examples as
arithmetic divisor orders.

Write the irreducible characters in the order
$1,\chi_3,\chi'_3,\chi_4,\chi_5$, and put
\[
 \Theta=a+b\chi_3+c\chi'_3+d\chi_4+e\chi_5,
 \qquad a,b,c,d,e\in\mathbb Z.
\]
Then the inequalities
$\langle\Theta,\operatorname{Ind}_H^{A_5}\eta\rangle\ge0$ for
every subgroup $H$ and every linear character $\eta$ are equivalent to
the following seven inequalities:
\[
 a\ge0,\quad e\ge0,\quad a+d\ge0,\quad b+c\ge0,\quad
 b+d+e\ge0,\quad c+d+e\ge0,\quad b+c+d\ge0.       \tag{MT-A5-cone}
\]
If $r=\Theta(1)=a+3b+3c+4d+5e$, then $d<0$ forces $r\ge3$,
and $b<0$ or $c<0$ forces $r\ge4$. These thresholds are sharp for
this abstract cone: respectively
\[
       1+\chi_3+\chi'_3-\chi_4,
       \qquad -\chi_3+\chi'_3+\chi_4               \tag{MT-A5-ex}
\]
have degrees three and four and satisfy every inequality.

\begin{proof}
We give the character and subgroup details so the finite assertion has
no unreported character-table oracle. The five conjugacy classes have
orders $1,2,3,5,5$ and sizes $1,15,20,12,12$. The natural four-dimensional
augmentation character has values $(4,0,1,-1,-1)$ and norm one, as in
MT-ART04. Its real representation is oriented because all permutations
are even. The Hodge-star decomposition of its exterior square has two
invariant real three-dimensional spaces. The exterior-square character
has values $(6,-2,0,1,1)$ by
\[
 \chi_{\wedge^2 V}(g)=\tfrac12(\chi_V(g)^2-\chi_V(g^2)).
\]
Its norm squared is two and it has no invariant vectors, so the two
three-dimensional summands are distinct irreducibles. Simplicity of
$A_5$ makes each faithful and its determinant trivial. An order-two or
order-three element in either corresponding $\mathrm{SO}_3$ representation
has trace $-1$ or $0$. At order five the two traces are
$\varphi=(1+\sqrt5)/2$ and $\bar\varphi=(1-\sqrt5)/2$; their sum is
one and each is the trace of a nontrivial order-five rotation.

There are six Sylow-five subgroups. Their normalizers have order ten
and are dihedral, since the centralizer of a five-cycle is its own
cyclic subgroup. The permutation action on the six subgroups has
values $(6,2,0,1,1)$: an order-five element fixes its unique Sylow
subgroup, an order-three element fixes none, and double counting gives
$6\cdot5/15=2$ fixed subgroups for an involution. Removing the constants
gives an irreducible character $\chi_5=(5,1,-1,0,0)$, of norm one.
The squared dimensions $1+9+9+16+25=60$ exhaust the character table.

Every proper subgroup is contained in an $A_4$, a $D_5$ of order ten,
or an $S_3$ of order six. Here is an elementary verification. Simplicity
excludes proper indices two, three and four. The remaining possible
proper subgroup orders divide sixty and belong to
$\{1,2,3,4,5,6,10,12\}$. There are no elements of order four, six or
ten. Thus subgroups of orders four, six and ten have types
$V_4,S_3,D_5$. The centralizer of an involution is its unique $V_4$,
whose normalizer is $A_4$. In a subgroup of order twelve, a normal
Sylow-three subgroup is impossible since the ambient normalizer of
$C_3$ has order six. There are therefore four Sylow-three subgroups,
contributing eight nonidentity elements. Three Sylow-two subgroups
would contribute nine further involutions, because distinct $V_4$'s
have disjoint nonidentity sets. Hence the Sylow-two subgroup is normal
and the order-twelve subgroup is its $A_4$ normalizer. The smaller
cyclic and Klein subgroups are contained in these listed groups.

Let $u,\bar u$ be the nontrivial linear characters of $A_4$,
$v$ its irreducible character of degree three, and let $\epsilon$
denote a sign character of either $D_5$ or $S_3$. The restriction
decompositions obtained from the displayed character values are
\[
\begin{array}{c|ccc}
 & A_4 & D_5 & S_3\\\hline
1&1&1&1\\
\chi_3&v&\epsilon+v_1&\epsilon+v_2\\
\chi'_3&v&\epsilon+v'_1&\epsilon+v_2\\
\chi_4&1+v&v_1+v'_1&1+\epsilon+v_2\\
\chi_5&u+\bar u+v&1+v_1+v'_1&1+2v_2
\end{array}                                                   \tag{MT-A5-res}
\]
Here $v_1,v'_1$ are the two dihedral characters of degree two and
$v_2$ is the standard $S_3$ character. The nonlinear characters in
these three groups are themselves induced from linear characters:
$v$ from a nontrivial character of $V_4$, $v_1,v'_1$ from $C_5$,
and $v_2$ from $C_3$. Transitivity of induction and Frobenius reciprocity
therefore show that all induced-linear tests imply nonnegative
multiplicities in every column of (MT-A5-res). Together with the
global trivial-character test $a\ge0$, these are exactly the seven
inequalities (MT-A5-cone), after deleting redundant inequalities.
Conversely those seven inequalities make every column an actual
character. Restricting further gives actual characters on all proper
subgroups, hence every required linear inner product is nonnegative.
The global linear test is again $a\ge0$. This proves the equivalence.

For the degree bounds, first put $d=-k<0$. Then $a\ge k$,
$b+c\ge k$, and $b,c\ge k-e$. If $e=0$, these imply
$r\ge k+6k-4k=3k\ge3$. If $e\ge1$, the first two inequalities give
$r\ge k+3k-4k+5e\ge5$. Next put $b=-k<0$. If $d\ge0$, we have
$c\ge k$, $d+e\ge k$, whence
$r\ge4d+5e\ge4k$. If $d=-j<0$, then
$a\ge j$, $c\ge k+j$, and $e\ge k+j$, giving
$r\ge j+3j-4j+5(k+j)\ge4k$. The argument for $c<0$ is identical.
Finally direct substitution verifies both examples (MT-A5-ex).
\end{proof}

\subsection*{MT-ART06. Pole multiplicity gates for actual \texorpdfstring{$A_5$}{A5} extensions}
Let $L/K$ be an $A_5$ extension and $s_0\ne1$. Its Heilbronn character
has coefficients equal to the local Artin orders. Hecke holomorphy and
induction give all the tests in MT-ART05, unconditionally. Hence every
Artin factor is holomorphic at a zero of $\zeta_L$ of order at most two.
A pole of the four-dimensional factor requires order at least three
and a zero of $\zeta_K$ at the same point. A pole of either
three-dimensional factor requires order at least four. Conditional on
$M_H$, all remaining nontrivial candidate poles lie on the critical line.

At the central point the lower bounds strengthen to six and eight,
respectively. To prove this final refinement, the permutation actions
of degrees five and six give
\[
 L(s,\chi_4)=\frac{\zeta_{L^{A_4}}(s)}{\zeta_K(s)},\qquad
 L(s,\chi_5)=\frac{\zeta_{L^{D_5}}(s)}{\zeta_K(s)}.
\]
The completed Dedekind zeta functional equation has sign $+1$, and
its completion factors are finite and nonzero at $1/2$. Consequently
the central orders $a,d,e$ are even integers, with $a,e\ge0$ by
(MT-A5-cone). If $d=-k<0$, then $k\ge2$; the preceding proof gives
$r\ge3k\ge6$ when $e=0$ and $r\ge5e\ge10$ otherwise.
If $b=-k<0$ and $d\ge0$, the nonnegative even integers $d,e$
satisfy $d+e\ge k\ge1$, so $r\ge4d+5e\ge8$.
If $d=-j<0$, then $j\ge2$ and the preceding bound
$r\ge5(k+j)$ is stronger. The same holds for $c<0$.
This proves all assertions; no assumption on simplicity of zeros has
been introduced. The degree-three and degree-four virtual examples
only show the limits of the indicated character tests away from the
additional central parity constraints.

These refinements put the residual Artin question into explicit local
inequalities on a finite Galois quotient, within the SGA 1 V \S7
framework. They sharpen ORDER item 42 and ZETA\_TWO ZT02 without
claiming general Artin holomorphy or constructing actual negative
Artin divisor orders.

\subsection*{MT-H01. A unitary Hecke family outside the algebraic-CM normalization}
Let $K=\mathbb Q(\sqrt5)$, $\epsilon=(1+\sqrt5)/2$, and let $\sigma$
be its nontrivial automorphism. Put $t_0=\pi/\log\epsilon$.
For each nonzero integer $k$ there is a conductor-one unitary Hecke
character defined on fractional ideals by
\[
       \eta_k((\alpha))=
       \left|\frac{\alpha}{\sigma(\alpha)}\right|^{ik t_0}.
                                                               \tag{MT-H1}
\]
It has infinite order, $\eta_k^\sigma=\eta_{-k}$, and is trivial
on the rational ideles. No positive power of $\eta_k$ becomes an
algebraic Hecke character after multiplying by a complex norm power
and a finite-order character. Thus the full unitary-Hecke input $M_H$
has an explicit scope beyond the algebraic Hecke characters used in
the potentially CM motive transfer.

\begin{proof}
The ring of integers is $\mathbb Z[\epsilon]$, with discriminant five.
Minkowski's ideal-class bound is $\sqrt5/2<2$, so every fractional
ideal is principal. An exact classical locator for the bound is
Milne, \emph{Algebraic Number Theory}, Theorem 4.3 (v3.08), PDF p.~72.
The units are $\{\pm\epsilon^j:j\in\mathbb Z\}$. For a direct
verification, a positive unit $u$ with $1<u<\epsilon$ would have
integer trace $u+u^{-1}$ in $(2,\sqrt5)$ if its norm were one, or
$u-u^{-1}$ in $(0,1)$ if its norm were minus one; neither interval
contains an integer. Multiplying an arbitrary positive unit by a
power of $\epsilon$ reduces it to $[1,\epsilon)$, so it is a power
of $\epsilon$, and signs give the assertion.

Multiplying $\alpha$ by a unit changes the logarithm of the absolute
ratio in (MT-H1) by $2j\log\epsilon$. Its phase therefore changes by
$2\pi i k j$, so (MT-H1) is well-defined and multiplicative. It is
unitary. The classical ideal/infinity-type description of Hecke
characters makes it an unramified finite-place character, with
archimedean exponents $\pm ikt_0$ (their simultaneous sign depends
on the ideal-to-idele convention). The unit compatibility just checked
is exactly what is required for this description. Rational ideles
have equal archimedean components and equal valuations at conjugate
places; their images are trivial. For a ramified rational place this
also follows here from the generator $\sqrt5$, whose absolute ratio
is one. Conjugating (MT-H1) inverts the ratio.

On the positive archimedean subgroup the exponents of any algebraic
Hecke character are integers, up to an overall convention sign.
A norm power adds the same complex number to the two exponents;
a finite-order character is trivial on this connected subgroup.
If a positive power $m$ of (MT-H1) had the asserted algebraic
normalization, subtraction of its two exponents would give
$2mikt_0$ equal to the difference of two integers. This is impossible.
The same nonconstant archimedean restriction proves infinite order.
\end{proof}

This does not assert a transcendence theorem for individual Hecke
values, or rule out every conceivable generalized motive. It precisely
separates the character class covered by CM01 from the broader unitary
character class quantified in Masullo's endpoint. A proposed transport
of the CM theta-value and algebraic-height steps must therefore address
this extra archimedean data; norm twisting alone cannot do so.

\subsection*{MT-H02. An explicit nonholomorphic \texorpdfstring{$\mathrm{GL}_2$}{GL2} family}
For $k\ge1$, the classical Hecke--Maa\ss\ correspondence associates to
$\eta_k$ a cuspidal dihedral Maa\ss\ newform $g_k$ of weight zero,
level five and nebentypus $\chi_5$. Its spectral parameter is $kt_0$,
its Laplace eigenvalue is $1/4+k^2t_0^2$, and
\[
       L(s,g_k)=L_K(s,\eta_k).
                                                               \tag{MT-H2}
\]
Consequently $M_H$ gives GRH for every member of this family. Its
entireness is already the classical Hecke theorem. The exact primary
locator for the correspondence, level, spectral parameter and coefficient
identity is Humphries--Khan, \emph{On the Random Wave Conjecture for
Dihedral Maa\ss\ Forms}, GAFA 30 (2020), Appendix A.1, printed pp.~106--107 (PDF pp.~73--74), the displayed
formula for $\psi((\alpha))$ and the subsequent formulas for
$t_g$ and $\lambda_{g_\psi}(n)$.

For clarity, the transfer proof is short and complete. MT-H01 gives
a non-norm character: its two archimedean exponents are unequal.
The cited correspondence is therefore cuspidal, and its formula
$\lambda_{g_k}(n)=\sum_{N\mathfrak a=n}\eta_k(\mathfrak a)$ identifies
the absolutely convergent Dirichlet series with (MT-H2), including all
finite primes. Analytic continuation gives the equality everywhere.
The character is unitary, so its nontrivial critical line is
$\operatorname{Re}s=1/2$, which is exactly the assertion supplied by
$M_H$. This proves the stated implication without a generic
non-dihedral Maa\ss\-form hypothesis.

These forms live on the analytic modular curve of level five and in
the automorphic-induction part of the adelic framework. They are an
explicit additional $\mathrm{GL}_2$ branch beside the holomorphic
CM and finite-image weight-one branches, not the missing general
non-CM $\mathrm{GL}_2$ theorem in ORDER item 42.

\subsection*{MT-H03. Exact tensor and symmetric-power products in that family}
Write $F_j(s)=L_K(s,\eta_j)$ for $j\ne0$ and
$F_0(s)=\zeta_K(s)=\zeta(s)L(s,\chi_5)$. Then $F_{-j}=F_j$.
For the actual local two-dimensional parameters $V_j$ of $g_j$, the
following finite Euler-product identities hold:
\[
 L(s,V_m\otimes V_n)=F_{m+n}(s)F_{m-n}(s),\qquad m,n\ge1,
                                                               \tag{MT-H3}
\]
and, for $n\ge0$ and $k\ge1$,
\[
 L(s,\operatorname{Sym}^n V_k)=
 \prod_{j=0}^{\lfloor(n-1)/2\rfloor}F_{(n-2j)k}(s)
 \begin{cases}\zeta(s),&n\text{ even},\\1,&n\text{ odd}.
 \end{cases}                                                   \tag{MT-H4}
\]
The product is empty when $n=0$. Conditional on $M_H$, all nontrivial
zeros of these products lie on $\operatorname{Re}s=1/2$. They are
entire except for a simple pole at one when $m=n$ in (MT-H3), or
when $n$ is even in (MT-H4). This is a covered arbitrarily large
degree family, not a use of an unknown general symmetric-power lift.

\begin{proof}
Conjugating ideals proves $F_{-j}=F_j$. At a split rational prime
the two local parameters of $V_k$ are $z^k,z^{-k}$, where
$z=\eta_1(\mathfrak p)$. The tensor weights are
$z^{m+n},z^{m-n},z^{-m+n},z^{-m-n}$, which gives (MT-H3).
The symmetric-power weights are $z^{(n-2j)k}$ for $0\le j\le n$;
pairing opposite weights leaves one middle weight $1$ precisely for
even $n$, proving (MT-H4) there.

At an inert prime the unique prime ideal is $(p)$ and has character
value one. The two induced Frobenius eigenvalues are $1,-1$ for
every nonzero $k$. Tensoring gives two of each sign. On the
symmetric power, the even-index monomials give eigenvalue $1$ and
the odd-index monomials give $-1$. If $n=2r$, these multiplicities
are $r+1,r$; if $n=2r+1$, they are $r+1,r+1$. These are exactly
the factors on the right sides. If $m=n$, the factor $F_0$ at an
inert prime is $(1-p^{-2s})^{-1}$, so the same check includes that
case.

The only ramified prime is five. The ideal $(\sqrt5)$ has
$\eta_k$-value one and the character is unramified at every finite
place. Thus its local character at that quadratic completion is
trivial, and its induced Weil parameter is $1\oplus\chi_{5,5}$.
The inertia-fixed dimension is one. A tensor product has two
trivial constituents and two $\chi_{5,5}$ constituents; a symmetric
power has respectively $r+1,r$ or $r+1,r+1$ of them. The latter
constituents have no inertia invariants. Hence the ramified Euler
factors also agree exactly, rather than only away from five.
Equivalently the proof is the weight-pair decomposition for diagonal
matrices $\operatorname{diag}(z,z^{-1})$ and an outside coset
acting by interchanging the two basis vectors.

Absolute convergence initially gives the identities for
$\operatorname{Re}s>1$. The right sides give their meromorphic
continuations. All $F_j$, $j\ne0$, are entire Hecke functions with
nonzero boundary value at one; the latter is Ramakrishnan--Valenza,
Theorem 7-28, printed p.~289. The same holds for $L(1,\chi_5)$.
Thus the displayed zeta factors give exactly the asserted simple
poles, with no cancellation at one. Finally $M_H$ locates all the
nontrivial zeros of each unitary Hecke factor. The identities contain
no negative exponents, so they introduce no hidden poles on the
critical line. This proves every assertion.
\end{proof}

The middle summand in (MT-H4) is the trivial character for this
specific family. In a differently normalized dihedral family it
depends on the central character; replacing it universally by either
$1$ or $\chi_K$ would be incorrect. This is the representation-ring
and lambda-operation part of the adelic framework, beside SGA 6 V
\S\S2,7, with the explicit analytic parameters of MT-H01 retained.

\subsection*{MT-ART07. The entire non-real $A_5$ Artin category}
Let $L/\mathbb Q$ be Galois with group $A_5$, and suppose $L$ is not
totally real. Every nontrivial irreducible Artin L-function for this
extension is entire, unconditionally. Under $M_H$ each satisfies GRH.
This closes the non-real $A_5$ specialization of the holomorphy gate;
MT-ART05/06 remain useful for totally real extensions and other base fields.

The holomorphy assertion is classical, not a new finding. An exact primary
locator is Calegari, \emph{The Artin conjecture for some $S_5$-extensions},
\S1, publisher-author PDF p.~2, the paragraph immediately before
Theorem 1.1: it states the stronger automorphy assertion for $A_5$ when
complex conjugation is nontrivial, citing Khare--Wintenberger and Kim's
icosahedral functorial constructions. The two conditions are equivalent
here because a Galois extension has a real embedding exactly when every
embedding is real, equivalently complex conjugation acts trivially.
We cite that classical result rather than repeat its proof.

Here is the complete GRH transfer and its tensor-category form.
For every finite-dimensional representation $V$ of $A_5$, put
$m=\dim V^{A_5}$. Complete reducibility gives
\[
 L(s,V)=\zeta(s)^m
          \prod_{\chi\ne1}L(s,\chi)^{\langle V,\chi\rangle},
 \qquad \langle V,\chi\rangle\in\mathbb Z_{\ge0}.
\]
The product after removing $\zeta^m$ is entire by the cited theorem.
For each irreducible factor, Brauer induction expresses its L-function
as a finite product of integer powers of finite-order Hecke L-functions.
At a point of the open critical strip off its middle line, $M_H$ makes
every such factor finite and nonzero. The product is consequently nonzero
there, including when some Brauer exponents are negative. This proves
the claimed GRH without making a positivity assumption on those exponents.

At $s=1$ every nontrivial irreducible factor is also nonzero. Indeed in a
Brauer expression $\chi=\sum_j a_j\operatorname{Ind}_{H_j}^{A_5}\eta_j$,
the order there is $-\sum_{\eta_j=1}a_j$, by the simple Dedekind pole and
Hecke boundary nonvanishing. Frobenius reciprocity identifies the sum with
$\langle\chi,1\rangle=0$. Thus $L(s,V)$ has a pole of order exactly $m$
at $1$, and has no other poles. Tensor products, duals, exterior and
symmetric powers, and every Schur functor produce representations of the
same finite group, so the same conclusions hold for them after removing
their explicitly determined invariant summands. This is a statement about
the whole Artin tensor category attached to this cover, not about a
general positive-weight motive. The companions in Grothendieck's language
are SGA~1, Expos\'e~V, \S7 (finite monodromy), and SGA~6, Expos\'e~V,
\S\S2,7 (representation-ring operations), together with AG-LAMBDA.

\subsection*{MT-ART08. A finite exceptional set is enough for a holomorphy transfer}
Here is a useful precise comparison lemma. Let $V$ be a self-dual,
finite-image complex Galois representation over a number field $K$.
Suppose an Euler product $F(s)$ converges in a right half-plane and
continues to an entire function, with the same finite local factors as $L(s,V)$ outside a
finite set $S$. Suppose also that, for each $v\in S$, the reciprocal of
the local factor of $F$ is a polynomial $Q_v((Nv)^{-s})$.
Then $L(s,V)$ is entire. No equality of ramified local parameters is
needed for this conclusion.

To prove this, write the Artin factor at $v$ as
$P_v((Nv)^{-s})^{-1}$. Every eigenvalue of Frobenius on the inertia
invariants is a root of unity, since $V$ has finite image. Hence the zeros
of $P_v((Nv)^{-s})$ all lie on $\operatorname{Re}s=0$.
Equality in a right half-plane, followed by meromorphic continuation,
gives
\[
 L(s,V)=F(s)\prod_{v\in S}
             \frac{Q_v((Nv)^{-s})}{P_v((Nv)^{-s})}.
\]
The right side is holomorphic on $\operatorname{Re}s>0$.
The classical completed Artin functional equation is
\[
 A^{s/2}\Gamma_V(s)L(s,V)
   =W A^{(1-s)/2}\Gamma_V(1-s)L(1-s,V),
\]
with $A>0$ and $W\ne0$; the gamma factor is a product of
$\Gamma_{\mathbb R}(s)$, $\Gamma_{\mathbb R}(s+1)$ and
$\Gamma_{\mathbb C}(s)$ with nonnegative integral multiplicities.
An exact locator for these formulas is Prasad--Yogananda,
\emph{A Report on Artin's holomorphy conjecture}, \S3.3,
PDF pp.~8--9, with Brauer induction in Theorem 3.3.
For $\operatorname{Re}s<1$, the numerator
$\Gamma_V(1-s)L(1-s,V)$ is holomorphic, since
$\operatorname{Re}(1-s)>0$. The reciprocal $1/\Gamma_V(s)$ is entire.
The functional equation therefore makes $L(s,V)$ holomorphic on this
second half-plane too. The two half-planes cover $\mathbb C$, proving
the assertion. For a non-self-dual $V$ the identical argument works if
the hypotheses are also supplied for $V^\vee$.

The finiteness of $S$ is an actual hypothesis. Agreement at a set of
places of density one does not give the finite product in the proof.
In particular this lemma cannot by itself upgrade the weak
five-dimensional correspondence in Calegari's Theorem 1.1.
In geometric terms this separates modification at finitely many closed
points of an arithmetic base from a density statement about infinitely
many Frobenius fibres.

\subsection*{MT-S501 and MT-S502. The $S_5$ transfer and its one-value gate}
Let $L/\mathbb Q$ be an $S_5$ extension, with complex conjugation of cycle
type $(12)(34)$. Assume $5$ is unramified and its Frobenius also has this
cycle type. Then the irreducible factors of dimensions $4$ and $6$ are
entire and, under $M_H$, satisfy GRH (MT-S501).

Let $E=L^{A_5}$, let $F\subset L$ have degree $6$ over $\mathbb Q$, and
put $H=EF$. Under $M_H$, the additional single condition
\[
                         \zeta_H(1/2)\ne0
\]
implies that every nontrivial irreducible factor is entire and satisfies
GRH (MT-S502). This is a proved reduction with the displayed extra input,
not a proof of that central nonvanishing.

The classical inputs are Calegari's Theorems 1.1 and 1.2,
publisher-author PDF pp.~2--3; the definition $H=EF$ is explicit in
Theorem 1.2 of arXiv:1112.1152v1, p.~2. The former supplies the
dimension-$4$ and dimension-$6$ holomorphy, and the latter supplies the
remaining case if $\zeta_H$ has no zero in the real interval $(0,1)$.
Under $M_H$ the only possible real zero in that interval is $1/2$.
Thus its interval condition is exactly equivalent to the displayed
single-value condition. Once holomorphy is available, the Brauer
zero-location argument of MT-ART07 proves GRH for every factor.
This is the full proof of both transfers. The one-dimensional factors
are Dirichlet functions and already fall under $A_D$.

The totally real quadratic field $E$ is the base for the odd projective
two-dimensional representations used in the cited theorem. Accordingly
this is a higher-rank consequence of a specific automorphic descent
theorem. It is not obtained by treating every $S_5$ factor as a single
Hecke character. In the Artin-motive programme, the remaining gate in
this family is one central Dedekind value. AG-SIGN controls a nonzero
leading coefficient after a central zero is removed; it does not prove
that the central order is zero.

\subsection*{MT-S503. One unconditional holomorphy example becomes a full GRH example}
Take the splitting field $L$ of
\[
                        x^5-x^3-x^2+x+1.
\]
Calegari, Theorem 5.1, publisher-author PDF p.~15, proves automorphy of
every representation of $\operatorname{Gal}(L/\mathbb Q)=S_5$;
the quintic field has discriminant $1609$. Consequently under $M_H$
every nontrivial irreducible Artin L-function of this explicit field
satisfies GRH, by the same signed Brauer argument. This example uses
that cited theorem, not the local-at-$5$ hypotheses of MT-S501, which
the source explicitly says this polynomial does not satisfy.
The historical bounded-height computation in Theorem 5.3 is part of
the classical proof's evidence; it is not a computation performed in
this run. Geometrically this supplies a concrete finite \(S_5\)-cover
of the arithmetic base whose complete Artin category is covered.

\subsection*{MT-M05. The stated Schauder quotient loses more than its zero evaluations see}
Lemma 9.3.2.5 on original PDF pp.~425--426 asserts a basis for the
quotient by the actual convolution ideal $\varphi*\mathcal S_E(G)$,
with coefficient functionals the Mellin jets at the zeros of
$\mathcal M\varphi$. We inspected the original p.~426 image: there
is no closure over this ideal. The assertion is false even when the
Mellin divisor consists of one simple zero. Thus the counterexample
does not rely on the zero-count defect in MT-M01.

Use $G=\mathbb R_{>0}$, trivial character and $u=\log t$; convolution
becomes ordinary additive convolution. Write
\[
 g_a(u)=(4\pi a)^{-1/2}e^{-u^2/(4a)},\qquad a>0.
\]
Choose $0<b<a$ and set $\varphi(e^u)=-g_a'(u)$,
$h(e^u)=-g_b'(u)$. Both functions belong to the exponential-Schwartz
space, and
\[
 \mathcal M\varphi(s)=s e^{as^2},\qquad
 \mathcal Mh(s)=s e^{bs^2}.
\]
The first function has just the simple zero $0$, lying in the interior
of the strip for $I=[-1,1]$. The only claimed coefficient functional is
$\delta_0$, and $\delta_0(h)=0$.
But $h\notin\varphi*\mathcal S_E$: if $h=\varphi*f$, multiplicativity
of the absolutely convergent Mellin transform gives
$\mathcal Mf(s)=e^{(b-a)s^2}$ for $s\ne0$, and then everywhere by
analyticity. On the imaginary axis this has modulus
$e^{(a-b)\tau^2}$. In contrast, $f(e^u)$ is integrable and therefore
$|\mathcal Mf(i\tau)|\le\int|f(e^u)|\,du$, a contradiction.
Thus the nonzero quotient class of $h$ has every asserted coefficient
zero, which is incompatible with uniqueness in a Schauder expansion.
The still simpler choice $\varphi=g_a$, $h=g_b$ has no Mellin zeros
and a nonzero quotient, so even the empty-divisor case fails.

\subsection*{MT-M06. The closed-ideal repair for every polynomial Gaussian divisor}
In logarithmic coordinates put
\[
 \mathcal E=\{f\in C^\infty(\mathbb R,\mathbb C):
   p_{N,j}(f):=\sup_x e^{N|x|}|f^{(j)}(x)|<\infty
   \text{ for every integer }N,j\ge0\}.
\]
This is $\mathcal S_E(\mathbb R_{>0})$ with its exponential-Schwartz
Fr\'echet topology. Polynomial factors in the usual Schwartz seminorms
are absorbed by increasing $N$. Let
\[
 P(s)=\prod_{\nu=1}^r(s-\rho_\nu)^{m_\nu},\quad
 d=\sum_\nu m_\nu,\qquad
 \varphi=P(-D)g_a,\quad D=d/dx.
\]
Distinct $\rho_\nu\in\mathbb C$, positive integral $m_\nu$, and
$a>0$ are arbitrary; $P=1,d=0$ is allowed. Then
\[
 \overline{\varphi*\mathcal E}
   =J_P:=\{h\in\mathcal E:
          (\mathcal Mh)^{(j)}(\rho_\nu)=0,
          \ 0\le j<m_\nu\},                         \tag{MT-M6}
\]
and the quotient by this \emph{closed} ideal is topologically
$\mathbb C^d$, with exactly the displayed Mellin jets as coordinates.
In particular its coordinate basis is absolute. This repairs the
finite-divisor case after making the closure operation explicit.

\begin{proof}
First, $g_a*\mathcal E$ is dense in $\mathcal E$. We give the full
argument, since density is distinct from the false surjectivity in
MT-M05. Convolution is continuous on $\mathcal E$: the inequality
$e^{N|x|}\le e^{N|x-y|}e^{N|y|}$ bounds each differentiated
convolution by an exponential supremum seminorm times an exponential
$L^1$ seminorm; the latter is bounded by a constant times
$p_{N+1,0}$.

Suppose a continuous linear functional $T$ vanishes on
$g_a*\mathcal E$. For every real $y$ and $\epsilon>0$ the function
$g_{a+\epsilon}(x-y)=g_a*g_\epsilon(\,\cdot-y)$ lies in this range.
It converges to $g_a(x-y)$ in every $p_{N,j}$ as
$\epsilon\downarrow0$, by uniform Gaussian estimates with the
parameter in a compact subinterval of $(0,\infty)$. Hence
$T(g_a(\,\cdot-y))=0$. The function
\[
 B(z)=T\bigl(e^{-x^2/(4a)}e^{zx/(2a)}\bigr)
\]
is entire: on every compact set of $z$, all its proposed complex
derivatives converge in every $p_{N,j}$ by Gaussian domination.
Completing the square gives $B(y)=0$ for real $y$, so the identity
theorem gives $B=0$.

For any Schwartz function $\psi$, Fourier inversion now gives
\[
 T(e^{-x^2/(4a)}\psi(x))
  =\frac1{2\pi}\int_{\mathbb R}\widehat\psi(\xi)
        B(2ai\xi)\,d\xi=0.
\]
This passage under $T$ is justified in the Fr\'echet space itself:
the $p_{N,j}$ seminorm of
$e^{-x^2/(4a)}e^{i\xi x}$ is at most
$C_{N,j,a}(1+|\xi|)^j$, whereas $\widehat\psi$ is rapidly
decreasing. For $h\in C_c^\infty$, take
$\psi=e^{x^2/(4a)}h$, which is again compactly supported and smooth.
Then $T(h)=0$. Finally $C_c^\infty$ is dense in $\mathcal E$:
multiply by a smooth cutoff equal to one on $[-R,R]$ and supported
in $[-2R,2R]$; the product rule and the bound from $p_{N+1,j}$ make
every $p_{N,j}$ of the discarded tail tend to zero. Thus $T=0$.
Hahn--Banach separation proves the asserted density.

Second, the range of $P(-D)$ on $\mathcal E$ is exactly $J_P$.
Integration by parts gives $\mathcal M(-Df)(s)=s\mathcal Mf(s)$,
so inclusion in $J_P$ follows. For the converse, if
$\mathcal Mh(\rho)=0$, define
\[
 f(x)=e^{-\rho x}\int_{-\infty}^x e^{\rho t}h(t)\,dt
     =-e^{-\rho x}\int_x^\infty e^{\rho t}h(t)\,dt.
\]
The equality is precisely the vanishing of the Mellin value.
The first expression controls the negative tail and the second the
positive tail. For any prescribed $N$, using an exponential bound
on $h$ with exponent larger than $N+|\operatorname{Re}\rho|+1$
proves exponential decay of $f$ on both tails; the equation
$f'=h-\rho f$ then gives the same bounds for every derivative.
Therefore $f\in\mathcal E$, and $(D+\rho)f=h$.
Equivalently $-f$ solves $(-D-\rho)(-f)=h$, and its Mellin transform
is $\mathcal Mh(s)/(s-\rho)$. Repeating this division through all
roots, including multiplicities, proves $P(-D)\mathcal E=J_P$.

The Mellin jet functionals are continuous, so $J_P$ is closed.
Since $g_a*\mathcal E$ is dense and $P(-D)$ is continuous,
$P(-D)(g_a*\mathcal E)$ is dense in $P(-D)\mathcal E=J_P$.
Differentiation commutes with convolution, identifying that range
with $\varphi*\mathcal E$. This proves (MT-M6).

For $d=0$ the quotient is zero and no interpolation is required.
For $d>0$, it remains to verify that no coordinate is missing. Given arbitrary
$d$ target jets, ordinary polynomial Hermite interpolation constructs
a polynomial $Q$ of degree below $d$ for which
$Q(s)e^{bs^2}$ has exactly those jets, with any fixed $b>0$.
Indeed multiplication by the nowhere-zero entire function $e^{bs^2}$
is invertible on each finite jet space, so interpolate the resulting
jets of $Q$. The function $Q(-D)g_b$ belongs to $\mathcal E$ and
realizes them. Choosing the targets to be the coordinate vectors
gives a continuous linear section from the finite-dimensional space
$\mathbb C^d$ to $\mathcal E$. The jet map consequently induces
a continuous isomorphism of the quotient with $\mathbb C^d$,
with continuous inverse. Finite-dimensional coordinate expansions are
absolute, completing the proof.
\end{proof}

For $P=s$, MT-M05 therefore exhibits a proper, nonclosed convolution
ideal whose closure is $\ker\delta_0$. The closed quotient really
does have the one coordinate the draft intended. For $P=1$ the
principal ideal is dense but proper, and the closed quotient is zero.
This is a concrete repair rather than a certification of the general
infinite-divisor statement. For arithmetic Hecke divisors the remaining
gate is controlled infinite interpolation and the claimed absolute
basis/nuclear estimates on the correct separated quotient. A local
$O(\log T)$ zero count alone does not prove those estimates.
In cote~1 trace language, Hausdorff completion and the identification
of the spectral quotient precede taking its nuclear trace.

\subsection*{MT-M07. An infinite closed quotient with an absolute Mellin basis}
The closed-ideal repair also admits a complete infinite-divisor example.
Fix $a,\ell>0$ in the logarithmic space $\mathcal E$ of MT-M06, and put
\[
 \varphi(x)=g_a(x-\ell)-g_a(x),\qquad
 \rho_n=\frac{2\pi i n}{\ell}\quad(n\in\mathbb Z).
\]
Then $\mathcal M\varphi(s)=e^{as^2}(e^{\ell s}-1)$, with precisely
the simple zeros $\rho_n$. Let
\[
 (\mathcal P_\ell f)(t)=\sum_{k\in\mathbb Z}f(t+k\ell).
\]
The quotient $\mathcal E/\overline{\varphi*\mathcal E}$ is
topologically $C^\infty(\mathbb R/\ell\mathbb Z)$ under
$\mathcal P_\ell$, and has an absolute Schauder basis whose coefficient
functionals are exactly $f\mapsto\mathcal Mf(\rho_n)$.

\begin{proof}
The periodization series and every derivative converge uniformly on a
fundamental interval, by exponential decay. This gives a continuous map
to the indicated smooth periodic space. Choose
$\chi\in C_c^\infty(\mathbb R)$ nonnegative and positive on
$[-\ell,\ell]$, and set
$\theta=\chi/\mathcal P_\ell\chi$. The denominator is positive,
smooth and periodic, so $\theta$ is smooth, compactly supported and
$\mathcal P_\ell\theta=1$. The map
$R(F)(x)=\theta(x)F(x)$ is a continuous section of periodization,
proving surjectivity.

Let $\Delta_\ell h(x)=h(x-\ell)-h(x)$. Telescoping shows that its
range lies in $\ker\mathcal P_\ell$. Conversely, if
$\mathcal P_\ell f=0$, define
\[
 h(x)=\sum_{k\ge1}f(x+k\ell)
     =-\sum_{k\le0}f(x+k\ell).
\]
These locally normally convergent series define a smooth function.
For $x\ge0$ use the first expression and for $x\le0$ the second.
For every positive integer $N$ and every $j$, summing the geometric
series of bounds $p_{N,j}(f)e^{-N|x+k\ell|}$ proves
$p_{N,j}(h)<\infty$; the $N=0$ case follows from $N=1$.
Thus $h\in\mathcal E$, and direct index shifting gives
$\Delta_\ell h=f$. The range of $\Delta_\ell$ is therefore the
closed space $\ker\mathcal P_\ell$.

By MT-M06, $g_a*\mathcal E$ is dense in $\mathcal E$.
Continuity of $\Delta_\ell$ and commutation with convolution imply
\[
 \overline{\varphi*\mathcal E}
 =\overline{\Delta_\ell(g_a*\mathcal E)}
 =\ker\mathcal P_\ell.
\]
The continuous section $R$ identifies the quotient topology with
the usual smooth periodic topology, rather than merely giving a
vector-space identification.

The Fourier expansion in the convention used here is
\[
 \mathcal P_\ell f(t)=\frac1\ell
       \sum_{n\in\mathbb Z}\mathcal Mf(\rho_n)e^{-\rho_n t}.
\]
Its coefficients are obtained by integrating over a fundamental
interval and summing translates, justified by absolute convergence.
They decrease faster than every power of $|n|$: integrate by parts
arbitrarily many times and use the integrability of every derivative
of $f$. The Fourier series consequently converges absolutely in
every smooth seminorm and equals $\mathcal P_\ell f$ by Fourier
uniqueness. Representatives for the basis are
\[
                     b_n(x)=\ell^{-1}\theta(x)e^{-\rho_n x}.
\]
Indeed $\mathcal Mb_n(\rho_m)=\delta_{m,n}$, by
$\mathcal P_\ell\theta=1$ and integration on one period. Moreover,
$p_{N,j}(b_n)\le C_{N,j}(1+|n|)^j$ since $\theta$ has compact
support. Hence the expansion
$\sum_n\mathcal Mf(\rho_n)b_n$ is absolutely convergent even in
$\mathcal E$, and represents the class of $f$ in the quotient.
The jets determine that class uniquely by the Fourier identification.
Ordering the indices as $0,1,-1,2,-2,\ldots$ gives an absolute
Schauder basis, as asserted.
\end{proof}

The actual ideal still cannot replace its closure: with $0<b<a$,
$h=\Delta_\ell g_b$ has all these Mellin values zero but is not
$\varphi*f$ for any $f\in\mathcal E$, by exactly the Fourier-growth
division argument of MT-M05. Thus this example tests an infinite
divisor with controlled spacing, not just the empty or finite case.

\subsection*{MT-M08. The corresponding trace and positive convolution formula}
For $f\in\mathcal E$, convolution by $f$ descends to the quotient of
MT-M07. On $L^2(\mathbb R/\ell\mathbb Z,dt)$ it is trace class and
\[
 \operatorname{Tr}C_f
   =\sum_{n\in\mathbb Z}\mathcal Mf(\rho_n)
   =\ell\sum_{k\in\mathbb Z}f(k\ell).                 \tag{MT-M8}
\]
All sums converge absolutely. If
$f=h*\widetilde h$ with $h\in\mathcal E$ and
$\widetilde h(x)=\overline{h(-x)}$, this operator is positive and
\[
                  \operatorname{Tr}C_f
                   =\sum_n|\mathcal Mh(\rho_n)|^2\ge0.
\]

To prove the assertions, first use absolute convergence and changes of
variables to periodize a convolution:
\[
 \mathcal P_\ell(f*g)(t)
  =\int_0^\ell(\mathcal P_\ell f)(t-u)
                     (\mathcal P_\ell g)(u)\,du.
\]
This proves descent and identifies the smooth kernel
$K_f(t,u)=(\mathcal P_\ell f)(t-u)$.
The orthonormal Fourier vectors
$\ell^{-1/2}e^{-\rho_n t}$ are eigenvectors with eigenvalues
$\mathcal Mf(\rho_n)$. These eigenvalues decrease faster than every
power, as proved in MT-M07, so their absolute sum is finite.
The diagonal expansion is thus a trace-class expansion and gives
the first trace sum. Evaluating the absolutely convergent Fourier
series of $\mathcal P_\ell f$ at zero gives the second equality
in (MT-M8), including its factor $\ell$. Finally the Mellin transform
of $\widetilde h$ at a purely imaginary argument is the conjugate
of that of $h$; the eigenvalues are therefore the displayed squares.

There is also a direct nuclear realization on the smooth quotient.
Explicitly, put $c_n(F)=\ell^{-1}\int_0^\ell F(t)e^{\rho_nt}\,dt$
and $e_n(t)=e^{-\rho_nt}$. The $c_n$ are an equicontinuous family,
bounded by the uniform norm. In
\[
 C_f=\sum_n(1+|n|)^{-2}\,
       c_n\otimes\bigl((1+|n|)^2\mathcal Mf(\rho_n)e_n\bigr),
\]
the scalars are summable and the vector family is bounded in every
$C^\infty$ seminorm. This is a nuclear decomposition. Fourier partial
sums converge uniformly on bounded sets of $C^\infty$ (apply a
uniform higher-derivative bound to the coefficient tail), so the
finite-rank trace limit is the same absolute sum above.

This supplies a fully checked spectral-quotient, basis and trace
model beside cote~1's nuclear trace language. It does not identify
an arithmetic prime distribution with this lattice or prove the
corresponding basis theorem for the actual Hecke zero set. In that
case the exact infinite interpolation and operator estimates remain
to be supplied on the correct closed quotient.

\subsection*{MT-M09. Even a bounded local zero count does not give the asserted point basis}
Passing to the closed ideal is necessary, but does not prove the
general basis claim. There exists $\varphi\in\mathcal E$ whose Mellin
transform has only simple zeros on the imaginary axis and only
$O(1)$ zeros in every unit-height interval, for which
\[
        Q=\mathcal E/\overline{\varphi*\mathcal E}
\]
has no Schauder basis with the individual zero evaluations as its
coefficient functionals. This holds for every enumeration of the
zeros, and therefore also rules out an absolute Schauder basis of
that specified kind.

\begin{proof}
Put $\epsilon_n=e^{-n^4}$ for $n\ge1$ and define
\[
 H(s)=\prod_{n=1}^{\infty}
       \left(1+\frac{s^2}{n^2}\right)
       \left(1+\frac{s^2}{(n+\epsilon_n)^2}\right),
 \qquad F(s)=e^{as^2}H(s),\quad a>0.
\]
The product converges uniformly on compact sets because the sum of
the absolute deviations of its factors from one does. Away from the
listed zeros the tail logarithms converge locally uniformly, so
there are no additional zeros. Its zeros are exactly
$\pm in,\ \pm i(n+\epsilon_n)$, all simple. Since
$0<\epsilon_n<1/2$, each unit-height interval contains a bounded
number of them.

For $r=|s|$ the product satisfies
\[
 \log|H(s)|\le2\sum_{n\ge1}\log(1+r^2/n^2)\le2\pi r.
\]
For the last bound compare the decreasing function of $x$ with its
integral on $(0,\infty)$; integration by parts gives
$\int_0^\infty\log(1+r^2/x^2)\,dx=\pi r$.
Consequently, on every strip $|\sigma|\le A$,
\[
 |F(\sigma+it)|
 \le \exp\{aA^2-at^2+2\pi\sqrt{A^2+t^2}\}.
\]
Define
\[
 \varphi(x)=\frac1{2\pi}\int_{\mathbb R}F(it)e^{-itx}\,dt.
\]
The Gaussian bound permits every derivative in $x$ under the integral.
Shifting the contour to $\operatorname{Re}s=\sigma$ gives, for
every real $\sigma$ and every $j\ge0$,
\[
 \varphi^{(j)}(x)
   =\frac{e^{-\sigma x}}{2\pi}
      \int_{\mathbb R}(-\sigma-it)^jF(\sigma+it)e^{-itx}\,dt.
\]
The horizontal sides of the contour tend to zero by the displayed
Gaussian estimate, uniformly for $\sigma$ in the finite interval
used in that shift. Choosing $\sigma=N+1$ for $x\ge0$ and
$\sigma=-N-1$ for $x\le0$ proves every seminorm $p_{N,j}$ is finite.
Thus $\varphi\in\mathcal E$. Cauchy estimates applied on slightly
larger strips show that $F(it)$ and all its derivatives are Schwartz.
Fourier inversion gives $\mathcal M\varphi(it)=F(it)$;
both sides are entire, so $\mathcal M\varphi=F$ everywhere.

Every evaluation at a listed zero annihilates the convolution ideal
and its closure, and therefore descends continuously to $Q$.
Let $p=p_{1,0}$ and let
$\overline p([f])=\inf_{j\in\overline{\varphi*\mathcal E}}p(f+j)$
be the corresponding quotient seminorm. For any $f\in\mathcal E$,
\[
 |(\delta_{in}-\delta_{i(n+\epsilon_n)})(f)|
  \le\epsilon_n\int_{\mathbb R}|x||f(x)|\,dx
  \le2\epsilon_n p(f).
\]
The first inequality follows from
$|e^{i\epsilon_n x}-1|\le\epsilon_n|x|$.
Because the difference functional annihilates the closed ideal,
the same bound holds with $p$ replaced by $\overline p$ on $Q$.

Suppose the claimed Schauder basis existed, and denote the basis
vector associated to $\delta_{in}$ by $b_n$. Biorthogonality of
basis vectors and coefficient functionals gives
$\delta_{in}(b_n)=1$ and
$\delta_{i(n+\epsilon_n)}(b_n)=0$. Hence
\[
                 \overline p(b_n)\ge(2\epsilon_n)^{-1}.
\]
Take the single test function $g_b\in\mathcal E$ for any fixed
$b>0$. Its coefficient at $\delta_{in}$ is $e^{-bn^2}$.
The corresponding terms of its proposed basis expansion satisfy
\[
 \overline p\bigl(e^{-bn^2}b_n\bigr)
        \ge\frac12 e^{n^4-bn^2}\longrightarrow\infty.
\]
Terms of a convergent series in a locally convex space must tend
to zero. In every enumeration these infinitely many distinct terms
occur arbitrarily far along the series, giving a contradiction.
This proves the assertion.
\end{proof}

This is a counterexample for general exponential-Schwartz Mellin
divisors, not a claim that such clustered zeros occur for an actual
Hecke L-function. It pinpoints why NT-Z10/MT-M04 cannot by themselves
establish the arithmetic Schauder step: even an $O(1)$ local count
does not control the norm of individual interpolation vectors.
The next possible repair must incorporate quantitative separation,
or group nearby evaluations into divided-difference blocks and prove
convergence of the resulting block expansion. Neither property is
supplied by a zero-location hypothesis alone in the present argument.

% END root
\part{Unchecked claims and unresolved gates}
The complete proofs above establish only their stated implications.
The rows below are retained as research gates or as carried claims whose
complete source proofs were unavailable. They are not additional theorems.
The manuscript faults MT-M01 and MT-M02 have proved counterexamples;
the associated repairs do not certify every downstream operator or trace.

\paragraph{NT-Q06} Which pair-correlation and low-zero statistics are actually established?

\textbf{State:} proved/\allowbreak{}classical NT-Z01–\allowbreak{}NT-Z05 with support/\allowbreak{}moment restrictions; full GUE/\allowbreak{}correlations remain gate NT-Z06

\textbf{Source:} Montgomery1973 §1; HR2003 Thm3.\allowbreak{}1/\allowbreak{}5.\allowbreak{}1.

\paragraph{NT-Q07} Does rank-one Hecke GRH discharge the GL2 family hypothesis?

\textbf{State:} no such implication established; exact additional GL2-GRH gate NT-Z07 retained

\textbf{Source:} ILS2000 p56 and Thm1.\allowbreak{}1.

\paragraph{NT-Q10} Can carried class-group generator theorem be freshly source-inspected?

\textbf{State:} current gate:\allowbreak{} original AMS PDF access; inherited CX-N07 statement retained honestly

\textbf{Source:} Bach1990 Thm4 pp375–\allowbreak{}376.

\paragraph{NT-Q19} Which prime-race statements follow from A alone and which still require ordinate independence?

\textbf{State:} NT-R01/\allowbreak{}R02 proved/\allowbreak{}classical under A; NT-R03/\allowbreak{}R04 retain explicit additional GSH gate

\textbf{Source:} Rubinstein–\allowbreak{}Sarnak1994 Thm1.\allowbreak{}1/\allowbreak{}1.\allowbreak{}2/\allowbreak{}1.\allowbreak{}4/\allowbreak{}1.\allowbreak{}5.

\paragraph{AG-Q01} Which finite-type Z-scheme zeta functions have the adopted zero-support conclusion?

\textbf{State:} proved inherited exact Artin–\allowbreak{}Tate transfers; M adds identified potentiallyCM factors; residual nonCM continuation/\allowbreak{}RH and exceptional fibres remain explicit gates

\textbf{Source:} BRIDGES Q01; ZETA\_\allowbreak{}TWO ZT01–\allowbreak{}ZT07.

\paragraph{AG-Q08} Can the positive squared-variable determinant extend beyond \ensuremath{\zeta}?

\textbf{State:} proved AG-DET/\allowbreak{}MOMENT/\allowbreak{}HANKEL and AG-SIGNED finite pole witness; prime-defined cohomological realization and moment positivity for arbitrary Artin quotients remain gates

\textbf{Source:} ALPOGE\_\allowbreak{}BRIDGES BR-TR-01/\allowbreak{}02; SA23 geo-H-lidskii/\allowbreak{}geo-H-genus; cote1/\allowbreak{}79–\allowbreak{}80,111–\allowbreak{}112.

\paragraph{AG-Q11} Does finite-field purity imply openness of closure[pi0,H] in pi-prime under the original hypotheses?

\textbf{State:} literal ambientpi-prime refuted AG-C34-TYPE; intrinsic repair proved AG-C34-REPAIR/\allowbreak{}INDEX; with a specified section and finite derived length the raw action commutators are openH by AG-C34-SECTION; omitting solvability is refuted as an abstract inference by AG-C34-NONSOLV

\textbf{Source:} BRIDGES Q03; cote34/\allowbreak{}64 Rem1.\allowbreak{}17; definitions53–\allowbreak{}63.

\paragraph{AG-Q13} Do zero-location inputs construct integral Lefschetz correspondences?

\textbf{State:} current gate:\allowbreak{} cycle/\allowbreak{}rank comparison on X×X plus denominator control; no such construction follows from displayed analytic reductions

\textbf{Source:} Hodge sketch§§3.\allowbreak{}2–\allowbreak{}3.\allowbreak{}3; cote16/\allowbreak{}14–\allowbreak{}18 no7,44–\allowbreak{}50.

\paragraph{AG-Q14} Does normalized L-sign prove arbitrary-cycle arithmetic Hodge index?

\textbf{State:} current gate:\allowbreak{} exact height identity for the given cycles and covered L-function; conditional sign deduction AG-HEIGHT proved

\textbf{Source:} Hodge sketch§3.\allowbreak{}5; cote44 Theorems4–\allowbreak{}5.

\paragraph{AG-Q16} Which finite/\allowbreak{}profinite GT gates move?

\textbf{State:} finite quotient detection/\allowbreak{}generation effective AG-FROB,CX-GEO-004/\allowbreak{}008; full profinite kernel/\allowbreak{}associator gate remains

\textbf{Source:} ALPOGE\_\allowbreak{}BRIDGES R-GTHAT; SA22 GT lines.

\paragraph{AG-Q17} Does global Hecke GRH give general weight-monodromy?

\textbf{State:} current gate:\allowbreak{} all-power trace bounds for finite exterior constructions and duals; prior reduction retained

\textbf{Source:} ALPOGE\_\allowbreak{}BRIDGES D-WM,BR-WM-02,CC-WM01–\allowbreak{}03.

\paragraph{AG-Q18} What K-group statement is actually exported by A?

\textbf{State:} endpoint A does not itself export an indexed K-group theorem; inherited2026K8 claim not independently certified; integral comparison/\allowbreak{}class-group components remain gates

\textbf{Source:} BRIDGES Q06; Alpoge step17/\allowbreak{}fn6; ORDER31.

\paragraph{AG-Q23} What exact extra positivity would eliminate poles after Hecke/\allowbreak{}Brauer divisor localization?

\textbf{State:} proved equivalence with all finite signed Hankel PSD and an explicit finite witness for any negative multiplicity; proving these inequalities from independent arithmetic data remains open

\textbf{Source:} AG-SIGNED; rootMT-ART02; cote1Fredholm framework.

\paragraph{AG-Q24} Does a specified arithmetic section and solvable H make its action-commutators open in H?

\textbf{State:} proved AG-C34-SECTION using all relevant finite covers; section existence not inferred

\textbf{Source:} Cote34PDF64/\allowbreak{}footer63 Rem1.\allowbreak{}17; AG-C34-SECTION.

\paragraph{AG-Q29} Does a Soulé-image assertion alone extend the regular-prime kernel-field proof to irregular primes?

\textbf{State:} exact freeness/\allowbreak{}generator-relation gate identified; under Greenberg canonical generators fail integrally despite Brown rational freeness; field-equality question remains unresolved

\textbf{Source:} LM6A44–\allowbreak{}A45; Fable§5; SharifiThm1.\allowbreak{}1/\allowbreak{}1.\allowbreak{}3.

\paragraph{AG-Q32} What remains to transport the formal gain to theta values or higher-dimensional CM formal factors?

\textbf{State:} formal bound and integral-product version proved; remaining gate is actual integral theta trivialization/\allowbreak{}operator normalization and integral formal-factor decomposition, not merely a rational isogeny

\textbf{Source:} Fable§2 higher-height CM proposal; AG-FG-PRODUCT/\allowbreak{}THETA.

\paragraph{CX-Q04} What is the strongest available general finite-field factorization transfer?

\textbf{State:} proved conditional on M:\allowbreak{} CX-ALG-009; A\_\allowbreak{}D-only extension gate CX-ALG-010 (high valuation coefficient unit or alternative descent)

\textbf{Source:} CX\_\allowbreak{}GRH CX-N15/\allowbreak{}CX-U02.

\paragraph{CX-Q05} What does GRH give for class groups and what computation remains?

\textbf{State:} proved conditional generator transfer CX-ALG-011/\allowbreak{}013; full polynomial relation-search bound remains gate CX-ALG-012

\textbf{Source:} Bach1990T4; CX\_\allowbreak{}GRH CX-N07/\allowbreak{}N08.

\paragraph{CX-Q15} Do the cited GRH transfers close general derandomization or GCT hardness gates?

\textbf{State:} current gate:\allowbreak{} CX-LIM-001, uniform hitting sets or independent border-hardness theorem; no logical nonimplication claimed

\textbf{Source:} GCT V T4.\allowbreak{}7/\allowbreak{}L4.\allowbreak{}8; PIT/\allowbreak{}BPP questions.

\paragraph{CX-Q16} Does positivity/\allowbreak{}purity remove efficient multiplicity-construction gates?

\textbf{State:} specified earlier counterexamples carried CX-LIM-002; actual quotient ranks and uniform data construction remain

\textbf{Source:} ALPOGE\_\allowbreak{}GCT P-F7/\allowbreak{}P-F10 and PLETH PLE-F5/\allowbreak{}F7.

\paragraph{CX-Q32} What is the smallest sufficient family established here for the full factoring composition?

\textbf{State:} proved CX-ALG-025:\allowbreak{} pure Q(p\textasciicircum{}(1/\allowbreak{}r)) only for encountered prime r\ensuremath{\le}d with r∤(p\ensuremath{-}1)m and higher valuation; general A\_\allowbreak{}D construction of these prime-degree fields remains open

\textbf{Source:} CX\_\allowbreak{}GRH CX-N15/\allowbreak{}CX-U02; IKRS2011 Theorem5.\allowbreak{}8, Proposition4.\allowbreak{}6, Lemma4.\allowbreak{}2.

\paragraph{CX-Q33} Does exact extension construction share the same smaller analytic and unconditional promise gates?

\textbf{State:} proved CX-ALG-026 in poly(n log q), with n charged to output size; primes r|pm or valuation-one cases are unconditional

\textbf{Source:} Lenstra1991 Theorem9.\allowbreak{}1 pp344–\allowbreak{}346; CX-ALG-023/\allowbreak{}024.

\paragraph{Q-ARTIN-POS} Can positivity remove critical-line Artin poles?

\textbf{State:} current gate MT-ART02:\allowbreak{} positive semidefinite local Heilbronn matrix at multiple zeta zeros; no proof of that positivity

\textbf{Source:} ORDER42; Alpoge steps4,7.

\paragraph{Q-A5-LOCAL} What sharper local Artin-pole gates follow from all proper-subgroup Hecke tests for A5?

\textbf{State:} proved MT-ART05/\allowbreak{}06:\allowbreak{} exact cone, degree3/\allowbreak{}4 thresholds and central6/\allowbreak{}8 thresholds; actual divisor positivity remains an arithmetic gate

\textbf{Source:} MT-ART04; ORDER42; ZETA\_\allowbreak{}TWO ZT02.

\paragraph{Q-A5-NONREAL} Does the non-real A5-over-Q specialization retain an Artin holomorphy gate?

\textbf{State:} proved closed by classical automorphy; M transfer and tensor-category form MT-ART07; totally real case remains

\textbf{Source:} ORDER42; Calegari2013 §1 PDF2.

\paragraph{Q-S5-CENTRE} What is the exact remaining nonvanishing gate under M in the stated S5 family?

\textbf{State:} proved reduction MT-S502 to \ensuremath{\zeta}\_\allowbreak{}H(1/\allowbreak{}2)\ensuremath{\ne}0; dimensions4/\allowbreak{}6 already covered by MT-S501; central value unproved in general

\textbf{Source:} Calegari2013 Theorems1.\allowbreak{}1/\allowbreak{}1.\allowbreak{}2; ORDER42.

\paragraph{Q-MAS-INFINITE-BASIS} Can the required arithmetic infinite-divisor quotient be identified and given its asserted absolute basis and nuclear estimates?

\textbf{State:} current gate:\allowbreak{} work on the separated closed-ideal quotient, prove controlled infinite Mellin interpolation and basis/\allowbreak{}nuclear bounds; local zero counts and finite-divisor repair alone do not establish it

\textbf{Source:} Masullo9.\allowbreak{}3.\allowbreak{}2.\allowbreak{}5; MT-M03/\allowbreak{}M04/\allowbreak{}M06; ORDER42.
\section*{Statement index}
\addcontentsline{toc}{section}{Statement index}
The complete statement, exact input, proof location and work placement for
every ID below is recorded in \path{CONSEQUENCES.tsv}. This index keeps
the paper navigable without duplicating the entire machine-readable table.
\paragraph{number theory} NT-P01, NT-P02, NT-P03, NT-P04, NT-P05, NT-P06, NT-P07, NT-P08, NT-P09, NT-P10, NT-C01, NT-C02, NT-C03, NT-C04, NT-C05, NT-H01, NT-H02, NT-H03, NT-H04, NT-A01, NT-A02, NT-A03, NT-F01, NT-F02, NT-F03, NT-F04, NT-L01, NT-L02, NT-L03, NT-L04, NT-L05, NT-L06, NT-L07, NT-L08, NT-Z01, NT-Z02, NT-Z03, NT-Z04, NT-Z05, NT-Z06, NT-Z07, NT-Z08, NT-Z09, NT-L09, NT-L10, NT-Z10, NT-P11, NT-P12, NT-Z11, NT-Z12, NT-H05, NT-H06, NT-H07, NT-H08, NT-R01, NT-R02, NT-R03, NT-R04, NT-Z13, NT-Z14, NT-Z15, NT-Z16, NT-P13, NT-P14, NT-P15, NT-P16, NT-L11, NT-L12, NT-H09, NT-H10.
\paragraph{geometry programme} AG-AT, AG-ARTIN, AG-DED, AG-LOCAL, AG-MODEL, AG-CM, AG-WEIGHT, AG-SCHEMES, AG-NONCM, AG-NONCM-W0, AG-FROB, AG-LAMBDA, AG-SIGN, AG-HEIGHT, AG-SIGN-NORM, AG-CRITICAL, AG-DET, AG-MOMENT, AG-HANKEL, AG-EXT, AG-EXT-K, AG-C1, AG-C16, AG-C34-74, AG-C34-64, AG-C44, AG-C149, AG-NODE, AG-INERTIA, AG-GT, AG-PURITY, AG-WM, AG-STANDARD, AG-COVER, AG-KUMMER, AG-KUMMER-2, AG-CM-REC, AG-CM-PERIOD, AG-SIGNED, AG-C34-TYPE, AG-C34-REPAIR, AG-C34-INDEX, AG-C34-SECTION, AG-C34-NONSOLV, AG-C149-COMP, AG-C149-LIFT, AG-C149-ELLIPTIC, AG-LM6-REG, AG-LM6-IRREG, AG-LM6-NORM, AG-FG-EXP, AG-FG-DER, AG-FG-PPOWER, AG-FG-PRODUCT, AG-FG-THETA.
\paragraph{complexity} CX-ALG-001, CX-ALG-002, CX-ALG-003, CX-ALG-004, CX-ALG-005, CX-ALG-006, CX-ALG-007, CX-ALG-008, CX-ALG-009, CX-ALG-010, CX-ALG-011, CX-ALG-012, CX-ALG-013, CX-ALG-014, CX-HN-001, CX-HN-002, CX-HN-003, CX-HN-004, CX-GEO-001, CX-GEO-002, CX-GEO-003, CX-GEO-004, CX-GEO-005, CX-GEO-006, CX-NN-001, CX-NN-002, CX-NN-003, CX-NN-004, CX-LIM-001, CX-LIM-002, CX-LIM-003, CX-GEO-007, CX-GEO-008, CX-ALG-015, CX-ALG-016, CX-ALG-017, CX-ALG-018, CX-GEO-009, CX-ALG-019, CX-ALG-020, CX-ALG-021, CX-ALG-022, CX-ALG-023, CX-ALG-024, CX-ALG-025, CX-ALG-026.
\paragraph{root} MT-A01, MT-M01, MT-M02, MT-M03, MT-M04, MT-ART01, MT-ART02, MT-ART03, MT-L01, MT-L02, MT-ART04, MT-Q01, MT-Q02, MT-ART05, MT-ART06, MT-H01, MT-H02, MT-H03, MT-ART07, MT-ART08, MT-S501, MT-S502, MT-S503, MT-M05, MT-M06, MT-M07, MT-M08, MT-M09.
\section*{Source registry and exact locators}
\addcontentsline{toc}{section}{Source registry and exact locators}
The inspection descriptions below record their actual scope. A located
classical theorem is not labelled as a fresh whole-proof audit unless
its event explicitly records that reading.

\paragraph{S-NT-LLS} Lamzouri–\allowbreak{}Li–\allowbreak{}Soundararajan, Conditional bounds for the least quadratic non-residue and related problems, corrected arXiv v3 (8 Dec 2016).

\textbf{Locator:} Theorem1.\allowbreak{}1 and Corollary1.\allowbreak{}1 p2; Corollary1.\allowbreak{}2 p3; Theorem1.\allowbreak{}5 pp3–\allowbreak{}4; Corollary1.\allowbreak{}3 and (1.\allowbreak{}3) p4.\allowbreak{} Author-page numbering.\allowbreak{}.

\textbf{Inspection:} Primary PDF text read for full statements, thresholds, U/\allowbreak{}V normalization and correction-version choice.\allowbreak{} These results were already background IA02; no fresh reproof of the long literature argument.\allowbreak{}

\url{https://arxiv.org/pdf/1309.3595v3}.

\paragraph{S-NT-GM} Grenié–\allowbreak{}Molteni, An explicit Chebotarev density theorem under GRH, J.\allowbreak{}Number Theory200(2019),441–\allowbreak{}485.

\textbf{Locator:} Theorem1.\allowbreak{}1 author p2; Corollary1.\allowbreak{}2 author p3; AppendixA TheoremA.\allowbreak{}1 author p26.\allowbreak{}.

\textbf{Inspection:} Original author PDF pp2–\allowbreak{}3 rendered and visually inspected for both explicit error polynomials and unramified convention.\allowbreak{} AppendixA finite-order Hecke zero-count statement read; that appendix is unconditional.\allowbreak{} The local O(log T) deduction uses consecutive heights.\allowbreak{}

\url{https://sites.unimi.it/molteni/research/papers-pdf/44-molteni-An_explicit_Chebotarev_density_theorem_under_GRH.pdf}.

\paragraph{S-NT-BS} Bach–\allowbreak{}Sorenson, Explicit bounds for primes in residue classes, Math.\allowbreak{}Comp.\allowbreak{}65(1996),1717–\allowbreak{}1735.

\textbf{Locator:} §2 printed1720; Corollary3.\allowbreak{}3 printed1721; Theorem5.\allowbreak{}1 printed1729/\allowbreak{}PDF13.\allowbreak{}.

\textbf{Inspection:} Primary author-upload PDF opened in this agent.\allowbreak{} Geometry agent confirmed precise Thm5.\allowbreak{}1 statement; RH \ensuremath{\zeta}\_\allowbreak{}L and \ensuremath{\zeta}\_\allowbreak{}Q; degree-one prime; L\ensuremath{\ne}Q.\allowbreak{} Arbitrary extra-prime exclusion is not in this theorem.\allowbreak{}

\url{https://www.researchgate.net/profile/Jonathan-Sorenson-2/publication/220576951_Explicit_bounds_for_primes_in_residue_classes/links/55100ea70cf2ac2905afbceb/Explicit-bounds-for-primes-in-residue-classes.pdf}.

\paragraph{S-NT-SCHOENFELD} Schoenfeld, Sharper bounds for the Chebyshev functions \ensuremath{\theta}(x) and \ensuremath{\psi}(x).\allowbreak{}II, Math.\allowbreak{}Comp.\allowbreak{}30(1976),337–\allowbreak{}360.

\textbf{Locator:} Theorem10 equations(6.\allowbreak{}1)–\allowbreak{}(6.\allowbreak{}5), printed337; Corollary1 to Theorem10.\allowbreak{}.

\textbf{Inspection:} Primary indexed full statement of Theorem10 read, including 73.\allowbreak{}2/\allowbreak{}599 and 23·10\textasciicircum{}8.\allowbreak{} Direct AMS PDF403; mirror https:\allowbreak{}/\allowbreak{}/\allowbreak{}electronicsandbooks.\allowbreak{}com/\allowbreak{}edt/\allowbreak{}manual/\allowbreak{}Magazine/\allowbreak{}M/\allowbreak{}Mathematics\%20of\%20Computation/\allowbreak{}1960-2002/\allowbreak{}pdf/\allowbreak{}1976\_\allowbreak{}v030\_\allowbreak{}n134/\allowbreak{}2005976.\allowbreak{}pdf timed out via web.\allowbreak{} Corollary1 original page still to inspect; do not label full source proof checked.\allowbreak{}

\url{https://www.ams.org/mcom/1976-30-134/S0025-5718-1976-0457374-X/S0025-5718-1976-0457374-X.pdf}.

\paragraph{S-NT-DGM} Dudek–\allowbreak{}Grenié–\allowbreak{}Molteni, Primes in explicit short intervals on RH, Int.\allowbreak{}J.\allowbreak{}NumberTheory12(2016),1391–\allowbreak{}1407.

\textbf{Locator:} Theorem1.\allowbreak{}1 author p1.\allowbreak{}.

\textbf{Inspection:} Primary theorem statement, c=1/\allowbreak{}2+2/\allowbreak{}log x and x\ensuremath{\ge}2 read.\allowbreak{} Computational c=.\allowbreak{}55 refinement not used.\allowbreak{}

\url{https://sites.unimi.it/molteni/research/papers-pdf/36-molteni-Primes_in_explicit_short_intervals_on_RH.pdf}.

\paragraph{S-NT-GMP} Grenié–\allowbreak{}Molteni–\allowbreak{}Perelli, Primes and prime ideals in short intervals, arXiv:\allowbreak{}1602.\allowbreak{}02906v2.

\textbf{Locator:} Theorems1,2 and subsequent asymptotic statements, author p2; Proposition1 author p3.\allowbreak{}.

\textbf{Inspection:} Primary theorem statements and distinction Dirichlet-GRH versus RH \ensuremath{\zeta}\_\allowbreak{}K read.\allowbreak{} Absolute constants remain unspecified.\allowbreak{}

\url{https://arxiv.org/pdf/1602.02906}.

\paragraph{S-NT-SHOUP} Shoup, Searching for primitive roots in finite fields (1992).

\textbf{Locator:} Theorem3, author PDF p4.\allowbreak{}.

\textbf{Inspection:} Complexity agent downloaded original PDF with a standard browser user agent.\allowbreak{} This agent inspected the actual rendered p4 image at work/\allowbreak{}complexity/\allowbreak{}sources/\allowbreak{}Shoup\_\allowbreak{}Theorem3.\allowbreak{}png.\allowbreak{} Exact r\textasciicircum{}4(logr+1)\textasciicircum{}4log²p bound confirmed; no numerical implicit constant claimed.\allowbreak{}

\url{https://www.shoup.net/papers/primroots.pdf}.

\paragraph{S-NT-BACH} Bach, Explicit bounds for primality testing and related problems, Math.\allowbreak{}Comp.\allowbreak{}55(1990),355–\allowbreak{}380.

\textbf{Locator:} Theorem4 printed375 and corollary376, inherited exact locator.\allowbreak{}.

\textbf{Inspection:} Carried as background with no fault found; direct AMS requests inaccessible.\allowbreak{} Primary fresh inspection not claimed.\allowbreak{} Complexity agent owns the algorithmic application.\allowbreak{}

\url{https://www.ams.org/mcom/1990-55-191/S0025-5718-1990-1023756-8/S0025-5718-1990-1023756-8.pdf}.

\paragraph{S-NT-LENSTRA} Lenstra, On Artin’s conjecture and Euclid’s algorithm in global fields, Invent.\allowbreak{}Math.\allowbreak{}42(1977),201–\allowbreak{}224.

\textbf{Locator:} §2 pp202–\allowbreak{}206; (2.\allowbreak{}14),(2.\allowbreak{}15) pp205–\allowbreak{}206; Theorem(3.\allowbreak{}1)206; Theorem(4.\allowbreak{}1)207; Theorem(9.\allowbreak{}1),(9.\allowbreak{}2)218.\allowbreak{}.

\textbf{Inspection:} Primary definitions and theorem statements read; original p218 rendered and inspected, confirming \#S\ensuremath{\ge}2.\allowbreak{} Full elementary reduction in Theorem3.\allowbreak{}1/\allowbreak{}lemma3.\allowbreak{}2 read, but Cooke–\allowbreak{}Weinberger base theorem not independently reproved.\allowbreak{}

\url{https://pub.math.leidenuniv.nl/~lenstrahw/PUBLICATIONS/1977d/fulltext-12.pdf}.

\paragraph{S-NT-HOOLEY} Hooley, On Artin’s conjecture, J.\allowbreak{}reine angew.\allowbreak{}Math.\allowbreak{}225(1967),209–\allowbreak{}220.

\textbf{Locator:} Unnumbered final theorem pp219–\allowbreak{}220; original internal locator corroborated by later primary research.\allowbreak{}.

\textbf{Inspection:} Original publication metadata checked; original fulltext endpoint blocked.\allowbreak{} For the exact proof used here, cite inspected Lenstra Theorem3.\allowbreak{}1 instead.\allowbreak{} Original theorem internal-page check remains a source gate, not a mathematical gate on the classical corollary.\allowbreak{}

\url{https://doi.org/10.1515/crll.1967.225.209}.

\paragraph{S-NT-TITCHMARSH} Titchmarsh, The Theory of the Riemann Zeta-function, 2nd ed.\allowbreak{}, revised by Heath-Brown (1986).

\textbf{Locator:} Lemma3.\allowbreak{}12 and proof printed61–\allowbreak{}62; Theorem14.\allowbreak{}2 and (14.\allowbreak{}2.\allowbreak{}5),(14.\allowbreak{}2.\allowbreak{}6) printed336–\allowbreak{}337; Theorem14.\allowbreak{}14(A); Theorem14.\allowbreak{}25(A)–\allowbreak{}(C) printed369–\allowbreak{}370.\allowbreak{}.

\textbf{Inspection:} Downloaded local PDF and text.\allowbreak{} Relevant theorem statements and the whole printed 14.\allowbreak{}25 argument read; Perron kernel locator checked for the new transfer proof.\allowbreak{} Not a check of the entire textbook.\allowbreak{} Two printed pages per PDF page.\allowbreak{} Full reference text excluded from deliverables.\allowbreak{}

\url{https://sites.math.rutgers.edu/~zeilberg/EM18/TitchmarshZeta.pdf}.

\paragraph{S-NT-SIMONIC} Simonič, Estimates for L-functions in the Critical Strip Under GRH with Effective Applications, Mediterr.\allowbreak{}J.\allowbreak{}Math.\allowbreak{}20(2023),87.

\textbf{Locator:} Theorem1 equations(4)–\allowbreak{}(6); Corollaries2–\allowbreak{}3; §2.\allowbreak{}1 SP axioms; Theorem2 equations(7),(8); proof §3.\allowbreak{}.

\textbf{Inspection:} Primary full-text theorem and axioms read.\allowbreak{} Fixed-character consequences derived in full in section.\allowbreak{}tex; effective Mertens constants checked in original HTML.\allowbreak{} No uniform hybrid-conductor Lindelöf estimate claimed.\allowbreak{}

\url{https://link.springer.com/article/10.1007/s00009-023-02289-2}.

\paragraph{S-NT-HARPER} Harper, Sharp conditional bounds for moments of the Riemann zeta function, arXiv:\allowbreak{}1305.\allowbreak{}4618.

\textbf{Locator:} Theorem1 authorp1; Theorem2 authorp15, proof§5 pp15–\allowbreak{}17, Propositions3/\allowbreak{}4.\allowbreak{}.

\textbf{Inspection:} Primary theorem statement read with k\ensuremath{\ge}0 fixed and implicit constant depending on k.\allowbreak{} Source introduction distinguishes order bounds from precise moments.\allowbreak{} No complete reproof of the paper.\allowbreak{} Theorem2 original p15 rendered and inspected for k\ensuremath{\ge}0, both signs of fundamental d, discriminant shell and exponent k(k+1)/\allowbreak{}2.\allowbreak{} New Markov deductions in NT-L12 are proved in full; the paper’s full moment argument is not independently recertified.\allowbreak{}

\url{https://arxiv.org/pdf/1305.4618}.

\paragraph{S-NT-MONTGOMERY} Montgomery, The pair correlation of zeros of the zeta function, Proc.\allowbreak{}Sympos.\allowbreak{}PureMath.\allowbreak{}24(1973),181–\allowbreak{}193.

\textbf{Locator:} §1 definition and unnumbered Theorem p181; Corollary2 p182; conjecture following restricted theorem.\allowbreak{}.

\textbf{Inspection:} Primary web PDF text read for normalization, \ensuremath{\alpha} range, uniformity and simplicity proportion; direct local download403.\allowbreak{} Full pair-correlation conjecture carefully separated.\allowbreak{}

\url{https://public.websites.umich.edu/hlm/paircor1.pdf}.

\paragraph{S-NT-HR} Hughes–\allowbreak{}Rudnick, Linear statistics of low-lying zeros of L-functions, Q.\allowbreak{}J.\allowbreak{}Math.\allowbreak{}54(2003),309–\allowbreak{}333.

\textbf{Locator:} v2:\allowbreak{} ordering p1; admissibility §2; Theorem3.\allowbreak{}1p5; Theorem4.\allowbreak{}1p7; Theorem5.\allowbreak{}1p9; Corollary8.\allowbreak{}2p20; Theorem8.\allowbreak{}3pp20–\allowbreak{}21.\allowbreak{}.

\textbf{Inspection:} Primary text read for unconditional complex-zero formulation, Fourier support, m<2/\allowbreak{}\ensuremath{\alpha} condition, and first NONNEGATIVE zero convention.\allowbreak{} Source statements checked, not full moment proof recertified.\allowbreak{}

\url{https://arxiv.org/pdf/math/0208230}.

\paragraph{S-NT-ILS} Iwaniec–\allowbreak{}Luo–\allowbreak{}Sarnak, Low lying zeros of families of L-functions, Publ.\allowbreak{}Math.\allowbreak{}IHES91(2000),55–\allowbreak{}131.

\textbf{Locator:} Global RH assumption printed56; Theorem1.\allowbreak{}1 printed58–\allowbreak{}59; proof Theorem7.\allowbreak{}2.\allowbreak{}.

\textbf{Inspection:} Primary published PDF text inspected.\allowbreak{} The family GL2-RH hypothesis remains separate from auxiliary Dirichlet-GRH; not silently discharged by M.\allowbreak{}

\url{https://www.numdam.org/article/PMIHES_2000__91__55_0.pdf}.

\paragraph{S-NT-RT} Rodgers–\allowbreak{}Tao, The de Bruijn–\allowbreak{}Newman constant is non-negative, arXiv:\allowbreak{}1801.\allowbreak{}05914v5.

\textbf{Locator:} §1 definition/\allowbreak{}equivalence pp1–\allowbreak{}2; Theorem1p3.\allowbreak{}.

\textbf{Inspection:} Primary version5 dated3Jul2021; theorem and classical RH equivalence read.\allowbreak{} Combination with A gives Λ=0, with normalization retained.\allowbreak{}

\url{https://arxiv.org/pdf/1801.05914}.

\paragraph{S-NT-LAGARIAS} Lagarias, An elementary problem equivalent to the Riemann hypothesis, arXiv:\allowbreak{}math/\allowbreak{}0008177v2.

\textbf{Locator:} Theorem1.\allowbreak{}1 author p1.\allowbreak{}.

\textbf{Inspection:} Primary theorem statement read, including equality only at n=1; no reproduction of entire classical proof.\allowbreak{}

\url{https://arxiv.org/pdf/math/0008177}.

\paragraph{S-NT-MILNE} Milne, Class Field Theory v4.\allowbreak{}03 (2020).

\textbf{Locator:} ChapterVI Corollary2.\allowbreak{}12 and Example2.\allowbreak{}13, printed191/\allowbreak{}PDF200.\allowbreak{}.

\textbf{Inspection:} Primary formula read for residue, full regulator, real-quadratic normalization, roots-of-unity factor.\allowbreak{} Source is author-hosted course notes.\allowbreak{}

\url{https://www.jmilne.org/math/CourseNotes/CFT.pdf}.

\paragraph{S-NT-JOHNSTON} Johnston, Improving bounds on prime counting functions by partial verification of the Riemann hypothesis, RamanujanJ59(2022),1307–\allowbreak{}1321.

\textbf{Locator:} Theorem3.\allowbreak{}2 equations(3.\allowbreak{}3)–\allowbreak{}(3.\allowbreak{}6); §2 paragraph after(2.\allowbreak{}1), ordinary-count normalization.\allowbreak{}.

\textbf{Inspection:} Primary full-text theorem and normalization statement read.\allowbreak{} Under A the finite-height hypothesis can be met for every x, yielding \ensuremath{\psi} threshold59 and \ensuremath{\pi} threshold2657.\allowbreak{} This supplies an exact directly checked alternative proof locator to inaccessible Schoenfeld Corollary1.\allowbreak{}

\url{https://link.springer.com/article/10.1007/s11139-022-00616-x}.

\paragraph{S-NT-RV} Ramakrishnan–\allowbreak{}Valenza, Fourier Analysis on Number Fields, GTM186(1999).

\textbf{Locator:} §7.\allowbreak{}4 Theorem7-19 printed279; §7.\allowbreak{}7 Theorem7-28 printed289.\allowbreak{}.

\textbf{Inspection:} Original PDF downloaded and relevant full statements read directly in extracted text:\allowbreak{} completed unitary-Hecke meromorphy/\allowbreak{}functional equation and pure-norm pole exception; nonvanishing on Re s=1 for every unitary idele-class character.\allowbreak{} Web fetch exceeds size limit; local download succeeded.\allowbreak{} No page-image check or whole-book proof check claimed.\allowbreak{}

\url{https://webusers.imj-prg.fr/~julien.marche/Ref_GL_GT/Fourier_Analysis_on_Number_Fields.pdf}.

\paragraph{S-NT-FPS} Fiorilli–\allowbreak{}Parks–\allowbreak{}Södergren, Low-lying zeros of quadratic Dirichlet L-functions:\allowbreak{} Lower order terms for extended support, CompositioMath153(2017),1196–\allowbreak{}1216; arXiv1601.\allowbreak{}06833.

\textbf{Locator:} Definitions(1.\allowbreak{}1)–\allowbreak{}(1.\allowbreak{}2) and Theorem1.\allowbreak{}1 authorp2; Lemma2.\allowbreak{}5p7; Theorem3.\allowbreak{}5p12; final proof of Theorem1.\allowbreak{}1p14.\allowbreak{}.

\textbf{Inspection:} Primary PDF downloaded; exact family \ensuremath{\chi}\_\allowbreak{}\{8d\} for odd squarefree d of both signs, even nonnegative Schwartz weights, Fourier convention and support (\ensuremath{-}2,2) read.\allowbreak{} Checked Dirichlet-only GRH in the prime-sum step.\allowbreak{} Full literature proof not independently revalidated; the new nonvanishing deduction in section.\allowbreak{}tex is complete original prose.\allowbreak{}

\url{https://arxiv.org/pdf/1601.06833}.

\paragraph{S-NT-TV} Tsfasman–\allowbreak{}Vlăduţ, Infinite Global Fields and the Generalized Brauer–\allowbreak{}Siegel Theorem, MoscowMathJ2(2002),329–\allowbreak{}402; arXivmath/\allowbreak{}0205129v1.

\textbf{Locator:} Definition2.\allowbreak{}1 and Lemma2.\allowbreak{}2 authorp15; GRH Theorem3.\allowbreak{}1p19, proofpp19–\allowbreak{}21; GRH Theorem7.\allowbreak{}2p44, proofpp44–\allowbreak{}49; Theorem7.\allowbreak{}3p44 gives a distinct unconditional almost-normal-tower result.\allowbreak{}.

\textbf{Inspection:} Original PDF downloaded; p19 basic-inequality constants and p44 generalized Brauer–\allowbreak{}Siegel statement inspected as rendered images.\allowbreak{} Exact-family definition, Dedekind-only RH and absence of normality/\allowbreak{}tower assumptions in Thm7.\allowbreak{}2 checked.\allowbreak{} Full literature proofs not independently revalidated.\allowbreak{} Our derived splitting/\allowbreak{}tower arguments are complete in section.\allowbreak{}tex.\allowbreak{}

\url{https://arxiv.org/pdf/math/0205129}.

\paragraph{S-NT-RS} Rubinstein–\allowbreak{}Sarnak, Chebyshev’s Bias, ExperimentalMath3(1994),173–\allowbreak{}197.

\textbf{Locator:} Theorem1.\allowbreak{}1p174, proof§2.\allowbreak{}1; Theorem1.\allowbreak{}2p175, proof§2.\allowbreak{}2; GSH/\allowbreak{}formula(1.\allowbreak{}2)/\allowbreak{}Remark1.\allowbreak{}3p176; Theorems1.\allowbreak{}4/\allowbreak{}1.\allowbreak{}5p177, proofsProposition3.\allowbreak{}1/\allowbreak{}§3.\allowbreak{}2.\allowbreak{}.

\textbf{Inspection:} Author-hosted original PDF downloaded; text encoding is corrupt, so printedpp174–\allowbreak{}178 were rendered and inspected visually.\allowbreak{} Checked exact \ensuremath{\pi}(x) normalization, continuous-test versus ordering-indicator distinction, lower logarithmic density under GRH alone, GSH collection including central zeros, law-symmetry classification and fixed-r large-q quantifier.\allowbreak{} Full classical proofs not independently reread.\allowbreak{}

\url{https://www.math.uwaterloo.ca/~mrubinst/publications/Chebyshev.pdf}.

\paragraph{S-NT-CCM} Carneiro–\allowbreak{}Chandee–\allowbreak{}Milinovich, A note on the zeros of zeta and L-functions, Math.\allowbreak{}Z.\allowbreak{}281(2015),315–\allowbreak{}332; arXiv1503.\allowbreak{}00955v1.

\textbf{Locator:} Theorem1 authorp2, proof§3 pp3–\allowbreak{}5; §4 axioms p6; Theorem5 equations(4.\allowbreak{}4),(4.\allowbreak{}5) p7, proofpp7–\allowbreak{}12; Theorem7 and proofp13; Examples8–\allowbreak{}10p14.\allowbreak{}.

\textbf{Inspection:} Original PDF downloaded.\allowbreak{} Axioms p6 and full Theorem5 p7 rendered and visually inspected.\allowbreak{} Read Theorem1, exact conductor normalization and universal-constant assertion, Theorem7 and its entire p13 proof.\allowbreak{} Critical restriction:\allowbreak{} gamma parameters real or in conjugate pairs.\allowbreak{} Applied to finite-order Hecke characters with explicit axiom verification; did not assume this applies verbatim to arbitrary unitary infinity types.\allowbreak{} Full Theorem5 classical proof not independently recertified; NT-Z14 deduction is proved and checked in full.\allowbreak{}

\url{https://arxiv.org/pdf/1503.00955}.

\paragraph{S-NT-CMS} Carneiro–\allowbreak{}Milinovich–\allowbreak{}Soundararajan, Fourier optimization and prime gaps, Comment.\allowbreak{}Math.\allowbreak{}Helv.\allowbreak{}94(2019),533–\allowbreak{}568; arXiv1708.\allowbreak{}04122v2.

\textbf{Locator:} Theorem3, Corollary4, Theorem5 authorp5; proofs§5 and §6 pp21–\allowbreak{}26.\allowbreak{}.

\textbf{Inspection:} Downloaded original author v2.\allowbreak{} Page5 rendered and visually inspected for strict21/\allowbreak{}25 limsup, exact22/\allowbreak{}25 all-x interval, x\ensuremath{\ge}4 and p\_\allowbreak{}n>3.\allowbreak{} Full long Fourier/\allowbreak{}explicit-estimate proof not independently recertified; numerical routines not rerun.\allowbreak{}

\url{https://arxiv.org/pdf/1708.04122v2}.

\paragraph{S-NT-CMQR} Carneiro–\allowbreak{}Milinovich–\allowbreak{}Quesada-Herrera–\allowbreak{}Ramos, Fourier optimization, the least quadratic non-residue, and the least prime in an arithmetic progression; arXiv2404.\allowbreak{}08380v2 (2025), Math.\allowbreak{}Comp.\allowbreak{} DOI10.\allowbreak{}1090/\allowbreak{}mcom/\allowbreak{}4154.

\textbf{Locator:} Theorems1/\allowbreak{}2p2, Theorem3p3, definition(1.\allowbreak{}8)/\allowbreak{}Theorem6p5, Corollary7p6; proof§4pp13–\allowbreak{}17, §4.\allowbreak{}4pp16–\allowbreak{}17, §5pp18–\allowbreak{}22; numerical proof§6.\allowbreak{}.

\textbf{Inspection:} Downloaded original v2 (dated August2025).\allowbreak{} Theorem pages2,3,6 rendered and inspected.\allowbreak{} Checked exact least-prime-residue convention, character zeros excluded, fixed character order, 0.\allowbreak{}7615/\allowbreak{}0.\allowbreak{}6707/\allowbreak{}0.\allowbreak{}6131/\allowbreak{}0.\allowbreak{}5765 constants, \ensuremath{\ell}\ensuremath{\ge}6 formula and strict8/\allowbreak{}9 AP asymptotic; proof conclusion(5.\allowbreak{}15)–\allowbreak{}(5.\allowbreak{}16) checked for a-uniformity.\allowbreak{} The paper’s computer-assisted optimization proof is cited, not rerun.\allowbreak{}

\url{https://arxiv.org/pdf/2404.08380v2}.

\paragraph{S-AG-ZETA} Supplied ZETA\_\allowbreak{}TWO return.

\textbf{Locator:} RETURN\_\allowbreak{}ZETA\_\allowbreak{}TWO\_\allowbreak{}final\_\allowbreak{}IMPLICATIONS.\allowbreak{}md §§1–\allowbreak{}4, ZT01–\allowbreak{}ZT08, CM01, SV01–\allowbreak{}SV02; matching LEDGER and QUESTIONS.

\textbf{Inspection:} Full implication text read and relevant reductions rechecked algebraically; original earlier full proof file not included.\allowbreak{} Earlier findings retain their IDs and are not claimed new.\allowbreak{}


\paragraph{S-AG-BRIDGES} Supplied ALPOGE\_\allowbreak{}BRIDGES and original BRIDGES returns.

\textbf{Locator:} RETURN\_\allowbreak{}ALPOGE\_\allowbreak{}BRIDGES\_\allowbreak{}cp02\_\allowbreak{}IMPLICATIONS.\allowbreak{}md and QUESTIONS:\allowbreak{} BR-TR-01/\allowbreak{}02, BR-MIX-01, BR-WM-02, AQ-FINITE, R-GTHAT; OURS\_\allowbreak{}BRIDGES\_\allowbreak{}BRIDGES\_\allowbreak{}20261120.\allowbreak{}tex F1–\allowbreak{}F4,Q01–\allowbreak{}Q08.

\textbf{Inspection:} Implication/\allowbreak{}ledger/\allowbreak{}question records read; reused background is distinguished from the original proof package, which is not attached.\allowbreak{}


\paragraph{S-AG-BS} Bach–\allowbreak{}Sorenson, Explicit bounds for primes in residue classes, Math.\allowbreak{}Comp.\allowbreak{}65(1996),1717–\allowbreak{}1735.

\textbf{Locator:} Theorem5.\allowbreak{}1 printed1729/\allowbreak{}PDF13; §2 printed1720 (RH \ensuremath{\zeta}E,\ensuremath{\zeta}Q); Corollary3.\allowbreak{}3 printed1721.

\textbf{Inspection:} Original author-upload PDF theorem and hypothesis text read; no reproof of full classical explicit estimate.\allowbreak{}

\url{https://www.researchgate.net/profile/Jonathan-Sorenson-2/publication/220576951_Explicit_bounds_for_primes_in_residue_classes/links/55100ea70cf2ac2905afbceb/Explicit-bounds-for-primes-in-residue-classes.pdf}.

\paragraph{S-AG-DG} Deligne–\allowbreak{}Goncharov, Groupes fondamentaux motiviques de Tate mixte.

\textbf{Locator:} arXiv math/\allowbreak{}0302267v2,27Feb2003; (2.\allowbreak{}1.\allowbreak{}3), Propositions2.\allowbreak{}2–\allowbreak{}2.\allowbreak{}3, §2.\allowbreak{}4, printed14/\allowbreak{}PDF16.

\textbf{Inspection:} Primary formula and surrounding hypotheses read, distinguishing rational Ext and free Lie presentation from canonical generators and integral torsion.\allowbreak{}

\url{https://arxiv.org/pdf/math/0302267}.

\paragraph{S-AG-MCMULLEN} C.\allowbreak{}McMullen, Advanced Complex Analysis, Harvard Math213a course notes.

\textbf{Locator:} Version16Nov2026; Theorem3.\allowbreak{}16 printed92/\allowbreak{}PDF93; Theorems3.\allowbreak{}10–\allowbreak{}3.\allowbreak{}12 zero exponent.

\textbf{Inspection:} Primary author notes used for exact classical Hadamard locator; statement and applicability read.\allowbreak{} The theorem is classical, not presented as a2026 innovation.\allowbreak{}

\url{https://people.math.harvard.edu/~ctm/home/text/class/harvard/101/22/html/www/math213a/home/course/course.pdf}.

\paragraph{S-AG-KUMMER} Stacks Project, Kummer theory.

\textbf{Locator:} Lemma59.\allowbreak{}28.\allowbreak{}1 and immediately following long exact cohomology sequence; fppf variant Lemma59.\allowbreak{}28.\allowbreak{}3 kept separate.

\textbf{Inspection:} Entire proof of Lemma59.\allowbreak{}28.\allowbreak{}1 and displayed sequence read; ell-invertibility and inverse-limit transition conventions checked.\allowbreak{}

\url{https://stacks.math.columbia.edu/tag/03PK}.

\paragraph{S-AG-MILNE-ANT} J.\allowbreak{}S.\allowbreak{}Milne, Algebraic Number Theory, v3.\allowbreak{}08,19July2020.

\textbf{Locator:} Theorem4.\allowbreak{}3 PDF72 (Minkowski bound); Theorem4.\allowbreak{}4 class group finiteness.

\textbf{Inspection:} Exact bound and hypotheses read; explicit -35 ideal-class example proved in this run, not copied from source.\allowbreak{}

\url{https://www.jmilne.org/math/CourseNotes/ANT.pdf}.

\paragraph{S-AG-149} Alexander Grothendieck, cote149, Autour de la Longue Marche à travers la théorie de Galois.

\textbf{Locator:} 136PDFpages including cover.\allowbreak{} PDF80/\allowbreak{}footer79 eqs13–\allowbreak{}17; PDF81/\allowbreak{}footer80 eqs18–\allowbreak{}19; PDF82/\allowbreak{}footer81 eqs20–\allowbreak{}21; PDF85/\allowbreak{}footer84 eq22.\allowbreak{} Images PDF79–\allowbreak{}86 inspected.\allowbreak{}.

\textbf{Inspection:} Original official PDF downloaded and selected original page images read.\allowbreak{} The content is normalizer/\allowbreak{}inertia extensions, normal line bundles and Kummer cohomology.\allowbreak{} ORDER31 association to LM6 A44–\allowbreak{}A45 is separate inherited evidence; no K8 theorem attributed to sheets.\allowbreak{}

\url{https://grothendieck.umontpellier.fr/149.pdf}.

\paragraph{S-AG-MASULLO} Supplied Masullo The Music of the Primes draft.

\textbf{Locator:} Original PDF425 Lemma9.\allowbreak{}3.\allowbreak{}2.\allowbreak{}4;426 Theorem9.\allowbreak{}3.\allowbreak{}2.\allowbreak{}6 eq9.\allowbreak{}3.\allowbreak{}2.\allowbreak{}10;804 DefinitionD.\allowbreak{}3.\allowbreak{}3.\allowbreak{}1;11 norm/\allowbreak{}SE definition.

\textbf{Inspection:} Original images425,426,804 inspected, norm text11 read.\allowbreak{} Independent whole-proof checks MT-M01–\allowbreak{}03 recorded in events.\allowbreak{}tsv.\allowbreak{} No claim that M itself is disproved by auxiliary counterexamples.\allowbreak{}


\paragraph{S-AG-HODGE} Supplied Hodge bridge sketch.

\textbf{Locator:} OURS\_\allowbreak{}HODGE\_\allowbreak{}BRIDGE\_\allowbreak{}SKETCH\_\allowbreak{}20261121.\allowbreak{}md §§1–\allowbreak{}5, archive map§5; cote16 scans14–\allowbreak{}18,44–\allowbreak{}50; cote44 scans2,4,6,8–\allowbreak{}9; cote47 scans3,5–\allowbreak{}9.

\textbf{Inspection:} Relevant arguments and reported archive locators read.\allowbreak{} Offset in cote16 scan convention retained.\allowbreak{} Separate2026CM-Hodge claims not independently certified from this sketch.\allowbreak{}


\paragraph{S-AG-SA} Supplied SA22 and SA23 intakes.

\textbf{Locator:} OURS\_\allowbreak{}INTAKE\_\allowbreak{}RESEARCH\_\allowbreak{}SA22\_\allowbreak{}20261120.\allowbreak{}md §5.\allowbreak{}1 A6-node-character/\allowbreak{}A:\allowbreak{}inertia-maximality,§5.\allowbreak{}3Soulé/\allowbreak{}GT,§5.\allowbreak{}4–\allowbreak{}5.\allowbreak{}8; SA23§4.\allowbreak{}1 geo-H-lidskii/\allowbreak{}geo-H-genus.

\textbf{Inspection:} Relevant statements and intake check records read; original proofs/\allowbreak{}scans absent except149 separately retrieved.\allowbreak{} Cote1 scans79–\allowbreak{}80,111–\allowbreak{}112 are inherited locators.\allowbreak{}


\paragraph{S-AG-ANDERSON} G.\allowbreak{}W.\allowbreak{}Anderson, Cyclotomy and an extension of the Taniyama group, Compositio Mathematica57(1986),153–\allowbreak{}217.

\textbf{Locator:} §5.\allowbreak{}3 printed178/\allowbreak{}PDF27; Theorem6 §5.\allowbreak{}6 printed179–\allowbreak{}180/\allowbreak{}PDF28–\allowbreak{}29; Cor5.\allowbreak{}7.\allowbreak{}2/\allowbreak{}Rmk5.\allowbreak{}7.\allowbreak{}3/\allowbreak{}Cor5.\allowbreak{}7.\allowbreak{}4 printed181/\allowbreak{}PDF30; Prop5.\allowbreak{}7.\allowbreak{}5 printed181–\allowbreak{}182/\allowbreak{}PDF30–\allowbreak{}31; §5.\allowbreak{}8 Theorem7 printed182/\allowbreak{}PDF31.

\textbf{Inspection:} Original primary PDF downloaded; images27–\allowbreak{}31 and complete proofs Theorem6, Prop5.\allowbreak{}7.\allowbreak{}5 and Theorem7 read.\allowbreak{} Deeper cited Serre/\allowbreak{}Langlands/\allowbreak{}Deligne/\allowbreak{}Blasius results are classical inputs, not re-proved.\allowbreak{} Printed Rmk5.\allowbreak{}7.\allowbreak{}3 refers to Cor5.\allowbreak{}8.\allowbreak{}2, an internal numbering error noted explicitly.\allowbreak{}

\url{https://www.numdam.org/item/CM_1986__57_2_153_0.pdf}.

\paragraph{S-AG-FRAMEWORK} Classical geometric framework locators carried in ZETA\_\allowbreak{}TWO.

\textbf{Locator:} SGA4 XVII§§5.\allowbreak{}1–\allowbreak{}5.\allowbreak{}3; SGA5 XV§3 no2 unnumbered theorem eq1; EGA IV3(8.\allowbreak{}8.\allowbreak{}2); WeilI Thm1.\allowbreak{}6; WeilII Thm3.\allowbreak{}3.\allowbreak{}1; SGA6 V§§2,7(7.\allowbreak{}1.\allowbreak{}1); Deligne1979§§1.\allowbreak{}1–\allowbreak{}1.\allowbreak{}3,(3.\allowbreak{}1.\allowbreak{}2).

\textbf{Inspection:} Carried exact locators, not a claimed new full reading of the original SGA/\allowbreak{}EGA/\allowbreak{}Deligne works.\allowbreak{}


\paragraph{S-AG-34} Alexander Grothendieck, cote34, ExposéII.

\textbf{Locator:} 210PDFpages includingcover; Thm1.\allowbreak{}1 PDF53/\allowbreak{}footer52; struckRem1.\allowbreak{}17 PDF64/\allowbreak{}footer63; Thm1.\allowbreak{}1 proof §§2–\allowbreak{}3 PDF70–\allowbreak{}78/\allowbreak{}footer69–\allowbreak{}77; (2.\allowbreak{}5.\allowbreak{}5) PDF74/\allowbreak{}footer73.

\textbf{Inspection:} Official84,513,436bytePDF downloaded; pageimages53–\allowbreak{}78 read, with enlarged crops53/\allowbreak{}64.\allowbreak{} Exact normal/\allowbreak{}locallyfinite/\allowbreak{}geoconnected hypotheses and pro-ell solvable quotient wording recovered.\allowbreak{} Full theorem1.\allowbreak{}1 proof route §§2–\allowbreak{}3 read and geometric reductions checked; foundational SGA citations not separately re-proved.\allowbreak{} Webopen failed because filelarge; original local download succeeded.\allowbreak{}

\url{https://grothendieck.umontpellier.fr/34.pdf}.

\paragraph{S-AG-RV} D.\allowbreak{}Ramakrishnan and R.\allowbreak{}J.\allowbreak{}Valenza, Fourier Analysis on Number Fields(1999).

\textbf{Locator:} §7.\allowbreak{}4 Theorem7-19 printed279; §7.\allowbreak{}7 Theorem7-28 printed289, whole proof289–\allowbreak{}292.

\textbf{Inspection:} Primary PDF downloaded by number\_\allowbreak{}theory; theorem/\allowbreak{}hypothesis text and complete boundary proof289–\allowbreak{}292 directly read here.\allowbreak{} Boundary includes all unitary characters and norm twists; completed norm-character pole exceptions retained.\allowbreak{} Full proof of Theorem7-19 not separately re-refereed.\allowbreak{}

\url{https://webusers.imj-prg.fr/~julien.marche/Ref_GL_GT/Fourier_Analysis_on_Number_Fields.pdf}.

\paragraph{S-AG-MILNE-JAV} J.\allowbreak{}S.\allowbreak{}Milne, Jacobian Varieties and Abelian Varieties, revised notes of1986 chapters.

\textbf{Locator:} JacobianVarieties version12June2021 Prop9.\allowbreak{}1 pp30–\allowbreak{}32; AbelianVarieties version2Jan2022 Remarks15.\allowbreak{}3/\allowbreak{}15.\allowbreak{}5 pp28–\allowbreak{}29.

\textbf{Inspection:} Primary statements and comparison proof read; the Galois-compatible Tate-module/\allowbreak{}abelianized-fundamental-group identification is used, not an asserted cycle comparison.\allowbreak{}

\url{https://www.jmilne.org/math/xnotes/JVs.pdf}.

\paragraph{S-AG-BES} N.\allowbreak{}Borne, M.\allowbreak{}Emsalem and J.\allowbreak{}Stix, Lifting Galois sections along torsors, Math.\allowbreak{}Nachr.\allowbreak{}288(2015),1939–\allowbreak{}1951.

\textbf{Locator:} arXiv1301.\allowbreak{}4429v2:\allowbreak{} algebraicK(pi,1) curves §3.\allowbreak{}1PDF4; Prop3.\allowbreak{}1/\allowbreak{}3.\allowbreak{}4PDF5–\allowbreak{}6; continuous-cohomology Def4.\allowbreak{}1/\allowbreak{}Prop4.\allowbreak{}2PDF7; Prop4.\allowbreak{}4PDF8–\allowbreak{}9; Lemmas4.\allowbreak{}5–\allowbreak{}4.\allowbreak{}6PDF9.

\textbf{Inspection:} PrimaryPDF4–\allowbreak{}9 fully read, including whole proofs of exactness, kernel action, continuous-cohomology limits, class/\allowbreak{}transgression comparison and section criterion; images5,8,9 inspected.\allowbreak{} Bibliographic journal details confirmed from publisher.\allowbreak{} Foundational étale-homotopy and Jannsen spectral-sequence references are invoked, not re-proved.\allowbreak{}

\url{https://arxiv.org/pdf/1301.4429}.

\paragraph{S-AG-SHARIFI} R.\allowbreak{}T.\allowbreak{}Sharifi, Relationships between conjectures on the structure of pro-p Galois groups unramified outside p, Proc.\allowbreak{}Sympos.\allowbreak{}PureMath.\allowbreak{}70(2002),275–\allowbreak{}284.

\textbf{Locator:} 13-page author version:\allowbreak{} definitionsPDF1–\allowbreak{}2; Thm1.\allowbreak{}1PDF2/\allowbreak{}proof8–\allowbreak{}9; Thm1.\allowbreak{}3PDF3/\allowbreak{}proof10; Eq2.\allowbreak{}10PDF7; full supporting §§2–\allowbreak{}3PDF3–\allowbreak{}10; Thm4.\allowbreak{}1PDF10–\allowbreak{}11.

\textbf{Inspection:} Primary author text read through the complete §§2–\allowbreak{}3 proofs, including maximal-image element construction, pro-p group lemma, regular-prime theorem and Greenberg obstruction.\allowbreak{} Exact domain and factorial normalization checked.\allowbreak{} Underlying Ihara, Iwasawa and Brown inputs cited rather than independently re-proved.\allowbreak{} Screenshot tool returned text markers only; no visual-inspection claim.\allowbreak{}

\url{https://www.math.ucla.edu/~sharifi/publ/galstr.pdf}.

\paragraph{S-AG-HOSHI} Y.\allowbreak{}Hoshi, On the Kernels of the Pro-l Outer Galois Representations Associated to Hyperbolic Curves over Number Fields, revised March2014 manuscript.

\textbf{Locator:} Definition1.\allowbreak{}8, Remark1.\allowbreak{}8.\allowbreak{}1 and Theorem1.\allowbreak{}9PDF11; references[4]/\allowbreak{}[22]PDF27.

\textbf{Inspection:} Primary exact field definition and regular-prime statement read, including its proof citation to Brown and Sharifi.\allowbreak{} Other Hoshi theorems not being certified here.\allowbreak{}

\url{https://www.kurims.kyoto-u.ac.jp/~yuichiro/rims1782revised.pdf}.

\paragraph{S-AG-BROWN} F.\allowbreak{}Brown, Mixed Tate motives over Z, Ann.\allowbreak{}Math.\allowbreak{}175(2012),949–\allowbreak{}976.

\textbf{Locator:} IntroductionConjecture1printed949/\allowbreak{}PDF1; Theorem1.\allowbreak{}1 and implication to Conjecture1 printed950/\allowbreak{}PDF2.

\textbf{Inspection:} Primary exact statement, category and Deligne–\allowbreak{}Ihara consequence in introduction read.\allowbreak{} The main 28-page proof was not re-refereed; treated as a classical cited theorem.\allowbreak{}

\url{https://annals.math.princeton.edu/wp-content/uploads/annals-v175-n2-p10-p.pdf}.

\paragraph{S-AG-ALPOGE-FABLE} Supplied Alpoge Remark and Fable programme explanation.

\textbf{Locator:} ALPOGE\_\allowbreak{}remark.\allowbreak{}pdf step2PDF3, step17/\allowbreak{}fn6PDF4; OURS\_\allowbreak{}FABLE\_\allowbreak{}PROGRAM\_\allowbreak{}SUGGESTIONS\_\allowbreak{}20261120\_\allowbreak{}EXPLAINED.\allowbreak{}md §§2,5,6; ALPOGE\_\allowbreak{}LEAN\_\allowbreak{}MAP\_\allowbreak{}20261120.\allowbreak{}md step2 row.

\textbf{Inspection:} Local notePDF3–\allowbreak{}4 text and exact step2 formula read; Fable§2 full inherited proof sketch read.\allowbreak{} New AG-FG proofs replace its special-normal-form passage by a direct recursion.\allowbreak{} Added archive locators are explicitly inherited from Fable, not visually rechecked.\allowbreak{} No exact Lean lemma claimed beyond supplied low-confidence map.\allowbreak{}

\url{https://alpo.ge/remark.pdf}.

\paragraph{CX-S-KOIRAN} Pascal Koiran, Hilbert’s Nullstellensatz is in the Polynomial Hierarchy, DIMACS 96-27.

\textbf{Locator:} Section 2 Theorem 2 pp.\allowbreak{}5–\allowbreak{}6; Section 5 Theorem 9 and Corollary 1 printed p.\allowbreak{}12 (PDF13); 8 August 1997 erratum p.\allowbreak{}1.

\textbf{Inspection:} Downloaded both author PS files, converted to PDF, read theorem statements, protocol proof, number-field input, entire erratum.\allowbreak{}

\url{https://perso.ens-lyon.fr/pascal.koiran/Publis/DIMACS.96-27.ps}.

\paragraph{CX-S-BACH} Eric Bach, Explicit bounds for primality testing and related problems, Math.\allowbreak{} Comp.\allowbreak{} 55 (1990), 355–\allowbreak{}380.

\textbf{Locator:} Section 6 Theorems 2–\allowbreak{}3 p.\allowbreak{}372, Miller consequence p.\allowbreak{}373; Section 7 Theorem 4 p.\allowbreak{}375, corollary p.\allowbreak{}376.

\textbf{Inspection:} Original AMS PDF returned 403.\allowbreak{} Exact locators and statements carried from CX\_\allowbreak{}GRH/\allowbreak{}ALPOGE\_\allowbreak{}GCT; not represented as freshly checked original proof.\allowbreak{}

\url{https://doi.org/10.1090/S0025-5718-1990-1023756-8}.

\paragraph{CX-S-SCHOOF} René Schoof, Four primality testing algorithms, MSRI Publications 44 (2008).

\textbf{Locator:} Theorem 2.\allowbreak{}2 pp.\allowbreak{}105–\allowbreak{}106, including full proof specifying quadratic Dirichlet characters.

\textbf{Inspection:} Read statement and full proof via primary PDF text.\allowbreak{} PDF text loses minus signs; transfer uses standard Miller test and the theorem’s stated cutoff.\allowbreak{}

\url{https://pub.math.leidenuniv.nl/~stevenhagenp/ANTproc/05rene.pdf}.

\paragraph{CX-S-AKS} Agrawal, Kayal and Saxena, PRIMES is in P, Ann.\allowbreak{} Math.\allowbreak{} 160 (2004), 781–\allowbreak{}793.

\textbf{Locator:} Theorems 4.\allowbreak{}1 and 5.\allowbreak{}1; 2019 erratum Ann.\allowbreak{} Math.\allowbreak{} 189, 317–\allowbreak{}318.

\textbf{Inspection:} Primary article and erratum publication located; used only for unconditional P classification.\allowbreak{}

\url{https://annals.math.princeton.edu/wp-content/uploads/annals-v160-n2-p12.pdf}.

\paragraph{CX-S-SHOUP} Victor Shoup, Searching for Primitive Roots in Finite Fields, Math.\allowbreak{} Comp.\allowbreak{} 58 (1992), 369–\allowbreak{}380.

\textbf{Locator:} Author version Theorem 3 p.\allowbreak{}4, Section 6 p.\allowbreak{}13; Proposition 9 p.\allowbreak{}10.

\textbf{Inspection:} Downloaded author PDF and visually inspected p.\allowbreak{}4:\allowbreak{} Theorem3 is O(r\textasciicircum{}4(log r+1)\textasciicircum{}4(log p)\textasciicircum{}2), r=omega(p\ensuremath{-}1).\allowbreak{} Text encoding is corrupted; image is authoritative.\allowbreak{} No full original proof check claimed.\allowbreak{}

\url{https://shoup.net/papers/primroots.pdf}.

\paragraph{CX-S-HAC} Menezes, van Oorschot and Vanstone, Handbook of Applied Cryptography (1997).

\textbf{Locator:} Algorithm 3.\allowbreak{}34 and Fact 3.\allowbreak{}35 p.\allowbreak{}100; Fact 3.\allowbreak{}46 (composite square-root/\allowbreak{}factoring relation).

\textbf{Inspection:} Primary author-hosted chapter, exact square-root algorithm and its random step checked.\allowbreak{}

\url{https://cacr.uwaterloo.ca/hac/about/chap3.pdf}.

\paragraph{CX-S-GCTV} Ketan Mulmuley, Geometric Complexity Theory V:\allowbreak{} Efficient Algorithms for Noether Normalization, arXiv:\allowbreak{}1209.\allowbreak{}5993v5.

\textbf{Locator:} Theorem 2.\allowbreak{}10 p.\allowbreak{}25; Theorem 3.\allowbreak{}4 pp.\allowbreak{}27–\allowbreak{}28; Theorem 4.\allowbreak{}1 p.\allowbreak{}30; Theorem 4.\allowbreak{}7/\allowbreak{}Lemma 4.\allowbreak{}8 p.\allowbreak{}33; Theorem 4.\allowbreak{}10 p.\allowbreak{}34.

\textbf{Inspection:} Primary PDF downloaded and relevant statements/\allowbreak{}proof of 3.\allowbreak{}4 read; output and equation-access qualifications checked.\allowbreak{}

\url{https://arxiv.org/pdf/1209.5993}.

\paragraph{CX-S-FS} Michael Forbes and Amir Shpilka, A PSPACE Construction of a Hitting Set for the Closure of Small Algebraic Circuits, ECCC TR17-163.

\textbf{Locator:} Theorem 7.\allowbreak{}1 p.\allowbreak{}21; introduction Theorems 1.\allowbreak{}2–\allowbreak{}1.\allowbreak{}3.

\textbf{Inspection:} Primary theorem and unary-encoding condition checked.\allowbreak{}

\url{https://eccc.weizmann.ac.il/report/2017/163/download/}.

\paragraph{CX-S-GSS} Guo, Saxena and Sinhababu, Algebraic Dependencies and PSPACE Algorithms in Approximative Complexity over Any Field, Theory of Computing 15(16) (2019).

\textbf{Locator:} Theorem1.\allowbreak{}3 p5; Lemma4.\allowbreak{}1/\allowbreak{}Theorem4.\allowbreak{}2 pp15–\allowbreak{}16; Theorems1.\allowbreak{}5–\allowbreak{}1.\allowbreak{}6 p6.

\textbf{Inspection:} Primary journal PDF freshly retrieved via theoryofcomputing.\allowbreak{}org; checked APS statement/\allowbreak{}closure characterization and precise separating-e.\allowbreak{}s.\allowbreak{}o.\allowbreak{}p.\allowbreak{} statements in Theorems1.\allowbreak{}5–\allowbreak{}1.\allowbreak{}6.\allowbreak{}

\url{https://theoryofcomputing.org/articles/v015a016/v015a016.pdf}.

\paragraph{CX-S-IKRS} Ivanyos, Karpinski, Rónyai and Saxena, Trading GRH for algebra:\allowbreak{} algorithms for factoring polynomials and related structures.

\textbf{Locator:} 2008 author version Main Theorem p3, Lemma4.\allowbreak{}2 p14, Proposition4.\allowbreak{}6 pp16–\allowbreak{}17; published authorPDF Theorems1.\allowbreak{}1/\allowbreak{}1.\allowbreak{}3 p3, Lemma2.\allowbreak{}3 pp9–\allowbreak{}10, Lemma4.\allowbreak{}2 p16, Proposition4.\allowbreak{}6 pp18–\allowbreak{}19, Theorem5.\allowbreak{}8 pp28–\allowbreak{}29.

\textbf{Inspection:} Published primary statements freshly checked.\allowbreak{} Complete proofs of Lemma2.\allowbreak{}3, Lemma4.\allowbreak{}2 and Proposition4.\allowbreak{}6 read for the new composition; Theorem5.\allowbreak{}8 exact statement/\allowbreak{}interface and runtime checked, not a fresh whole-paper proof audit.\allowbreak{} Older-version locators retained separately.\allowbreak{}

\url{https://www.cse.iitk.ac.in/users/nitin/papers/nogrh.pdf}.

\paragraph{CX-S-BS} Eric Bach and Jonathan Sorenson, Explicit Bounds for Primes in Residue Classes, Math.\allowbreak{} Comp.\allowbreak{}65 (1996),1717–\allowbreak{}1735.

\textbf{Locator:} Section2 p.\allowbreak{}1720; Theorem5.\allowbreak{}1 p.\allowbreak{}1729.

\textbf{Inspection:} Primary author-uploaded PDF retrieved; statement and precise analytic premise read; geometry agent independently verified same theorem.\allowbreak{}

\url{https://www.researchgate.net/profile/Jonathan-Sorenson-2/publication/220576951_Explicit_bounds_for_primes_in_residue_classes/links/55100ea70cf2ac2905afbceb/Explicit-bounds-for-primes-in-residue-classes.pdf}.

\paragraph{CX-S-MILLER} Gary L.\allowbreak{} Miller, Riemann’s Hypothesis and Tests for Primality, JCSS13(3) (1976),300–\allowbreak{}317.

\textbf{Locator:} Lemma5 and Theorem3 p310; Theorem4 pp311–\allowbreak{}312.

\textbf{Inspection:} Author-hosted primary PDF downloaded; p310 visually inspected and entire Lemma5 proof read.\allowbreak{} \ensuremath{\lambda}\ensuremath{{}^{\prime}} multiple and polynomial bit-growth conditions confirmed; explicit-input version proved here.\allowbreak{}

\url{https://www.cs.cmu.edu/~glmiller/Publications/Papers/Mi76.pdf}.

\paragraph{CX-S-BIMS} Bhargava, Ivanyos, Mittal and Saxena, Irreducibility and deterministic r-th root finding over finite fields, ISSAC2017.

\textbf{Locator:} Corollary1.\allowbreak{}3 authorPDF p2; Theorem3.\allowbreak{}1 and full proof p3; Corollary3.\allowbreak{}2 p4.

\textbf{Inspection:} Statement and complete Theorem3.\allowbreak{}1 proof read; original p3 rendered and inspected.\allowbreak{} Literal nonvanishing step fails for the explicit F7 cubic example CX-ALG-016.\allowbreak{} Corrected basis-search proof CX-ALG-017; even-degree quadratic case independently re-proved CX-ALG-018.\allowbreak{} Corollary conclusions not refuted.\allowbreak{}

\url{https://www.cse.iitk.ac.in/users/nitin/papers/irred-and-rth-root.pdf}.

\paragraph{CX-S-BE} Jean-François Biasse and Muhammed Rashad Erukulangara, A proof of the conjectured run time of the Hafner–\allowbreak{}McCurley class group algorithm, Advances in Mathematics of Communications17(6) (2023),1388–\allowbreak{}1405.

\textbf{Locator:} Author PDF MainTheorem and Conjecture1 p4; Algorithm4 p15; Proposition6 pp15–\allowbreak{}16; Corollary1 p16.

\textbf{Inspection:} Direct PDF request returned503 but web source recovered through indexed reference.\allowbreak{} Main theorem, explicit Hecke-GRH definition, complete Algorithm4, Proposition6 and Corollary1 with its proof read.\allowbreak{} Entire 18-page proof not claimed re-audited; prior CX-N12 detailed transfer retained.\allowbreak{}

\url{https://par.nsf.gov/servlets/purl/10324557}.

\paragraph{CX-S-JMV} David Jao, Stephen D.\allowbreak{} Miller and Ramarathnam Venkatesan, Expander graphs based on GRH with an application to elliptic curve cryptography, J.\allowbreak{} Number Theory129 (2009),1491–\allowbreak{}1504.

\textbf{Locator:} Author arXiv0811.\allowbreak{}0647v2:\allowbreak{} Theorem1.\allowbreak{}1 p3; Remark1.\allowbreak{}2(a)/\allowbreak{}Corollary1.\allowbreak{}3 p4; Theorem1.\allowbreak{}5 p5; Theorem1.\allowbreak{}6 pp6–\allowbreak{}7; Lemma2.\allowbreak{}1 pp7–\allowbreak{}8; proof of Theorem1.\allowbreak{}1 p9 and Cor1.\allowbreak{}3 p10; Proposition3.\allowbreak{}1 p11, Theorem3.\allowbreak{}3 p12.

\textbf{Inspection:} Primary author-hosted PDF read for all listed statements.\allowbreak{} Full proofs of Theorem1.\allowbreak{}1, Cor1.\allowbreak{}3 and Lemma2.\allowbreak{}1 read; verified estimates are character-by-character, enabling the genus quotient refinement.\allowbreak{} Isogeny statements and supplied-solver scope checked; no full underlying CM/\allowbreak{}Deuring proof audit claimed.\allowbreak{}

\url{https://web.williams.edu/Mathematics/sjmiller/public_html/crypto/handouts/JaoMillerVenkatesan_ExpanderGraphsGRHEllCurveCrypto.pdf}.

\paragraph{CX-S-L91} Hendrik W.\allowbreak{} Lenstra Jr.\allowbreak{}, Finding isomorphisms between finite fields, Mathematics of Computation56 (1991),329–\allowbreak{}347.

\textbf{Locator:} Theorems1.\allowbreak{}1/\allowbreak{}1.\allowbreak{}2 p330; subfields §2p332; Theorem3.\allowbreak{}1 pp334–\allowbreak{}336; Proposition4.\allowbreak{}2 p337; Theorems5.\allowbreak{}1/\allowbreak{}5.\allowbreak{}2 pp338–\allowbreak{}340; Theorem7.\allowbreak{}1 p342; Theorem9.\allowbreak{}1 pp344–\allowbreak{}346, Lemma9.\allowbreak{}2 p345.

\textbf{Inspection:} Author-hosted primary PDF downloaded.\allowbreak{} Complete normal-basis proof, prime-degree Teichmüller construction and converse, Artin–\allowbreak{}Schreier proof and final field-construction proof read.\allowbreak{} pp338–\allowbreak{}339 rendered and visually read to recover formula exponents lost in text extraction.\allowbreak{} New application and rewritten full proof in CX-ALG-023/\allowbreak{}026; no claim of new priority for these classical algorithms.\allowbreak{}

\url{https://pub.math.leidenuniv.nl/~lenstrahw/PUBLICATIONS/1991b/art.pdf}.

\paragraph{CX-S-AL86} Leonard M.\allowbreak{} Adleman and Hendrik W.\allowbreak{} Lenstra Jr.\allowbreak{}, Finding irreducible polynomials over finite fields, STOC1986,350–\allowbreak{}355.

\textbf{Locator:} Theorem1 p350; AlgorithmA pp350–\allowbreak{}351; Proposition3 and polynomial-time analysis p351; Lemma6 p352 (Gaussian periods under inertness).

\textbf{Inspection:} Fresh primary author PDF read through AlgorithmA and its complete correctness proof pp350–\allowbreak{}352.\allowbreak{} For prime degree r\ensuremath{\ne}p only steps2–\allowbreak{}3 and Lemma6 are needed; least-prime existence is replaced by the exact CX-ALG-024 input and bound.\allowbreak{} No extra radical or Hecke oracle is hidden in the period construction.\allowbreak{}

\url{https://pub.math.leidenuniv.nl/~lenstrahw/PUBLICATIONS/1986a/art.pdf}.

\paragraph{A-SOURCE} Levent Alpoge, Remark.

\textbf{Locator:} PDF pp3–\allowbreak{}4 steps2–\allowbreak{}20; ALPOGE\_\allowbreak{}Main.\allowbreak{}lean five declarations; ALPOGE\_\allowbreak{}Statements.\allowbreak{}lean definitions.

\textbf{Inspection:} User-supplied PDF text pp1–\allowbreak{}4 read; declarations and definitions read; no independent Lean kernel run.\allowbreak{} Cited theorem input A\_\allowbreak{}D as instructed.\allowbreak{}

\url{https://alpo.ge/remark.pdf}.

\paragraph{M-SOURCE} A.\allowbreak{}M.\allowbreak{} Masullo, The Music of the Primes, supplied 940-page draft.

\textbf{Locator:} §1.\allowbreak{}2 pp12–\allowbreak{}19; Lemma9.\allowbreak{}3.\allowbreak{}2.\allowbreak{}4 p425; Lemma9.\allowbreak{}3.\allowbreak{}2.\allowbreak{}5 pp425–\allowbreak{}426; Theorem9.\allowbreak{}3.\allowbreak{}2.\allowbreak{}6 p426; DefD.\allowbreak{}3.\allowbreak{}3.\allowbreak{}1 p804; Theorem12.\allowbreak{}5.\allowbreak{}1.\allowbreak{}1 pp519–\allowbreak{}520;13.\allowbreak{}3.\allowbreak{}1.\allowbreak{}1 and placeholder13.\allowbreak{}3.\allowbreak{}1.\allowbreak{}2 p522.

\textbf{Inspection:} Original PDF425,426 images inspected; source definitions11,804 read.\allowbreak{} Claims M remain conditional, independent intermediate assertions not inferred from endpoint.\allowbreak{} Original p426 re-inspected for Lemma9.\allowbreak{}3.\allowbreak{}2.\allowbreak{}5 actual ideal versus closure; MT-M05 refutes its quotient claim independently of the local zero-count issue.\allowbreak{}


\paragraph{KW09} Chandrashekhar Khare and Jean-Pierre Wintenberger, Serre's modularity conjecture (I).

\textbf{Locator:} Theorem10.\allowbreak{}1(ii), author PDF20; Cor10.\allowbreak{}2(ii), author PDF21.

\textbf{Inspection:} Primary source theorem and Corollary read, including A5 scope and Artin-L-function discussion.\allowbreak{}

\url{https://www.math.ucla.edu/~shekhar/papers/results.pdf}.

\paragraph{STARK74} Harold M.\allowbreak{} Stark, Some effective cases of the Brauer-Siegel theorem, Invent.\allowbreak{}Math.\allowbreak{}23(1974),135–\allowbreak{}152.

\textbf{Locator:} Theorem3,p144.

\textbf{Inspection:} Exact source locator corroborated by primary later paper Gun-Hazra-Sahu2025 introduction; original volume web fetch timed out.\allowbreak{} Classical theorem statement used; no claim to have read original proof.\allowbreak{}

\url{https://gdz.sub.uni-goettingen.de/download/pdf/PPN356556735_0023/PPN356556735_0023.pdf}.

\paragraph{GHS25} Sanoli Gun, Suhita Hazra and Dhananjaya Sahu, On holomorphy and non-vanishing of Artin L-functions, Monatshefte fur Mathematik206(2025),853–\allowbreak{}883.

\textbf{Locator:} Theorem5 and Corollary6; §6; introduction for Stark statement.

\textbf{Inspection:} Primary statements and proof mechanism of integer-root Euler deformation checked.\allowbreak{} MT-L01/\allowbreak{}L02 gives own proof for general |c|\ensuremath{\le}1 and irrational c.\allowbreak{}

\url{https://link.springer.com/article/10.1007/s00605-025-02060-7}.

\paragraph{GM19-A1} Loic Grenie and Giuseppe Molteni, An explicit Chebotarev density theorem under GRH.

\textbf{Locator:} AppendixA TheoremA.\allowbreak{}1, author PDF26.

\textbf{Inspection:} Number\_\allowbreak{}theory worker verified primary finite-order Hecke zero count; root uses the result as a cited classical input and subtracts consecutive-height counts.\allowbreak{}

\url{https://sites.unimi.it/molteni/research/papers-pdf/44-molteni-An_explicit_Chebotarev_density_theorem_under_GRH.pdf}.

\paragraph{HK20-A1} Peter Humphries and Rizwanur Khan, On the Random Wave Conjecture for Dihedral Maaß Forms, GAFA30(2020),34–\allowbreak{}125.

\textbf{Locator:} AppendixA.\allowbreak{}1 printed106–\allowbreak{}107/\allowbreak{}PDF73–\allowbreak{}74:\allowbreak{} character/\allowbreak{}infinity-type formula, correspondence, spectral parameter, coefficient sum and all finite-place Satake cases; LemmaA.\allowbreak{}1 printed107.

\textbf{Inspection:} Primary publisher HTML and PDF appendix read directly; exact character formula and normalization checked.\allowbreak{} Classical correspondence cited, not re-proved; explicit MT-H01/\allowbreak{}03 proved separately.\allowbreak{}

\url{https://link.springer.com/article/10.1007/s00039-020-00526-4}.

\paragraph{CALEGARI13} Frank Calegari, The Artin conjecture for some S5-extensions, Math.\allowbreak{}Ann.\allowbreak{}356(2013),191–\allowbreak{}207.

\textbf{Locator:} §1 publisher-author PDF2 paragraph before Theorem1.\allowbreak{}1; Theorems1.\allowbreak{}1/\allowbreak{}1.\allowbreak{}2 PDF2–\allowbreak{}3; Theorem5.\allowbreak{}1 PDF15 and Theorem5.\allowbreak{}3 PDF16; arXiv1112.\allowbreak{}1152v1 Theorem1.\allowbreak{}2 p2 for H=EF.

\textbf{Inspection:} Primary author-hosted publisher PDF statements and §1 read; arXiv v1 read to recover explicitly the H=EF definition missing from extracted publisher statement.\allowbreak{} Known automorphy cited, not re-proved; computations not re-run.\allowbreak{}

\url{https://math.uchicago.edu/~fcale/papers/S5.pdf}.

\paragraph{PY07} Dipendra Prasad and C.\allowbreak{}S.\allowbreak{}Yogananda, A Report on Artin’s holomorphy conjecture, February23 2007.

\textbf{Locator:} Theorem3.\allowbreak{}3 PDF5; §3.\allowbreak{}3 PDF8–\allowbreak{}9 completed functional equation and gamma factors.

\textbf{Inspection:} Primary author-hosted exposition read for exact finite Euler definition, Brauer theorem and functional-equation formulas.\allowbreak{} Its pre-2009 status discussion is not used for current modularity scope.\allowbreak{}

\url{https://www.math.iitb.ac.in/~dprasad/artin.pdf}.

\end{document}
